% 本文件由 LaTeX 排版 Agent 根据有效 translation.json 自动生成。
% author=Theodoret; work_id=2707; title=书信集

\chapter{书信集}

% structure: p0001 source=h1 target=omitted reason=page-title-already-rendered-by-parent

第一封信\newline{}第二封信\newline{}第三封信\newline{}第四封信\newline{}第五封信\newline{}第六封信\newline{}第七封信\newline{}第八封信\newline{}第九封信\newline{}第十封信\newline{}第十一封信\newline{}第十二封信\newline{}第十三封信\newline{}第十四封信\newline{}第十五封信\newline{}第十六封信\newline{}第十七封信\newline{}第十八封信\newline{}第十九封信\newline{}第二十封信\newline{}第二十一封信\newline{}第二十二封信\newline{}第二十三封信\newline{}第二十四封信\newline{}第二十五封信\newline{}第二十六封信\newline{}第二十七封信\newline{}第二十八封信\newline{}第二十九封信\newline{}第三十封信\newline{}第三十一封信\newline{}第三十二封信\newline{}第三十三封信\newline{}第三十四封信\newline{}第三十五封信\newline{}第三十六封信\newline{}第三十七封信\newline{}第三十八封信\newline{}第三十九封信\newline{}第四十封信\newline{}第四十一封信\newline{}第四十二封信\newline{}第四十三封信\newline{}第四十四封信\newline{}第四十五封信\newline{}第四十六封信\newline{}第四十七封信\newline{}第四十八封信\newline{}第四十九封信\newline{}第五十封信\newline{}第五十一封信\newline{}第五十二封信\newline{}第五十三封信\newline{}第五十四封信\newline{}第五十五封信\newline{}第五十六封信\newline{}第五十七封信\newline{}第五十八封信\newline{}第五十九封信\newline{}第六十封信\newline{}第六十一封信\newline{}第六十二封信\newline{}第六十三封信\newline{}第六十四封信\newline{}第六十五封信\newline{}第六十六封信\newline{}第六十七封信\newline{}第六十八封信\newline{}第六十九封信\newline{}第七十封信\newline{}第七十一封信\newline{}第七十二封信\newline{}第七十三封信\newline{}第七十四封信\newline{}第七十五封信\newline{}第七十六封信\newline{}第七十七封信\newline{}第七十八封信\newline{}第七十九封信\newline{}第八十封信\newline{}第八十一封信\newline{}第八十二封信\newline{}第八十三封信\newline{}第八十四封信\newline{}第八十五封信\newline{}第八十六封信\newline{}第八十七封信\newline{}第八十八封信\newline{}第八十九封信\newline{}第九十封信\newline{}第九十一封信\newline{}第九十二封信\newline{}第九十三封信\newline{}第九十四封信\newline{}第九十五封信\newline{}第九十六封信\newline{}第九十七封信\newline{}第九十八封信\newline{}第九十九封信\newline{}第一百封信\newline{}第一百零一封信\newline{}第一百零二封信\newline{}第一百零三封信\newline{}第一百零四封信\newline{}第一百零五封信\newline{}第一百零六封信\newline{}第一百零七封信\newline{}第一百零八封信\newline{}第一百零九封信\newline{}第一百一十封信\newline{}第一百一十一封信\newline{}第一百一十二封信\newline{}第一百一十三封信\newline{}第一百一十四封信\newline{}第一百一十五封信\newline{}第一百一十六封信\newline{}第一百一十七封信\newline{}第一百一十八封信\newline{}第一百一十九封信\newline{}第一百二十封信\newline{}第一百二十一封信\newline{}第一百二十二封信\newline{}第一百二十三封信\newline{}第一百二十四封信\newline{}第一百二十五封信\newline{}第一百二十六封信\newline{}第一百二十七封信\newline{}第一百二十八封信\newline{}第一百二十九封信\newline{}第一百三十封信\newline{}第一百三十一封信\newline{}第一百三十二封信\newline{}第一百三十三封信\newline{}第一百三十四封信\newline{}第一百三十五封信\newline{}第一百三十六封信\newline{}第一百三十七封信\newline{}第一百三十八封信\newline{}第一百三十九封信\newline{}第一百四十封信\newline{}第一百四十一封信\newline{}第一百四十二封信\newline{}第一百四十三封信\newline{}第一百四十四封信\newline{}第一百四十五封信\newline{}第一百四十六封信\newline{}第一百四十七封信\newline{}第一百四十八封信\newline{}第一百四十九封信\newline{}第一百五十封信\newline{}第一百五十一封信\newline{}第一百五十二封信\newline{}第一百五十三封信\newline{}第一百五十四封信\newline{}第一百五十五封信\newline{}第一百五十六封信\newline{}第一百五十七封信\newline{}第一百五十八封信\newline{}第一百五十九封信\newline{}第一百六十封信\newline{}第一百六十一封信\newline{}第一百六十二封信\newline{}第一百六十三封信\newline{}第一百六十四封信\newline{}第一百六十五封信\newline{}第一百六十六封信\newline{}第一百六十七封信\newline{}第一百六十八封信\newline{}第一百六十九封信\newline{}第一百七十封信\newline{}第一百七十一封信\newline{}第一百七十二封信\newline{}第一百七十三封信\newline{}第一百七十四封信\newline{}第一百七十五封信\newline{}第一百七十六封信\newline{}第一百七十七封信\newline{}第一百七十八封信\newline{}第一百七十九封信\newline{}第一百八十封信\newline{}第一百八十一封信\par

\SourceRecord{Translated by Blomfield Jackson.；Nicene and Post-Nicene Fathers, Second Series；Vol. 3.；Edited by Philip Schaff and Henry Wace.；Buffalo, NY: Christian Literature Publishing Co.,；1892.；<http://www.newadvent.org/fathers/2707.htm>.}

% structure: p0001 source=h1 target=section reason=page-title

\section{书信一}

\Emph{致一位不知名的收信人。}\par

先知的言语中，我们见到明智的听者与卓越的谋士一同被提及。而我，如今将我所著论及神圣宗徒的书卷呈上，不是奉献给一位明智的听者，而是呈交予一位公正而明察的审判官。当金匠欲检验黄金是否精纯无杂时，便将其置于试金石上；同样，我将书卷呈递于您尊前，为要知道它是否合乎标准，或尚需锤炼。您已阅读并归还，却对此未置一词。您的缄默使我猜想，审判官已作出定罪裁决，却因不愿伤我感受而未曾明言。恳请摒弃此念，勿迟疑，将您对书卷的意见告知于我。\par

\SourceRecord{Translated by Blomfield Jackson.；Nicene and Post-Nicene Fathers, Second Series；Vol. 3.；Edited by Philip Schaff and Henry Wace.；Buffalo, NY: Christian Literature Publishing Co.,；1892.；<http://www.newadvent.org/fathers/2707001.htm>.}

% structure: p0001 source=h1 target=section reason=page-title

\section{第二封信}

\Emph{致同一人。}\par

当人们热情地爱着时，我怀疑他们对于所爱之人的子女能否做到公正的评判。公义被情感所带走。父亲们以为自己的丑儿子是美的，儿子们则看不见父亲的丑陋。兄弟看待兄弟是出于情感而非出于本性。因此我担心，阁下您评判我所写的东西时，也是出于情感的温暖。因为爱的力量确实非常强大，它常常使我们看不见朋友身上的重大过失。正因为您有如此多的爱，亲爱的朋友，您才用善意的赞美来装点我所写的东西。我所做的一切，就是请求您的虔诚，恳求善良的主认可您的颂扬，使您所称赞的人能像其仰慕者言词中所描绘的那样。\par

\SourceRecord{Translated by Blomfield Jackson.；Nicene and Post-Nicene Fathers, Second Series；Vol. 3.；Edited by Philip Schaff and Henry Wace.；Buffalo, NY: Christian Literature Publishing Co.,；1892.；<http://www.newadvent.org/fathers/2707002.htm>.}

% structure: p0001 source=h1 target=section reason=page-title

\section{\WorkTitle{信函三}}

\Emph{致依勒内主教。}\par

神圣的宗徒禁止这种比较。他在\WorkTitle{罗马书}中写道：\Quote{所以，时候未到，什么都不要判断，只等到主来，祂要照亮暗中的隐情，显明人心的意念；那时，各人要从天主那里得著称赞。}他这样做完全正确；因为我们只能看见外在的行为，但万有的天主也知道行者的意念，当祂宣判时，所审判的不在于工作，而在于意志。因此，祂将给那位成为犹太人的犹太人、成为法律下之人的法律下之人、成为无法律之人无法律之人的圣宗徒戴上冠冕，他如此戴上演员的面具，目的是为了造福人类。他的生平并非见风使舵的行径。他所赚得的乃是亏损，但他确保了所教导之人的益处。正如我所说，圣保禄命我们等待天主的审判。但我们正探讨高深的主题；我们处理的是一种超越理解和言语的神学；我们并非像不虔敬的异端者那样寻求亵渎的立场，而是竭力驳斥他们的不虔敬，并尽我们所能赞美造物主；因此，我们尝试回答你的询问，并无不合理之处。\par

你提出了这样一个情况：一位不虔敬的法官向两位虔敬的斗士提供选择，要么向魔鬼祭献，要么投海自尽。你描述了其中一人选择了后者，毫不犹豫地跳入深海；而另一人两者都拒绝，他与同伴一样憎恶崇拜偶像，但拒绝投海，等待这一命运被强行施加。你提出了这些情况，并问这两人中哪一个采取了更好的道路。我想你会同意我的看法：后者更值得称赞。没有人应该擅自离开生命，而应等待自然的或暴力的死亡。我们的主给了我们这个教训，祂吩咐那些在一城受迫害的人逃到另一城，又命令他们离开那城，逃到另一城。神圣的宗徒遵循这一教导，逃脱了该城市总督的暴力，毫不迟疑地讲述他逃走的方式，而是提及篮子、墙壁和窗户，并以此夸耀和光荣。因为看似不光彩的事，因天主的命令而成为光荣的。同样，这位宗徒有时自称法利塞人，有时自称罗马人，并非因为他怕死，而是战斗中非常公正的行为。同样，当他得知犹太人的阴谋时，他向凯撒上诉，并打发他的外甥去千夫长那里报告针对他的计谋，并非因为他贪恋今世，而是服从天主的法律。因为我们的主确实不希望我们置身于明显的危险中；这既是言教也是身教，因为祂多次避开犹太人谋杀的暴力。伟大的伯多禄，宗徒之长，当他被解下锁链，从黑落德手中逃脱后，来到那别号马尔谷的若望的家中，用他的到访消除了朋友的焦虑，并吩咐他们保持沉默，然后躲到另一间屋子，试图通过离开更有效地隐藏自己。我们在\WorkTitle{旧约}中也发现同样的智慧，因为著名的梅瑟在与埃及人斗争中显出勇气，第二天发现凶杀已被知晓，就逃跑，长途跋涉，来到米德扬地。\BibleRef{出 2:11} 同样，伟大的厄里亚得知依则贝耳的威胁后，没有把自己交给那些想要杀他的人，而是离开世界，匆匆进入旷野。如果逃避敌人的暴力是正确而中悦天主的，那么当敌人命令人成为自己的凶手时，拒绝服从他们当然更加正确。我们的主没有顺从魔鬼，当魔鬼叫祂跳下去时；当魔鬼用鞭笞、荆棘和钉子装备犹太人的手来攻击祂时，受造物催促祂使邪恶的敌人遭彻底的毁灭，主自己却禁止了，因为祂知道祂的苦难将为世界带来救恩，正是为此，在祂受难前不久，祂对宗徒们说：\Quote{你们要祈祷，不要进入诱惑}，并教导我们祈祷：\Quote{不要让我们陷入诱惑}。现在让我们稍微转换一下角度，我们将更清楚地看到。让我们从论证中除去海洋，假设法官给了每位殉道者一把剑，命令拒绝祭献的人砍下自己的头；有谁会在理智下忍受用自己的血染红自己的手，成为自己的刽子手，举手攻击自己，服从法官的命令呢？\par

显然，你的第二位殉道者应得更高的赞美。前者确实因他的热心而应得称赞，但后者同样因正确的判断而受尊荣。\par

我照所赐给我的智慧之量答复了你；那位知道思念和作为的主，将在祂显现的日子显明谁是正确的。\par

\SourceRecord{Translated by Blomfield Jackson.；Nicene and Post-Nicene Fathers, Second Series；Vol. 3.；Edited by Philip Schaff and Henry Wace.；Buffalo, NY: Christian Literature Publishing Co.,；1892.；<http://www.newadvent.org/fathers/2707003.htm>.}

% structure: p0001 source=h1 target=section reason=page-title

\section{第四封书信}

\Emph{节日}\par

我们灵魂与肉身的造物主，将祂的恩赐同时赐给了两者，并以可触及心灵与感官的美善同时倾注于我们。在神圣庆节之时，祂赐给我们久盼的雨水，使我们的庆祝免于忧愁。我们已赞美了我们慷慨的主，如今按惯例写一封节日书信，致书于您的虔诚，恳求您以祈祷帮助我们。\par

\SourceRecord{Translated by Blomfield Jackson.；Nicene and Post-Nicene Fathers, Second Series；Vol. 3.；Edited by Philip Schaff and Henry Wace.；Buffalo, NY: Christian Literature Publishing Co.,；1892.；<http://www.newadvent.org/fathers/2707004.htm>.}

% structure: p0001 source=h1 target=section reason=page-title

\section{\WorkTitle{第五封信}}

\Emph{Festal.}\par

那创造我们的天主，在我们犯罪后，赐予我们忧虑与悲伤。但祂也藉着设立神圣的庆节，给我们提供了神圣的慰藉之机。这些庆节所引发的思绪，既提醒我们天主对我们的恩赐，又应许我们从一切苦难中获得完全的自由。享受这些美善，满怀喜乐，我们向阁下的尊威致辞，并依循庆节的惯例，偿还友谊之债。\par

\SourceRecord{Translated by Blomfield Jackson.；Nicene and Post-Nicene Fathers, Second Series；Vol. 3.；Edited by Philip Schaff and Henry Wace.；Buffalo, NY: Christian Literature Publishing Co.,；1892.；<http://www.newadvent.org/fathers/2707005.htm>.}

% structure: p0001 source=h1 target=section reason=page-title

\section{第六封信}

\Emph{节庆的}。\par

我们慈爱的主允许我们，以爱基督者的热忱，庆祝这神圣的救恩庆节，并享用从中流出的属神祝福的果实。既然我们知道阁下对我们怀有的情谊，我们写信告诉您这些。因为那些对他人怀有友善之心的人，总是乐于听到令人欣慰的关于他们的消息。\par

\SourceRecord{Translated by Blomfield Jackson.；Nicene and Post-Nicene Fathers, Second Series；Vol. 3.；Edited by Philip Schaff and Henry Wace.；Buffalo, NY: Christian Literature Publishing Co.,；1892.；<http://www.newadvent.org/fathers/2707006.htm>.}

% structure: p0001 source=h1 target=section reason=page-title

\section{第七封书信}

\Emph{致\Person{德奥尼拉}。}\par

倘若我早听闻您尊贵的夫君去世，我早已写信；如今我写信，并非要用安慰的话来平息您巨大的哀伤。它们并无必要。凡习得哲人智慧、深思此生为何物者，都会发现理性足以迎击并击碎哀伤升起的浪潮。即使您正追忆长久的相伴，理性也认明天主的旨意；为迎战哀伤之泪的力量，它同时调集了自然之途、天主的律法和复活的希望。既知如此，我无需多言。我只恳求您在需要时善用明悟。将逝者的离世视作一次远行，等候我们天主和救主的应许。因为那应许复活者绝不会说谎，祂是真理的泉源。\par

\SourceRecord{Translated by Blomfield Jackson.；Nicene and Post-Nicene Fathers, Second Series；Vol. 3.；Edited by Philip Schaff and Henry Wace.；Buffalo, NY: Christian Literature Publishing Co.,；1892.；<http://www.newadvent.org/fathers/2707007.htm>.}

% structure: p0001 source=h1 target=section reason=page-title

\section{第八封书信}

\Emph{致\Person{尤格拉菲亚}.}\par

我无需再次用精神的安慰来抚平你的悲伤。仅仅提及那为我们成就救恩的苦难，便足以驱散任何深重的哀愁。那一切苦难皆为人类而受。我们的主消灭死亡，并非为使一个身体战胜死亡，而是借着那同一个身体，实现我们共同的复活，并使我们对此复活的望德坚定确实。即使在神圣的庆典中，你的灵魂得到多重慰藉，你仍无法克服悲痛之感；那么，我敬爱的朋友，请允许我恳求你，阅读那份婚姻契约中紧随嫁妆之后的词句，看看婚礼之前是如何以死亡提醒为前奏。我们既知人皆有死，又思考生者的安宁，按惯例会订下所谓的条件，且在婚姻结合之前提及死亡时毫无迟疑。这些明文是：\Quote{若丈夫先死，约定如此这般；若此命运先降临妻子，亦如此这般。}这一切我们在婚礼前早已知道，可谓每日都等待着。为何还要介意呢？结合终必因丈夫的死亡或妻子的离世而破裂。生命的历程如此。你，我卓越的朋友，既深知天主的旨意，又明白人性；请驱散你的沮丧，耐心等候义人共同望德的实现。\par

\SourceRecord{Translated by Blomfield Jackson.；Nicene and Post-Nicene Fathers, Second Series；Vol. 3.；Edited by Philip Schaff and Henry Wace.；Buffalo, NY: Christian Literature Publishing Co.,；1892.；<http://www.newadvent.org/fathers/2707008.htm>.}

% structure: p0001 source=h1 target=section reason=page-title

\section{第九封信}

\Emph{致一位匿名通信者。}\par

你的孝爱之情为我所受的不公审判——未经审讯——感到烦恼与痛苦。你如此感受令我欣慰。若我被公正地定罪，我会因给了审判官们合理的理由而遗憾；但实际情况是，我的良心十分坦然，我反倒感到喜悦与欢欣，并盼望因这不公而获得其他罪的赦免。纳波特之所以活在人记忆中，只因为他遭受了那不义的死亡。只求你祈祷，使我们不被天主遗弃，任凭敌人继续作恶。天主的善愿足以令我非常振奋，若祂在我这一边，我便视一切烦恼为小事。\par

\SourceRecord{Translated by Blomfield Jackson.；Nicene and Post-Nicene Fathers, Second Series；Vol. 3.；Edited by Philip Schaff and Henry Wace.；Buffalo, NY: Christian Literature Publishing Co.,；1892.；<http://www.newadvent.org/fathers/2707009.htm>.}

% structure: p0001 source=h1 target=section reason=page-title

\section{第十封信}

\Emph{致博学之士\Person{埃利亚斯}。}\par

立法者为受压迫者制定了法律，辩护士为需要公正辩护的人运用了演说技艺。而你，我的朋友，既研习了雄辩术也研习了法律。现在，请将你的技艺付诸实践，借此压制压迫者，帮助被他们压制的人，并以法律作为盾牌捍卫他们。不要让任何有罪的当事人因你的辩护而受益，即使他是你的朋友。\par

如今，这些有罪者之一是那个恶棍亚巴郎。他在属于教会的一块田产上定居了相当长一段时间之后，又找了几个同伙一起行恶，并且毫不迟疑地承认了自己的所作所为。我把他连同他的所作所为、他侵害过的当事人以及可敬的副执事\Person{格隆提乌斯}一并交给了你。我不是要你将这个有罪之人送交官府，而是希望当他的受害者们向你诉说了他们不得不忍受的一切，并让你——我博学的朋友——对他们的境况感到同情时，你或许会被促使去强迫那个邪恶的家伙归还他所偷窃的东西。\par

\SourceRecord{Translated by Blomfield Jackson.；Nicene and Post-Nicene Fathers, Second Series；Vol. 3.；Edited by Philip Schaff and Henry Wace.；Buffalo, NY: Christian Literature Publishing Co.,；1892.；<http://www.newadvent.org/fathers/2707010.htm>.}

% structure: p0001 source=h1 target=section reason=page-title

\section{\WorkTitle{第十一書}}

\Emph{致\Place{君士坦丁堡}主教\Person{弗拉維亞諾}}\par

宇宙的創造者與引導者使祢成為世界的光明，將深沈無月的黑夜變為明朗的正午。正如海港旁的信標燈在夜間為水手指明港灣的入口，祢聖德的明亮光輝為所有因真道而受攻擊者帶來極大的安慰，並指示他們\OriginalTerm{宗徒信德}{Apostles’ faith}的安全港。那些已經認識這信德的人充滿安慰，而那些不認識的人則得以免於撞上礁石。我尤其應當讚美眾善的施予者，因為我找到了一位高貴的戰士，他以敬畏天主的力量驅散對人的恐懼，衷心地為福音的教義在前線奮戰，並欣然承受宗徒戰爭的衝擊。所以，今日眾口皆在讚頌祢的聖德，因為不僅真道的養育者欽佩祢信德的純潔，連真理的仇敵也歌頌祢的勇氣。謊言在真理的閃電下消失。\par

我這樣寫，因為我知道極可敬且虔誠的讀經員\Person{希帕提烏斯}既樂於遵從祢聖德的命令，又不斷地，我主，提及祢可讚揚的事蹟。我向祢致敬，聖善而極蒙主愛。我懇求祢以祈禱支持我們，使我們能在天主的律法下度過餘生。\par

\SourceRecord{Translated by Blomfield Jackson.；Nicene and Post-Nicene Fathers, Second Series；Vol. 3.；Edited by Philip Schaff and Henry Wace.；Buffalo, NY: Christian Literature Publishing Co.,；1892.；<http://www.newadvent.org/fathers/2707011.htm>.}

% structure: p0001 source=h1 target=section reason=page-title

\section{\WorkTitle{第十二封书信}}

\Emph{致\Person{依勒内}主教。}\par

\Person{约伯}，那座著名的金刚石之塔和美善的高贵斗士，甚至没有被各样各类连续苦难的打击所动摇，而是屹立不摇，坚固不移。然而，在他所有考验的末了，正义的立法者用这些话解释了它们的原因：\Quote{你认为我回答你是出于别的理由，而不是为了使你显得正义吗？} 我认为这些话为你的虔诚是众所周知的，你的虔诚能够承受众多各样麻烦和焦虑的攻击，并且非但没有退缩，反而显示出你管理的坚毅与稳固。因此，慷慨的主，看到你灵魂的英勇与圣善，不愿将一位可敬的斗士隐藏起来，而是将他带到战斗面前，用胜利之冠装饰你可敬的头，并将你的奋斗作为优良服务的崇高榜样赐予众人。所以，我亲爱的朋友，也在这场战斗中得胜，勇敢承受你女婿——我亲爱的朋友——的去世。以智慧战胜亲情的索取，以及对高贵慷慨品格的记忆，这份记忆总会回想起超越画家技艺或修辞家才能的东西。以那以智慧管理人类一切事务者的思想来击退悲伤的冲击，祂完全知晓未来，并正确地引导一切为我们获益。让我们与被解救脱离今生风暴的人一同喜乐。我们更应感恩，因为被和风护送，他已抛锚在无风之港，逃离了此生充斥的悲惨船难。但我何必对一位久经考验的美善角斗士说这一切？何必为坚忍而傅油一位其他运动员的教练？我仍然写信。写信对我自己是一种安慰。当我记起我如此珍视的情谊时，我实在由衷悲伤。我再次赞美万有的伟大向导，祂既知道什么对我们有益，也相应地引导我们的人生。我在写完前一封信后口述了这封信，当时我在\Place{安提约基雅}的一位朋友告诉我结局已经到来。\par

\SourceRecord{Translated by Blomfield Jackson.；Nicene and Post-Nicene Fathers, Second Series；Vol. 3.；Edited by Philip Schaff and Henry Wace.；Buffalo, NY: Christian Literature Publishing Co.,；1892.；<http://www.newadvent.org/fathers/2707012.htm>.}

% structure: p0001 source=h1 target=section reason=page-title

\section{第十三封书信}

致\Person{西流}。\par

我早已听说\Place{莱斯博斯}岛及其城市\Place{米蒂利尼}、\Place{麦西姆纳}等，但对其所产的葡萄之果实却一无所知。如今，因着您的勤勉，我得以了解它，并赞叹其色泽的洁白与味道的甘美。或许时间甚至能令它愈发醇厚，除非它反使其变酸；因为酒，如同身体、植物、建筑以及其他手工之物一样，会受时间损害。若如您所言，它能使饮者长寿，我恐怕它对我无甚益处，因为当生命的风暴如此众多而严酷时，我并无意愿活得长久。\par

然而，听闻隐修士的健康状况，我深感欣慰。我确实曾为他焦虑不堪，且错怪了医生们，因为他的病症正需要他们所施的治疗。我已托人给您送去一小罐蜜，那是\Place{基利家}的蜜蜂以苏合香花所酿制的。\par

\SourceRecord{Translated by Blomfield Jackson.；Nicene and Post-Nicene Fathers, Second Series；Vol. 3.；Edited by Philip Schaff and Henry Wace.；Buffalo, NY: Christian Literature Publishing Co.,；1892.；<http://www.newadvent.org/fathers/2707013.htm>.}

% structure: p0001 source=h1 target=section reason=page-title

\section{第十四封书信}

\Emph{\Person{亚历山德拉}}\par

如果我仅仅考虑你所遭受的损失的性质，我自己就本应需要安慰，不仅因为我认为凡是关乎你的事也关乎我，无论是喜是悲，而且还因为我那样深深地爱着那位可敬的、真正卓越的人。但是，天主的旨意已将他从我们这里带走，并把他迁移到更美好的生命中去。因此，我将忧伤的阴云从我的灵魂中驱散，并恳请你，我尊贵的朋友，以理性的力量战胜你悲伤的痛苦，并在这需要时刻，让你的灵魂受到天主圣言的感召。我们为何从摇篮起就如从乳房吸吮乳汁一样吸取圣经的教导？岂不是为了当患难降临时，我们能将圣神的训诲当作药膏来医治我们的疼痛？我知道，当一个人体验过所爱之物的价值后，突然被剥夺它，并在一瞬间从幸福跌入痛苦，是多么悲伤、多么痛苦啊。但对于那些具有良好理智并运用正确理性能力的人而言，没有什么是完全出乎意料的；人间之事没有什么是稳固的，没有什么是持久的——无论是美貌、财富、健康、尊严，还是人们视为最高的一切。有些人从富贵的顶峰跌入最底层的贫穷；有些人失去健康，与各种疾病搏斗；有些人为自己家世的光辉而骄傲，却被拖入奴役的沉重桎梏。美貌被疾病损坏，被衰老毁容，至高主宰明智地不让这些事物持久或长存，意在使拥有者因惧怕变化而降低傲慢的目光，并认识到这一切所有物都如潮水般涨落，从而不再信赖那些短促易逝之物，而将希望寄托于一切美善的施予者。我深知，我卓越的朋友，你懂得这一切，我恳请你思考人的本性：你会发现人是必死的，从一开始就接受了死亡的判决。正是天主对亚当说的: \BibleQuote{你是尘土，仍要归于尘土。} 赐予法律者是从不撒谎的，经验也证明了祂的真实。圣经告诉我们: \BibleQuote{众人进入生命的门径相同，离世的途径也一样。} 每个人出生后都等待着坟墓。而且并非所有人都活同样的时长：有些人过早地结束生命，有些人在壮年，有些人则在经历衰老的考验之后。同样，那些缔结婚姻羁绊的人也要从中解脱，必定是丈夫先离去，或是妻子先到达人生的终点。有些人刚刚进入洞房，就遭遇哭泣和哀悼；有些人只一起生活了短暂的时间。只要记着这悲伤是共同的，就足以给理性提供克服悲伤的理由。除此之外，即使被最痛苦的悲伤所征服的人，也可以因想到亡者留下了儿子而得到安慰；他留下了已成年的孩子；他曾达到非常高的地位，且在此地位上非但没有引起嫉妒，反而使人更加爱戴他；他身后留下了慷慨、嫉恶如仇、温和以及各种美德的声誉。\par

但是，如果我们记起天主自己的应许和基督徒的望德——我是说复活、永生、在天国里永存，以及\BibleQuote{眼所未见，耳所未闻，人心所未想到的，就是天主为爱祂的人所准备的}——那么我们还有什么理由沮丧呢？宗徒不是强调地说: \BibleQuote{弟兄们，关于长眠的人，我们不愿意你们不知道，免得你们忧伤，像那些没有望德的人一样。} 我认识许多人，他们即使没有希望，单凭理性的力量也能克服悲伤；那么，被如此希望支持的人，若竟比毫无希望的人更软弱，岂非反常？因此，我求你，将死亡看作一次长途旅行。当他外出旅行时，我们固然会难过，但我们等待他的归来。现在，让分离确实使我们有些悲伤——因为我并非劝你违背人性——但不要让我们像对着尸体一样哀号；不如让我们祝贺他的出发和离去，因为他如今已从充满无常的世界中获得自由，不再惧怕灵魂、身体或肉体状况的任何改变。战斗现已结束，他在等待奖赏。不要过于为孤寡而悲伤。我们有一位更大的保护者，祂的律法要求所有人都应好好照顾孤儿和寡妇，关于祂，圣达味说: \BibleQuote{上主扶持孤儿和寡妇，但将恶人的道路倾覆。} 只要我们生命的舵交在祂手中，我们就会遇到从不失落的眷顾。祂的保护比任何人的保护都更可靠，因为祂的话是: \BibleQuote{一个女人岂能忘记她吃奶的婴儿，而不怜悯她腹中的儿子？纵然她能忘记，我也不会忘记你。} 祂比父母更亲近我们，因为祂是我们的创造者和造物主。婚姻并不能造就父亲，而是按照祂的旨意使人成为父亲。\par

我现在这样写是出于不得已，因为我的桎梏不容我赶赴你处；但是，你最敬爱天主、至圣洁的主教，能独自以言语和行为、以见面和思想的交流，以及他那属灵的、天主赐予的智慧，给你那非常忠信的灵魂带来一切安慰——我相信，藉着这安慰，你悲伤的风暴将会平息。\par

\SourceRecord{Translated by Blomfield Jackson.；Nicene and Post-Nicene Fathers, Second Series；Vol. 3.；Edited by Philip Schaff and Henry Wace.；Buffalo, NY: Christian Literature Publishing Co.,；1892.；<http://www.newadvent.org/fathers/2707014.htm>.}

% structure: p0001 source=h1 target=section reason=page-title

\section{第十五封书信}

\Emph{致首席主教\Person{西尔瓦诺}。}\par

我知道，在安慰的话语上，我有些迟缓，但我并非无端拖延，因为我以为值得让您悲伤的剧烈情绪自然发泄。最高明的医生决不会在热病最炽烈时施以治疗，而是等待有利时机运用其医术。因此，在估量您的痛苦必定何等尖锐之后，我让这几天过去，因为如果我本人听到消息时都如此忧虑、充满悲伤，那么一位丈夫和伴侣——圣经上说，他们已成为一体——在因时间和爱情而紧密结合的联合被粗暴撕裂时，又会遭受何等的痛苦呢？这样的苦楚本是自然；但让理性以这样的提醒来施以安慰：人性脆弱，悲伤是普遍的，还有复活的希望和那明智地管理我们生活的天主的旨意。我们必须接受那不可估量的智慧的谕令，并承认它们对我们有益；因为那些如此虔诚地反思的人，将收获虔诚的赏报，从而脱离过度的哀叹，在平安中度日。另一方面，那些被悲伤奴役的人，哭泣毫无益处，反而会过疲惫的生活，并使我们的守护者忧伤。那么，我至为尊敬的朋友，请接受一个父亲般的劝勉：\BibleQuote{主赐予，主收回。祂做了祂所喜悦的一切。愿主的名受赞美。}\par

\SourceRecord{Translated by Blomfield Jackson.；Nicene and Post-Nicene Fathers, Second Series；Vol. 3.；Edited by Philip Schaff and Henry Wace.；Buffalo, NY: Christian Literature Publishing Co.,；1892.；<http://www.newadvent.org/fathers/2707015.htm>.}

% structure: p0001 source=h1 target=section reason=page-title

\section{书信第十六}

\Emph{致伊肋内主教。}\par

看来，我们面前毫无善事可言，远非风平浪静，扰乱教会的风暴反而日甚一日。会议召集人已抵达，并将传召信交给了数位都主教，包括我们的都主教，我已将信函副本寄送尊座，以便让您知道，正如诗人所言：\Quote{祸患由祸患焊接而成。} 我们唯赖主的良善方能平息风暴。平息风暴对祂而言轻而易举，但我们不配得到平静，然而祂忍耐的恩宠为我们足够，好使我们藉此或许胜过敌人。诚如神圣宗徒教导我们祈祷：\BibleQuote{祂必在试炼时也开一条出路，叫你们能忍受得住。}\BibleRef{格前 10:13} 但我恳求您的虔敬堵住反对者之口，使他们明白：那些正如俗语所说\Quote{站在射程之外}的人，无权嘲笑在队列中战斗、给予并承受打击的人；因为士兵用什么武器击倒敌手，这有什么关系？即使伟大的达味击杀外邦勇士时，也未使用全副武装；三松用驴腮骨一日击杀千人。没有人对胜利发怨言，也没有人指责征服者胆怯，只因他赢得胜利时没有挥舞长矛、没有用盾牌护身、没有投掷标枪、没有射箭。捍卫真宗教的人也应受到同样的评判，我们不可寻求挑起纷争的言辞，而应寻找清楚宣告真理、使胆敢反对者自觉羞愧的论证。\par

我们将至圣童贞同时称为人之母与天主之母，或者称她为其后裔的母亲和婢女，加上她是作为人的主耶稣基督的母亲、但作为天主而言是祂的婢女，从而既避免了那成为诽谤借口的术语，又用另一短语表达了相同的意见，这有何妨碍？此外还必须记住：前一个称谓是通用的，后一个则专属于童贞女；而所有争论都由此而起——但愿从未发生过。大多数古代教父已将更光荣的称谓应用于童贞女，正如尊座本人也曾在两三篇讲道中所做；您虔敬寄给我的几篇讲道，我亲自保存着，在这些讲道中您并未将\Quote{人之母}与\Quote{天主之母}连用，而是用其他词语解释了其含义。但既然您责备我在我的权威名单中遗漏了圣善可颂的教父狄奥多若和狄奥多勒，我认为有必要对此补充几句。\par

首先，亲爱的朋友，我遗漏了许多其他著名且卓越的教父。其次，必须记住这一事实：被指控的一方有责任提出无可指责的证人，连指控者也无法质疑他们的证词。但如果被告传唤被控方指控的权威，连法官本人也不会同意采纳。如果我在编纂教父颂词时遗漏了这些圣人，我承认我是错误的，并且证明我对我的老师忘恩负义。但如果我在受指控时提出了辩护，并提出了无可指责的证人，那么那些不愿看到任何这些证词的人为何要对我施加不合理的责备？我对这些作者的敬仰已由我为他们所写的辩护书充分证明，在那部书中，我驳斥了对他们的指控，毫不畏惧指控者的权势，甚至不畏惧对我自身的暗中攻击。这些如此喜爱空谈的人最好为他们的言辞诡计另找借口。我的目标不是使我的言行迎合这个或那个人的喜好，而是建立天主的教会，并使她的新郎和主喜悦。我呼求我的良心作证：我并非出于对物质事物的挂虑，也不是因为我贪恋那带有种种烦恼的尊荣——我甚至不愿称之为不幸的尊荣——才如此行事。若非害怕天主的审判，我早就自愿退隐了。现在你们清楚地知道，我在等待我的结局。而且我认为它正在临近，因为针对我的阴谋表明了这一点。\par

\SourceRecord{Translated by Blomfield Jackson.；Nicene and Post-Nicene Fathers, Second Series；Vol. 3.；Edited by Philip Schaff and Henry Wace.；Buffalo, NY: Christian Literature Publishing Co.,；1892.；<http://www.newadvent.org/fathers/2707016.htm>.}

% structure: p0001 source=h1 target=section reason=page-title

\section{第十七封书信}

\Emph{致女执事卡西安娜}\par

若我仅考虑你忧伤之深，本应暂缓写信，好让时间助我疗治此痛；但我深知你虔敬的明智，故敢冒昧献上几句安慰的话，这些话部分出自人的本性，部分出自神圣的圣经。因为我们的本性脆弱，而整个生命充满了这样的灾祸；统驭并治理世界的全能主宰——那明智安排我们万事的主——藉祂神圣的谕言赐予我们各种安慰，圣史的著作和蒙福先知们的神圣话语都充满了这些安慰。但我确信，无需为你摘引这些段落并呈于你的虔敬，因你自幼即受默感的圣言养育，依此治理生命，无需其他教导。然而我恳求你记念那些教训，命我们克制情感、应许永生、宣告死亡之毁灭并昭示众人共同复活的言语。除此一切，不，在这一切之上，我求你思量：命定这些事如此成就的是主，祂是全然智慧、全然良善的主，深知什么于我们最为有益，并为此引导我们全部的生命。有时死亡胜于生命，看似痛苦的事反比虚幻的快乐更愉悦。我恳求你的虔敬接纳我谦卑所提供的安慰，好使你以高贵的忍耐服侍万有的主，并为男女众人树立真智慧的榜样。因为众人都将钦佩那份坚强的心志，它勇敢承受了忧伤的攻击，并以决心的宽宏击溃其猛烈的冲击。我们亦非缺乏极大的安慰：在你离世之子活生生的肖像中，他留下了令人深爱的后裔，或能遏止我们过度的悲伤。\par

最后，我恳求你在忧伤中记念你身体的软弱所能承受的限度，避免因过度哀悼而加重痛苦；并恳求我们的主以祂无限的资源赐予你安慰的理由。\par

\SourceRecord{Translated by Blomfield Jackson.；Nicene and Post-Nicene Fathers, Second Series；Vol. 3.；Edited by Philip Schaff and Henry Wace.；Buffalo, NY: Christian Literature Publishing Co.,；1892.；<http://www.newadvent.org/fathers/2707017.htm>.}

% structure: p0001 source=h1 target=section reason=page-title

\section{第十八封信}

\Emph{致\Person{内奥普托勒摩斯}}\par

每当我将目光投于那称因婚姻结合者为\Quote{一体}的神圣法律时，我不知如何安慰那被分离的肢体，因为我深知那痛苦的巨大。但当我思及自然的历程，以及造物主在\Quote{你是尘土，仍要归于尘土}这句话中所设立的法则，并思及世上每日在陆地和海洋上所发生的一切——无论是丈夫先走近生命的终点，还是妻子先遭遇此命运——我从这些思考中找到许多慰藉的理由；而最首要的是我们的主和救主所赐予我们的望德。因为降生成人奥迹实现的理由，乃在于使我们既得知死亡的失败，便不再为所爱者的亡故而过度哀伤，而是期待那复活之望德的渴望实现。我恳请阁下思量这些事，并克服你悲伤的痛苦；尤其因为你们共同爱情的子女与你同在，给你一切安慰的理由。让我们赞美那以智慧治理我们生命的主，不要因无节制的哀悼而激怒祂，因为祂以智慧知道什么对我们有益，并以慈悲赐予我们。\par

\SourceRecord{Translated by Blomfield Jackson.；Nicene and Post-Nicene Fathers, Second Series；Vol. 3.；Edited by Philip Schaff and Henry Wace.；Buffalo, NY: Christian Literature Publishing Co.,；1892.；<http://www.newadvent.org/fathers/2707018.htm>.}

% structure: p0001 source=h1 target=section reason=page-title

\section{书信十九}

\Emph{致司铎\Person{巴西略}。}\par

我发觉那位十分雄辩的演说家\Person{阿塔纳修}，正如你的信中所描述的那样。他的口才因他的言辞而增色，他的言辞因他的品格而增色，他的一切因他丰盛的信德而光辉灿烂。至蒙主爱之友，请常赐予我们这样的礼物。你确知，因我与他交往，你给了我极大的喜乐。\par

\SourceRecord{Translated by Blomfield Jackson.；Nicene and Post-Nicene Fathers, Second Series；Vol. 3.；Edited by Philip Schaff and Henry Wace.；Buffalo, NY: Christian Literature Publishing Co.,；1892.；<http://www.newadvent.org/fathers/2707019.htm>.}

% structure: p0001 source=h1 target=section reason=page-title

\section{第二十封书信}

\Emph{致司铎\Person{玛尔提里}。}\par

在性格决意之前，自然性情已显现在我们面前，就此而言，性情占先；但性情被决意克服，极善言辞的演说家\Person{亚大纳削}清楚证明了这一点。他虽是埃及人，却没有埃及人那种缺乏自制，反而表现出一种为温和所调节的性格。他更是神圣事物的热心爱好者。因此，他与我共度多日，期望从停留中获益。但如你所知，最敬天主的挚友，我羞于尝试这样从他人获益，且远不能将其给予寻求它的人，这不是因为我吝啬，而是因为我没有可资给予的东西。因此，愿你圣善祈祷，使人们对我的称赞能被事实证实，不仅言语上报告我的好事，而且行动上证明。\par

\SourceRecord{Translated by Blomfield Jackson.；Nicene and Post-Nicene Fathers, Second Series；Vol. 3.；Edited by Philip Schaff and Henry Wace.；Buffalo, NY: Christian Literature Publishing Co.,；1892.；<http://www.newadvent.org/fathers/2707020.htm>.}

% structure: p0001 source=h1 target=section reason=page-title

\section{第二十一書信}

\Emph{致博學的\Person{厄色伯}}。\par

傳播這重大消息的人，以爲這會使我極其不快，幻想他們可用這種方式煩擾我。然而靠天主的恩寵，我歡迎這消息，並愉快地等待結果。事實上，任何因神聖教義而臨到我身上的困苦，都令我萬分感激。因爲，如果我們真正信賴主的應許：\BibleQuote{現今的苦楚，與將來在我們身上要顯示的光榮，是不能較量的。}\par

何況我何必談論那所盼望的美好事物的享受呢？因爲，即使沒有賞報提供給爲真道而奮鬥的人，僅真理本身以其自有的力量，就足以說服那些愛慕她的人樂意爲她承受一切危險。神聖的宗徒見證了我所說的，他呼喊說：\BibleQuote{誰能使我們與基督的愛隔絕呢？是困苦嗎？是窘迫嗎？是迫害嗎？是飢餓嗎？是赤身露體嗎？是危險嗎？是刀劍嗎？正如經上所記載：\InnerQuote{我們爲了你的緣故終日被殺；人看我們如同將宰的羊。}}\par

然後，爲教導我們他並不尋求賞報，而只愛他的救主，他立刻加上說：\BibleQuote{然而，靠著那愛我們的主，我們在這一切事上大獲全勝。}\par

他繼續更清楚地顯示自己的愛：\BibleQuote{因爲我深信：無論是死，是生，是天使，是掌權者，是現今的事，是將來的事，是權能者，是崇高，是深淵，或是其他任何受造之物，都不能使我們與天主的愛隔絕，這愛是在我們的主基督耶穌內的。}\par

請看，我的朋友，宗徒愛情的火焰；請看愛德的火炬。\par

他說：我不貪圖祂的什麼，我只渴慕祂本人；我的這份愛是無法熄滅的愛，我樂意放棄一切現在和將來的幸福，是的，寧可受苦並再次忍受各種痛苦，也要使我心中這火焰保持全部力量。這在神聖的作者身上，不僅以言語，而且以行動，在陸地和海洋各處留下了他受苦的紀念。因此，當我注目於他，以及其他族長、先知、宗徒、殉道者、司鐸時，凡通常被視爲悲慘之事，我不得不視爲喜樂。我承認感到羞愧，因爲想起連那些從未學過我們所學的道理，而只以人性爲指導的人，在德行的競賽中竟然佔據了顯著的位置。著名的\Person{蘇格拉底}，\Person{索弗羅尼斯庫斯}之子，在面對誹謗的指控時，不僅蔑視控告者的謊言，而且在苦難中表達了他的愉快，他說：\Quote{\InnerQuote{\Person{阿尼托斯}和\Person{梅勒托斯}可以殺死我，但不能傷害我。}}而\Place{佩尼亞}的演說家，他既智慧又雄辯，用以下話語豐富了他當代和後世的人：\Quote{\InnerQuote{對全人類而言，生命的終結就是死亡，即使一個人爲安全而將自己關在小室內；所以善人應當永遠致力於每一件光榮的事，以美好的希望如同盾牌保衛自己，並勇敢地承受天主所賜予的任何命運。}}\par

此外，比\Person{德摩斯梯尼}更早的作者，即\Person{奧洛羅斯}之子，寫了許多高尚的語句，其中一句是：\Quote{\InnerQuote{我們必須勇敢地忍受諸神所注定的和戰爭的命運。}}何需引用哲學家、歷史學家和演說家呢？因爲連那些更重視神話而非真理的人，也在他們的故事中插入許多有益的勸誡；正如\Person{荷馬}在他的詩篇中介紹希臘人中最有智慧者準備自己從事英勇的行爲，他說道：\par

他責備自己憤怒的心靈，捶著胸膛，\newline{}說道：\newline{}\Quote{忍住吧，我的心靈，想著這一點：\newline{}曾有過更痛苦的煎熬測試過你的忍耐。}\par

類似段落可輕易從詩人、演說家和哲學家那裏收集，但對我們而言，聖經已經足夠。\par

我引用這些，是爲了證明：如果僅僅追隨本性的人勝過了我們這些受過先知和宗徒教導、信賴救主的苦難、期待身體的復活、脫離敗壞、獲得不朽的恩賜和天國的人，那是多麼可恥。\par

所以，我親愛的朋友，安慰那些因流傳的謠言而氣餒的人；如果有人對它們感到高興，就告訴他們我們也快樂，我們正歡欣踴躍，他們所謂的懲罰，我們卻視爲天國本身。\par

为使不了解我们心意的人得知，最卓越的朋友，请放心，我们相信，正如我们所受教的，相信圣父、圣子和圣神。有些人诽谤说我们受教相信、或受洗、或相信、或教导别人相信两个儿子，这绝非事实。正如我们知道一位圣父和一位圣神，同样我们知道一位圣子，我们的主耶稣基督，天主的独生子，降生成人的天主圣言。然而，我们不否认各性体的属性。我们认为那些将一位主耶稣基督分裂为两个儿子的人是错误的，我们也称那些企图混淆性体的人为真理的敌人。我们相信一种无混淆的结合，并认为某些属性属于人性，某些属于天主性；因为正如人——我指一般的人——是有理智和会死的存有，有灵魂和身体，被认为是一个存有，同样，两个性体之间的区别并不将一个人分成两个位格，而是在一个人身上同时承认灵魂的不朽和身体的必死，并承认不可见的灵魂和可见的身体，但如前所述，是一个同时有理性且会死的存有；同样，我们承认我们的主和天主，即我们的主基督、天主子，即使在降生成人之后，仍是一位儿子；因为结合是不可分的，我们深知它是无混淆的。我们也承认，天主性是无始的，而人性是晚近才有的；因为一个性体来自亚巴郎和达味的后裔，童贞玛利亚就是从他们而出，但天主性是由天主圣父在万世之前、无时间、无情感、无分离地生的。但是，假如肉体与天主性之间的区别被摧毁，我们将用什么武器与亚略和欧诺米作战？我们如何解除他们对独生子的亵渎？实际上，我们将谦卑的言辞应用于人，将尊崇和天主的言辞应用于天主，真理的阐述对我们而言非常容易。\par

但这一关于信仰的讨论已超出信函的限度。然而，即使这几句话也足以显示宗徒信仰的特征。\par

\SourceRecord{Translated by Blomfield Jackson.；Nicene and Post-Nicene Fathers, Second Series；Vol. 3.；Edited by Philip Schaff and Henry Wace.；Buffalo, NY: Christian Literature Publishing Co.,；1892.；<http://www.newadvent.org/fathers/2707021.htm>.}

% structure: p0001 source=h1 target=section reason=page-title

\section{第二十二书函}

致乌尔皮亚努斯伯爵\par

俗话说，人行为中的缺失可用言辞加以纠正和改善。然而我认为，天性优美之人，其自身即使无需言辞，亦能使言语悦目，正如天生丽质无需脂粉。这些品质在口才出众的演说家阿塔纳修斯身上显而易见，我对他尤为满意，因为他热切敬爱阁下，且不断称扬您的德范。然而在此我与他争胜，历数您的高贵品质，胜过了他，因为我对您的善行所知比他更多。但我仍感烦恼，不能一一赞扬，且自认对您美德的概述有负于您当受的称誉；不过，若蒙天主恩赐，即使仅能接近真相，您也将在所有同时代人中，在每一种美德上卓然超群。\par

\SourceRecord{Translated by Blomfield Jackson.；Nicene and Post-Nicene Fathers, Second Series；Vol. 3.；Edited by Philip Schaff and Henry Wace.；Buffalo, NY: Christian Literature Publishing Co.,；1892.；<http://www.newadvent.org/fathers/2707022.htm>.}

% structure: p0001 source=h1 target=section reason=page-title

\section{书信第廿三封}

\newline{}\newline{}\Emph{致贵族阿雷奥宾达斯}\par

\newline{}\newline{}万有的创造者与治理者在人中间分配财富与贫穷时，并未作出不公的审判，而是将穷人的贫穷赐给富人，作为行善的手段。同样，祂施惩戒于人，不仅是为惩罚他们的过犯，也是为给富人提供向人类显示仁慈的机会。今年，主赐予我们的鞭挞，远不及我们的罪过所应得的，却足以使农夫困苦；关于他们的苦难，我近日已通过您自己的仆从告知了您的尊贵。我恳求您怜悯这些田地耕作者，他们辛苦劳作却收获甚微。愿这歉收之年成为属灵丰盈的提示，你们通过施行怜悯，聚集天主的怜悯之收获。为此，卓越的狄约尼削已赶往您的尊前，告知这困境，以期获得解救。他携带此信，如同求援的橄榄枝，希望藉此能得到更仁慈的对待。\par

\SourceRecord{Translated by Blomfield Jackson.；Nicene and Post-Nicene Fathers, Second Series；Vol. 3.；Edited by Philip Schaff and Henry Wace.；Buffalo, NY: Christian Literature Publishing Co.,；1892.；<http://www.newadvent.org/fathers/2707023.htm>.}

% structure: p0001 source=h1 target=section reason=page-title

\section{第二十四函}

\Emph{致\Place{萨莫撒塔}主教\Person{安德肋}书}\par

阁下之虔诚，天主之爱的乳儿，我确信渴望与我来往。但我更渴慕与您来往，因为我知道由此我将获得更多益处。匮乏在某种程度上自然使我们的愿望更加强烈，但是万有的主能够赐予我们所渴望的。祂亲自治理万有；知道什么必定对我们有益，并且从不停止将这一恩赐赐予每个人。我实在无法告诉您，您的来信令我多么欣喜，而极其可敬虔诚的执事\Person{塔拉修斯}通过告诉我我非常想知道的事情，更增加了我的喜乐，因为还有什么比得知您一切安好的消息更令我欢迎的呢？又有什么比我们中间伟大人物的节制更能增加您的福祉呢？您行事如同一位明智而活跃的医生，不等病人来请，便主动前来照顾那些需要他照料的人。这令我极为欣喜，并且我通过自身经验领悟了诗人所说的\Quote{含泪欢笑}的含义。愿丰厚赐予万善的天主赐予您的圣德在这些事上卓越超群，并使我们效法所有善人身上可赞美之处。那么，我亲爱的朋友，请帮助我们，并说服那位有能力者允准我们的请求。\par

\SourceRecord{Translated by Blomfield Jackson.；Nicene and Post-Nicene Fathers, Second Series；Vol. 3.；Edited by Philip Schaff and Henry Wace.；Buffalo, NY: Christian Literature Publishing Co.,；1892.；<http://www.newadvent.org/fathers/2707024.htm>.}

% structure: p0001 source=h1 target=section reason=page-title

\section{书信二十五}

\Emph{节日书信}\par

当\OriginalTerm{独生子天主}{only begotten God}降生成人，并为我们完成了救恩时，那些日子里亲眼看见这些恩惠之源的人并未举行庆节。但在我们这个时代，陆地与海洋、城镇与村庄，尽管无法用肉眼看见它们的恩主，却因纪念祂为他们所做的一切而举行庆节；从这些庆典中涌出的喜乐如此之大，以致属神喜乐的溪流奔涌四方。因此我们现在向您的虔诚致敬，既为表明这庆节在我们心中引起的欢愉，也为恳求您的祈祷，使我们能持守这庆节直到终末。\par

\SourceRecord{Translated by Blomfield Jackson.；Nicene and Post-Nicene Fathers, Second Series；Vol. 3.；Edited by Philip Schaff and Henry Wace.；Buffalo, NY: Christian Literature Publishing Co.,；1892.；<http://www.newadvent.org/fathers/2707025.htm>.}

% structure: p0001 source=h1 target=section reason=page-title

\section{第二十六封信}

\Emph{节庆。}\par

主的慈爱的泉源不断涌流出美善给相信的人；但某些额外的益处是通过那些保存最大恩惠记忆的庆典，赐予那些以更大热忱守庆节的人。我们刚才举行了礼仪，享受了它们的祝福，并以此向你的虔敬致意，因为庆节的传统和爱德之律如此规定。\par

\SourceRecord{Translated by Blomfield Jackson.；Nicene and Post-Nicene Fathers, Second Series；Vol. 3.；Edited by Philip Schaff and Henry Wace.；Buffalo, NY: Christian Literature Publishing Co.,；1892.；<http://www.newadvent.org/fathers/2707026.htm>.}

% structure: p0001 source=h1 target=section reason=page-title

\section{第二十七封信}

\Emph{给执事兼隐修院院长\Person{阿基利诺}}\par

凡获得了天主义子身份的人，不会为孤儿状态哭泣，因为有什么监护比我们在高天的父亲更强大呢？因为祂，地上的父亲才成为父亲。因祂的旨意，有些人因本性成为父亲，有些人因恩宠成为父亲。因此，让我们紧紧持守祂，并保持对已亡者的纪念。因为纪念那些善生者，激励我们效法他们，将使我们获益。\par

\SourceRecord{Translated by Blomfield Jackson.；Nicene and Post-Nicene Fathers, Second Series；Vol. 3.；Edited by Philip Schaff and Henry Wace.；Buffalo, NY: Christian Literature Publishing Co.,；1892.；<http://www.newadvent.org/fathers/2707027.htm>.}

% structure: p0001 source=h1 target=section reason=page-title

\section{第二十八封信}

\Emph{致\Person{雅各伯}，司铎兼隐修士。}\par

那些以德行勤勉使其壮年之力辉煌的人，幸福地奔向老年，因回忆往昔的胜利而欣喜，并因年老而免于进一步的奋斗。我认为您自己的虔诚也拥有这种喜悦，而且您因回忆青年时代的辛劳而更能承受老年。\par

\SourceRecord{Translated by Blomfield Jackson.；Nicene and Post-Nicene Fathers, Second Series；Vol. 3.；Edited by Philip Schaff and Henry Wace.；Buffalo, NY: Christian Literature Publishing Co.,；1892.；<http://www.newadvent.org/fathers/2707028.htm>.}

% structure: p0001 source=h1 target=section reason=page-title

\section{书信 29}

\Emph{致\Person{Apellion}}\par

\Place{迦太基}人所受的苦难，就其伟大而言，或许会要求甚至超出\Person{埃斯库罗斯}或\Person{索福克勒斯}的悲剧语言所能表达的力量。昔日的\Place{迦太基}被\Place{罗马}人艰难攻克。她一次又一次地与\Place{罗马}争夺世界霸权，并使\Place{罗马}濒临毁灭的危险。如今，它的毁灭却只是野蛮人的一场余兴罢了。如今，她赫赫有名的元老院中的尊贵成员们流落世界各地，依靠慷慨的陌生人的施舍维持生计，令旁观者落泪，并教导人们命运的无常与不稳。\par

我见过许多从那里来的人，心中感到恐惧，因为我不知道，正如圣经所说，\BibleQuote{明日将带来什么}。我尤其钦佩令人敬佩而最可敬的\Person{塞莱斯蒂尼亚努斯}，他如此勇敢地承受自己的不幸，并将失去幸福作为修身的机会，赞颂万有的主宰，并认为凡天主所命定或允许的，都是好的。因为神圣上智安排是不可言喻的。他与妻儿一同旅行，我恳请阁下以\Person{亚伯拉罕}般的款待接待他。我完全信赖你的仁善，已应允将他引荐于你，并正告诉他你的右手是何等慷慨。\par

\SourceRecord{Translated by Blomfield Jackson.；Nicene and Post-Nicene Fathers, Second Series；Vol. 3.；Edited by Philip Schaff and Henry Wace.；Buffalo, NY: Christian Literature Publishing Co.,；1892.；<http://www.newadvent.org/fathers/2707029.htm>.}

% structure: p0001 source=h1 target=section reason=page-title

\section{信函三十}

\Emph{致诡辩家\Person{艾里乌}}。\par

如今正是你的学园证明其讨论之用的时刻。我得知有一群杰出的人士聚集在你家中，成员既出身显赫，又言辞精妙，你们辩论德行、灵魂不朽以及其他类似主题。现在请适时展示你灵魂的高贵和德行的丰盛，并以那些深知人间福祉迅速变迁之人的精神，接纳最可敬、最尊贵的\Person{塞莱斯蒂尼亚努斯}。他昔日曾是\Place{迦太基}城的荣耀，在那里他向许多司铎敞开家门，从未想过需要陌生人的善举。请作他的代言人，我的朋友，在他需要你的嗓音时帮助他，因为他无法忍受那位诗人的劝告——诗人规劝需要发言的人即使羞愧也要说话。\par

我恳求你，说服你团体中能够如此行的人，效法\Person{阿尔基努斯}的好客精神，除去他意外遭遇的贫困，将厄运变为好运。让他们赞美我们仁慈的主，因祂用他人的灾祸使我们智慧，而没有将我们送到陌生人的家中，却将陌生人带到了我们的门前。祂向善待他人的人应许，赐予言语无法表达、理智无法理解的事物。\par

\SourceRecord{Translated by Blomfield Jackson.；Nicene and Post-Nicene Fathers, Second Series；Vol. 3.；Edited by Philip Schaff and Henry Wace.；Buffalo, NY: Christian Literature Publishing Co.,；1892.；<http://www.newadvent.org/fathers/2707030.htm>.}

% structure: p0001 source=h1 target=section reason=page-title

\section{第三十一封书信}

\Emph{致\Place{安提约基雅}主教\Person{多姆诺斯}。}\par

最令人敬佩和尊贵的\Person{塞来斯提尼安努斯}生于著名的\Place{迦太基}，出身于该城的一个显赫家族。如今他已从那里被流放。他流落异乡，不得不仰赖热爱天主之人的善心。他肩负着一个无法逃避的重担，这重担增加他的忧虑——我指的是他的妻子、孩子和仆役，他为这些人花费巨大。我惊叹他的精神。因为他赞美伟大的舵手，仿佛正被和风托举，对那可怕的风暴毫不在意。他从自己的灾难中收获了虔诚的果实，这至福的收益正是由他的不幸带来的；因为当他顺遂时，他从未接受这教导，而当恶日使他赤贫如洗时，在他的其他损失中，他也失去了自己的不敬，如今他拥有了信德的财富，并为此而轻视自己的破产。\par

因此，我恳求您的圣德，让他在这些异乡之地找到祖国，并嘱咐那些富裕之人去安慰一位曾经如他们一样富有的人，驱散他灾难的乌云。这是完全正当合理的：在秉性相同的人中间，既然所有人都曾有过失，那么逃脱了惩罚的人应当安慰那些遭遇恶运的人，并通过同情后者来平息天主的仁慈。\par

\SourceRecord{Translated by Blomfield Jackson.；Nicene and Post-Nicene Fathers, Second Series；Vol. 3.；Edited by Philip Schaff and Henry Wace.；Buffalo, NY: Christian Literature Publishing Co.,；1892.；<http://www.newadvent.org/fathers/2707031.htm>.}

% structure: p0001 source=h1 target=section reason=page-title

\section{第三十二书}

\Emph{致\Person{迪奥克提图斯}主教。}\par

如果万有的天主立即惩罚一切犯错的人，祂必早已毁灭了全人类。但祂宽容，祂是仁慈的审判者；因此，祂惩罚一些人，又借着受惩罚者的惩戒来教导其他人。在我们这个时代，就有一个这样仁慈对待的例子。从昔日称为\Place{利比亚}、今称\Place{阿非利加}的地方来的流亡者，被祂带到了我们的门前，祂让我们看见他们的痛苦，以激动我们恐惧，又借恐惧唤起我们的同情；这样祂就同时完成了两个目的：一方面借着他们的惩罚使我们获益，另一方面借着我们给他们带来安慰。我现在恳求你，将这安慰赐给那极为可敬和可贵的\Person{塞莱斯蒂尼亚努斯}，他曾是阿非利加主要城市的荣耀，如今却既无城邦，也无家宅，更无一生活所需。如今，在你圣德管辖下，受托管理灵魂牧职的人，理应将有益于同胞的事向他们提出，因为他们确实需要这样的教导。为此，我们知道，神圣的宗徒在他的\BibleBook{弟铎}{书}中写道：\Quote{我们的人也该学习行善，以备必要的用途}，因为即使我们的城，孤零零，人口稀少，又贫乏，尚且援助这些陌生人，那么那在真宗教中养育的\Place{贝洛雅}，更当如此行，尤其是在你圣德的领导下。\par

\SourceRecord{Translated by Blomfield Jackson.；Nicene and Post-Nicene Fathers, Second Series；Vol. 3.；Edited by Philip Schaff and Henry Wace.；Buffalo, NY: Christian Literature Publishing Co.,；1892.；<http://www.newadvent.org/fathers/2707032.htm>.}

% structure: p0001 source=h1 target=section reason=page-title

\section{第三十三封信}

\Emph{致伯爵与首席主教\Person{斯塔西穆斯}。}\par

要叙述最尊贵和显赫的\Person{塞莱斯蒂尼亚努斯}所受的苦难，需要悲剧式的雄辩。悲剧作家详尽陈述人类的灾祸，但我只能以一言向阁下禀告：他的祖国是利比亚——长久以来载于众人之口的土地，他的城市是闻名遐迩的迦太基，他世袭的品级是那座著名议会中的一个席位，他的家产丰裕。然而这一切如今已成故事，只是剥离了现实的话语。蛮族战争剥夺了他所有这一切。但命运便是如此；她不愿永远与同一人同在，而是急于变换住所，与他人同住。我恳求向阁下引荐这位客人，并恳求他得以享受您远近闻名的恩惠。我也恳求通过阁下，他得以被所有身居官位和富有的人认识，为使您既能成为他们得利的媒介，亦能从我们仁慈的天主那里赢得更高的赏报。\par

\SourceRecord{Translated by Blomfield Jackson.；Nicene and Post-Nicene Fathers, Second Series；Vol. 3.；Edited by Philip Schaff and Henry Wace.；Buffalo, NY: Christian Literature Publishing Co.,；1892.；<http://www.newadvent.org/fathers/2707033.htm>.}

% structure: p0001 source=h1 target=section reason=page-title

\section{第三十四封书信}

\Emph{致\Person{帕特里基乌斯}伯爵。}\par

各种美善都是可称赞的，但仁爱使它们更加优美。我们恳切地向万有的天主祈求仁爱；唯独藉着仁爱，我们在犯错时获得宽恕；仁爱使富者俯就穷人；因我知道阁下富于仁爱，我满怀信心地向您推荐令人钦佩且卓越的\Person{塞莱斯蒂尼亚努斯}，他一度拥有巨大的财富和产业，突然间被剥夺一切，但他承受贫穷之轻松，如同很少有人承受财富那般。那场导致他命运跌落之悲剧的主题是\Place{利比亚}和\Place{迦太基}遭受蛮族入侵。我已将他引荐于阁下；请向他人提及他的境况，并感动他们心生怜悯。您将因施予多人仁爱的教训而赢得更大的利益：\par

\SourceRecord{Translated by Blomfield Jackson.；Nicene and Post-Nicene Fathers, Second Series；Vol. 3.；Edited by Philip Schaff and Henry Wace.；Buffalo, NY: Christian Literature Publishing Co.,；1892.；<http://www.newadvent.org/fathers/2707034.htm>.}

% structure: p0001 source=h1 target=section reason=page-title

\section{第三十五封信}

\Emph{致\Person{依勒内}主教。}\par

主，您在多种善行上显赫，您的圣德更因仁爱、轻看财富以及为了帮助需要者而涌流的慷慨而特别美丽。我也知道，您认为那些在顺境中长大、却陷入困境的人值得特别关注。我深知这一点，所以冒昧向您介绍那位非常可敬和卓越的\Person{塞莱斯提尼亚努斯}。他曾在\Place{迦太基}因财富和地位而闻名，如今失去这些后，却因他的虔敬和哲学而闻名，因为他以顺服承受人们所谓的不幸，因为这不幸为他带来了灵魂的救恩。他带着一封描述他昔日繁荣的信来找我，他与我共度数日后，我亲身体验了人们对他所说的话的真实性。因此，我毫不犹豫地将他推荐给您的圣座，并恳求您把他介绍给城中的富裕人士。很可能，当他们得知他的遭遇后，因害怕同样的命运降临到自己身上，会努力通过施以怜悯来逃避审判。他别无生计，只能四处乞讨，因为他带着妻子和儿女以及与他一起逃脱蛮族暴行的仆人们，这使他开销更大。\par

\SourceRecord{Translated by Blomfield Jackson.；Nicene and Post-Nicene Fathers, Second Series；Vol. 3.；Edited by Philip Schaff and Henry Wace.；Buffalo, NY: Christian Literature Publishing Co.,；1892.；<http://www.newadvent.org/fathers/2707035.htm>.}

% structure: p0001 source=h1 target=section reason=page-title

\section{第三十六封书信}

\Emph{致\Person{庞皮安}，\Place{埃梅萨}主教。}\par

我深知你的资财微薄而心怀广大，在你身上，慷慨并不因资源有限而受阻。因此，我向你的圣德引荐那位可敬而卓越的\Person{塞勒斯提尼亚努斯}，他从前享有许多财富与顺遂，如今却仅带着自由之身从蛮族人手中逃脱，除了像你这样虔诚之人的慈悲外，别无生计。而且，忧虑环绕着他，因为与他同行的有他的妻子、儿女和仆人，他带他们出来，并非出于其他动机，而是出于人性考量，因为他认为，当他们不肯离弃他时，弃之不顾是不对的。我恳求你的善心，让他为我们富有的市民所认识，因为我想，若他们由你的圣德告知，并看到顺遂何等容易消逝，便会思念我们共同的人性，效法你的宽宏，尽己所能地帮助他。\par

\SourceRecord{Translated by Blomfield Jackson.；Nicene and Post-Nicene Fathers, Second Series；Vol. 3.；Edited by Philip Schaff and Henry Wace.；Buffalo, NY: Christian Literature Publishing Co.,；1892.；<http://www.newadvent.org/fathers/2707036.htm>.}

% structure: p0001 source=h1 target=section reason=page-title

\section{信函三十七}

\Emph{致总督萨路斯提乌斯。}\par

当统治者使公义的天平保持正直，并使其均匀平衡时，他们便为臣民带来种种益处；若他们还赋有明智，并对有需要者显示慈惠，那么他们的统治便会为治下之人带来多方面的利益。因着尊驾，百姓已享有这些善事，并在您先前的治理中体验了它们；如今，他们获悉治理的舵柄已托付给您的慷慨，心中充满喜乐。我祈求他们能获得更大的善，尊驾能赢得更高的赞誉，而颂扬者们的赞美词，能因在您所有其他可敬的荣衔之上，加上这善事的圆满——即真宗教，而得以印证。由于我被迫在\Place{希埃拉波利斯}停留数日，本希望有幸与尊驾会面，并不断向新来者询问官印是否已送达给您。然而，我因着神圣的救恩庆节，不得不匆匆赶回托付给我的城市。如今，我收到了尊驾的信函，怀着极大的喜悦回复您的问候，并已按您的要求，毫不迟延地派遣了那位可敬且虔敬的执事——他靠着天主的恩宠，是一位能寻找水源的人。愿主以祂的慈惠，使他能为该城提供良好的服务，并增加尊驾的光荣。\par

\SourceRecord{Translated by Blomfield Jackson.；Nicene and Post-Nicene Fathers, Second Series；Vol. 3.；Edited by Philip Schaff and Henry Wace.；Buffalo, NY: Christian Literature Publishing Co.,；1892.；<http://www.newadvent.org/fathers/2707037.htm>.}

% structure: p0001 source=h1 target=section reason=page-title

\section{第38封信}

\Emph{节庆}\par

救恩的神圣盛宴为我们带来了天主美善恩赐的源泉、十字架的祝福，以及从我们主的死亡而来的不朽，即我们主耶稣基督的复活，这复活预许了我们众人的复活。这些既是盛宴的恩赐，也彰显了神圣恩宠的丰裕，它使我们充满属灵的喜乐。但是，我们既被众多巨大的灾祸四面围困，盛宴的光彩便黯淡了，哀叹与哭泣与我们的圣咏混杂在一起。这样的忧苦乃是罪所产生。是罪使我们的生命充满痛苦；是因罪，死亡对我们来说比生命更可爱；是因罪，当我们在想象中思考那不朽的审判台时，我们甚至对来生也感到战栗。因此，愿你们的虔诚祈祷，使天主的仁慈降在我们身上，使这阴郁可怕的云雾消散，使阳光再次迅速赐予我们喜乐。\par

\SourceRecord{Translated by Blomfield Jackson.；Nicene and Post-Nicene Fathers, Second Series；Vol. 3.；Edited by Philip Schaff and Henry Wace.；Buffalo, NY: Christian Literature Publishing Co.,；1892.；<http://www.newadvent.org/fathers/2707038.htm>.}

% structure: p0001 source=h1 target=section reason=page-title

\section{第三十九封书信}

\Emph{节期的。}\par

我本意想以欢愉的言辞书写，奏响庆节属神喜乐的旋律，但因我们的众罪众多而受阻，这些罪正为我们招致天主的审判。因为，究竟谁能如此麻木，以至于察觉不到天主的义怒？那么，愿你们的虔诚祈祷事务能有所改善，好使我们也能改变书信的风格，以欢欣的话语取代哀哭之辞。\par

\SourceRecord{Translated by Blomfield Jackson.；Nicene and Post-Nicene Fathers, Second Series；Vol. 3.；Edited by Philip Schaff and Henry Wace.；Buffalo, NY: Christian Literature Publishing Co.,；1892.；<http://www.newadvent.org/fathers/2707039.htm>.}

% structure: p0001 source=h1 target=section reason=page-title

\section{第四十封书函}

\Emph{致代理主教\Person{狄奥多若}}\par

节庆的惯例促使我写一封节庆书函，但灾难的阴云却不允许我从其中收获通常的喜乐果实。谁的心能硬如石，不为主的忿怒与忧伤感到震惊和恐惧？谁不被触动而追忆过犯？谁不期待公义的判决？这一切都使节日的荣光黯淡，但主满有慈爱，我们相信祂不会真正执行祂的警告，反而会以慈悲眷顾我们，驱散我们的忧伤，开启慈悲的泉源，并显示祂惯常的忍耐。我向阁下致敬，并恳请您赐知我您的安康消息，我衷心相信您正享此安康。\par

\SourceRecord{Translated by Blomfield Jackson.；Nicene and Post-Nicene Fathers, Second Series；Vol. 3.；Edited by Philip Schaff and Henry Wace.；Buffalo, NY: Christian Literature Publishing Co.,；1892.；<http://www.newadvent.org/fathers/2707040.htm>.}

% structure: p0001 source=h1 target=section reason=page-title

\section{第四十一封信}

\Emph{致\Person{克劳狄安}。}\par

神圣庆典如常赐予我们属灵的恩惠；然而罪的苦果不许我们喜乐地享用。它们产生惯常的后果：起初使荆棘、蒺藜、汗水、辛劳和痛苦萌芽；此刻，罪使大地向我们震动，使万国从四方起来攻击我们。我们哀叹，因为我们迫使那愿意施恩于我们的善主，使我们受苦，并迫使他施加惩罚。\par

然而，当我们思念他慈悲的不可测深度时，便得安慰，并相信主必不弃绝他的子民，也不离弃他的产业。在问候您的尊贵之余，我恳求您告知我您久盼的健康消息。\par

\SourceRecord{Translated by Blomfield Jackson.；Nicene and Post-Nicene Fathers, Second Series；Vol. 3.；Edited by Philip Schaff and Henry Wace.；Buffalo, NY: Christian Literature Publishing Co.,；1892.；<http://www.newadvent.org/fathers/2707041.htm>.}

% structure: p0001 source=h1 target=section reason=page-title

\section{第四十二封信}

\Emph{致\Person{君士坦提乌斯}长官。}\par

若非有必不得已之事迫使余致函于尊驾，余或可被指为僭妄，因既未量己之度，亦未识尊权之赫赫。然今者天主所托付于余治理之城邑与地区，所余者已濒于尽灭，且有人竟敢以诽谤之词加于近日之巡察，余深信阁下若察余作书之必要与余书写之用意，必能宽恕此书之唐突。余叹息哀恸，因不得不笔伐一人，其过误本应遮掩，因其属于圣职之列。然余仍书写以辩护其所欺压之穷人之事。此人被控多端，被排除于共融之外，待神圣主教会议召集；彼为逃避主教会议之裁决，已由此地逃脱，以此践踏教会法律，且因蔑视绝罚之判决而暴露其动机。彼所提出的控诉，甚至不适合于卑贱之工匠，并因对显赫的斐理伯之敌意，转而加害于不幸的纳税人。余以为无须再提其品性、其生平之初以及其恶行之重大，惟有一事恳请大人，勿信其谎言，而应批准巡察，并宽恕不幸的纳税人。是的，请宽恕三倍不幸的市议员，他们无法征收所要求之款项。有谁不知我等境内田亩征税之严苛？因此我等之大部分地主已逃离，农夫已逃亡，我等之大部分土地已荒芜。讨论土地时，用几何术语并无不妥。我乡之长度四十英里，宽度亦然。其内多高山，有者全无草木，有者覆盖不毛之植被。此区域内有五万自由犹格，另有一万属于帝国国库。惟愿阁下之智慧思量此不公之大。若乡间无荒地，且皆为农夫提供易耕之田，彼等亦必因无法忍受征税之严苛而陷于贡赋。此有证据：在荣福纪念的依西多禄时，一万五千亩土地以金币纳税，但科米提安评估之征收者因无法承受损失而屡次抱怨，并以奉献恳请尊驾减免二千五百亩不毛之地之税，且阁下之前任下令从不孝的市议员中除去不毛之地，并以等量之数替代科米提安之数；然即使如此，彼等亦未能完成其数目。\par

余以此多言恳求阁下之恩惠，并恳请尊驾摒弃对不幸纳税人之诬告，遏止此悲惨地区之苦难，使其得以再度抬头。如此，阁下将为后世留下不朽之荣誉纪念。余之祈求得我区所有圣人相助，尤其是至圣而虔诚的天主之人——主\Person{雅各伯}，彼极其珍视沉默，以致不能动笔，但彼祈祷，愿我城——因有彼为邻而增辉，并为彼之祈祷所护佑——能获得余所求之恩惠。\par

\SourceRecord{Translated by Blomfield Jackson.；Nicene and Post-Nicene Fathers, Second Series；Vol. 3.；Edited by Philip Schaff and Henry Wace.；Buffalo, NY: Christian Literature Publishing Co.,；1892.；<http://www.newadvent.org/fathers/2707042.htm>.}

% structure: p0001 source=h1 target=section reason=page-title

\section{第43封书信}

\Emph{致奥古斯塔·\Person{普尔喀利亚}}\par

因为您以虔诚装饰帝国，以信德使紫袍更加明亮，我们才敢冒昧致信于您，不再顾及我们的卑微，因您始终对圣职人员施予应有之尊敬。怀着这些情感，我恳求陛下纡尊降贵，对我这不幸的乡里显示仁慈，下令批准已多次进行的视察，不接受某些人对它所提的诬告。我恳求您不要信任那位名义上是主教、但行为连体面的奴隶都配不上的人。他本人曾受到严重指控，且在召集主教会议调查对他的指控之前，受至圣而敬爱天主的\Place{安提约基雅}总主教\OriginalTerm{多姆努斯大人}{Lord Domnus}所绝罚。如今他逃走了，前往\Place{君士坦丁堡}，干起告密者的勾当，攻击他的祖国——那里有成千上万的穷人，并且因对一个人的仇恨，就向所有人鼓舌。出于对我身份合宜的考虑，我对他的品格和教育不予置评；而且他目前的所作所为已昭然若揭。但关于该地区，我要说：当整个行省的负担都减轻时，这一部分虽然承担了极重的分担，却从未享受到舒缓的益处。结果是许多庄园失去了农夫；更甚者，许多庄园被其主人完全抛弃，而可怜的\OriginalTerm{市议员}{decurions}却要为这些产业缴税，实在无法忍受勒索，于是有的去乞讨，有的则逃走。这座城市似乎只剩下一个人，若无陛下的虔诚提供补救，他将无法支撑。但我希望陛下会治愈这座城的创伤，并为您众多的善行再添这一件。\par

\SourceRecord{Translated by Blomfield Jackson.；Nicene and Post-Nicene Fathers, Second Series；Vol. 3.；Edited by Philip Schaff and Henry Wace.；Buffalo, NY: Christian Literature Publishing Co.,；1892.；<http://www.newadvent.org/fathers/2707043.htm>.}

% structure: p0001 source=h1 target=section reason=page-title

\section{\WorkTitle{第四十四封信}}

\Emph{致贵族} \Emph{元老。}\par

感谢世界的救主，因为祂不断为您尊驾的威仪增添尊荣和名誉。我至今没有写信表达我对您尊荣圆满的喜悦，是因为我不想打扰您阁下的尊贵。此刻我写信之时，天主眷顾托付于我管理的地区，正如谚语所说，悬于刀锋之上。您还记得当我们最初蒙受您莅临的恩惠时所进行的视察吗？它是在最卓越的总督\Person{弗洛伦提乌斯主}时期艰难建立起来的，并由现任官员确认。有一个称之为主教、但行为甚至不如戏子的人，在受到绝罚判决之时从主教会议中逃脱，并企图诽谤和诋毁这次视察，同时因他对杰出的\Person{斐理伯}的仇恨而攻击真理。因此，我恳请阁下使他的谎言无效，使依法确认的视察不受干扰。确实，您的威仪理应在其他善行中收获这份善果，接受受惠者的欢呼，并以此同时向万有的天主和他真正的仆人、真正属天主的人\Person{主雅各伯}致敬，他与我一同向您呈上这份恳求。倘使他惯于写信，他会亲自写信的。\par

\SourceRecord{Translated by Blomfield Jackson.；Nicene and Post-Nicene Fathers, Second Series；Vol. 3.；Edited by Philip Schaff and Henry Wace.；Buffalo, NY: Christian Literature Publishing Co.,；1892.；<http://www.newadvent.org/fathers/2707044.htm>.}

% structure: p0001 source=h1 target=section reason=page-title

\section{第四十五封信}

\Emph{致贵族 \Person{阿纳托利}。}\par

阁下深知东方的全体居民对您的尊威怀有何种情感，如子女对慈父一般。为何您却向爱您的人显示憎恨，剥夺了您的仁慈关怀，并以私利为先而将他人的服务置之度外，使他们全都陷入哭泣与哀叹？的确，我认为他们中没有一个敬畏主的人不为失去您的治理而深感悲痛，而且我认为即使其余的人，虽然对神性之事没有正确的认识，但当他们想到您所施的恩惠时，也都分担这些痛苦之感。就我而言，当我想到您的尊严与淳朴品格时，我尤其悲伤，并祈求万有的天主常以祂不可战胜的右手作为您的壁垒，并赐予您各种丰厚的祝福。我们恳求阁下，无论您在座与否，都一如既往地向我们施以惯常的保护，并平息那位不配之主教（他的意图阁下知之甚详）的愤怒。据我所知，他正试图彻底摧毁我们的地区，并充当告密者以诽谤近期的巡视，而此时众人皆知，我们地区的赋税非常沉重，以致许多田庄已被农夫遗弃。但此人藐视自己被逐出教会的处分，并从 \OriginalTerm{神圣公会}{holy synod} 逃逸，竟张口攻击不幸的穷人。愿阁下允许留意，使真理不被谎言战胜。关于\Place{基利家人}，我也呈上同样的恳求。因为我们不停地哀哭，直到不义被消除。主许诺连一滴水也要奖赏，祂必因这辛劳而报答您。\par

\SourceRecord{Translated by Blomfield Jackson.；Nicene and Post-Nicene Fathers, Second Series；Vol. 3.；Edited by Philip Schaff and Henry Wace.；Buffalo, NY: Christian Literature Publishing Co.,；1892.；<http://www.newadvent.org/fathers/2707045.htm>.}

% structure: p0001 source=h1 target=section reason=page-title

\section{第四十六封信}

\Emph{致博学的\Person{伯多禄}}。\par

任何事物都不能阻止那些高度珍视正义之人的可赞美的志向。这一点由阁下您最近得知消息时所流露的悲痛，以及您决不容忍正义所受的攻击所证实。您适时地放下了您的忧愁，并正义地堵住了真理之仇敌的口。我们一听说此事，发现真正的哲学与修辞技巧如此结合，便对您的卓越更加热诚。现在我们更加恳切地求您抵制这好人的谎言，并确认给予不幸穷人的安慰。\par

\SourceRecord{Translated by Blomfield Jackson.；Nicene and Post-Nicene Fathers, Second Series；Vol. 3.；Edited by Philip Schaff and Henry Wace.；Buffalo, NY: Christian Literature Publishing Co.,；1892.；<http://www.newadvent.org/fathers/2707046.htm>.}

% structure: p0001 source=h1 target=section reason=page-title

\section{第四十七封书信}

\newline{}\Emph{致普罗克鲁斯}，\Emph{君士坦丁堡主教。}\newline{}\par

\newline{}一年前，多亏了你的圣德，本城杰出的总督菲利普得以摆脱严重危险。他享受了你仁慈所赐的安全之后，满口称赞你的美德。然而，某个极其虔诚的人物正竭力使你的全部辛劳化为乌有。他诽谤十二年前数次进行的巡视，并采用了一种连体面的奴隶都不应有的诽谤手法。现在，我恳求你的圣德制止他的谎言，并促使杰出的总督们确认他们公正而仁慈地作出的决定。事实上，本城税收比省内所有城市都重，在其他城市都已减轻之后，本城至今仍被评估为超过六万两千英亩。最后，那些尊位上的官员好不容易才同意派遣巡视员到该地区；他们的报告首先由可敬的伊西多尔接收，并由荣耀而爱基督的弗洛伦提乌斯主人确认，随后我们现任统治者——他的公正为宝座增光——对整个事件进行了非常仔细的调查，并以御旨确认了评估。但这位热爱真理的人，只因为对杰出的菲利普一人的仇恨，就向穷人宣战。在这种情况下，我恳求你的圣德以你公义的言辞之力来对抗他不义的言辞，庇护受攻击的真理，同时证明真理的力量和谎言的虚妄。\newline{}\par

\SourceRecord{Translated by Blomfield Jackson.；Nicene and Post-Nicene Fathers, Second Series；Vol. 3.；Edited by Philip Schaff and Henry Wace.；Buffalo, NY: Christian Literature Publishing Co.,；1892.；<http://www.newadvent.org/fathers/2707047.htm>.}

% structure: p0001 source=h1 target=section reason=page-title

\section{第四十八封信}

\Emph{致\Person{尤斯塔修斯}，\Place{贝里图斯}主教。}\par

我很高兴地接受了这项指控，尽管我不难驳斥这份诉状。我写的信件不是三封，而是四封；我怀疑有两种可能：要么是那些承诺送信的人没有把信送到，要么是阁下虽然收到了信，但仍渴望更多，于是对我提出了懒惰的指控。我如前所述，并不因这指控而苦恼，因为这清楚地证明了你对我的深情厚谊。那么，继续施展你的计谋吧，不要停止你的抱怨，这样就能使我感到欣慰。\par

\SourceRecord{Translated by Blomfield Jackson.；Nicene and Post-Nicene Fathers, Second Series；Vol. 3.；Edited by Philip Schaff and Henry Wace.；Buffalo, NY: Christian Literature Publishing Co.,；1892.；<http://www.newadvent.org/fathers/2707048.htm>.}

% structure: p0001 source=h1 target=section reason=page-title

\section{第四十九封书信}

\Emph{致\Person{达米亚努斯}}，\Emph{\Place{西顿}主教}。\par

镜子的本性是反映注视者的面容，因此任何人看镜子都看到自己的形象。眼睛的瞳孔也是如此，它们映出他人面容的相似。您的圣德提供了这样一个实例，因为您没有看到我的丑陋，反而欣赏了自己的美丽。我实际上没有您所列举的任何品质。尽管如此，我祈愿您的话语能被事实证实，并恳求您的虔诚通过您的祈祷使之实现，这样您的赞美就不会因为缺乏相应的现实而落空。\par

\SourceRecord{Translated by Blomfield Jackson.；Nicene and Post-Nicene Fathers, Second Series；Vol. 3.；Edited by Philip Schaff and Henry Wace.；Buffalo, NY: Christian Literature Publishing Co.,；1892.；<http://www.newadvent.org/fathers/2707049.htm>.}

% structure: p0001 source=h1 target=section reason=page-title

\section{第五十封信}

\Emph{致院长格伦提乌斯。}\par

心灵的特征常藉言语描绘，其无形之态亦得以显明；如今，尊驾的书信正展示了您圣洁灵魂的虔诚。您对那判决的等待、焦虑，寻求辩护者并预备答辩，这一切都清楚表明您灵魂对天主之事的热忱。相反的，我们却近乎怠惰沉睡，游手好闲，极需祈祷的助佑。请为我们赐予这些祈祷吧，天主所爱的人，好叫我们至少能醒悟过来，并稍加关心灵魂。\par

\SourceRecord{Translated by Blomfield Jackson.；Nicene and Post-Nicene Fathers, Second Series；Vol. 3.；Edited by Philip Schaff and Henry Wace.；Buffalo, NY: Christian Literature Publishing Co.,；1892.；<http://www.newadvent.org/fathers/2707050.htm>.}

% structure: p0001 source=h1 target=section reason=page-title

\section{书信第五十一}

\Emph{致司铎\Person{阿加比乌斯}}。\par

善工本身是可钦可佩的，但如果有一位能将其恰当叙述的雄辩之才，它们就更显得可钦可佩。在备受天主爱戴的主教\Person{托马斯}大人的事例中，这两样长处都未曾缺乏：因为他本人为虔诚奉献了自己的辛劳，也找到了你的圣德这舌唇来适当地赞扬那些辛劳。他带着这样有利的见证而来，我们见到他更加喜悦，并在享受了一段短暂的共融之后，送他回到他的职责。\par

\SourceRecord{Translated by Blomfield Jackson.；Nicene and Post-Nicene Fathers, Second Series；Vol. 3.；Edited by Philip Schaff and Henry Wace.；Buffalo, NY: Christian Literature Publishing Co.,；1892.；<http://www.newadvent.org/fathers/2707051.htm>.}

% structure: p0001 source=h1 target=section reason=page-title

\section{书信第五十二}

\Emph{致厄得撒主教依巴斯}。\par

我想，万有的天主为了我们共同救恩的眷顾，使某些人遭遇灾祸，为的是这惩戒对犯错误者可成为治疗的良药，对德行的运动员成为坚持的鼓励，对旁观者成为有益的榜样。因为当我们看见别人受罚，自己也会充满畏惧。基于这些考量，我把非洲的患难视为普遍利益。首先，当我想到他们昔日的繁荣，现在看到他们突然覆灭，我看到人事多么变幻无常，并学会双重教训：既不在幸福中欢喜，好像它永不结束；也不在灾难中痛苦，好像难以忍受。然后，我回忆起过去的错误，战栗害怕自己遭受同样的痛苦。我现在写信给您的动机，是向您圣德引见非常蒙天主所爱的西彼廉主教，他从著名的非洲出发，如今因蛮族的残暴被迫在异乡旅行。\par

他从极圣善的、明智管理加拉太人的欧瑟比主教那里带了一封信给我们。当您的敬德以惯常的仁慈接待他后，我请求您写一封信，把他送到您认为合适的任何虔敬主教那里，好使他在享受他们仁慈安慰的同时，也能成为他们获得天上永恒恩惠的媒介。\par

\SourceRecord{Translated by Blomfield Jackson.；Nicene and Post-Nicene Fathers, Second Series；Vol. 3.；Edited by Philip Schaff and Henry Wace.；Buffalo, NY: Christian Literature Publishing Co.,；1892.；<http://www.newadvent.org/fathers/2707052.htm>.}

% structure: p0001 source=h1 target=section reason=page-title

\section{书信 53}

\Emph{致\Person{索弗罗纽斯}，\Place{君士坦丁纳}的主教。}\par

既然我知道，敬爱天主的人，你的右手是多么慷慨和丰裕，我便将一件令人渴望的好处放在你伸手可及之处；因为正如贪求世俗利益的人见到需要金钱援助的人便感到烦恼，同样，慷慨的人却感到喜悦，因为他们所追求的是天上的财富。提供这一绝佳机会的人，是敬爱天主的\Person{西彼廉}主教，他以前在服务他人的人中为人所知，但现在，当他讲述非洲灾难的悲惨情况时，却不得不仰赖他人的仁爱，依靠虔诚灵魂的慷慨。我希望他也能享受你的兄弟般的善意，并被转介到其他避风港。\par

\SourceRecord{Translated by Blomfield Jackson.；Nicene and Post-Nicene Fathers, Second Series；Vol. 3.；Edited by Philip Schaff and Henry Wace.；Buffalo, NY: Christian Literature Publishing Co.,；1892.；<http://www.newadvent.org/fathers/2707053.htm>.}

% structure: p0001 source=h1 target=section reason=page-title

\section{第五十四封书信}

\Emph{庆节书信。}\par

依照我们神圣且带来救恩的庆节，心灰意冷者得到鼓舞，喜乐者更加喜乐。这是我凭经验学到的，因为当我被绝望的波涛淹没时，一看到庆节这港口，就凌驾于浪涛之上。愿你的虔敬祈求，使我完全脱离这场风暴，并使爱我们的主赐我忘却忧伤。\par

\SourceRecord{Translated by Blomfield Jackson.；Nicene and Post-Nicene Fathers, Second Series；Vol. 3.；Edited by Philip Schaff and Henry Wace.；Buffalo, NY: Christian Literature Publishing Co.,；1892.；<http://www.newadvent.org/fathers/2707054.htm>.}

% structure: p0001 source=h1 target=section reason=page-title

\section{第五十五封书信}

\Emph{庆节的}.\par

我们深感痛苦，因为我们本非岩石，而是具有人的本性；然而，对主显节的回忆对我而言是一剂极为有效的良药；因此我立刻遵照节日的惯例写信，并以祈祷问候您的尊贵，愿您生活昌盛、享有美誉。\par

\SourceRecord{Translated by Blomfield Jackson.；Nicene and Post-Nicene Fathers, Second Series；Vol. 3.；Edited by Philip Schaff and Henry Wace.；Buffalo, NY: Christian Literature Publishing Co.,；1892.；<http://www.newadvent.org/fathers/2707055.htm>.}

% structure: p0001 source=h1 target=section reason=page-title

\section{第五十六封信}

\Emph{节日}。\par

我之悲伤已至极点，心神为之重创；然思及履行节庆之礼，实属应当，故今执笔向阁下致敬，并偿情爱之债。\par

\SourceRecord{Translated by Blomfield Jackson.；Nicene and Post-Nicene Fathers, Second Series；Vol. 3.；Edited by Philip Schaff and Henry Wace.；Buffalo, NY: Christian Literature Publishing Co.,；1892.；<http://www.newadvent.org/fathers/2707056.htm>.}

% structure: p0001 source=h1 target=section reason=page-title

\section{书信五十七}

\Emph{致\Person{尤特雷修斯}总督。}\par

除了其他恩惠，宇宙的主宰还恩赐我们听闻阁下获得荣耀，并同时为您的高升和您的臣民以如此温和的统治而庆贺。我认为若不表达我的满意并藉书信显明它，便是错误的。您的显赫完全知道我们对您怀有多深的爱——一种被最深切地回报的爱。充满这样的爱，我们恳求一切美善的赐予者，不断赐予您祂丰厚的恩赐。\par

\SourceRecord{Translated by Blomfield Jackson.；Nicene and Post-Nicene Fathers, Second Series；Vol. 3.；Edited by Philip Schaff and Henry Wace.；Buffalo, NY: Christian Literature Publishing Co.,；1892.；<http://www.newadvent.org/fathers/2707057.htm>.}

% structure: p0001 source=h1 target=section reason=page-title

\section{书信第五十八}

\Emph{致执政官诺穆斯。}\par

想到要寄信给尊驾，我心存犹豫。一方面，我深知万事均仰仗您的判断；目睹您身负重担，为国事操劳，故以为保持沉默更为妥当。另一方面，我深知您才智广博，不忍一言不发，唯恐落得疏忽之责。此外，与您短暂的交游令我心怀眷恋，更促我提笔。当时因那位极有福之人的疾病与离世，我未能尽享其乐；如今我想，写信或可聊以慰藉。我祈求万有的主宰引导您的人生，常乘顺风，使我们能得享您仁爱关照的益处。\par

\SourceRecord{Translated by Blomfield Jackson.；Nicene and Post-Nicene Fathers, Second Series；Vol. 3.；Edited by Philip Schaff and Henry Wace.；Buffalo, NY: Christian Literature Publishing Co.,；1892.；<http://www.newadvent.org/fathers/2707058.htm>.}

% structure: p0001 source=h1 target=section reason=page-title

\section{\WorkTitle{Letter 59}}

\Emph{致克劳狄安。}\par

真诚的友谊既不会因地域的阻隔而消散，也不会因时间的流逝而减弱。时间确实给我们的身体带来侮辱，夺去其青春的光彩，招致衰老；但友谊却使它的美丽更加灿烂，不断点燃它的火焰，使之更加温暖和明亮。因此，我虽与阁下相隔多日的路程，却在友谊的鞭策下，写下这封问候的信函。它由掌旗官\Person{帕特若伊诺斯}转交，此人因品格高尚，值得一切尊敬，因他殷勤努力遵守天主的法律。至卓越的阁下，请借由他告知我们您珍贵的健康状况，以及您所期望的诺言的实现。\par

\SourceRecord{Translated by Blomfield Jackson.；Nicene and Post-Nicene Fathers, Second Series；Vol. 3.；Edited by Philip Schaff and Henry Wace.；Buffalo, NY: Christian Literature Publishing Co.,；1892.；<http://www.newadvent.org/fathers/2707059.htm>.}

% structure: p0001 source=h1 target=section reason=page-title

\section{第六十封书信}

\Emph{致亚历山大主教\Person{迪奥斯哥鲁}}\par

在您圣德所装饰的许多德行中（因为所有人耳中都充满了您荣耀的飞鸿之誉，它向四方传播），一致称赞的是您的谦逊，这是我们的主在祂的法律中以身作则的德行，祂说：\BibleQuote{你们跟我学，因为我是良善心谦的}；因为天主是崇高的，或者更确切说是至高者，祂在降生成人时尊重了谦逊温和的精神。注视祂，阁下，您不看您属下的人数，也不看您宝座的崇高，而是看见人性，以及生命的迅速变化，并遵守神圣的法律，遵守法律给我们带来天国。听到您圣德方面的这种谦逊，我鼓起勇气在这封信中向一位神圣而亲爱的天主之人致敬，并献上以救恩为果实的祈祷。写作的机会是由非常虔诚的司铎\Person{欧瑟伯}给我的，因为我听说他要去您那里，我便立刻写了这封信，邀请您的圣德以祈祷支持我们，并以您的回信赐予我们一场神宴，把您话语的蒙福盛宴送给我们这些饥饿的人。\par

\SourceRecord{Translated by Blomfield Jackson.；Nicene and Post-Nicene Fathers, Second Series；Vol. 3.；Edited by Philip Schaff and Henry Wace.；Buffalo, NY: Christian Literature Publishing Co.,；1892.；<http://www.newadvent.org/fathers/2707060.htm>.}

% structure: p0001 source=h1 target=section reason=page-title

\section{第六十一封书信}

\Emph{致司铎阿尔奇比乌}\par

我并未忽视方才从你处收到的两封信，而是即刻写了回信，并将我的信交给了那位非常虔诚的司铎欧西比乌斯。但因某些耽搁，此事暂时延后，因为天气使船只滞留港内——海面上预示着一场风暴即将来临，水手与领航员都需稍作等待。所以我暂时偿还了这笔债务，并非就此不再欠债，而是要使债务增加。因为这项义务在偿还之后反而变得更多倍：凡试图遵守友谊律法者，会增强其爱的力量；好似将火星吹成火焰，燃起更炽热的情感；而所有被这火焰点燃的人，都力图在爱上超越彼此。请接受我的辩护吧，我可敬的朋友；请宽恕我；并请来信告知你的近况。\par

\SourceRecord{Translated by Blomfield Jackson.；Nicene and Post-Nicene Fathers, Second Series；Vol. 3.；Edited by Philip Schaff and Henry Wace.；Buffalo, NY: Christian Literature Publishing Co.,；1892.；<http://www.newadvent.org/fathers/2707061.htm>.}

% structure: p0001 source=h1 target=section reason=page-title

\section{书信第六十二}

\Emph{致司铎若望}\par

从前一个被称为明智的人说过：\Quote{隐姓埋名地生活。}我赞赏这句话，并决心以行动来证实这话，因为我觉得，从别人那里汲取善美并无不当，正如人们所说，蜜蜂采集蜂蜜，从苦草和可食的甜草中同样吸取甘汁。我自己也曾见过它们停在一块贫瘠的岩石上，吮吸那稀少的湿气。更合理的是，那些被认为有理性的人，应当从各种来源中收获善美。因此，如我所说，我努力隐姓埋名地生活，尤其喜爱安宁与平静。那位极虔敬的司铎欧瑟比最近从你们那里回来时，宣布你们曾举行过某次会议，在交谈中提到了我，而你的虔敬称赞了我这微不足道的人。因此，我认为，如果那位对我如此善意地称赞的人不能得到同样的回报，那就是忘恩负义，也是不公平的；因为虽然我们没有什么值得称赞的，但我们仍然钦佩那些称赞我们的人的意图，因为这种称赞是情感的产物。因此，我问候你的尊威，用那位将你们关于我的口传之言带给我的人作为传递信函的工具。最虔敬的先生，当你收到我的信后，请回信。你先发言；我先写信；我用信函来回答言语。现在，你只须以信函来回答信函。\par

\SourceRecord{Translated by Blomfield Jackson.；Nicene and Post-Nicene Fathers, Second Series；Vol. 3.；Edited by Philip Schaff and Henry Wace.；Buffalo, NY: Christian Literature Publishing Co.,；1892.；<http://www.newadvent.org/fathers/2707062.htm>.}

% structure: p0001 source=h1 target=section reason=page-title

\section{第六十三封书信}

\Emph{节日的}\par

我们已享受了庆节慣有的恩惠。我们举行了救恩苦难的纪念庆节；藉着主的复活，我们听到了众人复活的喜讯，并颂扬了我们天主和救主不可言喻的慈爱。但摇撼教会的风暴不容我们分享纯粹的喜乐。如果当一个肢体痛苦时，全身都分担痛苦，那么当整个身体都痛苦时，我们怎能不哀叹呢？更令我们气馁的是，想到这些事是普遍背道的前奏。愿您的虔诚祈祷，使我们在这困境中获得天主的援助，正如神圣宗徒所说的，我们可能\Quote{能抵禦邪惡的日子}。但如果还有时间处理此生的事务，请祈祷风暴过去，教会恢复往日的平静，真理的敌人不再因我们的不幸而得意。\par

\SourceRecord{Translated by Blomfield Jackson.；Nicene and Post-Nicene Fathers, Second Series；Vol. 3.；Edited by Philip Schaff and Henry Wace.；Buffalo, NY: Christian Literature Publishing Co.,；1892.；<http://www.newadvent.org/fathers/2707063.htm>.}

% structure: p0001 source=h1 target=section reason=page-title

\section{第六十四封书信}

\Emph{节庆信函}。\par

当主为了人类而承受救恩的受难时，神圣宗徒们的团体十分沮丧，因为他们不清楚这受难的果实将是什么。但当他们知道由此而生的救恩时，他们便称这受难的宣告为喜讯，并热切地将它献给全人类。而那些信了的人，由于心思受光照，便欣然接受它，并谨守这节日以纪念受难，把死亡时刻当作欢乐和庆典的机会。因为与它紧密相连的复活消除了死亡的悲伤，并成为众人复活的保证。在刚刚参与这庆典之后，我们向你们发送这节日的消息，如同芬芳的香水，并问候你们的虔敬。\par

\SourceRecord{Translated by Blomfield Jackson.；Nicene and Post-Nicene Fathers, Second Series；Vol. 3.；Edited by Philip Schaff and Henry Wace.；Buffalo, NY: Christian Literature Publishing Co.,；1892.；<http://www.newadvent.org/fathers/2707064.htm>.}

% structure: p0001 source=h1 target=section reason=page-title

\section{第六十五封信}

\Emph{致泽诺将军。}\par

遭受人间疾苦是众人皆有的命运；勇敢地承受苦难并超越其攻击则不再常见。前者属于人性，后者则取决于决心。正因如此，我们惊叹于那些哲学家如何选择了最高尚的生活方式，并以智慧战胜了他们的灾祸。哲学由我们理性的力量所产生，这力量主宰我们的情欲，而不被情欲左右。悲伤是人间疾苦之一，我们恳请阁下战胜它；您若思考人性，并深思悲伤的无益，便不难战胜这种情感。因为，我们哀号痛哭，对逝者有何益处呢？然而，当我们想到共同的出生、多年的交往、战场上的辉煌服务以及赫赫功绩时，让我们想到那被这些所装饰的人是一个服从死亡法则的人；而且，万事万物都由天主安排，祂按照对万事有益的圣智引导人间事务。至此，我在书信的篇幅内尽我所能写下这些，恳请尊驾为我们所有人保重身体——身体常因喜乐而健康，因沮丧而毁坏。因此，出于我对我们所有人利益的关心，我写下了这封信。\par

\SourceRecord{Translated by Blomfield Jackson.；Nicene and Post-Nicene Fathers, Second Series；Vol. 3.；Edited by Philip Schaff and Henry Wace.；Buffalo, NY: Christian Literature Publishing Co.,；1892.；<http://www.newadvent.org/fathers/2707065.htm>.}

% structure: p0001 source=h1 target=section reason=page-title

\section{第六十六书}

\Emph{致智者\Person{阿厄略}。}\par

那生养你的，邀你赴渴望的盛宴。圣所被屋顶加冕；它装饰合宜；它渴望那些为它所建造的居住者。这些是宗徒和先知，旧盟约和新盟约的高声宣报者。因此，请以你的光临装点这盛宴；领受从那盛宴洋溢的降福，并使这盛宴对我们更为欢悦。\par

\SourceRecord{Translated by Blomfield Jackson.；Nicene and Post-Nicene Fathers, Second Series；Vol. 3.；Edited by Philip Schaff and Henry Wace.；Buffalo, NY: Christian Literature Publishing Co.,；1892.；<http://www.newadvent.org/fathers/2707066.htm>.}

% structure: p0001 source=h1 target=section reason=page-title

\section{第六十七封书信}

\Emph{致\Person{马拉纳斯}。}\par

这曾是您的工作，我的好先生，去召唤其余的人也来参加奉献节。藉着您的热忱和精力，圣殿已被建起，真理的洪亮传报者已来居住其中，并守护那些凭信德前来的人。然而，我写信指明这节期的时节。\par

\SourceRecord{Translated by Blomfield Jackson.；Nicene and Post-Nicene Fathers, Second Series；Vol. 3.；Edited by Philip Schaff and Henry Wace.；Buffalo, NY: Christian Literature Publishing Co.,；1892.；<http://www.newadvent.org/fathers/2707067.htm>.}

% structure: p0001 source=h1 target=section reason=page-title

\section{第六十八封书信}

\Emph{\Person{厄丕法尼}。}\par

我曾希望邀请你参加圣宗徒与先知的庆节，不仅因为你是一位同胞，而且因为你是与我共享同一信德和同一家室的人。但我因你的意见状况而受阻。因此，除了我们祖国的名义之外，我不提出其他要求，我邀请你共融于圣宗徒与先知的宝贵祝福。这种共融是不受意见分歧阻碍的。\par

\SourceRecord{Translated by Blomfield Jackson.；Nicene and Post-Nicene Fathers, Second Series；Vol. 3.；Edited by Philip Schaff and Henry Wace.；Buffalo, NY: Christian Literature Publishing Co.,；1892.；<http://www.newadvent.org/fathers/2707068.htm>.}

% structure: p0001 source=h1 target=section reason=page-title

\section{第六十九封书信}

\Emph{致\Person{欧格拉菲亚}。}\par

若非因不可避免的阻碍，我一听到你那伟大而光荣的丈夫安眠，便会立刻赶到你身边。我曾从你手中领受了许多种多样的荣耀，我向你满心感激。当我因极不情愿的阻碍而无法偿还债务时，我认为在你悲痛至极的时刻给你寄信是欠妥的；那时我的使者不可能接近你的尊驾，而悲痛也使你无法阅读我所写的内容。但现在你的理智已从悲痛的迷醉中苏醒，能够抑制你的情感，并管束悲恸的放纵，我便冒昧写信，恳求你的尊驾思念人的本性，反思你所哀悼的损失何其普遍，尤其要接受神圣的教导，不让你的忧苦超越信德的界限。因你那最杰出的丈夫，正如主自己所说：\Quote{不是死了，而是睡着了}——这睡眠比他平日稍长而已。这希望是主赐予我们的；这应许我们是从神圣的谕示中领受的。我确实知道分离是多么令人痛苦，极其痛苦；尤其是当情感因性格的相投和时间的久长而更加牢固时。但让你的悲伤是为去远方旅行，而非为生命的终结。这种哲学特别适合于在虔敬中成长的人，而我所恳求你，我尊敬的友人，寻求的正是这种哲学的装饰。我并非以一个自身缺乏感觉的人的身份给你这劝告；事实上，当我得知我所深爱者的离去时，我的心也感到忧伤。但我追忆世界的统治者及其不可言喻的智慧，祂为我们的益处安排一切。我恳求你的圣洁深思这些道理，超越你的悲伤，并赞美那作为我们众人的主宰的天主。祂以不可言喻的眷顾引导人的生命。\par

\SourceRecord{Translated by Blomfield Jackson.；Nicene and Post-Nicene Fathers, Second Series；Vol. 3.；Edited by Philip Schaff and Henry Wace.；Buffalo, NY: Christian Literature Publishing Co.,；1892.；<http://www.newadvent.org/fathers/2707069.htm>.}

% structure: p0001 source=h1 target=section reason=page-title

\section{\WorkTitle{书信第七十}}

\newline{}\Emph{致\Person{犹斯塔修}，\Place{埃盖}主教。}\newline{}\par

高贵的\Person{玛利亚}的故事，堪称一出悲剧。正如她自己所言，并经多位人士证实，她是尊贵的\Person{欧德蒙}之女。在\Place{利比亚}遭遇的灾祸中，她失去了父亲的自由身份，沦为奴隶。一些商人从蛮族手中买下她，又将她卖给了我们的同乡。与她一同被卖的，还有一位曾是她家仆的女子；于是，奴隶制的痛苦轭同时落在主仆二人身上。然而，那仆人并未忽视她们之间的区别，也没有忘记昔日的尊卑：在患难中，她保持了善意，在伺候共同的主人之后，还服侍那被视为与她同奴的女主人，为她洗脚、铺床，并尽心履行其他类似职责。此事被买主得知。于是，全城都传扬着女主人原先的自由身份和仆人的美善行为。这些情况被驻扎在我城（当时我正外出）的忠诚士兵们知晓，他们便向买主付了赎金，将女子从奴役中解救出来。我回来后，得知这悲惨的境遇和士兵们可赞的用心，便为他们祝福，将这位高贵的女子托付给一位可敬的执事照料，并吩咐为她提供充足的生活所需。十个月过去了，她听说父亲仍在世，并在\Place{西方}担任高官，自然渴望回到他身边。据报，许多来自\Place{西方}的使者正前往你处正在举办的集市。她请求允许带着我的信出发。因此，我写了这封信，恳请你的敬爱照管这位高贵的女子，并委派一位可靠之人与船夫、领航员和商人沟通，将她托付给能把她安全送回父亲身边的可靠之人。毫无疑问，那些在一切复原希望都已失去时，将女儿带回父亲身边的人，将获得丰厚的赏报。\par

\SourceRecord{Translated by Blomfield Jackson.；Nicene and Post-Nicene Fathers, Second Series；Vol. 3.；Edited by Philip Schaff and Henry Wace.；Buffalo, NY: Christian Literature Publishing Co.,；1892.；<http://www.newadvent.org/fathers/2707070.htm>.}

% structure: p0001 source=h1 target=section reason=page-title

\section{第七十一书}

\Emph{致\Person{Zeno}}，\Emph{将军与执政官。}\par

你的刚毅激发普遍的钦佩，它由温和与温良所调和，在你家眷中表现为仁慈，在敌人面前表现为勇敢。这些品质指明了一位令人钦佩的将军。在军人的品格中，主要装饰是勇敢，但在指挥官身上，谨慎优于勇敢；之后是自制和公正，由此聚集了德行的财富。这样的财富是追求善的灵魂的赏报，它双眼凝视着果实的甘甜，视劳苦为十分愉快。因为万有的天主如同一位伟大的竞赛赐予者，向德行的竞技者提供了奖品，一些在今世，一些在无穷的未来之世。阁下已经享有了今世的那些，并已达到最高的荣誉。愿尊驾也获得那些永恒持久的福祉，不仅接受执政官的长袍，也接受那不可言喻且神圣的衣裳。在所有理解那恩赐伟大的人中，这是共同的祈求。\par

\SourceRecord{Translated by Blomfield Jackson.；Nicene and Post-Nicene Fathers, Second Series；Vol. 3.；Edited by Philip Schaff and Henry Wace.；Buffalo, NY: Christian Literature Publishing Co.,；1892.；<http://www.newadvent.org/fathers/2707071.htm>.}

% structure: p0001 source=h1 target=section reason=page-title

\section{第七十二封书信}

\Emph{致助理法官\Person{赫莫西涅斯}}\par

当人们被无知黑暗所淹没的时候，并非所有人都遵守同样的庆节，而是在不同的城市举行不同的礼仪。在\Place{埃利斯}有奥林匹亚竞技会，在\Place{德尔斐}有皮提亚竞技会，在\Place{斯巴达}有海阿辛斯节，在\Place{雅典}有泛雅典娜节、地母节和狄俄尼索斯节。这些是最著名的；此外，有些人庆祝某些邪神的狂欢节，另一些人则庆祝别的。但现在，这些迷雾已被理智之光驱散，在每一片陆地和海洋上，大陆居民和岛民共同庆祝我们天主和救主的庆节。无论任何人想要外出旅行，无论是向日出之地还是向日落之地，处处都会发现同一庆节在同一天被遵守。不再需要为了服从那适应犹太人软弱的梅瑟法律而聚集到一座城中庆祝我们福分的纪念，而是每个城镇、每个村庄、乡间和最远边疆，都充满了天主的恩宠，在每一处都为万有的天主祝圣了神圣的神殿和圣地。因此，在每个城镇，我们各自庆祝我们的庆节，并在庆节中彼此相通。同一的天主和主在我们的赞歌中被尊崇，我们的奥迹祭献也是向祂奉献的。有鉴于此，众所周知，我们邻人之间彼此以书信致意，宣示在庆节中临到我们的喜乐。现在，我也如此向您致意，并献上节日的问候。您无疑会回复，并尊重庆节的惯例。\par

\SourceRecord{Translated by Blomfield Jackson.；Nicene and Post-Nicene Fathers, Second Series；Vol. 3.；Edited by Philip Schaff and Henry Wace.；Buffalo, NY: Christian Literature Publishing Co.,；1892.；<http://www.newadvent.org/fathers/2707072.htm>.}

% structure: p0001 source=h1 target=section reason=page-title

\section{第七十三封书函}

\Emph{致\Person{阿波洛尼乌斯}。}\par

那位令人赞叹的史家用赞美的笔触描述\Person{尼奥克利斯}之子\Person{特米斯托克利}，这位远近闻名的杰出将领，说他天生就具备卓越的德性。至于\Person{桑提波斯}之子\Person{伯里克利}，据说他同样从教育中获得了才能，能以雄辩的口才取悦听众，并且天赋异禀，既能知晓应采行之策，又能以言辞使之生效。我在论述他时，借用他自己的言辞并无不妥。这些事彰显了您的卓越，因为我们的造物主\OriginalTerm{天主}{God}赐予您天赋，而您的教育又使这光芒更加璀璨。因此，您的高贵品质已无所欠缺，唯独缺少对赐予者本身的认识；只要加上这一点，我们所拥有的诸德就将臻于圆满。当我得知您到来的消息时，便写信给您，恳求一切美善的施与者赐予您灵魂之眼一道光明，使您看见祂恩赐的伟大，点燃您对那拥有的爱慕，并将那渴望已久的心愿赐予渴望者。\par

\SourceRecord{Translated by Blomfield Jackson.；Nicene and Post-Nicene Fathers, Second Series；Vol. 3.；Edited by Philip Schaff and Henry Wace.；Buffalo, NY: Christian Literature Publishing Co.,；1892.；<http://www.newadvent.org/fathers/2707073.htm>.}

% structure: p0001 source=h1 target=section reason=page-title

\section{第七十四书函}

\Emph{致\Person{乌尔巴努斯}。}\par

我们慷慨的主再次赐予我们享受庆节，并向尊驾致以节日的问候。我们祈求您安康、昌盛，并分享那不可言喻的天主恩赐，这恩赐为接近的人提供所盼望之福的种籽，并赐予无尽生命与国度的标记。我们恳求仁爱的主将这些恩赐分施给您，因为朋友请求朋友蒙福是理所当然的。\par

\SourceRecord{Translated by Blomfield Jackson.；Nicene and Post-Nicene Fathers, Second Series；Vol. 3.；Edited by Philip Schaff and Henry Wace.；Buffalo, NY: Christian Literature Publishing Co.,；1892.；<http://www.newadvent.org/fathers/2707074.htm>.}

% structure: p0001 source=h1 target=section reason=page-title

\section{第七十五封书信}

\Emph{致\Place{贝罗亚}的圣职人员}\par

我深知，我对诸位尊长的好感是有理由的，因为你们亲切的信函使我相信我的感情得到了回报。我对你们的这份感情有多种原因。首先，你们的父亲，那位伟大而充满宗徒精神的人，也是我的父亲。其次，我看待现今治理你们教会的这位真正虔诚的主教，如同看待一位在血缘和心意上都与我相合的兄弟。第三，我们城邦的邻近，第四，我们之间的频繁交往，这自然产生友谊，并使之增长。如果你们愿意，我还可以提出第五点，那就是我们与你们的关系正如舌头与耳朵的关系，前者发出言语，后者接收言语；因为你们最乐意聆听我的话，而我也乐于向你们倾吐我的一点心得。但我们合一的圆满在于我们在信仰上的和谐；在于我们拒绝接受任何伪劣的学说；在于我们保存了那古老的、宗徒的教导，这教导已由古老的智慧带给了你们，并由德行的艰苦劳作所滋养。因此，我恳求你们更加关心羊群，将它完好无损地保存给牧者，并勇敢地发出族长的那句著名的话：\Quote{那从野兽生的，我没有献给祢。}\par

\SourceRecord{Translated by Blomfield Jackson.；Nicene and Post-Nicene Fathers, Second Series；Vol. 3.；Edited by Philip Schaff and Henry Wace.；Buffalo, NY: Christian Literature Publishing Co.,；1892.；<http://www.newadvent.org/fathers/2707075.htm>.}

% structure: p0001 source=h1 target=section reason=page-title

\section{书信七十六}

致塞浦路斯总督乌拉纽斯。\par

真正的友谊因交往而加强，但分离并不能切断它，因为它的纽带是坚固的。这一真理本可用许多其他例证轻易证明，但我们自己的情况足以验证我所说的。在我与你之间确实隔着许多事物——山脉、城市和海洋——然而，没有任何东西能抹去我对你阁下的记忆。每当我们看到任何从那些沿海城镇抵达的人，谈话便转向塞浦路斯及其尊贵的总督，我们很高兴听到你的崇高声誉。最近，我们更是格外欣慰地得知了最令人喜悦的消息：因为，最卓越的先生，有什么比看到你那高贵的灵魂被知识之光所照亮更令我们欢喜呢？我们认为，那位以多种德行装饰自己的人，应当再加上它们的\OriginalTerm{顶峰}{colophon}，并且我们相信我们将看到我们所渴望的。因为你的高贵无疑会热切地抓住天主所赐的恩惠，为此而被真正的朋友们所感动，他们清楚地理解其价值，并被引导至慷慨的天主，祂\BibleQuote{愿意所有的人都得救，并得以认识真理}，藉由人作工具将人网向救恩，并将祂所捕获的人带入永生。渔夫确实夺去猎物的生命，但我们的\Emph{渔者}却将祂所活捉的一切从死亡的痛苦束缚中释放，因此，祂\BibleQuote{曾在世上出现，并与人往来}，将祂的生命带给人，藉由可见的人性传授教训，并给予有理性者合宜的生活与交往的法则。这法则祂以奇迹证实，并藉肉身的死亡摧毁了死亡。藉着复活肉身，祂向我们所有人赐下了复活的应许，以祂自己宝贵身体的复活作为我们复活的可靠保证。祂如此爱世人，即使世人恨祂，以至于这\OriginalTerm{救恩计划}{\GreekText{οἰκονομία}}的奥迹因祂苦难的剧烈程度而难以被一些人相信，而祂慈爱的深度足以显明：祂至今仍日日呼唤那些不信的人。祂这样做，并非因为祂需要人的服务——宇宙的创造者会缺少什么呢？——而是因为祂渴望每一个人的救恩。那么，我卓越的朋友，请抓住祂的恩赐；赞美施恩者，并为我们预备一场非常盛大而美好的盛宴。\par

\SourceRecord{Translated by Blomfield Jackson.；Nicene and Post-Nicene Fathers, Second Series；Vol. 3.；Edited by Philip Schaff and Henry Wace.；Buffalo, NY: Christian Literature Publishing Co.,；1892.；<http://www.newadvent.org/fathers/2707076.htm>.}

% structure: p0001 source=h1 target=section reason=page-title

\section{第七十七封书信}

\Emph{致\Place{波斯属亚美尼亚}的主教\Person{欧拉里乌斯}。}\par

我知道撒殚想要筛你们像筛麦子一样，主允许他如此行，为要显明麦子，考验金子，加冕竞赛者，宣布得胜者的名字。尽管如此，我仍恐惧战栗，并非为你们这些真理的高贵战士而忧愁，而是因为我知道会发生有些人心志软弱。如果在十二位宗徒中有一个成了叛徒，那么毫无疑问，在数倍于其数目的人群中，任何人都不难发现许多人未能达至完美。如此思考，我深感困惑与沮丧，因为，正如神圣宗徒所说：\BibleQuote{若一个肢体受苦，所有肢体都一同受苦。}我们互为肢体，组成一个身体，以主基督为头。然而在我焦虑中尚有一安慰，就是想到你的圣德。因为你既在天主圣言中长大，并由总牧人教导何为善牧的标记，毫无疑问你将为羊群舍命。因为主说：\BibleQuote{雇工看见狼来，就逃跑，因为他是雇工，不关心羊；但善牧为羊舍命。}同样，最好的将军不是在和平时期显示其天生的勇武，而是在战争时期，既要激励他人，自己也置身险境为部属。因为他若享有指挥官的尊荣，却在需要时躲避危险，岂非荒谬？三倍蒙福的先知们就是这样行的，他们轻看身体的安全，为了那些恨恶离弃他们的犹太人，忍受各种危险和劳苦。关于他们，神圣宗徒说：\BibleQuote{他们被石头砸死，被锯锯断，被试探，被刀剑杀死；披着绵羊山羊皮漂泊，困苦、患难、受折磨，世界不配拥有他们；在旷野、山岭、山洞、地穴中漂流。}同样，神圣宗徒们走遍世界传道，没有家、床、铺盖、饭食或任何生活必需品，却受鞭打、刑架、监禁，经历无数种死亡。他们经历这一切，不是为朋友，而是自愿面对这些危险，为了那些迫害他们的人。如今你更有义务接受当前的危险，为了同信仰的弟兄和子女们。甚至无理性的动物也显示这种感情：可见麻雀竭力为幼雏争斗，拼尽全力保护它们；其他鸟类也为幼雏冒危险。但我何必言及鸟类？熊、豹、狼、狮，也都为保护幼崽甘受任何痛苦，因为它们不躲避猎人，反而等待攻击，为幼崽战斗。\par

我列举这些例子，并非以野兽的榜样来膏抹你的虔诚以忍耐和勇敢，而是为了安慰我的沮丧，并确信你不会在狼群进攻时使基督的羊群没有牧人，反而会呼求羊群之主帮助，并全心为羊群争战。这样的危机证明谁是牧人，谁是雇工；谁殷勤牧养羊群，谁反而喝羊奶，不关心羊的安全。\BibleQuote{但天主是信实的，必不容许你们受试探超过你们所能承担的；在受试探时，也必开一条出路，使你们能承受得住。}但我恳求你的尊威，要更加留意软弱者；不仅要坚固动摇者，也要扶起跌倒者，因为牧人绝不忽略羊群中患病的，而是把它们从群体中分开，用尽一切方法医治它们，我们也必须如此行。我们必须使那些跌倒的站起来，伸出援手，说鼓励的话。当他们被咬伤时，我们必须医治；不可放弃拯救他们的努力，也不可让他们留在魔鬼的口中。神圣宗徒\Person{保禄}就是这样行的；当\Place{迦拉达}人领受了救恩的圣洗和圣神的恩赐后，却堕入犹太主义的病症，接受了割礼，他比最慈爱的母亲更悲痛哀哭，照料他们，使他们摆脱那软弱。我们可以听见他呼喊：\BibleQuote{我的孩子们！我为你们再受生产之苦，直到基督在你们内形成。}同样，对\Place{格林多}人中那犯了可憎淫乱的人，他既像父亲一样责罚，又巧妙医治；在第一封书信中将他开除后，在第二封中重新接纳，并说：\BibleQuote{所以你们倒不如宽恕他、安慰他，免得这样的人被过度的忧伤吞没。}又说：\BibleQuote{免得撒殚占了我们的便宜，因为我们不是不晓得他的计谋。}同样，对于那些吃祭偶像之物的人，他恰当地斥责，适当地劝诫，并使他们脱离严重的错误。\par

所以，我们的主 \Person{耶稣基督} 允许第一位宗徒（祂曾将他的宣认作为教会的一种基础和根基）摇摆不定并否认祂，然后又再次扶起他。这样，祂给了我们两个教训：不要自信于自己的力量，并坚固动摇者。因此，我恳求你，向跌倒的人伸出援手，\Quote{将他们从可怕的坑中、从污泥中拉上来，把他们的脚立在磐石上}，并\Quote{将新歌放入他们口中，赞美我们的天主}，好使他们的生活榜样成为得救的榜样，使\Quote{许多人看见而畏惧，并信赖主}。要禁止他们参与神圣奥迹，但不要阻止他们参加慕道者的祈祷，也不要阻止他们聆听圣经和教师的劝勉；要禁止他们领受神圣奥迹，不是直到死亡，而是在一段特定时间内，直到他们认识自己的疾病，渴望健康，并为离弃真君王而投靠暴君、离开恩主而投奔仇敌而真诚痛悔。\par

同样的教训也由神圣而蒙福的教父们的训导赐给我们。我这样写，并非要教导你虔敬，而是如同兄弟般提醒你，深知即使最优秀的舵手在风暴时刻也需要来自船员的提醒。同样，那位伟大的、闻名的\Person{梅瑟}，举世闻名，行了那些伟大的奇事，并未拒绝\Person{耶特洛}的建议，此人仍陷于拜偶像的错谬中；因为他并未看重他的不敬虔，而承认了他建议的得当。此外，我恳请你的虔敬为我向天主恳切祈祷，使我在余生能依照祂的律法生活。\par

我藉着最可敬、最虔诚的长老\Person{斯德法诺}写了此信，因他品格的良善，我见到他十分欢喜。\par

\SourceRecord{Translated by Blomfield Jackson.；Nicene and Post-Nicene Fathers, Second Series；Vol. 3.；Edited by Philip Schaff and Henry Wace.；Buffalo, NY: Christian Literature Publishing Co.,；1892.；<http://www.newadvent.org/fathers/2707077.htm>.}

% structure: p0001 source=h1 target=section reason=page-title

\section{\WorkTitle{书函 七八}}

\Emph{致亚美尼亚波斯的主教\Person{欧瑟伯}。}\par

每当舵手发生什么事，要么在前舷指挥的军官，要么级别最高的水手接替他的位置，并非因为他成了自命的舵手，而是因为他为船的安全着想。同样，在战争中，当指挥官倒下时，首席护民官接过指挥权，并非企图以暴力夺取权位，而是因为他关心他的士兵。同样，那三位有福的\Person{弟茂德}，当被神圣的\Person{保禄}派遣时，也接替了他的位置。因此，你的虔诚应当接受舵手、船长、牧人的责任，甘愿为基督的羊群冒一切风险，不让祂的受造物被遗弃和孤独。你反而应当包扎破碎的，扶起跌倒的，使迷途者从错误中转回，保持全体健康，并跟随那些站在羊圈前与豺狼争战的好牧人。让我们也记住先祖\Person{雅各伯}的话：\BibleQuote{白日受尽炎热，黑夜受尽寒霜，睡眠也离我而去。你羊群中的公羊，我没有吃过。被野兽撕裂的，我没有带给你；我自己赔偿了损失。无论是白日被偷的，或是黑夜被偷的，你都向我索讨。}这些是牧人的标志；这些都是牧放羊群的法则。如果著名的先祖对没有理性的牲畜如此关心，并向托付给他的人如此辩护，那么我们这些受托管理理性羊群、并从万有的天主领受这托付的人，当记起主为它们舍了生命，我们岂不应做更多？谁不恐惧战兢，当他听见天主藉\Person{厄则克耳}所说的这句话：\BibleQuote{我要在牧人和羊之间施行审判，因为你们吃脂油，穿羊毛，却不牧放羊群。}又说：\BibleQuote{我立了你作以色列家的警醒者；你若没有去警戒那恶人离开他的恶道，那恶人必死在他的罪孽之中，我却要向你追讨他的血债。}主在比喻中所说的话与此一致：\BibleQuote{你这又恶又懒的仆人……你应当把我的银子交给兑换银钱的人，到我回来时，我就可以连本带利收回。}所以，起来吧，我恳求你，让我们为天主的羊群而战。他们的主就在近处。祂必定显现，驱散豺狼，荣耀牧人。\BibleQuote{上主善待等候祂的人，善待寻求祂的灵魂。}让我们不要抱怨已经兴起的风暴，因为万有的主知道什么对我们有益。因此，当宗徒请求脱离他的考验时，主不肯允准他的恳求，却说：\BibleQuote{我的恩宠为你足够了，因为我的德能在软弱中才得以成全。}让我们勇敢地承受临到我们的灾祸；正是在战争中才能识别英雄；在竞技中运动员才获得冠冕；在海浪的冲击中才显出舵手的技艺；在火中金子才被试验。我恳求你，不要只顾自己，倒要多为其他人着想，尤其是为患病的人多于为健康的人，因为宗徒的训言大声疾呼：\BibleQuote{安慰那怯懦的，扶持那软弱的。}让我们向那些躺卧的人伸出双手，让我们包扎他们的伤口，并让他们重新站立，去与魔鬼争战。没有什么比看见他们争战并还击更能使魔鬼烦恼。我们的主满有慈爱。祂接纳罪人的悔改。让我们听听祂自己的话：\BibleQuote{我指着我的永生起誓——主上主说——我决不喜欢恶人死亡，而喜欢恶人离开他的道而生存。}所以祂以誓言作前言，而祂是禁止别人起誓的，自己却发誓，为使我们相信祂多么渴望我们的悔改和得救。神圣的经卷，无论是旧约还是新约，都充满了这样的教训，并且圣教父们的训导也教导同样的道理。\par

但我写信给你，并非因为你无知；而是提醒你所知道的，就像那些安全站在岸上的人救助被风暴颠簸的人，给他们指出一块岩石，或警告一个隐藏的浅滩，或接住并拉回抛出的绳索。\BibleQuote{赐平安的天主快要将撒殚践踏在你们脚下。}并要因你们从风暴转向平静的消息而喜乐我们的耳朵，那时祂对海浪说：\BibleQuote{静默吧，安静吧。}\par

你也要为我们祈祷，因为你曾为祂的缘故经历危险，所以可以更加大胆地说话。\par

\SourceRecord{Translated by Blomfield Jackson.；Nicene and Post-Nicene Fathers, Second Series；Vol. 3.；Edited by Philip Schaff and Henry Wace.；Buffalo, NY: Christian Literature Publishing Co.,；1892.；<http://www.newadvent.org/fathers/2707078.htm>.}

% structure: p0001 source=h1 target=section reason=page-title

\section{第七十九封书信}

\Emph{致贵族\Person{阿纳托利}}\par

主天主于现今赐予我们阁下来为莫大的安慰，并在风暴中为我们提供了合宜的避风港。故此，我们满怀信心地将我们的苦楚禀告主上。不久前，我们曾告知阁下，可敬的伯爵\Person{鲁福}向我们出示了一份皇帝亲笔书写的谕令，命令这位英勇的将军谨慎而勤勉地安排我们在\Place{居鲁士}的居所，不许我们前往另一座城市，借口是我们企图在\Place{安提约基雅}召集主教会议，扰乱正统信仰者。如今，我向您表明，我已服从皇帝诏令来到\Place{居鲁士}。过了六七天，他们派遣虔诚的统领\Person{优弗洛尼乌斯}，携带一封信函，求我书面承认皇帝谕令已向我出示。因此，我承诺留在\Place{居鲁士}及其邻近地区，牧养托付给我的羊群。为此，我恳求阁下确切查明，这些命令是否确实已颁布，以及出于何故。我确实自知还有许多其他罪过，但我并不知自己曾得罪过天主的教会或公共秩序。我如此写信，并非因为我厌恶住在\Place{居鲁士}，因为实际上，她对我来说比任何最著名的城市都更珍贵，因为我在那里的职务是由天主赐予的。但我并非出于自愿，而是被迫被束缚于她，这似乎有些令人痛心；此外，这确实给了恶人以把柄，使他们胆大妄为，拒绝听从我们的劝诫。\par

在此情形下，我恳求主上：若此类命令并未真正发出，请告知我；但若那封信函确来自得胜的皇帝，请禀告他的虔诚陛下，不可轻信诽谤，也不可只听控告者的话，而应要求被控告者作出解释。然而，事实上，仅凭事实本身便足以说服他的虔诚之心，指控我的罪名是虚假的。因为何时我曾对陛下或他的主要官员冒犯过任何事？或者何时我曾令此地诸多显赫的业主感到厌恶？相反，阁下深知，我曾将教会收入的相当一部分用于建造柱廊、浴场，修建桥梁，并为公共事业进一步提供经费。但若有人不满我为\Place{腓尼基}教会的倾覆而哀恸，那么请主上须知，当我看见犹太人的角高扬，而基督徒垂泪悲伤、被驱逐到地极时，我不可能不悲痛。我们无法对抗宗徒的法令，因为我们记得宗徒的话说：\BibleQuote{应当服从天主而不是人}，且对我们来说，比今生任何痛苦更可怕的，是主\Person{基督}的\BibleQuote{审判台}，我们众人将站立在它面前，为我们的言行交账。为了那审判台，必须忍受此生的艰辛。对于受屈者而言，对将来的盼望已是足够的安慰，但仁慈的主更在您，最卓越的先生，您的一生以虔诚和信德闪耀光芒，之中赐予了我们额外的安慰。\par

\SourceRecord{Translated by Blomfield Jackson.；Nicene and Post-Nicene Fathers, Second Series；Vol. 3.；Edited by Philip Schaff and Henry Wace.；Buffalo, NY: Christian Literature Publishing Co.,；1892.；<http://www.newadvent.org/fathers/2707079.htm>.}

% structure: p0001 source=h1 target=section reason=page-title

\section{Letter 80}

\Emph{致长官欧特罗基乌斯书}\par

我大为惊讶，阁下竟未向我通报针对我的阴谋。要挫败这些阴谋，对任何无力揭穿其煽动者谎言的人而言，很可能是一件困难的事；但通报所发生之事，所需的不在于能力，而在于友善。我们曾希望，当阁下奉召至帝都，并被选任装饰长官的高位时，教会的一切风暴将会平息。然而，我们遭遇的纷扰，甚至比争论初期所见更为严重。腓尼基的教会处于困境；巴勒斯坦的教会也处于困境，众口一词如此报告；最虔诚的主教们的书信也证实了这种痛苦。我们中间的众圣人均在叹息，每一个虔诚的会众都在哀悼。正当期盼从前苦难的止息时，我们却遭受了新的苦难。我自己已被禁止离开居鲁士的海岸——如果给我看的那份文书是真实的，它据说是我们得胜皇帝的手谕。内容如下：\Quote{既然此城的主教某某不断召集会议，这对正教徒造成了困扰，你当以适当的勤勉和智慧留意，使他常驻居鲁士，不得离开前往另一城市。}我接受了判决，仍留在这裡。阁下可以见证我的情感，因为你知道我到达安提约基雅后，因那些想留我的人而匆忙离开。那些偏听诽谤我的人，不肯留一只耳朵给我的人，无疑是错误的。即使对杀人犯和玷污他人床榻的人，也给予他们辩护的机会，他们不会在未经当面定罪或对所控罪状认罪之前就受判决。然而，一位大司祭，在隐修院度过前半生后担任主教已二十五年，从未搅扰法庭，也从未被任何人起诉，却像被当作诽谤的玩物，甚至不被给予盗墓者常有的特权——询问指控是否属实。然而他们行了恶；我未行恶。但我准备承受更严重的麻烦。即使他们对我哀悼腓尼基的灾祸感到恼火，只要我看见这些灾祸，我就不会停止哀悼。对我而言，唯一可畏的审判是天主的审判。然而，我为他们祈祷，愿他们从万有的天主那里获得宽恕；为阁下祈祷，愿你永远尊荣地生活，在一切美善上卓越，勇敢地反对谎言，为真理而战。愿这些阴谋的策划者知道，即使我前往天涯海角，天主也决不容许不敬神的教义得到确认，祂必点头并消灭那些向可憎教义屈膝的人。\par

\SourceRecord{Translated by Blomfield Jackson.；Nicene and Post-Nicene Fathers, Second Series；Vol. 3.；Edited by Philip Schaff and Henry Wace.；Buffalo, NY: Christian Literature Publishing Co.,；1892.；<http://www.newadvent.org/fathers/2707080.htm>.}

% structure: p0001 source=h1 target=section reason=page-title

\section{第八十一封书信}

致执政官\Person{诺穆斯}。\par

我只享受了与阁下交往的短暂时光，因为我被无法避免的环境剥夺了我如此渴望的东西。我曾希望我们的短暂会面能点燃善意和友好交往，但我失望了。我现在已写了两次信给你，但没有收到任何回音；根据皇帝谕旨，我被禁止离开\Place{居鲁士}的疆界。这种表面惩罚的原因并不存在，除了我召集了主教会议这一事实。没有公布起诉书；没有出现控告者；被告未被定罪；但判决已经下达。我们服从，因为我们知道受冤者的赏报。然而我知道，当犹太人要求处死神圣的\Person{保禄}时，受托管理犹太人的总督\Person{斐斯托}公开回答说：\BibleQuote{我们罗马人不能将任何人定罪，除非被告与原告当面对质，并获准就所控告的罪行进行申辩。}这些话出自一个不相信我们主人基督、而是多神论谬误的奴隶之口。从未有人询问我是否在召集会议，或者我为何召集会议，或者这可能给教会或政府带来什么不快；然而，好像我是一个非常有罪的罪犯，我被禁止访问其他城市；而对其他人而言，每座城市都是开放的，不仅对\OriginalTerm{亚略派}{Arians}和\OriginalTerm{欧诺弥派}{Eunomians}，也对\OriginalTerm{摩尼派}{Manichees}和\OriginalTerm{马西翁派}{Marcionists}，对那些患有\OriginalTerm{华伦提努派}{Valentinus}和\OriginalTerm{孟他努派}{Montanus}病症的人，甚至对异教徒和犹太人，而我，福音教义的首要捍卫者，却被从每座城市驱逐。然而有人认为我并未持守福音。那么，召开一次公议会吧！让那些有判断力的虔诚主教聚集在那里；然后让那些在神圣学识中受过教育的官员和显贵聚集起来。让我陈述我的主张，让裁判们宣告哪种意见符合宗徒的教导。我这样写，并非出于想见大城的渴望，也不是试图前往其他城市。事实上，我更喜欢那些希望以隐修方式管理教会之人的宁静。我希望殿下知道，无论在真福且已列圣品的\Person{狄奥多托}时代，还是在有福纪念的\Person{若望}时代，抑或在至圣的主教\Person{多米诺斯}大人时代，我都未曾自愿进入\Place{安提约基雅}；我曾被邀请五六次，但勉强同意；我同意时，也是遵守教会的教规，该教规命令被召集到会议却拒绝出席的人被视为有罪。当我出现时，我做了什么令天主不悦的事呢？难道是我将某些犯有不可言说之恶者的名字从神圣名单中删除？难道是我祝圣了一些品德高尚、生活可敬的人为司铎？难道是我向百姓宣讲了福音？如果这些事情应受控告和惩罚，我乐意为此迎接更严厉的惩罚。我的控告者迫使我发言。在我受孕之前，我的父母就应许将我奉献给天主；从襁褓之中，他们照此应许奉献了我，并相应地培养我；在我担任主教之前的时光，我在隐修院度过，然后不情愿地被祝圣为主教。二十五年中，我如此生活，从未被任何人传唤受审，也从未控告过任何人。我手下虔诚的圣职人员中，没有一人经常出入法庭。在这些年中，我从未从任何人那里收取过一枚奥博尔或一件衣服。我的仆人中也没有一人收过一片面包或一个蛋。我无法忍受拥有任何东西，除了我身穿的破衣。我用自己的教区收入建造了公共柱廊；修建了两座大桥；照管了公共浴场。发现城市没有被流经的河流供水时，我修建了水道，为干涸的城镇供水。但且不提这些事，我带领了八个\OriginalTerm{马西翁派}{Marcionists}的村庄及其周边地区走上真理之路；又带领了另一个充满\OriginalTerm{欧诺弥派}{Eunomians}的村庄和另一个\OriginalTerm{亚略派}{Arians}的村庄归向神圣知识的真光，并且，靠着天主的恩宠，在我们中间没有留下一棵异端的莠子。这一切我并非未受惩罚而行；多次我流了血；多次我被他们用石头砸，被带到死亡的门口。但我是在自夸中的愚人，然而我的话是出于必要，而非自愿。有一次，三次有福的\Person{保禄}被迫同样行事，以堵住控告者的口。然而我忍受表面的耻辱，视为极大光荣，因为我听见宗徒的声音呼喊说：\BibleQuote{凡愿意在基督耶稣内虔敬度日的，都要遭受迫害。}\par

但我恳请阁下关注教会的事务，平息已兴起的风暴，因为事实上即使在争论之初，教会也未曾被如此混乱困扰。没有人告知您危险之大，\Place{腓尼基}基督徒的悲叹，以及我们最圣洁的隐修士们的哀号。因此我较长地写信给您，希望了解教会的动荡后，阁下能予以制止，并收获此善举带来的果实。\par

\SourceRecord{Translated by Blomfield Jackson.；Nicene and Post-Nicene Fathers, Second Series；Vol. 3.；Edited by Philip Schaff and Henry Wace.；Buffalo, NY: Christian Literature Publishing Co.,；1892.；<http://www.newadvent.org/fathers/2707081.htm>.}

% structure: p0001 source=h1 target=section reason=page-title

\section{第八十二封信}

\Emph{致\Place{安西拉}主教\Person{欧瑟比}。}\par

我此前曾希望时常听到尊驾的消息。我蒙受纯属诽谤的指控，亟需兄弟般的安慰。因为那些如今重新鼓吹马尔西翁、瓦伦提努斯、摩尼及其他幻象说者异端的人，恼恨我公开揭露他们的异端，便竭力欺骗皇帝圣听，称我为异端者，并诬告我将我们的主耶稣基督——降生成人的天主圣言——分为两个儿子。他们的言论并未如预期得逞。于是便致函尊贵的统军与执政官，信中虽无异端指控，却另有同样毫无根据的责难。他们声称我企图在安提约基雅频繁召集主教会议；某些人因此不悦；为此我应停止此类行动，管理好托付给我的各教会。当此函出示于我时，我抓住其中一句作为行善的机会。首先，我获得了久盼的安息；此外，我深信，因真理仇敌对我图谋的恶事，我诸多过犯的污点将被涂抹。即使在此生，至高的统治者极其清楚地表明祂多么眷顾受冤屈的人。当我安息在家乡的边界内，当东方众人皆忧愁痛泣，却因降临的恐慌而被迫沉默（因发生在我身上的事令众人心惊胆战），主从天垂顾，证实了诬告者的谎言，揭露了他们不敬的意图。他们甚至策动亚历山大里亚反对我，借其得力工具向众人鼓吹我宣讲两个儿子而非一个。\par

与此相反，我远未持有此种可憎见解，以致当我发现尼西亚大公会议的一些圣教父在其著作中反对亚略的疯狂，并在与对手斗争中被迫作出过分的区分时，我曾提出异议，拒绝承认这种区分，因为我知道区分的需要会导致夸张。\par

唯恐有人以为我因恐惧而如此言语，任何人若愿意，可取我在厄弗所大公会议前写的古老著作，以及十二年前会后所写的作品。因天主的恩宠，我诠释了全部\WorkTitle{先知书}、\WorkTitle{圣咏集}和\WorkTitle{宗徒著作}；我早已写下反对亚略派、马其顿派、阿波利纳里的诡辩和马尔西翁疯狂的作品；在我的每一本书中，因天主的恩宠，教会的思想清晰闪耀。此外，我写了\WorkTitle{论奥迹}、\WorkTitle{论天主眷顾}、\WorkTitle{论贤士的问题}、\WorkTitle{圣人列传}，此外还有许多，不再一一细述。\par

我列举这些并非出于野心，而是为了挑战我的控告者和审判者，请他们任意检验我的作品。他们将会发现，因天主的恩宠，我所持的见解无非是从圣经所领受的。\par

尊驾既闻此，我恳请告知无知者，并劝说那些毁谤我及受其欺骗之人的放肆口舌，不要相信他们从诬告者处听来的关于我的话。请他们相信立法者的呼喊：\BibleQuote{不可接受虚假的报告。}请他们等待事实得以证实。\par

我祈求各教会得享平静，这场漫长痛苦的暴风雨消散。但若我们罪孽众多，不容此事成就；若因我们的罪，我们被交付于筛子者，我们祈求能分担为信德所受的危险，以便我们既无此生的自信，至少因持守完整的信德，在主显现之日得蒙怜悯和宽恕。为此我们恳请尊驾与我们同祷。\par

\SourceRecord{Translated by Blomfield Jackson.；Nicene and Post-Nicene Fathers, Second Series；Vol. 3.；Edited by Philip Schaff and Henry Wace.；Buffalo, NY: Christian Literature Publishing Co.,；1892.；<http://www.newadvent.org/fathers/2707082.htm>.}

% structure: p0001 source=h1 target=section reason=page-title

\section{第八十三封信}

\Emph{\Person{泰奥多勒图斯}，\Place{居鲁士}主教，致\Person{狄奥斯哥鲁斯}，\Place{亚历山大}的总主教。}\par

对于遭受诬告的人来说，圣经的话语给了极大的安慰。当这样一个受苦者被无缰之舌的谎言所伤，感到痛苦的尖刺时，他想起可敬的\Person{若瑟}的故事；当他看到那个贞洁的模范、一切美德的典范，因诬告侵犯他人床第而被囚禁、戴上镣铐，在牢狱中度过漫长岁月时，这故事所提供的良药减轻了他的痛苦。同样，当他发现温良的\Person{达味}被\Person{撒乌尔}像暴君一样追捕，随后又抓住敌人并安然放走时，他的痛苦也得到了止痛剂。但当他看到主基督自己——万世的创造者、万物的造主、真天主、真天主子——被邪恶的\Person{犹太人}称为贪食者和酗酒者时，他得到的不仅是安慰，更是极大的喜乐，因为他被认为配得分享主的苦难。\par

因此，当我读到您的圣座致最虔诚神圣的总主教\Person{多姆努斯}的信时，我不得不写信，因为信中所述，有人来到由您的圣座管理的荣城，指控我在\Place{安提约基雅}讲道时，将唯一的主耶稣基督分为两个儿子，而那里有无数的听众聚集成群。我为那些胆敢对我捏造虚空诽谤的人哭泣。但我感到悲伤，我的主，请原谅我，因为痛苦迫使我说话：您虔诚的卓越没有为保留一只无偏见的耳朵给我，反而相信了控告者的谎言。然而他们不过三、四人或大约十二人，而我却有无数听众可以证明我教导的正统。我在\Place{安提约基雅}已故主教、神圣记忆的\Person{泰奥多图斯}时期连续教导了六年，他以卓越的事迹和对神圣教义的学识而闻名。在神圣受福记忆的\Person{若望}主教时期，我教导了十三年，他对我讲道如此喜悦，以至于高举起双手，一次又一次地站起来；您的圣座在自己的信中作证，他从小就在神圣的圣言中成长，对神圣教义的认识极为精确。除此之外，这是最虔诚的主总主教\Person{多姆努斯}的第七年。直到今日，过了那么长的时间，没有一位虔诚的主教，没有一位虔诚的圣职人员曾对我的言论提出任何指责。而基督徒民众是多么欣喜地聆听我们的讲道，您虔诚的卓越从那些从那里旅行到此地的人，以及从我们这里到达您城市的人那里，都容易得知。\par

我这样说并非为了夸耀，而是因为我被迫自卫。我不是在引起对我讲道名声的注意，而只是关注它们的正统性。就连那位伟大的世界导师，他惯于自称圣徒中最末的、罪人中的首恶，为了堵住说谎者的口，也被迫列举自己的辛劳；为了表明这番关于他受苦的叙述是出于必要，而非出于自愿，他补充说：\BibleQuote{我成了愚妄的人，是你们逼我的。}我承认自己可怜——是的，三倍可怜。我犯了许多过错。我只靠信德指望在主显现的日子找到一些怜悯。我渴望并祈求能跟随圣父们的足迹，我热切盼望保持福音教导的无玷，这教导总括而言是由在\Place{比提尼亚}的\Place{尼凯亚}举行的神圣公议中的圣父们所传授给我们的。我相信有一个天主圣父，和一位发自圣父的圣神；同样，也有一位主耶稣基督，天主的独生子，在万世之前由圣父所生，是祂光荣的辉耀，圣父本体的确切形像，为人类的救恩而降生成人、成为人，并由童贞\Person{玛利亚}在肉身中诞生。因为聪明的\Person{保禄}这样教导我们：\BibleQuote{圣祖也是他们的，并且按肉身说，基督也出自他们；祂是万有之上，永受赞美的天主。阿们。}以及\BibleQuote{论到祂的圣子我们的主耶稣基督：按肉身说，是从\Person{达味}的后裔生的；按圣德的神性说，由于从死者中复活，被立为具有大能的天主子。}为此，我们也称圣童贞者为\OriginalTerm{天主之母}{\GreekText{Θεοτόκος}}，并认为反对这一称号的人已远离真宗教。\par

同样，我们称那些将我们唯一的主耶稣基督分为两个位格或两个儿子或两个主的人为败坏者，并将他们排除在基督徒的集会之外，因为我们听到非常神圣的\Person{保禄}说：\BibleQuote{只有一个主，一个信德，一个洗礼}，又说：\BibleQuote{只有一个主耶稣基督，万物都是藉着祂而有的}，又说：\BibleQuote{耶稣基督昨天、今天、直到永远，常是一样的}，又在另一处：\BibleQuote{那降下的，就是那升到诸天之上以上的。}在宗徒的著作中可以找到无数其他此类段落，宣明一个主。\par

同样，神圣的福音宣讲者也大声说：\BibleQuote{圣言成了血肉，住在我们中间；我们看见了他的光荣，正如父独生者的光荣，满溢恩宠和真理。}\par

而与他同名的（\Person{若翰}）宣告说：\BibleQuote{在我以后来的，成了在我以前的，因祂原先我而有。}当他显示了一个位格时，他同时表达了神性和人性，因为\Quote{人}和\Quote{来}这些词是属人的，而\Quote{祂原先我而有}这句话表达了神性。但尽管如此，他并没有承认那后来者和那在前者之间的区别，而是承认同一存在作为天主是永恒的，但作为人，由童贞女所生，是在他自己之后。\par

同样，三倍有福的\Person{多默}，当他将手放在主的肉身上时，称祂为主和天主，说：\BibleQuote{我的主，我的天主！}因为通过可见的本性，他辨识了那不可见的。\par

因此，我们不知道同样的肉身和神性之间有何区别，而是承认成为人的天主圣言是一位圣子。\par

这些教训我们都从圣经和解释圣经的圣教父们学得，即\Person{亚历山大}、\Person{亚大纳削}——真理的洪亮传声筒，你们宗座的光荣——以及\Person{巴西略}、\Person{额我略}和其余世界的明灯。而且，我们在努力封住胆敢反对真福\Person{德奥斐禄}和\Person{济利禄}之人的口时，我们使用了他们的著作，我们自己的作品可为证。因为我们极其渴望用至圣之人提供的药物来医治那些否认主的肉躯与天主性之间的区别，并且时而主张天主性被改变为肉躯，时而主张肉躯被转变为天主性之本性的人。\par

因为他们清楚地教导我们两性之间的区别，宣告天主性不可改变，称主的肉躯为天主性的，因为它是天主圣言所成的肉躯；但将其转变为天主性之本性的教义，他们斥为不虔。\par

我想大人您很清楚，真福\Person{济利禄}常写信给我，当他将反对\Person{儒利安}的著作以及类似地论替罪羊的书寄到\Place{安提约基雅}时，他请求安提约基雅主教真福\Person{若望}将这些书展示给东方的大师们；应此请求，真福\Person{若望}将书寄给了我。我阅读后赞叹不已，并写信给真福\Person{济利禄}；他回信称赞我的精确与友善。这封信我保存至今。\par

我曾两次签署真福\Person{若望}关于\Person{聂斯多略}的著作，这是我的亲笔所证；但那些试图通过诽谤我来掩盖自己错误的人，就是这样在背后议论我的。\par

因此我恳请尊驾远离说谎者；关注教会的安宁，要么用有益的药物医治那些试图摧毁真理教义的人，要么若他们拒绝接受您的治疗，就将他们逐出羊栈，使羊群免于传染。我请求您赐予我惯常的问候。我向您袒露了真实的心声，这由我关于圣经以及反对\OriginalTerm{亚略派}{Arians}和\OriginalTerm{欧诺米派}{Eunomians}的著作所证明。\par

我将再写几句简短的话。若有人拒绝承认圣童贞为\OriginalTerm{天主之母}{Theotokos}，或称我们的主耶稣基督为纯粹的凡人，或将唯一受生者和首生者分裂为两个儿子，我愿他丧失对基督的望德，愿全体人民说：阿们，阿们。\par

既然我已如此说，请我主赐予我神圣的祈祷，并以一封回信使我欣慰，告知我尊驾已不再理会我的控告者。\par

我和我全家问候你们在基督内虔敬的众弟兄。\par

\SourceRecord{Translated by Blomfield Jackson.；Nicene and Post-Nicene Fathers, Second Series；Vol. 3.；Edited by Philip Schaff and Henry Wace.；Buffalo, NY: Christian Literature Publishing Co.,；1892.；<http://www.newadvent.org/fathers/2707083.htm>.}

% structure: p0001 source=h1 target=section reason=page-title

\section{第八十四封书信}

\Emph{致基里基雅诸位主教}\par

你们的虔敬已听闻对我的诽谤。真理的反对者声称，我将我们唯一的主耶稣基督，天主的独生子，分割为两个儿子；有人说，他们诽谤的依据，是从你们中间一小撮持此意见的人而来，他们将降生成人的天主圣言分割为两个儿子。他们本该听从宗徒的话，他公开宣告：\BibleQuote{只有一个主耶稣基督，万物藉祂而有}，又说：\BibleQuote{只有一个主，一个信德，一个圣洗}。他们本该追随主的教导，因为主亲自说：\BibleQuote{没有人升过天，除了那从天上降下来而仍在天上的人子}。又说：\BibleQuote{如果你们看见人子升到他先前所在之处，你们将如何呢}。而圣洗的传统教导我们，只有一个圣子，正如只有一个圣父和一个圣神。因此，我希望你们的虔敬能俯允，若真有（虽然我无法相信）不服从宗徒训诲的人，照教会的法规斥责他们，封住他们的口，并教导他们追随圣教父的足迹，保存由圣洁可敬的教父在\Place{比提尼亚}的\Place{尼西亚}所订立、并总结福音作者和宗徒教导的信德，纯洁无瑕。因为爱天主的人，理当留意天主的荣耀和我们共同的声誉，不可忽视这少数人因无知或好争而对我们全体发起的攻击——若他们果真有过失；若他们并非无辜，而像我们一样，遭受到诬告者尖刻舌头的伤害。\par

请垂念我们，在你们向天主的祈祷中纪念我们，因为爱德的法律如此规定。\par

\SourceRecord{Translated by Blomfield Jackson.；Nicene and Post-Nicene Fathers, Second Series；Vol. 3.；Edited by Philip Schaff and Henry Wace.；Buffalo, NY: Christian Literature Publishing Co.,；1892.；<http://www.newadvent.org/fathers/2707084.htm>.}

% structure: p0001 source=h1 target=section reason=page-title

\section{第八十五封信}

\Emph{致\Person{巴西略}主教。}\par

神圣的\Person{保禄}称首要的美善是爱，并以爱吩咐信德的乳儿得到喂养。你的虔诚拥有丰厚的爱，因此告知了合宜的事，并带来了令人愉悦的消息。因为对于敬畏主的人，还有什么比真理教义的健康与和谐更令人愉悦的呢？最虔诚的先生，请确信，我们听到我们共同朋友的消息时非常喜乐；而且，正如我们先前因听闻他将肉身与天主性的本性描述为同一，并公开将救恩的苦难归于不可受苦的天主性而深感悲痛，同样，当我们读到你的圣德的信函，得知他完整地保存了诸本性的特性，既否认天主圣言变化为肉身，也否认肉身变化为天主性的本性，反而坚持在独生子、我们的主耶稣基督——降生成人的天主圣言——中，每一本性的特性毫不混淆地持存时，我们都极其欢喜。我们为万有的天主赞美神圣信德的和睦。然而，虽然我们的情报不完整，但我们仍写信给\Place{基里基雅}两地，询问是否真有真理的反对者，并嘱咐虔诚的主教们搜查并查明是否有将独一的主耶稣基督分为两个儿子的人，或通过劝诫使他们醒悟，或将他们从兄弟名册中剔除。因为事实上，我们同等地弃绝那些胆敢主张肉身与天主性为同一本性的人，以及那些将独一的主耶稣基督分为两个儿子并试图超越宗徒定义的人。\par

但请你的圣德确信，我们倾向于和平。因为如果先知说：\Quote{与憎恨和平的人，我保持和平，}那么我们就更乐意欢迎天主的和平。\par

那些以谎言为食的人中，有些匆忙赶到\Place{亚历山大里亚}，捏造了对我的诽谤，结果那座城市的虔诚主教被他们的陈述所误导，尽管他已从我的信函中充分了解情况，还是派了一位虔诚的主教到帝国都城。因此我恳请你向他展示你惯常的仁慈，以真理对抗谎言。\par

\SourceRecord{Translated by Blomfield Jackson.；Nicene and Post-Nicene Fathers, Second Series；Vol. 3.；Edited by Philip Schaff and Henry Wace.；Buffalo, NY: Christian Literature Publishing Co.,；1892.；<http://www.newadvent.org/fathers/2707085.htm>.}

% structure: p0001 source=h1 target=section reason=page-title

\section{\WorkTitle{第八十六封书信}}

\Emph{\Person{弗拉维亚诺}，\Place{君士坦丁堡}主教。}\par

现今，最敬爱天主的吾主，我遭受了波涛的多次冲击，但我呼求了伟大的舵手，得以在风暴中站稳；然而，如今对我发动的攻击超越了任何悲剧故事。关于那些针对宗徒信德所策划的攻击，我以为会在最敬神的\Place{亚历山大里亚}主教，吾主\Person{狄约斯古禄}身上找到同盟和同工，于是派了一位我们虔诚的司铎，一位极其谨慎的人，携带一封主教团书信，通知他的敬爱：我们持守有福记忆的\Person{济利禄}时代所达成的协议，并接受他所写的信函，以及极有福且圣德的\Person{亚大纳削}写给有福的\Person{厄比克德多}的信函，而且在此之前，也接受由圣而有福的教父们在\Place{彼提尼雅}的\Place{尼西亚}所制定的信德声明。我们劝他促使那些不愿遵守这些文件的人立即遵守。但是反对派中的一人，维持着这些骚乱，通过欺骗一些在场的人，并制造无数针对我的诽谤，掀起了一场不义的反对我的骚动。\par

但是那最敬神的主教\Person{狄约斯古禄}给我们写了一封信，这样一封信本不该由一位从万有的天主那里学到不听虚言的人所写。他相信了那些对我的指控，仿佛亲自调查了每一项，并在审讯后得出了真相，就这样定了我的罪。然而，我勇敢地承受了这诽谤性的指控，并给他回了一封谦恭的信，向他的敬爱表明这项指控完全是假的，东方的敬神主教中没有一个人持有违反宗徒谕令的意见。而且，他派来作为信使的虔诚圣职人员，已被事实的实际证据所说服。然而，他却不加理会地遣散了他们，转而倾听我的诽谤者，其行为若非整个教会都为之作证，简直难以置信。他容忍那些对我呼喊\Quote{弃绝}的人；不仅如此，他还从座位上站起来，以声援他们的话来证实他们的话。除此之外，他还派了一些敬神的主教到帝都，据我们所知，希望加剧对我的骚动。我首先以鉴察万有的那一位为护佑者，因为我是在为天主的谕令而奋斗——其次，我恳请您的圣德为受到攻击的信德而战，并为正在被践踏的教规而战。当有福的教父们在帝都聚会，与曾在\Place{尼西亚}开会的教父们和谐一致时，他们区分了教区，并指定每个教区管理自己的事务，明确命令不得有人从一个教区侵入另一个教区。他们命令\Place{亚历山大里亚}主教独自管理\Place{埃及}的治理，每个教区管理自己的事务。\par

然而，\Person{狄约斯古禄}拒绝遵守这些决定；他正在颠覆有福的\Person{玛尔谷}的宗座；这些事他做，尽管他完全知道\Place{安提约基雅}大都会拥有伟大的\Person{伯多禄}的宗座，\Person{伯多禄}是有福的\Person{玛尔谷}的老师，也是宗徒唱诗班的首位和\OriginalTerm{领袖}{\GreekText{κορυφαῖος}}。\par

然而，我知道这宗座的威严，我也知道并衡量自己。我从起初就学到了宗徒的谦逊。我恳求您的圣德不要忽视圣教规被践踏，而要热心地挺身而出，作天主信德的护佑者，因为在这信德中我们有得救的望德，并因此确信我们将蒙受怜悯。\par

但是，为使您的圣德不至于对此无知，吾主，请知道，自从我为了服从圣教父们的教规而赞同在有福记忆的\Person{普罗克路斯}时代由您的宗座发出的主教团书信以来，他就对我显出了恶意；关于这一点，他屡次责备我，说我违反了\Place{安提约基雅}教会以及，如他所说，\Place{亚历山大里亚}教会的权利。他记着这一点，并且，如他所想，找到了机会，就显示了他的敌意。但没有什么比真理更强。真理甚至以少言就能获胜。我恳求您的圣德在向主的祈祷中记念我，使我能够战胜四面冲击我的波浪。\par

\SourceRecord{Translated by Blomfield Jackson.；Nicene and Post-Nicene Fathers, Second Series；Vol. 3.；Edited by Philip Schaff and Henry Wace.；Buffalo, NY: Christian Literature Publishing Co.,；1892.；<http://www.newadvent.org/fathers/2707086.htm>.}

% structure: p0001 source=h1 target=section reason=page-title

\section{第八十七封书信}

\Emph{致\Place{阿帕梅亚}主教\Person{多姆努斯}}。\par

兄弟友爱的律法要求，此时我应当收到许多来自您圣德的信函。因为神圣的宗徒嘱咐我们要与哭泣的人一同哭泣，与喜乐的人一同喜乐。我连一封也没有收到，尽管不久前，您的隐修院中几位虔诚的隐修士与虔诚的司铎\Person{厄利亚}曾来探望我。尽管如此，我已写信，并向您的圣德致意；并让您知道，主的安慰已代替了一切其他，因为事实上，即使我有如头发般多的口，也不能相称地赞美祂，因我得以因宣认祂而受苦，为这表面的耻辱，我视之胜过一切光荣。即使我被放逐到地极，我更要赞美祂，因我被算为配得更大的福分。不过，我期望您的圣德能为圣教会的安宁祈祷。正是因为这风暴袭击她们，我才哀哭、叹息和悲伤。据我所知，那份安宁被\Place{奥斯洛厄内}的圣职人员赶走了，他们向我倾泻了无数话语，尽管我没有参与对他们的谴责，也没有参与对他们的判决；相反，正如您的圣德所知，我恳求在复活节时将共融赐予他们。但毁谤者可能毫不费力地说他们喜欢的话。我的安慰在于主的祝福，祂说：\BibleQuote{几时人为了我而辱骂迫害你们，捏造一切坏话毁谤你们，你们是有福的。你们欢喜踊跃罢！因为你们在天上的赏报是丰厚的，因为他们在你们以前的先知，也曾这样迫害过。}\par

\SourceRecord{Translated by Blomfield Jackson.；Nicene and Post-Nicene Fathers, Second Series；Vol. 3.；Edited by Philip Schaff and Henry Wace.；Buffalo, NY: Christian Literature Publishing Co.,；1892.；<http://www.newadvent.org/fathers/2707087.htm>.}

% structure: p0001 source=h1 target=section reason=page-title

\section{第八十八封书信}

\Emph{致贵族\Person{陶鲁斯}。}\par

毁谤者迫使我超越了节制的界限，并强使我写信给您——您曾装饰最高官职，获得最显赫的荣誉。因此我恳求您宽恕我，因为我并非出于自满而写信，而是被需要所推动。并非因为我预料会不义地陷入患难与困苦——这是所有真诚侍奉天主之人的共同命运，而是因为我渴望说服阁下：那些控告我见解的人正对我提出不实的指控。我自幼从母胎中便受宗徒训诲的滋养，圣洁可敬的教父们在尼西亚所订立信经，我既学习也教导。凡持任何其他见解者，我皆斥为不敬；若有人执意声称我教导相反之道，愿他不要提出我无法辩护的控告，而是当面定我的罪。因为这合乎天主与人的法律，但谁比您，基督之友，更适合为受冤者辩护呢？您的家世光辉、地位崇高和在法律中的首要地位，赐予您直言无忌的胆量。\par

\SourceRecord{Translated by Blomfield Jackson.；Nicene and Post-Nicene Fathers, Second Series；Vol. 3.；Edited by Philip Schaff and Henry Wace.；Buffalo, NY: Christian Literature Publishing Co.,；1892.；<http://www.newadvent.org/fathers/2707088.htm>.}

% structure: p0001 source=h1 target=section reason=page-title

\section{第八十九封书信}

\Emph{致贵族弗洛伦提乌斯}\par

当我向您伟大的尊驾致书时，我正冒着我能力之外的危险；但我冒险的原因并非自信，而是那些诽谤我之人的中伤。我认为值得向您正义的耳朵说明，我意见的反对者是如何公开地诽谤我。我承认自己曾犯许多错误，但直到现在，我始终保存着宗徒们无玷的信德，并唯独因此期望在主显现的日子蒙受怜悯。为这信德，我继续与各种异端争辩；这信德，我始终传授给虔诚的乳婴；藉此信德，我将无数豺狼化为绵羊，引领他们归向那为我们众人总牧的救主。我不仅从宗徒和先知，而且从他们的著作诠释者——依纳爵、尤斯塔修斯、亚大纳削、巴西略、额我略、若望以及其余世界之光——学到这些；在我之前，还有尼西亚大公会议的神圣教父，我将他们的信仰宣认完整保存，如同祖传产业，称所有胆敢逾越其法令的人为败类与真理的仇敌。我恳请您的伟大，既已从我这里听到这些，请堵住我诽谤者的口。我认为全然不合理的是，在他人缺席时指控他们的话竟被当作真理；相反，合法且正当的是，那些想要充当控告者的人，应在被告在场时指控他们，并尽力当面定他们的罪。在这样的条件下，法官将能毫无困难地发现真理。\par

\SourceRecord{Translated by Blomfield Jackson.；Nicene and Post-Nicene Fathers, Second Series；Vol. 3.；Edited by Philip Schaff and Henry Wace.；Buffalo, NY: Christian Literature Publishing Co.,；1892.；<http://www.newadvent.org/fathers/2707089.htm>.}

% structure: p0001 source=h1 target=section reason=page-title

\section{书信第九十}

\Emph{致\Person{卢皮奇努斯}大人。}\par

我经历了年轻时的角力。我已看见老年近在咫尺，本期望作为长者能得到更多的尊敬。然而我却成了毁谤之矢的目标，被迫以辩护应对加诸于我的指控。在此情形下，我恳请大人不要相信控告者的谎言。若我一直过着默默无闻的生活，或许会有人怀疑我信仰不纯。但我不断在教会中宣讲，因此，靠天主的恩宠，我有无数的见证者证明我教导的正统。我遵循宗徒的规律和规则。我用圣洁可敬的教父们在\Place{尼西亚}所订立的信德，如同用尺子和规则一般，来衡量我的教导。若有人声称我持有相反的意见，请他当面控告我；不要在我背后毁谤我。即使被告也应有机会发言，以辩护应对对他的指控，然后，唯有然后，审判者才能合法地宣判。我愿藉着大人的协助求得这一恩惠。若有人不愿听我申辩就要定我的罪，我甚至甘愿接受他们不义的判决。因为我等待上主的审判，在那里我们既不需要证人也不需要控告者。在祂面前，正如神圣的宗徒所说：\BibleQuote{万物都是赤裸敞开的。}\par

\SourceRecord{Translated by Blomfield Jackson.；Nicene and Post-Nicene Fathers, Second Series；Vol. 3.；Edited by Philip Schaff and Henry Wace.；Buffalo, NY: Christian Literature Publishing Co.,；1892.；<http://www.newadvent.org/fathers/2707090.htm>.}

% structure: p0001 source=h1 target=section reason=page-title

\section{\WorkTitle{第九十一封书信}}

\Emph{致行政长官厄特雷基乌斯。}\par

我深知，无需言语，阁下如何看待我。行动比言语更清楚地表明了一切，但我一直渴望您能知晓那对我的指控之缘由。因为我正受一种极不寻常的责难，同时被攻击为未婚者，又曾两次结婚。如果现今的毁谤者断言我歪曲了宗徒的教导，那么，他们为何不在我面前，反而在我缺席时，试图定我的罪呢？仅此事实就足以彻底驳斥他们的谎言，正因为他们知道我教导的宗徒特性有无数的见证人，才对我提起了一项无辩护的控诉。相反，合法的法官必须为被告保留一只不偏不倚的耳朵。如果他们只听反对者的辩词，并作出合他们心意的判决，我将忍受这不义，视其使我更接近天国，并等待那公正的法庭，那里没有控告者，没有辩护人，没有证人，没有身份差别，只有对言行和正义报应的审判。\Quote{因为，}据说，\Quote{我们众人必须出现在基督的审判台前，为使各人因着身体所作的事，无论善恶，都领取相当的报应。}\par

\SourceRecord{Translated by Blomfield Jackson.；Nicene and Post-Nicene Fathers, Second Series；Vol. 3.；Edited by Philip Schaff and Henry Wace.；Buffalo, NY: Christian Literature Publishing Co.,；1892.；<http://www.newadvent.org/fathers/2707091.htm>.}

% structure: p0001 source=h1 target=section reason=page-title

\section{第九十二封书信}

\Emph{致贵族\Person{阿纳托利}。}\par

至圣的总主教\Person{多米诺斯}大人已安排最虔诚的主教们前往帝都，以便彻底驳斥对我们所有人提出的虚假控告。此时我们特别需要您的尊贵的援助，因为万有的主赐予您纯正信德、为此热忱、才智和能力，并且有能力执行您审慎的忠告。因此我恳请您为受冤者辩护，与谎言抗争，并捍卫如今受到攻击的宗徒教导。毫无疑问，教会的导师和主必会祝福您的努力，驱散阴云，并以晴朗天空祝福信德的子民。即使祂允许风暴肆虐，您的伟大将收获完美的赏报，而我们将在风暴前低头，准备无论被驱往何处都愉快地生活，并等待天主的审判和他真实而公义的判决。\par

\SourceRecord{Translated by Blomfield Jackson.；Nicene and Post-Nicene Fathers, Second Series；Vol. 3.；Edited by Philip Schaff and Henry Wace.；Buffalo, NY: Christian Literature Publishing Co.,；1892.；<http://www.newadvent.org/fathers/2707092.htm>.}

% structure: p0001 source=h1 target=section reason=page-title

\section{信函九十三}

\Emph{致\Person{元老院议员贵族}。}\par

我珍藏着对阁下尊贵不可磨灭的记忆，如今藉着非常虔诚圣善的主教们向您致敬。至圣的主教\Person{多姆努斯}大人安排他们前往\Place{帝都}，以终结针对我的不实指控。因为某些人对我罗织了明显的诽谤，并严重扰乱了教会——为了这些教会，主基督\BibleQuote{忍受了十字架，轻视了羞辱}；为了这些教会，神圣宗徒的队伍和得胜殉道者的团体交付于多种死亡。为了教会的平安，我恳请阁下奋力相助。万有的天主本可以点头驱散密布的乌云；但祂等待时机，藉此一面显示受攻击者的坚忍，一面赐给我们行善的机会。\par

\SourceRecord{Translated by Blomfield Jackson.；Nicene and Post-Nicene Fathers, Second Series；Vol. 3.；Edited by Philip Schaff and Henry Wace.；Buffalo, NY: Christian Literature Publishing Co.,；1892.；<http://www.newadvent.org/fathers/2707093.htm>.}

% structure: p0001 source=h1 target=section reason=page-title

\section{第九十四封信}

\Emph{致\Person{普罗托革内斯}} \Emph{\OriginalTerm{总督}{Præfect}}\par

主的慈爱已给你机会实现你的善愿。如今他给了你更大的机会，使你阁下能更容易地维护受攻击的真理，消灭谎言，并使教会获得它们所热切渴望的安宁。阁下已从许多其他来源得知东方教会遭受的巨浪有多么汹涌，但你将从那些非常虔诚的主教们那里获得更准确的消息——他们因此而在冬季长途跋涉，依靠天主的恩宠之后，也依赖你权威的眷顾。那么，基督徒啊，为我们驱散风暴，将无月之夜变为晴朗的阳光，并勒住那些妄加非议我们的舌头。我们靠天主的恩宠始终为宗徒的谕令而战，并保持\Place{尼西亚}所确立的信德纯洁无损，称所有胆敢违反其教义的人为不敬虔者。作为我所说之话的真实性证据，可以引用我的慕道者、那些不时由我施洗的人，以及在教会中听我讲道的人。如果他们想依法控告我，就必须当面定我的罪，而不是在背后诽谤我。阁下在其他案件审判时正是如此，从控辩双方的陈述中判断哪一方有理，然后宣判。\par

\SourceRecord{Translated by Blomfield Jackson.；Nicene and Post-Nicene Fathers, Second Series；Vol. 3.；Edited by Philip Schaff and Henry Wace.；Buffalo, NY: Christian Literature Publishing Co.,；1892.；<http://www.newadvent.org/fathers/2707094.htm>.}

% structure: p0001 source=h1 target=section reason=page-title

\section{第九十五封信}

\Emph{致安提约古长官}\par

您已放下那极为重要之政府的忧虑，但您的声望却在众人中兴盛；因为那些收获您仁慈果实的人——他们人数众多，遍布各处——不断地颂扬它，从各方宣扬您的佳誉，并激励听者齐声欢呼。当我看到那以其美丽装点其远扬名声的枝条的宝贵果实时，我满心欢喜。因此，我呼吁阁下您从事更伟大、更高尚的事业，并恳求您关注各教会的安宁。他们被那些编造诽谤我的人所引发的一场巨大风暴所淹没，在这种情况下，极其虔诚的主教们，不顾长途跋涉、体弱多病和年事已高，离开他们无人牧放的羊群，踏上这遥远的路途，急切地要驳斥那些针对我们所有人所说的谎言。我恳请阁下您给予他们保护，关怀受诽谤的东方，并为宗徒信仰的福祉未雨绸缪。您理应为您的其余善行增添这又一荣耀。\par

\SourceRecord{Translated by Blomfield Jackson.；Nicene and Post-Nicene Fathers, Second Series；Vol. 3.；Edited by Philip Schaff and Henry Wace.；Buffalo, NY: Christian Literature Publishing Co.,；1892.；<http://www.newadvent.org/fathers/2707095.htm>.}

% structure: p0001 source=h1 target=section reason=page-title

\section{第九十六封信}

\Emph{致诺穆斯贵族。}\par

我给您写过两封信，确切地说，我认为是三封，但没有得到任何答复。我曾希望不再多说，而是认识自己的位置和尊荣的伟大，并恳求您告知我沉默的原因。实在我不知道我能对阁下犯下什么过失。我们无论无意还是有意都会犯错，有时完全不知自己是如何过犯的。因此我恳请阁下，记念那些明确命令我们的神圣律法：\BibleQuote{如果你的兄弟得罪了你，去，要在你和他独处的时候，指出他的过错}（\BibleRef{玛 18:15}），请屈尊向我说明这烦恼的根源，好让我要么证明自己无辜，要么在意识到我错在哪里后，求您宽恕。我依凭良心的见证而自信，希望是前者。所有人都因宽宏大量而得装饰，尤其是那些效法阁下榜样的人，他们既受世俗教育培养，又受神圣原则教导，既听到宗徒律法大声疾呼：\BibleQuote{不可让太阳在你们的愤怒中落下}（\BibleRef{弗 4:26}），又记起\Person{荷马}的诗句。\par

\Quote{用恰如其分的界限约束你伟大的心灵；\newline{}仁爱为最善。}\par

我这样写，并非好像要告知你什么，而是提醒一位十分忙碌的人，我这样做是记念主的诫命，他说：\BibleQuote{所以，如果你在祭台前献上你的礼物，在那里想起你的兄弟有什么反对你的事，就把你的礼物留在祭台前，先去与你的兄弟和好，然后回来献上你的礼物。}（\BibleRef{玛 5:23-24}）听从这些话，我认为应当通过最虔诚的主教们向阁下致意，并劝请你关注教会的平安。它们确实被一场巨大的风暴所淹没。\par

\SourceRecord{Translated by Blomfield Jackson.；Nicene and Post-Nicene Fathers, Second Series；Vol. 3.；Edited by Philip Schaff and Henry Wace.；Buffalo, NY: Christian Literature Publishing Co.,；1892.；<http://www.newadvent.org/fathers/2707096.htm>.}

% structure: p0001 source=h1 target=section reason=page-title

\section{第九十七封信}

\Emph{致\Person{斯波拉基乌斯}伯爵。}\par

我非常高兴收到阁下的来信。那位极其虔诚的司铎兼隐修士\Person{杨布利科斯}告诉我您的温暖热忱、在宗教上的真诚，以及您对我的真正善意，这更增添了我的喜悦。听到这些，以及光荣而虔诚的\Person{帕特里基乌斯}大人为我所作的努力，我将蒙福的\Person{欧乃西佛若}从那神圣的口舌中获得的宗徒的祝福给予您；\BibleQuote{愿主赐怜悯于你的家，因为他屡次使我快慰，不以我的锁链为耻；}\BibleQuote{愿主赐予他，使他在那一日获得主的怜悯。}我这样为您祈祷，即使真理的仇敌如他们所想，加给我更大的痛苦；因为我们受教要顾及人的意向；但请确信，凭着真宗教，死亡对我非常愉快，流放到地极也是。尽管如此，我们仍为教会的风暴而忧虑，万有之主有能力驱散它。\par

\SourceRecord{Translated by Blomfield Jackson.；Nicene and Post-Nicene Fathers, Second Series；Vol. 3.；Edited by Philip Schaff and Henry Wace.；Buffalo, NY: Christian Literature Publishing Co.,；1892.；<http://www.newadvent.org/fathers/2707097.htm>.}

% structure: p0001 source=h1 target=section reason=page-title

\section{第九十八书}

\Emph{致\Person{潘卡里乌斯}}\par

我们看见教会受风暴困扰，深感痛苦；但教会的主人及统治者永远透过惊涛骇浪向人显示祂的智慧与能力。祂斥责风浪，带来平静，如同祂曾在宗徒的船上所做的一样。因此，虽然我痛苦，但我因认识我们救主的这种能力，也知晓祂为我们安排的一切，即使逆境临到我，我也感恩，并视之为天主的恩赐。我已学会不太关心现世，而期待那应许的福乐。但阁下亟应热心护卫宗徒信仰，好从万有的天主那里接受如此行为的报偿。\par

\SourceRecord{Translated by Blomfield Jackson.；Nicene and Post-Nicene Fathers, Second Series；Vol. 3.；Edited by Philip Schaff and Henry Wace.；Buffalo, NY: Christian Literature Publishing Co.,；1892.；<http://www.newadvent.org/fathers/2707098.htm>.}

% structure: p0001 source=h1 target=section reason=page-title

\section{第九十九封信}

\Emph{致\Person{克劳迪亚努斯}\OriginalTerm{书记员}{Antigrapharius}。}\par

虽然你尚未见过我，但我认为阁下知道那些公开散布的诽谤，因为你常听到我在教会中讲道，那时我宣扬主耶稣，并指出神性与人性的各自特性；因为我们并不将一位圣子分为两位，而是朝拜独生子，同时指明肉身与神性的区别。我认为，连亚略派也承认这一点，他们既不称肉身是神性，也不将神性呼作肉身。圣经清楚地教导我们两种本性。然而，尽管我向来如此宣讲，仍有人用谎言攻击我。但我依靠我的良心，并请鉴察人心者作我教导的见证。所以，如同先知所言，我把诽谤的阴谋视为\Quote{蜘蛛网}。我等待那无需言语的伟大审判，它要显明那隐匿未显的事。\par

我通过这几位极虔诚的主教寄出此信，认为值得借此问候阁下，并提醒你你的诺言。因为尽管我遭受攻击，我仍不停捕鱼，因为我知道神圣的宗徒们在遭受攻击时，也不停撒下圣神的网。\par

\SourceRecord{Translated by Blomfield Jackson.；Nicene and Post-Nicene Fathers, Second Series；Vol. 3.；Edited by Philip Schaff and Henry Wace.；Buffalo, NY: Christian Literature Publishing Co.,；1892.；<http://www.newadvent.org/fathers/2707099.htm>.}

% structure: p0001 source=h1 target=section reason=page-title

\section{第一百封书信}

\Emph{致\Person{亚力山德拉}。}\par

我最近收到了阁下的来信。为了你为我所表现的热忱，我感谢你，并祈求万有的天主守护你所拥有的恩物，以更多的恩赐增加它们，并赐你享受未来永久的祝福。我认为祂甚至垂听被判流放者的祈祷，尤其是当他们为了祂神圣的教义而忍受明显的耻辱时。我正通过非常虔诚的主教们写信，并恳求他们能得到你仁慈的关照。他们是为了福音的信德和教会的和平而承担了这次长途旅行。\par

\SourceRecord{Translated by Blomfield Jackson.；Nicene and Post-Nicene Fathers, Second Series；Vol. 3.；Edited by Philip Schaff and Henry Wace.；Buffalo, NY: Christian Literature Publishing Co.,；1892.；<http://www.newadvent.org/fathers/2707100.htm>.}

% structure: p0001 source=h1 target=section reason=page-title

\section{第一〇一封书信}

\Emph{致女执事\Person{塞拉里娜}。}\par

反对我们的战争烈焰再次燃起。人类的仇敌在退却片刻之后，再次武装起那些在谎言中长大的人来攻击我，他们公开诽谤我，说我将我们唯一的主耶稣基督分割为两个儿子。然而，我明辨天主性与人性之间的区别，并宣认独一圣子，即降生成人的天主圣言。我断言祂是永恒的天主，在末世成了人，这并非因天主性的改变，而是通过人性的摄取。然而，我无需向您的虔敬禀告我的见解，因为您确切知道我所宣讲的，以及我如何教导无知的人。因此，我恳求您，既然谎言的作工人将他们侮辱倾倒在东方所有虔诚的主教身上，并以风暴淹没教会，愿您的虔敬为福音的教义和教会的和平展现一切可能的热情。为此，那些极其虔诚的主教们已经离开了他们牧养的教会，不顾冬天的严酷，忍受了他们漫长旅途的劳苦，以便他们能够平息已兴起的风暴。我相信您的神圣尊贵将视他们为虔诚的斗士和教会的治理者。\par

\SourceRecord{Translated by Blomfield Jackson.；Nicene and Post-Nicene Fathers, Second Series；Vol. 3.；Edited by Philip Schaff and Henry Wace.；Buffalo, NY: Christian Literature Publishing Co.,；1892.；<http://www.newadvent.org/fathers/2707101.htm>.}

% structure: p0001 source=h1 target=section reason=page-title

\section{第一〇二封书信}

\Emph{致\Person{巴西略}主教。}\par

那些不认识我的人默默听着对我的指责，这并不奇怪；但您的圣德竟不驳斥诽谤者的谎言，或至少只是在一定程度上、且不很热心地这样做，这令任何了解您品格和行径的人难以置信。我说这话，不是因为友谊应当优先于真理，而是因为真理的见证站在友谊一边。您阁下曾多次听到我在教堂讲道，也在其他集会上听我谈论教义问题；您听过我所说的话，我不知道有哪一次您曾因我表达不合正统的意见而指责我。但如今的情况如何？究竟为什么，我亲爱的朋友，您对谎言一言不发，却容忍朋友被诽谤、真理受攻击？如果这是因为您轻视无依无靠和微不足道的人，请记住主明明白白的诫命：\BibleQuote{你们小心，不要轻视这些信我的小子中的一个，因为我告诉你们：他们在天上的天使，常常见到我天父的面。}如果却是我的诽谤者的势力使您缄默不语，那么您必须听从另一条法律，它说：\BibleQuote{不可尊重有权势的人}和\BibleQuote{施行公义的审判}和\BibleQuote{不可随众行恶}和\BibleQuote{那闭眼不看邪恶、塞耳不听流血之事的}。您可以在圣经中找到无数类似的段落，我觉得写信给一位在天主圣言中长大、并以教导灌溉基督徒的人，没有必要收集这些。但我要说这一点：我们都要站在基督的审判台前，并为我们的言行交账。我，因各种理由畏惧这法庭，如今被诽谤包围，却在这思想中找到我的主要安慰。\par

\SourceRecord{Translated by Blomfield Jackson.；Nicene and Post-Nicene Fathers, Second Series；Vol. 3.；Edited by Philip Schaff and Henry Wace.；Buffalo, NY: Christian Literature Publishing Co.,；1892.；<http://www.newadvent.org/fathers/2707102.htm>.}

% structure: p0001 source=h1 target=section reason=page-title

\section{\WorkTitle{第一〇三封信}}

\Emph{致阿颇罗纽伯爵。}\par

那些非常虔诚的主教们因针对我的毁谤而被引往帝都旅行，而我通过他们的圣德向阁下您致以问候，并偿还友谊之债，这并非为了消除珍视的义务，而是使之增长。因为事实上，友谊的义务因履行而增加。我现在正承受毁谤的果实并不出奇，因为作为人，我料想任何事情都可能发生。一切此类烦恼都应由学会智慧的人承受；只有一件事令人痛苦——即灵魂应遭受损害。\par

\SourceRecord{Translated by Blomfield Jackson.；Nicene and Post-Nicene Fathers, Second Series；Vol. 3.；Edited by Philip Schaff and Henry Wace.；Buffalo, NY: Christian Literature Publishing Co.,；1892.；<http://www.newadvent.org/fathers/2707103.htm>.}

% structure: p0001 source=h1 target=section reason=page-title

\section{第一〇四号书信}

\Person{福拉维亚诺}致\Place{君士坦丁堡}主教。\par

\Person{福拉维亚诺}致\Place{君士坦丁堡}主教。 我已另致一书，禀告您的圣德，我们教义的诽谤者如何公开诋毁我们。如今同样藉着极为虔诚的主教们，我也这样行，不仅以此为我教义正统的见证，还有无数东境各教会中听闻我讲道的人。超越这一切，我有我的良心，以及那洞察我良心的那一位。我也知道，神圣的宗徒常常诉诸良心的见证，因为\Quote{我们的夸耀就是良心的见证}又\Quote{我在基督内说实话，并不说谎，我的良心也藉着圣神为我作证}。因此，圣善的阁下，请知：从未有人听过我们宣讲两位子；事实上，这一教义在我看来可憎且不虔，因为只有一个主耶稣基督，万物藉祂而有。我承认祂既是永生的天主，又在末世成为人，并只以独生子的敬礼敬拜祂。然而，我学习了肉体与神性的区别，因为合而不混。如此列阵对抗\Person{亚略}和\Person{欧诺米}的疯狂，我们很容易驳斥他们针对独生子所妄加的亵渎，即：把关于主谦卑所言、且适合祂所取本性的，归给人；而另一方面，把属于神性且表明神性本质的，归于天主；不将祂分成两位，而是教导那前后的属性皆属于独生子，后者属于祂——作为创造万有的天主和主，前者属于祂——为了我们而成为人。因为神圣的圣经说祂成为人，不是藉着神性的变化，而是藉着取人性，出自\Person{亚巴郎}的后裔。神圣的宗徒明确地在经上说：\Quote{因为祂实在没有承取天使的本性，而是承取了\Person{亚巴郎}的后裔，因此祂应当在各方面相似于弟兄们。}又说：\Quote{诺言原是向\Person{亚巴郎}和他的后裔所许的；没有说\InnerQuote{后裔们}，如同指着许多人，而是指着一个人，即\InnerQuote{你的后裔}，就是基督。}\par

这些和类似的章节已被\Person{西门}、\Person{巴西理德}、\Person{瓦伦提努斯}、\Person{巴尔德撒内斯}、\Person{马尔西翁}，以及那位以疯狂异端命名的人从圣经中删除。因而他们称主基督仅为天主，描述祂毫无人性，只是凭想象和幻象向人显现为人。另一方面，\Person{亚略}派和\Person{欧诺米}派则说天主圣言只取了身体，祂自己在身体中代替了灵魂的位置。\Person{阿波利拿里}描述主的身体有灵魂；但他从某处——我不知何处——得出灵魂与理智有分别，从而剥夺了理智在已完成的救恩中的分。相反，神圣宗徒的教导指出：一个既有理性又有智慧的灵魂连同肉体一同被取了，而期望得救的人所怀的救恩是完整的。\par

还有另一群异端者持不同意见。\Person{佛提诺}、\Person{玛尔塞鲁}和\Place{撒摩撒特}的\Person{保禄}断言我们的主和天主只是人。当与这些人辩论时，我们不得不提出天主性的证据，并表明主基督是永生的天主。另一方面，当我们与前一派（仅称我们的主耶稣基督为天主）争论时，我们不得不集结圣经的军队对抗它们，收集人性被取的证据。因为医生必须使用对症的疗法，使药物适合病情。\par

因此，我恳求您的圣德驱散针对我的诽谤，勒住如今徒然辱骂我的舌头。因为，在降生成人之后，我敬拜一位天主之子、一位主耶稣基督，并谴责所有持异议者为不虔。阁下，请也将您的神圣祈祷赐予我，使我因天主的恩宠，能渡越危险之海，将锚抛在主无风之港中。\par

\SourceRecord{Translated by Blomfield Jackson.；Nicene and Post-Nicene Fathers, Second Series；Vol. 3.；Edited by Philip Schaff and Henry Wace.；Buffalo, NY: Christian Literature Publishing Co.,；1892.；<http://www.newadvent.org/fathers/2707104.htm>.}

% structure: p0001 source=h1 target=section reason=page-title

\section{第一〇五号书信}

\Emph{致\Person{欧洛吉乌斯} \OriginalTerm{总管}{\GreekText{οἰκονόμος}}。}\par

我们听说了许多关于您虔诚为真正信仰所付出的努力。因此，您理所当然应当乐于援助那为同样原因而受诬蔑的人，并驳斥谤者的谎言。您，虔诚的阁下，知道我所持守和教导的是什么，也知道没有人曾听闻我宣讲\OriginalTerm{两个儿子}{two sons}。我恳求您，在此事上也发挥您神圣的热忱，堵住恶言者的口。在这种冲突中，不仅要帮助朋友，就连那些使我们痛苦的人也应施以援手。\par

\SourceRecord{Translated by Blomfield Jackson.；Nicene and Post-Nicene Fathers, Second Series；Vol. 3.；Edited by Philip Schaff and Henry Wace.；Buffalo, NY: Christian Literature Publishing Co.,；1892.；<http://www.newadvent.org/fathers/2707105.htm>.}

% structure: p0001 source=h1 target=section reason=page-title

\section{第一〇六函}

\Emph{致\OriginalTerm{管家}{Œconomus}\Person{亚巴郎}}\par

敬虔的主教们向您问安。我恳求您留心教会的平静，驱散诽谤的波涛。正如神圣宗徒所说：\BibleQuote{人种什么，也收什么}。无疑，那为宗徒道理而战的人，必将收获宗徒祝福的果实，并享有宗徒们的事奉。\par

\SourceRecord{Translated by Blomfield Jackson.；Nicene and Post-Nicene Fathers, Second Series；Vol. 3.；Edited by Philip Schaff and Henry Wace.；Buffalo, NY: Christian Literature Publishing Co.,；1892.；<http://www.newadvent.org/fathers/2707106.htm>.}

% structure: p0001 source=h1 target=section reason=page-title

\section{书信第一〇七}

\Emph{致司铎德奥多托。}\par

\Person{你}的虔诚在宗徒教义上所受的挣扎并非无人知晓，反而被那些亲历者以及从他们听闻者屡屡提及。我亲爱的先生，请继续努力，为教父们的教义而战。因为我也为这些教义而四面受击，承受巨浪的冲击，我恳求我们的治理者或点头驱散风暴，或藉祂的恩宠使风暴中的遇难者能显出勇毅。\par

\SourceRecord{Translated by Blomfield Jackson.；Nicene and Post-Nicene Fathers, Second Series；Vol. 3.；Edited by Philip Schaff and Henry Wace.；Buffalo, NY: Christian Literature Publishing Co.,；1892.；<http://www.newadvent.org/fathers/2707107.htm>.}

% structure: p0001 source=h1 target=section reason=page-title

\section{书信108}

\Emph{致\Person{阿卡西乌}司铎。}\par

达味圣咏的许诺确实是真实的，因为透过他，真理之神将此许诺赐予相信的人：\BibleQuote{将你的道路托付给主，也要信赖他；他必成就；他必使你的义德如光发出，你的公平如日中天。}我们发现这一点也在你的虔敬身上应验了。因为你赐予那些为孤苦而哭泣的人极大的关怀，以及你为宗徒道理所做的奋斗，众口皆碑，因此，正如先知们所说：\BibleQuote{隐藏的事被显明出来。}因为我既已听闻你虔敬的可贵努力，特以此信向你致意，最虔诚的先生，并恳请你通过增加劳苦来增添你的光荣，为福音道理而战斗，使我们能保持我们不朽的祖传遗产，并能将我们主人的塔冷通连本带利归还。\par

\SourceRecord{Translated by Blomfield Jackson.；Nicene and Post-Nicene Fathers, Second Series；Vol. 3.；Edited by Philip Schaff and Henry Wace.；Buffalo, NY: Christian Literature Publishing Co.,；1892.；<http://www.newadvent.org/fathers/2707108.htm>.}

% structure: p0001 source=h1 target=section reason=page-title

\section{第一〇九封信}

\Emph{致\Place{安奇拉}的主教\Person{欧瑟比}}\par

许多计谋暗中针对我，并透过我拼凑起来反对宗徒的信德。然而，圣人们、先知、宗徒、殉道者以及在恩宠的言语中教会里出名的人所受的苦难安慰了我；此外，我们的天主和救主的应许也安慰了我，因为在此生祂并没有应许我们任何愉快或喜悦的事，反倒是麻烦、劳苦、危险和仇敌的攻击。祂说：\BibleQuote{在世界上，你们将有苦难}，又\BibleQuote{他们若迫害了我，也要迫害你们}，以及\BibleQuote{他们若称家主为贝耳则步，何况他的家人呢？}，以及\BibleQuote{时候将到，凡杀你们的，还以为是事奉天主}，以及\BibleQuote{门是窄的，路是狭的，通往生命的}，以及\BibleQuote{他们在这城迫害你们，你们就逃到别的城去}，我还可以引述所有类似的经文。神圣的宗徒也说了同样的话：\BibleQuote{凡在基督耶稣内愿意敬虔生活的，都必受迫害；但恶人和骗子却越来越坏，迷惑人，也被人迷惑。}这些话在我这困苦中给了我最大的安慰。因着那些对我的诽谤可能已经传到您的圣座耳中，我恳求您的圣座不要相信我的毁谤者所说谎言。直到如今，我并不知道自己曾教导任何人相信两个儿子。我所受的教导是相信一位\OriginalTerm{独生子}{only begotten}，我们的主耶稣基督，\OriginalTerm{天主圣言}{God the Word}成了人。但我认识\OriginalTerm{肉体和神性}{flesh and Godhead}的区别，并且认为所有将我们的主耶稣基督分为两个儿子的人，以及那些走向相反方向，称基督主的神性和人性为\OriginalTerm{一性}{one nature}的人都是不敬虔的。因为这两种极端彼此对立，而在它们之间是福音教义的道路，由先知和宗徒，以及所有在他们之后因教导的恩赐而卓越的人的足迹所美化。我本想引用他们的意见，并指出他们如何为我作证，但一封信的篇幅不够；因此我简要地写下了我所受的教导，关于独生子的\OriginalTerm{降生成人}{incarnation}；我将我的陈述呈送给您的敬虔阁下。我写这信不是为教导别人，而是为辩护自己免受对我的指控，并向那些不知道我的看法的人解释我的观点。在您的圣座读完我所写的内容后，若发现它符合宗徒的教义，我希望您能藉回信确认我的意见；相反，若我所说的任何事与天主的教导有冲突，我请求您的圣座告知我。因为我虽然花了很长时间教导，但仍需要有人教导我。神圣的宗徒说：\BibleQuote{我们知道一部分}，他又说：\BibleQuote{若有人以为自己知道什么，他还没有按他所应当知道的知道什么。}因此，我希望我能从您的圣座听到真理，并希望您也留意教会的平安，为天主的教义而战。正是为了这些教义，极其敬虔的主教们不辞旅途的艰辛和冬季的严寒，出发前往皇城，试图为风暴带来一些终结。我请求您用祈祷送他们上路，也用您的祈祷加强我。\par

\SourceRecord{Translated by Blomfield Jackson.；Nicene and Post-Nicene Fathers, Second Series；Vol. 3.；Edited by Philip Schaff and Henry Wace.；Buffalo, NY: Christian Literature Publishing Co.,；1892.；<http://www.newadvent.org/fathers/2707109.htm>.}

% structure: p0001 source=h1 target=section reason=page-title

\section{书信第一一〇号}

\Emph{致安提约基亚主教多姆努斯。}\par

当我读到你的信时，我想起了那位极其蒙福的苏撒纳；她看见那著名的恶徒，并相信万有的天主临在，便发出那著名的呼喊：\Quote{我四面受困；}但她宁愿落入诽谤的罗网，也不愿蔑视公义的天主。我，先生，正如我常说，有两个选择：要么得罪天主、伤害我的良心，要么因人的不公判决而跌倒。我认为，最虔诚的皇帝对此一无所知。因为如果这真合他心意，有什么阻碍他写信并下令举行祝圣呢？到底为什么他们在外面发出威胁、引起恐慌，却不公开送信命令这事？两者必居其一：要么这位非常虔诚的皇帝没有被说服写信，要么他们试图使我们违法，然后被他们以非法之名起诉。我面前有蒙福的主教普林奇皮乌斯的榜样；在那件事上，他们通过信件下令，却因他服从而惩罚他。此外，在送信人到达的那天，我所读的信件内容相反。因为一位神圣的隐修士写信给某人说，他已收到极杰出的侍卫官和极荣耀的前任长官的信件，表示非常虔诚的主教依勒乃的事将更为顺利；为回报这番善意，他们请求为他们祈祷。因此，我认为应当给那些从帝都写信的圣职人员回信，表明：我服从了非常虔诚的腓尼基主教们的判决，且深知非常虔诚的主教依勒乃的热忱、宽宏、爱贫及其他德行，此外还有他观点的正统性，因而祝圣了他。我不知道他曾反对称童贞玛利亚为\OriginalTerm{天主之母}{Theotokos}，或持有任何与福音教义相悖的观点。关于再婚者的问题，我遵循了前任的做法；因为神圣而可纪念的亚历山大，这宗座之荣，以及非常蒙福的贝洛雅主教阿卡奇乌斯，祝圣了蒙福的\OriginalTerm{再婚者}{digamus}狄奥革涅斯；同样，蒙福的普拉伊里乌斯祝圣了凯撒勒雅的\OriginalTerm{再婚者}{digamus}多米努斯。因此，我们遵循了先例，并效法了这些在学识和品格上著名而杰出之人的榜样。君士坦丁堡主教、蒙福的普罗克鲁斯深知此事及许多其他事例，他不仅接受了祝圣，还写信称赞和钦佩它。本都教省的主要虔诚主教们和所有巴勒斯坦人也如此。\par

\Quote{关于此事没有提出疑问，我们认为谴责一个因多种高贵行为而显赫的人是不对的。}依我之见，这样写是合适的。如果您的圣德持有其他看法，就按您认为好的去做。我，如他们所料，已经受了一次惩罚，并准备依靠天主的帮助再受一次。即使是第三次、第四次，若他们愿意，靠着天主恩宠的扶持，我也将忍受，并赞美主。如果您的圣德认为合适，让我们看看从巴勒斯坦传来什么答复，并在更准确地考虑应采取的行动后，再这样写信给君士坦丁堡。\par

\SourceRecord{Translated by Blomfield Jackson.；Nicene and Post-Nicene Fathers, Second Series；Vol. 3.；Edited by Philip Schaff and Henry Wace.；Buffalo, NY: Christian Literature Publishing Co.,；1892.；<http://www.newadvent.org/fathers/2707110.htm>.}

% structure: p0001 source=h1 target=section reason=page-title

\section{\WorkTitle{致阿纳托利书第一一一}}

\Emph{致\Person{阿纳托利}总执事。}\par

尊驾所赐予我的恩惠，万有的天主必将酬报，因为凡为祂而行的事，必有其赏报。我嘲笑所有诽谤我的人。那些受最严厉鞭笞之人的身体并不感到疼痛，因为被鞭打的皮肉已经麻木。然而，我仍为那些人哀叹，因他们肆无忌惮的嘴巴竟说出如此谎言。指控虔诚主教\Person{依巴斯}的人，究竟在哪方面被我亏待，以致他们对我发出这样的诬蔑？首先，我甚至不是审判官之一，因为我遵奉帝国敕令，居住在\Place{居鲁士}。此外，正如我从许多人那里听到的，他们始终把我的缺席视为一种不满，因为我曾安排他们参加逾越节救恩盛宴的圣体圣事；他们常表示愿意见我，我便以善意接待他们，并就应采取的正当行动给予劝告。但为极虔诚的主教、\Person{多姆努斯}主教的辩护起见，我也要说：他应当采取什么正当行动呢？他公开受到攻击；他看到有人被主教会议判决撤职后，竟被送到另一教区，并违反教会法律恢复司铎职务；他看到神圣和天主的事物被教会仇敌嘲笑和戏弄；他该做什么？当他知晓此事时，便将案件交给他人处理，不仅交给极虔诚的主\Person{依巴斯}，也交给\Place{阿米达}城圣主教\Person{西默盎}，以便两省的教省主教能听取指控。指控同一个人既残忍又仁慈，这有什么公平可言？若我们施行绝罚，则陷入危险；若我们不施行绝罚，亦不能逃脱。我们在全世界中唯独成为攻击对象。其他教区平安无事。唯独我们暴露在诽谤者面前——尤其是我本人，尽管我并未参与审判，且对此事完全不负责任。\par

因此，我读罢尊驾的信函，得知由此缘故，人们对我——一个固守本教区、爱好和平、甚至未与本省虔诚主教们商议之人——掀起巨大骚动，遂被迫写此信。事实上，尽管我们省内已有两次主教祝圣，我却一次也未参与。若非受帝国敕令束缚，我早已离去，在某个偏远之处度我余生。我为针对我的阴谋而疲惫不堪。我确信那些\Place{厄德撒}人并非自愿编造对我的诽谤。他们是在真正诚实的邻居唆使下对我发动攻击的。我感谢我们的救主，祂虽不配，却仍视我为配享福音的真福。为此，我欣然接受流放的判决。我准备好被放逐，并且为了那\Quote{为我存留的希望}，甘心承受他们可能施加的任何命运。我不停地为尊驾祈祷，并恳请所有圣人同我一起祈祷。\par

\SourceRecord{Translated by Blomfield Jackson.；Nicene and Post-Nicene Fathers, Second Series；Vol. 3.；Edited by Philip Schaff and Henry Wace.；Buffalo, NY: Christian Literature Publishing Co.,；1892.；<http://www.newadvent.org/fathers/2707111.htm>.}

% structure: p0001 source=h1 target=section reason=page-title

\section{第一一二封信}

\Emph{致\Place{安提约基雅}主教\Person{多姆努斯}}\par

当得知得胜的皇帝的小气已被终止，他本人与非常虔诚的主教之间达成和解，会议召集被取消，以及教会的和平得以恢复时，我曾希望我们的麻烦已成过去。但我从圣座那里听到的令我深感痛苦。除非仁慈的主以其惯常的眷顾摧毁狂暴魔鬼的计谋，否则不可能期望从这个臭名昭著的会议得到任何好处。即使在伟大的会议中，我指的是\Place{尼西亚}大公会议，亚略派也与正统派一同投票，并在宗徒信理阐述上签字。但他们并未停止与真理作战，直到他们撕裂了教会的身体。三十年间，宗徒教义的支持者与感染亚略亵渎的人继续共融。但在\Place{安提约基雅}，当最近的会议结束后，他们曾把天主的人、伟大的\Person{梅利提乌}安置在宗徒之座上，但几天后却通过皇帝权威驱逐了他，而无疑感染了亚略瘟疫的\Person{欧佐伊}被推举出来，宗徒教义的捍卫者随即分离出去，此后分裂持续下去。\par

当我回顾那时发生的事，并展望未来类似的事件时，我可怜的精神叹息哀号，因为我看不到任何好的前景。其他教区的人不知道\OriginalTerm{十二章}{Twelve Chapters}中包含的毒素；由于顾及作者的声誉，他们不怀疑其中有危害；而他在教座上的继任者，我认为正采用一切手段在第二次会议上加以确认。因为假设他最近受命写了这些章节，并诅咒所有不愿遵守它们的人，那么如果他主持一个大公会议，他还有什么办不到的呢？请确信，我的主，任何知道它们所包含的异端的人，都不会容忍接受它们，即使更多这类人颁布它们。在此之前，尽管有更多数目的人轻率地确认了它们，我在\Place{厄弗所}抵制了，并拒绝与作者共融，直到他同意我提出的要点，并使他的教导与它们一致，而不提及这些章节。圣座可以毫无困难地查证这一点，如果命令调查会议记录；因为它们按照会议签署的惯例被保存，现存有五十多份会议文件，显示对十二章的指控。因为在前往\Place{厄弗所}之前，\Person{真福若望}曾给非常虔诚的主教们——\Place{提亚纳}的\Person{尤特里乌}、\Place{凯撒勒雅}的\Person{费尔穆}和\Place{安卡拉}的\Person{德奥多图}写信，谴责这些章节是阿波林派的。而在\Place{厄弗所}，这些章节的阐述和确认导致我们废黜了\Place{亚历山大}主教和\Place{厄弗所}主教。此外，在\Place{厄弗所}，写了许多会议书信给得胜的皇帝、高级官员、\Place{君士坦丁堡}的平信徒以及尊敬的圣职人员。而且当我们被召到\Place{君士坦丁堡}时，我们在皇帝面前进行了五次讨论，之后给皇帝送去了三份抗议书。并且我们给西方非常虔诚的主教们——即\Place{米兰}、\Place{阿奎莱亚}和\Place{拉文纳}的主教——写信谈论同一话题，抗议这些章节充满了阿波林派的创新。此外，它们的作者通过\Person{真福保禄}之手收到了\Person{真福若望}一封信，公开责备它们；同样也收到了有福记忆的\Person{阿卡修斯}的信。为了使圣座得到简明信息，我送去了\Person{真福阿卡修斯}的信，以及\Person{真福若望}写给\Person{真福济利禄}的信，这样您可以明白，尽管他们写信给他谈论和好的事，他们却责备这些章节。而\Person{真福济利禄}本人在他给\Person{真福阿卡修斯}的信中明确指出了这些章节的倾向，用以下话语：\Quote{我写这些反对他的创新，当和平实现时，它们将显明。}辩护本身证明了指控。我送去了他在和好时所写的副本，以便您可以看到，我的主，他没有提到它们，并且那些出席大公会议的人有义务提出在和好时所写的东西，清楚说明什么导致了分歧，以及在什么条件下分裂的部分得到了补偿。因为那些被召为真理而战的人必须不畏惧任何劳苦，并必须祈求天主的助佑，以便我们能够完好地保存我们先祖遗留给我们的遗产。\par

圣座必须在虔诚的主教中寻找志同道合的人，使他们成为旅途的同伴；同样地，在尊敬的圣职人员中寻找那些热心于真理的人，以免甚至被我们同侧的人出卖，我们被迫做任何令万有的天主不悦的事，或者在我们被抛弃时，轻易落入敌人的爪牙。\par

信德是我们救恩的希望所在，我们必须不遗余力地防止任何虚假的东西被引入其中，并防止宗徒教导被败坏。\par

我从远方写下这些话，带着叹息和呻吟，我恳求我们共同的主驱散这片乌云，再次赐予我们明亮阳光的恩赐。\par

\SourceRecord{Translated by Blomfield Jackson.；Nicene and Post-Nicene Fathers, Second Series；Vol. 3.；Edited by Philip Schaff and Henry Wace.；Buffalo, NY: Christian Literature Publishing Co.,；1892.；<http://www.newadvent.org/fathers/2707112.htm>.}

% structure: p0001 source=h1 target=section reason=page-title

\section{第一一三封信}

\Emph{致罗马主教良}\par

\Person{保禄}，真理的宣讲者、\Person{圣神}的号角，曾急忙赶往伟大的\Person{伯多禄}那里，为要向他领受关于遵守法律生活的疑难问题的满意解答；那么，我们这些渺小卑微的人，更当急忙赶往你的\OriginalTerm{宗座}{apostolic see}，以求从你那里获得教会创伤的医治。因为种种理由，你应当居首位，因为你的宗座享有许多特权。其他城市固然以它们的规模、美丽和人口而著称；有些在这些方面有所欠缺的，却因某些神恩而熠熠生辉。但你的城市，伟大的眷顾者已赐予丰富的恩赐。她是世界上最大、最辉煌、最显赫的城市，人口众多，应有尽有。此外，她已取得了目前的统治权，并以自己的名字赐予其臣民。她更因自己的信德而特别光彩，对此，那位神圣的宗徒曾高呼：\BibleQuote{你们的信德传遍了全世界}。即使在接受救恩信息的种子之后，她的枝条立刻结出了这些令人钦佩的果实，那么，有什么言语能恰当地赞美如今在她内所实践的虔诚呢？她保存着那些光照信友灵魂的坟墓，就是我们共同的父辈和真理导师\Person{伯多禄}和\Person{保禄}的坟墓。这对三倍有福而且神圣的一对，从日出之地兴起，向四方传播光芒。现在，从日落之地，他们甘愿接受了此生日落的区域，照亮了世界。他们使你的宗座最为光荣；这是你诸般好事的冠冕和完成；但在这些日子里，他们的天主以你的圣德装饰了他们的宝座，因为发射出正统信仰的光辉。对此，我可以提出许多证据，但只需提及你圣德最近对声名狼藉的\OriginalTerm{摩尼教徒}{Manichees}所表现的热忱，就足以证明你虔诚对神事的热心。你近来的著作，也足以表明你的宗徒性质。因为我们见到了你圣德关于我们的天主和救主降生成人的论述，并惊叹你措辞的精确。\par

因为两份著作都一致阐明了\OriginalTerm{独生子}{Only-begotten}出自永远圣父的永远神性，以及出自\Person{亚巴郎}和\Person{达味}后裔的人性；并且所取的本性在一切方面与我们相似，仅在这一点上与我们不同，即它保持无罪；因为罪不是出于本性，而是出于自由意志。\par

这些书信还包含这一点，即天主的\OriginalTerm{独生子}{Only-begotten}是独一的，他的神性是\OriginalTerm{不受苦难的}{impassible}、\OriginalTerm{不变的}{immutable}、\OriginalTerm{不可变的}{invariable}，如同生他的圣父和圣神一样；因此他取了\OriginalTerm{可受苦的本性}{passible nature}，因为神性不能受苦，为要藉他自己肉身的苦难，将不受苦难的恩赐赐予那些相信他的人。这些陈述及其他类似性质的内容都包含在你的书信中。我们钦慕你的属神智慧，赞美通过你发出的圣神的恩宠，并呼求、恳求、乞求、哀求你的尊贵，保护那些现今被风暴袭击的天主的教会。\par

我们曾期望，通过你圣德派往\Place{厄弗所}的代表，风暴会得以平息，但我们却遭遇更猛烈的风暴攻击。因为那位很正义的\Place{亚历山大里亚}主教，不满足于非法且极不正义地罢免了最神圣而虔诚的\Place{君士坦丁堡}主教，主\Person{弗拉维亚诺}，他的灵魂也不以类似地屠杀其他主教为满足，反而在我缺席时，用笔刺伤了我；他没有传唤我到庭，没有在我面前审讯我，没有询问我对我们的天主和救主降生成人的看法。即使是杀人犯、掘墓者、奸淫者，除非他们自己通过承认确认了对他们的指控，或者被其他人的证词明确证明有罪，否则法官也不会定他们的罪。然而我，受神圣法律养育的人，却被他随心所欲地定了罪，而我当时还在三十五天的行程之外。\par

他所做的还不止这些。就在去年，当两个沾染\Person{阿波利纳里}错误的人到了那里，捏造对我的诽谤时，他竟在教堂里站起来诅咒了我，而且此前我已经写信给他，向他解释了我的看法。\par

我为教会的混乱而哀恸，渴望和平。二十六年以来，我靠着你的祈祷，治理了万有者天主托付给我的教会。在可敬的\Person{狄奥多图斯}，东方首席主教的时代，从未；在他\Place{安提约基雅}宗座的继任者时代，也从未招致丝毫指责。靠着天主恩宠的帮助，与我合作，我从\OriginalTerm{马尔基昂}{Marcion}的瘟疫中拯救了一千多个灵魂；许多其他人从\OriginalTerm{亚略派}{Arian}和\OriginalTerm{欧诺弥派}{Eunomian}的派别中领回了我们的主基督。我在八百座教堂里履行了牧职，因为\Place{基禄斯}管辖着这么多堂区；并且在其中，靠着你的祈祷，连一棵莠子也没有留下，我的羊群脱离了所有异端和错误。那位洞悉万物的主知道，邪恶的异端者曾向我投掷多少石块，我在东方大多数城市中与外邦人、犹太人、每一种异端进行了多少战斗。在这一切考验和危险之后，我竟未经审判就被定罪了。\par

但我等候着贵宗座的判决。我恳求并哀请圣座援助我向您公正公义的法庭提出的上诉。请命我速速到您那里，向您证明我的教导追随宗徒的足迹。我手中有我二十年前、十八年前、十五年前、十二年前所写的著作：反对亚略异端者与欧诺米异端者，反对犹太人与外邦人；反对波斯术士；论天主眷顾；论神学；论降生成人。因天主的恩宠，我诠释了宗徒们的著作和先知们的预言。从这些著作中不难看出我是恪守了信德的正确准则，还是偏离了它的直路。恳求您，不要轻视我的祈祷；恳求您，顾念我在所有辛劳之后加在我这可怜白发上的侮辱。\par

尤其我恳求您告诉我，我是否应当忍受这不义的撤职，或者不；因为我等候您的决定。若您命令我服从定罪判决，我就服从；从今往后我将不再打扰任何人，只等候我们天主和救主的公义审判。天主是我的见证，我的主，我不在乎尊荣与光荣。我只在乎引起的恶表，因为许多较单纯的信友，尤其是那些我曾从种种异端中拯救出来的人，因着依附我审判者的权威，完全无法理解教义的准确真理，也许将以为我犯有异端。\par

整个东方都知道，在我任主教的所有期间，我没有购置一所房屋、一块地、一文钱、一座坟墓，而是自愿拥抱神贫，在父母去世后，我将从他们继承的全部财产都分施了。\par

尤其我恳求您，圣洁的、蒙主爱的人，赐我以您祈祷的帮助。我通过可敬且虔诚的\OriginalTerm{乡村主教}{chorepiscopi}\Person{伊帕提乌}、\Person{阿布拉米乌}以及我们隐修士的\OriginalTerm{督教}{exarch}\Person{阿利皮乌}，将这些事告诉您。我本应亲自赶赴您那里，若不是被帝国谕令的锁链所困，这锁链像困住别人一样困住了我。恳求您，对待我的使者如同父亲对待儿子；请仁慈而无偏颇地听他们；请赐予您对我被诽谤和徒然攻击的暮年的保护。尤其请尽您所能，顾念那被合谋攻击的信德；为各教会保存其祖产不受损害。如此，圣座将从一切美好恩赐的伟大施予者那里领受对这些行为的相应报酬。\par

\EditorialNote{在米涅的版本中，此处接续了\Person{圣良}致\Person{德奥多勒}的答复，即\Person{圣良}著作中的第120号书信。}\par

\SourceRecord{Translated by Blomfield Jackson.；Nicene and Post-Nicene Fathers, Second Series；Vol. 3.；Edited by Philip Schaff and Henry Wace.；Buffalo, NY: Christian Literature Publishing Co.,；1892.；<http://www.newadvent.org/fathers/2707113.htm>.}

% structure: p0001 source=h1 target=section reason=page-title

\section{第一一四封书信}

致\Person{安迪贝里斯}。\par

尊敬的司铎\Person{伯多禄}不仅在司铎品位上卓尔不凡，而且以其精湛的医术著称。他在我们这里长久居住期间，以其和解的态度赢得了所有人的心。得知我将离开，他现在决定离开\Place{居鲁士}；因此我将他推荐给您阁下，希望他——他在\Place{亚历山大}居住时曾行医——完全有能力为城市提供良好服务，并能在您手中得到善待。\par

\SourceRecord{Translated by Blomfield Jackson.；Nicene and Post-Nicene Fathers, Second Series；Vol. 3.；Edited by Philip Schaff and Henry Wace.；Buffalo, NY: Christian Literature Publishing Co.,；1892.；<http://www.newadvent.org/fathers/2707114.htm>.}

% structure: p0001 source=h1 target=section reason=page-title

\section{第一一五封信}

《致\Person{阿培拉}》\par

当我承担\Place{居鲁士}教区的治理时，我从各地为它招募了从事必要技艺的人，此外还劝使技艺精湛的医生居住在那里。其中一位是令人尊敬的司铎\Person{伯多禄}，他以智慧从事其职业，并以自己的品行装饰它。在我离开时，有几位离开了城市，\Person{伯多禄}也决定离开。在这种情况下，我恳求阁下给予他您的慈爱关怀。他完全能够照料病人，并与他们的疾病作战。\par

\SourceRecord{Translated by Blomfield Jackson.；Nicene and Post-Nicene Fathers, Second Series；Vol. 3.；Edited by Philip Schaff and Henry Wace.；Buffalo, NY: Christian Literature Publishing Co.,；1892.；<http://www.newadvent.org/fathers/2707115.htm>.}

% structure: p0001 source=h1 target=section reason=page-title

\section{书信一一六}

\Emph{司铎雷纳杜斯。}\par

我们已听说您圣善的热忱与义怒，以及您在谴责\Place{厄弗所}会议上那鲁莽的行为时所运用的公正而合法的勇气。此事不仅为我所知，您正统信仰的名声已传遍各地，众人都在颂扬您的正义、热忱、勇气以及您对我所受不公待遇的斥责。您圣善在目睹一次大屠杀后便采取了这一行动。若您见到您离开后所发生的其他事，您或许会效法那位著名的\Person{丕乃哈斯}的热情。我是后来被定罪的人之一，被皇帝命令禁止参加公会议，并在缺席的情况下被判刑。\par

我担任主教已有二十六年；经受了无数劳苦；为真理奋力拼搏；使成千上万的异端者脱离谬误，引领他们归向救主；而现在他们剥夺了我的司铎职；他们将我逐出城邑。他们不尊重我的年迈，不尊重我为真理而灰白的头发。因此，我恳求您的圣善，说服那至圣至敬的总主教命我速速前往你们的公会议。因为那圣座在全世界所有教会中居于首位，原因众多；尤其在于它从未沾染异端污点，从未有任何异端见解的主教坐过它的宝座，而它始终保持了宗徒的恩宠无玷。我信赖您的正义，将接受您的任何决定，并主张由我的著作来受审。我写了三十多部书反对\Person{亚略}、\Person{欧诺弥}、\Person{马西翁}、\Person{马其顿尼}，反对异教徒和犹太人；我诠释了圣经，任何人若愿意，都能轻易看出我追随宗徒的足迹，宣扬独一圣子、独一圣父、独一圣神；天主圣三的独一天主性、独一主权、独一能力、永恒、不变、不受苦难、独一旨意；主耶稣基督的天主性是完美的，所取的人性为我们的救恩也是完美的，并为我们被交于死亡。我不认识一个为人子、另一个为天主子的两位，而是同一位，天主子由天主所生、是天主，同时借仆人的形象、出于\Person{亚巴郎}和\Person{达味}的后裔而为人子。我继续教导这些以及类似的道理；我也在至圣至敬的总主教\Person{良}的著作中发现了这些道理，我赞美万有之主，因我认同他宗徒性的教导。请接受我的恳求，不要忽视我所受的冤屈。为此，我派遣了敬主的司铎\Person{希帕提乌斯}和\Person{阿布拉米乌斯}（乡村主教）以及我们隐修士的首领\Person{阿利庇乌斯}到您圣善处，他们以善行装饰自己，并能口述关于我卑微之事的准确信息。\par

\SourceRecord{Translated by Blomfield Jackson.；Nicene and Post-Nicene Fathers, Second Series；Vol. 3.；Edited by Philip Schaff and Henry Wace.；Buffalo, NY: Christian Literature Publishing Co.,；1892.；<http://www.newadvent.org/fathers/2707116.htm>.}

% structure: p0001 source=h1 target=section reason=page-title

\section{第一一七书}

\Emph{致弗洛伦提乌斯主教。}\par

诚然，我们的天主和救主的恩宠尚未抛弃人类，却在您的圣德中为我们留下了一个遗民，\Quote{免得我们变成索多玛，变得像哥摩辣。}这遗民不容我们完全沮丧，却吩咐我们等候那可怕风暴的过去；这使我们充满希望。\par

为此，我们将极其虔诚的司铎、\OriginalTerm{乡村主教}{chorepiscopi}希帕提乌斯和阿布拉米乌斯，以及我们隐修士的\OriginalTerm{修院院长}{exarch}阿利比乌斯派往您的圣德，以便您能结束降临东方教会的那场灾祸：首先，确认最初由圣宗徒传授给我们的信德，谴责那兴起的异端，并公开指控那些胆敢贬损\OriginalTerm{救恩经济}{Œconomy}宣讲的人；其次，为那些因真理而受攻击的人作斗士，保卫他们。因为正是为了宗徒信德，至圣的，我们遭受了那不义的屠杀，因为我们不肯放弃福音教义的真理。如今，您的圣德不应该忽视那些与您同心合意的人所受的不义迫害，而应通过您正义的援助制止不义，教导真理的袭击者：那些试图按自己的私欲任意妄为的人，不能让他们达成目的。\par

\SourceRecord{Translated by Blomfield Jackson.；Nicene and Post-Nicene Fathers, Second Series；Vol. 3.；Edited by Philip Schaff and Henry Wace.；Buffalo, NY: Christian Literature Publishing Co.,；1892.；<http://www.newadvent.org/fathers/2707117.htm>.}

% structure: p0001 source=h1 target=section reason=page-title

\section{\WorkTitle{书信一一八}}

\newline{}\Emph{致罗马总执事。}\par

一场可怕的暴风雨袭击了我们的教会，但宗徒信德的拥护者在您的圣德中找到了安全宁静的港湾。您不仅捍卫福音教义的事业，而且极其憎恶对我的不公。当我远在三十五天的路程之外时，就被那些最公义的法官凭他们的喜好定罪了。他们抛弃了从天主我们的救主降临直到今日在教会中持有的教导，引入了一种新奇而虚伪的教义，与宗徒的传统截然相反，并公开与持守古老训导的人为敌。因此，最敬神的主啊，请屈尊点燃那位极其神圣和圣的总主教的热情，使东方的教会也能享受您的仁慈关怀。尤其要为从起初由圣宗徒传授的信德而战；保全我们祖先的遗产不受损伤，驱散压迫我们的迷雾。赐给我们晴朗的阳光代替无月之夜，并谴责那不义地加害于我们的屠杀之恶。将这一热心之举添加到您其他的善行中，与您的圣德相称。\par

\SourceRecord{Translated by Blomfield Jackson.；Nicene and Post-Nicene Fathers, Second Series；Vol. 3.；Edited by Philip Schaff and Henry Wace.；Buffalo, NY: Christian Literature Publishing Co.,；1892.；<http://www.newadvent.org/fathers/2707118.htm>.}

% structure: p0001 source=h1 target=section reason=page-title

\section{\WorkTitle{第一一九封书信}}

\Emph{致贵族阿纳托利}\par

阁下方才已完全知悉在厄弗所最公义的审判官们所行的事，因为他们的声音已传遍全地，他们最公义的判决也传到地极。哪个教会没有感受到由此掀起的风暴？一方受了冤屈，另一方行了不义，但那些既未受害也未行不义的人，却与受冤屈者一同忧伤，并为那些如此野蛮、违反所有神人律法、残害自己肢体的人哀悼。就连入室行窃的盗贼，在作案现场被抓获，也要先受审讯，然后由审判官惩处；就连杀人犯、盗墓者、奸淫者，也要先被带到法庭，命原告提出控诉，并检验证人的动机，看他们是否为了讨好控方而作证，或是对被告抱有偏见；在这之后，才命他们对所受的指控进行辩护。这样进行两次、三次，有时甚至四次；然后，也只有在那时，在从原告和被告双方的陈述中寻求了真相之后，才做出判决。至于这些人在其他人的案件中如何审判，我什么也不说，免得我似乎多管闲事。我只被迫为自己辩护，因为那不义的暴力行为迫使我如此。皇帝的诏令将我拘禁在家中，不许我离开我牧养之城的地界。主教会议的判决对我不利，一个人被判了罪，而他却远在三十五天的路程之外。\par

如今，万有的天主对先祖亚巴郎论及索多玛和哈摩辣说：\BibleQuote{因为控告索多玛和哈摩辣的声音实在很大，他们的罪恶实在深重；我要下去看看他们所行的，是否确实达到我所听见的控诉；如果不然，我也知道。}祂十分清楚那些人的邪恶，却仍然说：\BibleQuote{我要下去看看，}如此教导我们要等待事实的证明。但这些人从未传唤我受审，从未听过我的声音，拒绝听我陈述我的见解，却将我作为待宰的牺牲，交到真理仇敌的狂怒之中。\par

然而，我欢迎我的安息，尤其是在此时，当宗徒的法令被许多人破坏，新的异端得以加强之时。但惟恐有不认识我的人相信对我的诽谤属实，并且因认为我持有与福音不同的见解而跌倒，我恳求阁下向得胜的君王请求恩准：让我前往西方，在最虔诚圣洁的主教们面前为自己辩护；如果发现我甚至在最小程度上违背了信德的准则，就让我被投入深海之中。若他不赐予您这个请求，至少请他命令我居住在我的隐修院中，它离基禄斯一百二十英里，离安提约基雅七十五英里，离阿帕梅亚三英里。\par

在这些请求中，如果可能，我恳求前者；若不可能，至少我恳求通过阁下的调停，赐予我第二个请求。我将永远把您的仁慈铭记于心、挂在嘴边，恳求万军的上主，以现世和将来的福泽回报阁下。我被迫如此写信给您，因为我听说有人正竭力设法将我从此地迁走。\par

\SourceRecord{Translated by Blomfield Jackson.；Nicene and Post-Nicene Fathers, Second Series；Vol. 3.；Edited by Philip Schaff and Henry Wace.；Buffalo, NY: Christian Literature Publishing Co.,；1892.；<http://www.newadvent.org/fathers/2707119.htm>.}

% structure: p0001 source=h1 target=section reason=page-title

\section{第一二〇书}

\Emph{致\Person{路比西乌斯}。}\par

我想，即使是真理的仇敌，也必然对我们所受暴力的不义与非法感到愤慨。蒙真理养育的人，在您阁下领导之下，更应因这新异而令人震惊的悲剧而悲伤。这是理所当然的；那些更为悲伤的人，理应表现出更大的热诚与努力，以对抗那些不敬且非法的事，并将教会这如今面临撕裂之险的身体，恢复至昔日的和睦。因此，我恳求您阁下，将当前的危机视为一个属灵互惠的机会；请您一方为真理付出热诚，并从我们慷慨的主那里，同时领受祂在今生慈祥的眷顾，以及在来世的天国。\par

\SourceRecord{Translated by Blomfield Jackson.；Nicene and Post-Nicene Fathers, Second Series；Vol. 3.；Edited by Philip Schaff and Henry Wace.；Buffalo, NY: Christian Literature Publishing Co.,；1892.；<http://www.newadvent.org/fathers/2707120.htm>.}

% structure: p0001 source=h1 target=section reason=page-title

\section{第一二一号书信}

\Emph{致贵族\Person{阿纳托利乌斯}}\par

主眷顾并治理万有，祂既显示了我道理中的宗徒真理，也显示了加于我身的诽谤的虚妄。因为从极其虔诚和圣善的\Person{良}主，\Place{大罗马}的总主教，寄给圣善记忆的\Person{弗拉维安}以及聚集在\Place{厄弗所}的其他人的书信，完全与我所写和常在教会中宣讲的和谐一致。因此，我一读到它们，就赞美主的仁慈，因主没有完全抛弃教会，而保护了正统信仰的火星；或者——我岂不更应该说？——不是火星，而是一个巨大的火炬，足以点燃并光照世界；因为他确实在他的著作中保持了宗徒的印记，我们在其中立刻找到了由圣洁和蒙福的先知、宗徒以及他们在福音宣讲中的继承者所传承的，以及\Place{尼西亚}会议上的圣教父们所传承的。我承认我坚持这些，并控告所有持有其他教义者为不敬之罪。我将这些我的著作与他寄往\Place{厄弗所}的一封信并列放置，以便当阁下阅读时，能记住我常在教会中说的话，能认识这些教义的和谐，并能憎恶那些说谎者以及那些建立新异端反对宗徒教义的人。\par

\SourceRecord{Translated by Blomfield Jackson.；Nicene and Post-Nicene Fathers, Second Series；Vol. 3.；Edited by Philip Schaff and Henry Wace.；Buffalo, NY: Christian Literature Publishing Co.,；1892.；<http://www.newadvent.org/fathers/2707121.htm>.}

% structure: p0001 source=h1 target=section reason=page-title

\section{第一二二书}

\Person{乌拉尼乌斯} \Place{厄梅萨}主教。\par

我非常高兴，我们这些在品格上相合的人，也得以书信往来。但我不太明白你说\Quote{这不正是我的话吗？}的意思。如果这只是为了问候而说，我并不介意；但如果是想提醒我那个建议缄默的劝告以及所谓的\OriginalTerm{经济}{œconomy}，我非常感激，但我不接受这个提示。因为神圣的宗徒吩咐我们采取完全相反的方针：\BibleQuote{不论时机是否方便，总要坚持} \BibleRef{弟后4:2}。而且主对这位代言人本身说：\BibleQuote{不要害怕，只管讲} \BibleRef{宗18:9}；对依撒意亚说：\BibleQuote{大声呼喊，不要顾惜} \BibleRef{依58:1}；对梅瑟说：\BibleQuote{下去，吩咐百姓} \BibleRef{出19:21?}；对厄则克耳说：\BibleQuote{我立了你作以色列家的守望者} \BibleRef{则3:17}，又说：\BibleQuote{你若不去警告恶人} \BibleRef{则3:18}，以及诸如此类。因为我认为对一个知情者无需长篇大论。因此，我们不仅不以坦言为苦，反而欢喜快乐，赞美那位认为我们配受这些苦难的天主；而且我还要呼召我的朋友们去面对同样的危险。\par

如果他们认为我们没有遵守宗徒信仰的准则，而是偏左偏右，就让他们憎恨我们；让他们加入对立的一方；让他们与我们作战的人同列。但如果他们见证我们持守福音信息的正确教导，我们就以呼声欢呼他们：\BibleQuote{你们也要\InnerQuote{站稳，用真理作带束腰，……用和平福音的预备作鞋穿在脚上}}，等等，因为据说德性不仅包含节制、正义和明智，还包含勇敢，而且藉着勇敢，其余组成部分得以保存。因为正义在对邪恶的战争中需要勇敢的联盟；节制藉着勇敢的帮助战胜不节制。因此，万有的天主对先知说：\BibleQuote{义人因信德而生活，如果他退缩，我的心就不会喜悦他} \BibleRef{哈2:4/希10:38}。退缩，祂称之为怯懦。所以，我亲爱的朋友，要坚守宗徒的教义，因为\BibleQuote{要来的那一位将要来到，绝不迟延} \BibleRef{希10:37}，并且\BibleQuote{祂要按照各人的行为报应各人} \BibleRef{罗2:6}，因为\BibleQuote{这世界的局面正在逝去} \BibleRef{格前7:31}，真理必将显现。\par

\SourceRecord{Translated by Blomfield Jackson.；Nicene and Post-Nicene Fathers, Second Series；Vol. 3.；Edited by Philip Schaff and Henry Wace.；Buffalo, NY: Christian Literature Publishing Co.,；1892.；<http://www.newadvent.org/fathers/2707122.htm>.}

% structure: p0001 source=h1 target=section reason=page-title

\section{第一百二十三封信}

\Emph{致同一人。}\par

你的信很长，令人愉快，它显示出你的情谊是多么炽热而真诚。我对此如此欢喜，以至对于错误猜测你上一封信的开头之意一点也不感到遗憾。因为我误解你信函的意图，反而揭示了你的兄弟之爱，显明了你信德的真诚，并表明了你对真宗教的热忱。我们确实共同分享了先知的话语和磨难：你的圣德使用了话语；我则被飓风和巨浪打击，对着船上的划桨者用他的话语呼喊：\BibleQuote{敬奉虚伪偶像的人，离弃了自己的恩主。}或许那位既是约纳的主也是我的主的那一位，会恩赐我也能起身并从巨兽中被释放。但如果波涛继续翻滚，我仍相信即使如此我也将享受天主的保护，并从我的经验中学习祂的力量如何\BibleQuote{在软弱中显得完全}，因为祂按我的软弱来衡量危险。我提到的那位神圣先知被他的船员们一起抛入海中，但我却得到了你圣德和其他虔敬之人的安慰。为他们、也为你的虔敬，我祈求那赐给卓越的\Person{敖尼息佛}的福分也归于你，因为你没有以我的嘲笑为耻；反倒为了信德的缘故，分担了我的苦难。\par

还有一事我希望你知道：虽然其他虔敬的主教们给我送来了他们的馈赠，我都谢绝了——并非出于对赠送者的不敬，上天禁止——而是因为迄今合宜的食物由那位甚至不吝惜赐给乌鸦者提供。但对于主教大人您，我的做法不同，因为您情谊的热切确实克服了我迄今的既定原则。请确实相信，我虔诚的朋友，自从我们之间孕育出友情以来，我们爱德的热火便被点燃得更旺了。\par

\SourceRecord{Translated by Blomfield Jackson.；Nicene and Post-Nicene Fathers, Second Series；Vol. 3.；Edited by Philip Schaff and Henry Wace.；Buffalo, NY: Christian Literature Publishing Co.,；1892.；<http://www.newadvent.org/fathers/2707123.htm>.}

% structure: p0001 source=h1 target=section reason=page-title

\section{第一二四封信}

\Emph{致博学的\Person{马拉纳斯}。}\par

我也为教会的灾难感到痛心，为那肆虐的风暴哀哭；就我自己而言，我很高兴能摆脱动荡，享受那令我愉悦的平静。至于您学识所提及的那些仍在作恶的人，他们为其现今狂妄的不法行为付出代价的日子已不远。万有的主以权衡和法则治理万有，每当有人堕入无边的邪恶，祂的容忍便告终结，那时祂便作为审判者施行惩罚。预见此事，我祈求他们停止放纵，以免我再次为他们哀哭，眼见他们遭受惩治。\par

您的卓越我永不忘记，我恳求我们共同的主使您的家充满祝福。\par

\SourceRecord{Translated by Blomfield Jackson.；Nicene and Post-Nicene Fathers, Second Series；Vol. 3.；Edited by Philip Schaff and Henry Wace.；Buffalo, NY: Christian Literature Publishing Co.,；1892.；<http://www.newadvent.org/fathers/2707124.htm>.}

% structure: p0001 source=h1 target=section reason=page-title

\section{第一二五封信}

\Emph{致\Person{阿弗托尼乌斯}、\Person{塞奥多利图斯}、\Person{农努斯}、\Person{斯库拉基乌斯}、\Person{阿普托尼乌斯}、\Person{若望}，\Place{泽乌格马地区}的长官们。}\par

我知道你们信德的坚强与稳固，心中充满莫大的喜乐。因为，我们这些崇拜永恒天主圣三的人构成一个身体，自然，当部分肢体健全时，其余肢体也一同喜乐。神圣的宗徒这样说：\BibleQuote{若一个肢体受尊荣，众肢体都一同喜乐。}因此，我为你们为宗徒教义而奋斗，以及追随那位著名的纳波特在更优越之事上的榜样，而与你们一同喜乐。纳波特为了他的葡萄园遭受最不义的屠杀，因为他不愿舍弃祖业。你们战斗不是为葡萄园，而是为神圣的教义，你们拒绝这新奇而虚假的异端，因其玷污福音教导的光明；你们不容许有福的天主圣三之数目被减少或增加。因为，那些将独生子的苦难归诸天主性的人，就减少了它；那些胆敢引入第二个儿子的人，就增加了它。你们信仰一位独生子，正如你们信仰一位圣父和一位圣神。在降生成人的独生子内，你们看见他所取的本性，这本性取自我们，又为我们献上。否认这本性，就将我们的救恩置于远处；因为，如果独生子的天主性是不可受苦的，正如圣三的本性是不可受苦的，而我们又不承认那本性上适宜受苦者，那么，宣扬从未发生的苦难便是虚空而无益的。因为，如果那受苦者不存在，苦难怎能发生呢？我们宣称天主性是不可受苦的——这是我们的对手与我们共同宣认的道理。那么，当没有能够受苦的主体时，苦难怎能发生呢？那伟大的\OriginalTerm{救恩计划}{œconomy}的奥迹将会显得是一个幻象，一个纯然的外表而非现实。这是\OriginalTerm{瓦伦提努}{Valentinus}、\OriginalTerm{巴尔德萨内斯}{Bardesanes}、\OriginalTerm{马西昂}{Marcion}和\OriginalTerm{摩尼}{Manes}所杜撰的寓言。然而，从起初传承给教会的教导，即使在降生成人之后，仍承认一位独生子，我们的主耶稣基督，并宣认他是永生的天主，以及在末后的日子成为人；他成为人并非由于天主性的变化，而是由于人性的被取。因为，假如天主性变成了人性，那么它就不再是它原有的；如果它不是它原有的，那么那些崇拜这些对象的人称他为天主便是虚假。我们相反地承认天主的独生子作为天主是不变的，且是真实天主的儿子。因为我们从圣经学到：他原具有天主的形体，却取了奴仆的形体；他取了亚巴郎的后裔，并非变成了亚巴郎的后裔；他同我们一样分享了血肉和灵魂，但灵魂是不死且无玷的。他将这些保存起来，为了我们罪恶的身体，他献上自己无罪的身体，并为我们的灵魂献上他那毫无瑕疵的灵魂。正因如此，我们怀有普遍复活的希望，因为人类必定与它的初果共享，正如我们与亚当同死，同样我们也要与基督我们的救主共享他的生命。神圣的宗徒已明确教导我们，因为\Quote{如今}他说：\BibleQuote{基督从死者中复活了，成为睡了之人初熟的果子。因为死既是因一人而来，死人复活也是因一人而来；在亚当内众人都死了，照样，在基督内众人也都要复活。}\par

我这样写，并非为了教导你们，而是为提醒你们。我本想简短，但恐怕已超出了书信的篇幅。然而，至可尊敬且虔诚的司铎兼总隐修院院长\Person{梅奇马斯}因遵循爱德的法律，跋涉如此长途，向我们报告了你们卓越的热忱，并恳求我们以书信激励之。于是，我应允了他的请求，写了这封信，并恳求万有的主使你们在信德中安全稳固，比那筛我们的人更为坚强。\par

\SourceRecord{Translated by Blomfield Jackson.；Nicene and Post-Nicene Fathers, Second Series；Vol. 3.；Edited by Philip Schaff and Henry Wace.；Buffalo, NY: Christian Literature Publishing Co.,；1892.；<http://www.newadvent.org/fathers/2707125.htm>.}

% structure: p0001 source=h1 target=section reason=page-title

\section{第一二六封信}

\Emph{致\Person{萨比尼亚诺}主教}\par

当您离开那令人羡慕的教区时，我曾赞扬您的圣德。它一度可敬，如今却荒唐可笑，因为我们使它成了可买卖之物。我惊愕地听说您向驱逐您的人上诉。您本应反其道而行之，当被邀请执掌船舵时，应当婉拒，因为您的同船者已变成您的仇敌。最虔诚的先生，您难道不知道我们的救主通过祂神圣的宗徒教导我们宣讲什么吗？您难道不知道宗徒教义的继承者刚刚确立了哪些敬拜的对象吗？因为从福音最初宣讲直到现今盛行的黑暗时期，古代导师中谁曾听过有人宣讲\OriginalTerm{肉性与天主性的单一本性}{one nature of flesh and Godhead}，或曾敢于称那\OriginalTerm{独生者的可受难本性}{nature of the only begotten passible}呢？这类教义在当今被一些人公开大胆地宣讲，而在另一些人中，他们的宣讲却被忽视，且因沉默，人们成了亵渎的参与者。那么，有人可能会问，憎恶这类教义的人应采取什么适当的行动？我应当回答，他们面前有两种选择：要么正面交锋，证明这些教义的虚假；要么拒绝与公开不敬的对手共融。\par

我确实将所受的不公视为天主的祝福。我并非说我已感谢那些亏待我的人；我怎能感谢杀弟兄者，以及那些成为\Person{加音}追随者的人呢？\par

但我赞美我的主，因祂认为我配得受不义者的份，将我从作恶者和亵渎者中分别出来，并赐予我最甘美的安息。\par

\SourceRecord{Translated by Blomfield Jackson.；Nicene and Post-Nicene Fathers, Second Series；Vol. 3.；Edited by Philip Schaff and Henry Wace.；Buffalo, NY: Christian Literature Publishing Co.,；1892.；<http://www.newadvent.org/fathers/2707126.htm>.}

% structure: p0001 source=h1 target=section reason=page-title

\section{\WorkTitle{第一二七封书信}}

\Emph{致\Person{约彼乌斯}司铎兼隐修院院长。}\par

先祖亚巴郎在晚年获得了胜利。伟大的梅瑟当时已是老人，当他举手祈祷时，就打败了阿玛肋克。神圣的撒慕尔年老时驱逐了外邦人。你可敬的晚年效法了这些人。在我们为真宗教而战的战争中，你表现出大丈夫气概，捍卫福音教义的事业，以你精神的活力使年轻人黯然失色。\par

我听到感到欣喜，非常高兴，渴望拥抱你可敬的白发。但此事我不能做到，因为你因年迈被留在家，而我因皇帝的敕令被拘禁在此。但我藉这封信聊慰我的爱意，并给你虔诚以最爱的拥抱。我恳求你在祈祷中帮助现今被风暴席卷的教会，并为我赢得天主的支持，因为我为了福音教义的缘故遭受攻击，极需要来自上天的帮助。\par

\SourceRecord{Translated by Blomfield Jackson.；Nicene and Post-Nicene Fathers, Second Series；Vol. 3.；Edited by Philip Schaff and Henry Wace.；Buffalo, NY: Christian Literature Publishing Co.,；1892.；<http://www.newadvent.org/fathers/2707127.htm>.}

% structure: p0001 source=h1 target=section reason=page-title

\section{\WorkTitle{第一二八封信}}

\Emph{致\Person{坎迪杜斯}司铎兼院院长。}\par

我怕你虔诚灵魂的力量已被年迈所胜，你不像往常那样举手祈祷了。所以\Person{阿玛肋克}试图取胜。但愿有人扶持你的软弱，如古时\Person{乌尔}和\Person{亚郎}扶助立法者的手一样，使你能够打倒\Person{阿玛肋克}，拯救\Place{以色列}。这些日子我们尤其需要更热切的祈祷，因为外邦人、犹太人和一切异端都处于和平，唯有教会遭受风暴打击，被汹涌的波涛包围。\par

我们尤其需要你祈祷的援助，因为我们原以为与我们并肩作战的人，如今却与我们的敌人作战。\par

\SourceRecord{Translated by Blomfield Jackson.；Nicene and Post-Nicene Fathers, Second Series；Vol. 3.；Edited by Philip Schaff and Henry Wace.；Buffalo, NY: Christian Literature Publishing Co.,；1892.；<http://www.newadvent.org/fathers/2707128.htm>.}

% structure: p0001 source=h1 target=section reason=page-title

\section{\WorkTitle{第一二九封書信}}

\Emph{致司鐸\Person{馬格努斯·安托宁}。}\par

夜間的水手因看見港口的燈光而振奮，同樣，為宗徒信仰而冒險的人，也因同道者的熱忱而振奮。我們因聽到你的虔誠在神聖教義方面所作的努力而大得安慰，因為這心志是由一切美善恩賜的賜予者賜給你的，而你為了保守這些教義承受了一切勞苦。如今，我因你的熱忱得安慰，作微不足道的回報，勸你堅持你的神聖勞苦，藐視你的敵人如同容易捕捉的獵物（因為誰比那些缺乏真理的人更軟弱呢？），並信靠那曾說：\BibleQuote{我決不離開你，決不拋棄你}，以及\BibleQuote{看！我同你們天天在一起，直到今世的終結}的那一位。也請用你的祈禱幫助我，使我能夠自信地說：\BibleQuote{主在我身邊，我無所畏懼；人能對我做什麼呢？}\par

\SourceRecord{Translated by Blomfield Jackson.；Nicene and Post-Nicene Fathers, Second Series；Vol. 3.；Edited by Philip Schaff and Henry Wace.；Buffalo, NY: Christian Literature Publishing Co.,；1892.；<http://www.newadvent.org/fathers/2707129.htm>.}

% structure: p0001 source=h1 target=section reason=page-title

\section{第一三〇书}

\Emph{致主教弟茂德书。}\par

至高统治者并非毫无目的地允许那敌对我们的恶灵掀起不敬的波澜。祂这样做，是要考验水手们的勇气，并在彰显一些人英勇的同时，揭露另一些人的怯懦，从一些表面虔诚的人脸上剥下面具，并宣告另一些人是真理行列中的先锋。我们在当今时代看到了一个实例。风暴高涨；一些人暴露出他们隐秘的不敬；另一些人放弃了他们曾持守的真理，投靠了敌人的方阵，如今与他们一起攻击那些他们曾称为首领的人。这些事的见证者憎恶敌人，怜悯叛逃者，却不敢援助那些受到攻击、忠于宗徒教义的人。不仅如此，倘若叛徒更坚决地逼迫他们，他们或许自己也会转向攻击者的一方，不会宽恕他们的同道，反而与那些他们指责的人并肩向同道射出利箭。他们虽然从圣经中学到：加害邻人必受惩罚，而忍受不义则带来伟大而持久的赏报，但依然会如此行。\par

你本人的虔诚、对信德的热忱以及对我的善意，已在这场动荡中得到证明。你藐视一切可能阻止你的因素，两次写信给我，从而显示了你的兄弟之爱。你也表明了你为宗徒教义所承受的争战。你请求我写信告诉你关于救赎苦难我们应当如何思考和宣讲。我欣然接受了你的请求，并且——并非为了给你新知，而只是提醒一位天主所爱的人——我将向你讲述我从圣经和解释圣经的教父那里所学到的东西。\par

那么，最虔诚的先生，要知道：首要的是必须注意术语的区别，此外还要注意天主降生成人的原因。一旦这些清楚明了，关于苦难就没有任何含混之处了。因此，我们将首先向那些试图反驳我们的人提出这个问题：在给天主的独生子所起的名字中，哪些是在降生之前，哪些是在降生之后，或者更确切地说，与降生救世工程的运作相关？他们会回答说：先前的术语是\Quote{圣言}、\Quote{独生子}、\Quote{全能者}、\Quote{万有的主}；而\Quote{耶稣基督}这个名字属于降生。因为降生之后，天主的圣言、天主的独生子才被称为耶稣基督；因为经上说：\BibleQuote{看，今天为你们诞生了基督主}；又因为其他人曾被称为基督——司祭、君王和先知——为避免有人以为祂与他们相似，天使将\Quote{主}这个称号与\Quote{基督}连在一起，以证明所生者的至高尊威。再者，加俾厄尔对荣福童贞说：\BibleQuote{看，你将怀孕生子，并要给他起名叫耶稣}，\BibleQuote{因为祂要把自己的子民从他们的罪恶中拯救出来。}然而，在降生之前，祂从未被称为基督或耶稣。因为确实，神圣的先知们在预言未来之事时，使用了这些词语，正如他们在这些事件尚未发生时，就预言了诞生、十字架和苦难。尽管如此，即使在降生之后，祂仍被称为天主的圣言、主、全能者、独生子、创造者、造物主。因为祂不是藉着变化而成为人，而是保持祂所是的一切，取了我们的人性，因为\BibleQuote{祂虽具有天主的形体}，借用神圣宗徒的话说，\BibleQuote{却取了奴仆的形体。}因此，即使在降生之后，祂也被冠以降生之前的称号，因为祂的本性是恒常不变的。但当圣经叙述苦难时，绝未使用\Quote{天主}一词，因为那是绝对本性的名称。凡听见\BibleQuote{在起初已有圣言，圣言与天主同在，圣言就是天主}以及类似语句的人，都不会认为肉体在万世之前就已存在，或与宇宙的天主同性同体，或是世界的创造者。人人都知道这些术语专属天主性。同样，凡读到圣玛窦的族谱的人，都不会认为达味和亚巴郎按本性是天主的祖先，因为所取的人性是从他们而来的。\par

既然这些论点即使在极端的异端者中也属明显且无可置疑，我们既承认那先于万有的本性，也承认那近时的本性，我们就必须同时承认肉身的\OriginalTerm{可受苦性}{passibility}与神性的\OriginalTerm{不可受苦性}{impassibility}，不分割那\OriginalTerm{结合}{union}，也不将独生子分作两位，而是在独生子中默观各\OriginalTerm{本性}{natures}的属性。以灵魂和身体为例，两者属于同时代、自然结合的本性，我们惯于作此区分：描述灵魂为单纯、有理、不死，而身体为复合、可受苦、可朽。我们不分割那结合，也不将一个人切成两半。那么，更当如此：对于那在万世之前由圣父所生的\OriginalTerm{神性}{Godhead}，以及从达味的后裔所取的人性，采用类似的方法，明确地认识一种本性的永恒、无始、单纯、无限、不死、不变之特性，以及另一种本性的近时、复杂、有限、变动之特性。我们承认肉身如今是不死不腐的，尽管在复活之前它曾易受死亡与苦难；否则它如何被钉在木头上，并被置于坟墓呢？虽然我们承认本性的区分，但我们必须钦崇一位圣子，并承认同一子既是天主子也是人子，天主的形象和仆人的形象，达味之子与达味之主，亚巴郎的后裔与亚巴郎的创造者。结合使名称成为共通，但名称的共通并不混淆本性。对心怀正直者而言，一些名称显然适用于天主，另一些则适用于人；如此，可受苦与不可受苦都正当地用于主基督，因为按祂的人性祂受了苦，而按天主性祂保持不可受苦。如果，按不虔敬者的论证，祂是在神性中受苦，那么我以为，取肉身便是多余的；因为假设神性能够受苦，那么祂就不需要那会受苦的人性了。但即使如他们自己的论证所争辩的，神性是不可受苦的，而苦难是真实的，那么让他们当心不要否认那受苦者，免得连同否认了苦难的真实性；因为如果受苦者不存在，那么苦难就是不实的。现在，任何愿意打开神圣福音者四人集的人，都容易看到圣经明确地宣告了身体的苦难，并从中得知：阿黎玛特雅的\Person{若瑟}如何来到\Person{比拉多}面前，求\Person{耶稣}的身体；比拉多如何下令交付耶稣的身体；若瑟如何从木头上取下耶稣的身体，用殓布包裹，安放在新坟墓中。这一切均由四位福音书作者描述，并频繁提及身体。但如果我们的对手引用天使对\Person{玛利亚}及其同伴说的话：\BibleQuote{来，看主安放的地方}\BibleRef{玛 28:6}，让他们参考\WorkTitle{宗徒大事录}中记载虔敬的人\BibleQuote{埋葬了斯德望}\BibleRef{宗 8:2}的经文，并注意：得胜的斯德望所受到的习惯礼仪的是他的身体，而非灵魂。直到今日，当我们走近得胜殉道者的圣龛时，我们常问那埋葬在墓中的人是谁，而了解情况的人或许回答说：\Quote{殉道者犹利安}，或\Quote{罗马努斯}，或\Quote{弟茂德}。\par

很多时候，埋葬的不是完整的身体，而是非常小的遗骸，但我们仍然用属于整个人的名称来称呼身体。正是在此意义上，天使称主的身体为\Quote{主}，因为那是宇宙之主的身体。再者，主自己应许为世界生命而赐予的，不是祂不可见的本性，而是祂的身体。因为祂说：\BibleQuote{我所要赐给的食粮，就是我的肉，是为世界的生命而赐给的}\BibleRef{若 6:51}，当祂拿起神圣奥迹的象征时，祂说：\BibleQuote{这是我的身体，为你们而舍弃的}\BibleRef{路 22:19}。或按宗徒的说法：\BibleQuote{擘开的}\BibleRef{格前 11:24}。在祂提及苦难的任何地方，祂从未提到那不可受苦的神性。\par

因此，首要之事是必须向那些试图反驳我们的人提出这个问题：他们是否承认天主圣言取了完满的人性，并主张结合是无混淆地完成的？一旦这些要点被接纳，其余自会随之而来，苦难将被归于那会受苦的本性。我现在已总结这些要点，且已超出了书信的篇幅。我还寄去了我最近在一位非常虔诚圣洁的天主之人——主——的建议下所写的东西，以简明指导的形式，旨在教导宗徒道理的真理。若找到好的抄写员，我将把以对话形式写的东西也寄给您的圣善，以延伸论证，并借教父的教导加强我的立场。我还寄去了古代导师的一些陈述，足以显示他们教导的旨趣。最虔诚的先生，请以您的祈祷帮助我，使我得以穿越可怕的风暴，抵达救主平静的港湾。\par

\SourceRecord{Translated by Blomfield Jackson.；Nicene and Post-Nicene Fathers, Second Series；Vol. 3.；Edited by Philip Schaff and Henry Wace.；Buffalo, NY: Christian Literature Publishing Co.,；1892.；<http://www.newadvent.org/fathers/2707130.htm>.}

% structure: p0001 source=h1 target=section reason=page-title

\section{第一三一封信}

\Emph{致多利刻隐修院院长隆吉努斯。}\par

你同时显示了对真宗教的热忱和对邻人的爱德，这两者在当前是显然相联的，因为我正是为了宗徒的法令而受攻击，因为我拒绝放弃祖传的遗产，宁愿忍受任何苦难，也不愿轻视对福音信德一丝一毫的剥夺。你已分担了我的苦难，不仅通过来信安慰我，而且派来了十分尊贵且虔诚的玛窦与依撒格。我确信，你将从公义的主口中听到：\BibleQuote{我在监里，你们来探望了我。}我们卑微不足道，背负着沉重的罪担，但主是宽厚而慷慨的。他记念小者而非大者，并说：\BibleQuote{你们对我这些小子中最小的一个所做的，就是对我做的。}我祈求你，既然你以正确的道理著称，又以相称的生活发光，因此在天主面前大有胆量，求你以祈祷帮助我，使我能\BibleQuote{站住}，借用宗徒的话，\BibleQuote{抵挡错谬的阴谋}，逃脱毁灭者的罪，并在公义审判者显现之日，虽只有少许胆量，却能站立得住。\par

\SourceRecord{Translated by Blomfield Jackson.；Nicene and Post-Nicene Fathers, Second Series；Vol. 3.；Edited by Philip Schaff and Henry Wace.；Buffalo, NY: Christian Literature Publishing Co.,；1892.；<http://www.newadvent.org/fathers/2707131.htm>.}

% structure: p0001 source=h1 target=section reason=page-title

\section{第一三二封信}

\Emph{致\Place{埃德萨}的主教\Person{依巴斯}。}\par

主教导那遭受冤屈的人，不要沮丧，反要喜乐，并从古时的榜样中汲取安慰。因为从最初的人直到我们的日子，我们都能找到这样的人：他们热切敬拜万有的天主，却因同代人的缘故受屈，并陷入许多严重患难。若非我致信给一位精通神圣圣经的人，我本可详述这全部名单。但既然你，天主所爱的人，自幼便在天主圣言中受培育，我以为这样做是多余的。我仅求你用心看这些事，用怜悯看待那些行不义的好心圣职人员，用同情看待那些轻忽恶事的人，并为教会的动荡而忧伤。我求你因我为了真宗教而共受苦难而喜乐欢欣，并不住赞美那将此命运加给我者。至于尊荣、安逸、主教座的尊严及可怜的声名，就让它们归给那些杀害者吧。\par

让我们只持守福音的教理，若必要，借此忍受任何极致的痛苦，并选择可敬的贫困，胜过那多忧的财富。\par

我写这些话，并非为了劝诫你，因为我知道你在苦难中的圣德之勇气。我的目的是向你的虔诚表明我的心思，并告诉你，你身边有那些甘愿为真理冒险的同袍。我渴望给你写信已有一些时日，但一直找不到人递送。如今我遇到了非常可敬的虔诚司铎\Person{奥泽亚斯}，他既为真理而战，又敬爱你的虔诚。因此我写信问候你的圣德，并恳求你以你的祈祷支持我，并以你的来信安慰我。\par

\SourceRecord{Translated by Blomfield Jackson.；Nicene and Post-Nicene Fathers, Second Series；Vol. 3.；Edited by Philip Schaff and Henry Wace.；Buffalo, NY: Christian Literature Publishing Co.,；1892.；<http://www.newadvent.org/fathers/2707132.htm>.}

% structure: p0001 source=h1 target=section reason=page-title

\section{书信第一三三}

\Emph{致日尔马尼奇亚主教若望。}\par

大人，我一直知道您并没有忘记我们的友谊。我也一直期望并祈祷您的虔敬能留意准确的真理，并远离那些背叛真宗教者的共融，将那至高治理者对我们的眷顾归功于祂。因为，当我沉默不语、无所作为时，祂已终结了我们极其剧烈而可怕的苦难，并以这明亮的平静取代了那可怕的风暴。如今，主的慈爱已将这恩赐赐予我们，我发现我退隐的宁静确实令人愉悦，因为我感到有责任说服那些被针对我的诽谤所误导的人，既要使他们相信福音教导的真理，又要驳斥谎言的攻击。一旦这次驳斥完成，真理的胜利得到确保，我便打算退出公众生活，退隐到我长久渴望的安息中。至于真理的仇敌，我同先知一起呼喊：\BibleQuote{他们的纪念随喧哗而消灭，但主却永远长存。}至于我们自己，我与圣咏作者一同歌唱：\BibleQuote{祂从高处遣使，祂接了我，祂从多水中拉出我，祂救我脱离了我的强敌。}\par

此信是对您圣善所寄来的两封信的回复：一封由贝洛阿的司铎阿纳斯塔西转交，另一封由掌旗官德奥多托转交。在您最近的来信中，您提到另一封信，但并未送达。至于我前往那里的行程，在我得知最虔敬的皇帝对我下达了何种命令之前，我无法说任何话。他的信还未到达。\par

\SourceRecord{Translated by Blomfield Jackson.；Nicene and Post-Nicene Fathers, Second Series；Vol. 3.；Edited by Philip Schaff and Henry Wace.；Buffalo, NY: Christian Literature Publishing Co.,；1892.；<http://www.newadvent.org/fathers/2707133.htm>.}

% structure: p0001 source=h1 target=section reason=page-title

\section{第一三四封信}

\Emph{致\Place{贝若亚}主教\Person{德奥克提斯图}。}\par

我们的救主、立法者和主曾被人问：\BibleQuote{什么是第一条诫命？}祂回答说：\BibleQuote{你应全心、全灵、全意爱主你的天主。}祂又加上说：\BibleQuote{这是第一条诫命：第二条与此相似：你应爱你的近人如你自己。}然后祂接着说：\BibleQuote{这两条诫命是全部法律和先知的基础。}\par

因此，按照主的定义，持守这些诫命的人就是圆满地履行了法律；违犯它们的人就是犯了违犯全部法律的罪。那么，让我们在我们良心的公正而确切的法庭前检查一下，我们是否履行了神圣的诫命。第一条诫命是藉着完整地保守天主所赐的信德，憎恨真理的仇敌如同憎恨敌人，全心憎恨那些恨所爱者的人而持守的；第二条诫命是藉着极其重视对近人的照顾，并且不仅在顺境中，也在明显的逆境中遵守友谊的法则而持守的。反之，那些只顾自己安全，自以为安全的人，为此轻视友谊的法则，在朋友遭受攻击和打击时不关心他们，就被算在恶人和外人的行列中。万有的主向祂的门徒要求更好的东西。祂说：\BibleQuote{爱你们的仇敌，因为如果你们爱那些爱你们的人，你们有什么赏报呢？因为罪人和税吏也这样做。}然而，我连税吏所得的恩待也没有得到。我说税吏吗？我甚至连杀人犯和巫术者在监牢里所得的安慰也没有得到。如果每个人都效法这种残酷，那么我在生前除了因贫困而耗尽之外，没有别的剩下，死后也不得入土，反而成为狗和野兽的食物。但是，我在那些不关心现世生命、等候永福的人那里找到了支持，他们给了我多种安慰。但是仁慈的主\BibleQuote{使审判从天上发出；大地恐惧而静默，当天主起来审判时。}但恶人必要灭亡。新异端的虚假已被禁止，神圣福音的真理已被公开宣告。我本人与有福的\Person{达味}一同呼喊：\BibleQuote{愿主天主受赞颂，惟独祂行奇事；愿祂荣耀的圣名受赞颂；愿全地充满祂的荣耀；阿们，阿们。}\par

\SourceRecord{Translated by Blomfield Jackson.；Nicene and Post-Nicene Fathers, Second Series；Vol. 3.；Edited by Philip Schaff and Henry Wace.；Buffalo, NY: Christian Literature Publishing Co.,；1892.；<http://www.newadvent.org/fathers/2707134.htm>.}

% structure: p0001 source=h1 target=section reason=page-title

\section{第一三五封信}

\Emph{致罗慕路主教。}\par

你提醒了我古代的史事，并指出来自叙利亚的国王如何想起以色列诸王的仁慈，遂以恳求者的姿态出现，而如愿以偿。所以，阁下，请记住天主的义怒。天主因阿哈布滥用仁慈而将他交付于完全的毁灭，并借先知的口宣布了判决说：\BibleQuote{你的性命必代替他的性命，你的百姓必代替他的百姓。}我们因此被命令以正义调和仁慈，因为并非每一种仁慈都中悦万有的天主。当前的局势特别需要审慎的商议；因为我们是在为神圣的教义而争战，其中存有我们救恩的希望。然而，在此也能看到人与人之间的巨大差别。有些人确实感染了普遍的邪恶；另一些人则不加区分，时而提倡一种教义，时而又提倡其反面。有些人明知真理，却将其隐藏在灵魂的密室中，而与众人一同宣讲邪恶；还有些人充满嫉妒，将私人的恶意作为向真理开战的借口，对真理的先知施加种种恶行。又有些人接受了宗徒教义的真理，却因害怕当权者的势力而不敢宣扬，他们虽为我们的不幸之多而哀叹，却仍与掀起巨浪的人站在一边。我们把您的尊敬归入这最后一类。我们相信您在神圣教义上是纯正的，并认为您对我保持友爱，而随波逐流只是出于怯懦。既然如此，尽管我没有写信给其余的人，我仍写信给您的圣德，并收到了您的回复。我看清了您的意图，并且在某种程度上原谅您的懦弱。但如今，仁慈的主已除去了所有怯懦的缘由，显明了新奇的邪恶，并揭示了福音的简明真理。我，即使有头发般多的口舌，也无法恰当地称扬主的仁慈，因祂迫使我最强的对手公开宣讲我所宣讲的。因为我听说，与您的圣德同住的那位，当他听闻在大城市中发布了绝罚令之后，便不再模仿螃蟹的弯曲步态，而是在某一集会中辩论教义后，走上了直路。我们绝不能使言语迎合时代，而应永远持守真理的不屈准则。\par

\SourceRecord{Translated by Blomfield Jackson.；Nicene and Post-Nicene Fathers, Second Series；Vol. 3.；Edited by Philip Schaff and Henry Wace.；Buffalo, NY: Christian Literature Publishing Co.,；1892.；<http://www.newadvent.org/fathers/2707135.htm>.}

% structure: p0001 source=h1 target=section reason=page-title

\section{第一三六书}

寄\Person{济利禄}．\OriginalTerm{侍从官}{Magistrianus}书\par

闻君所遭之患，我心甚忧。我岂能不忧？因我以君之利为己利，且记宗徒之律，命我侪不仅\BibleQuote{与喜乐者同乐，亦与哭泣者同哭}。苦难本身能使仇敌亦生同情。\par

何者如丧妻之痛？彼妇无可指摘地负起婚姻之轭，使夫生活愉快，共理家务，主持内事，参决诸务；常能进善言，顺夫意。然而，又有何哀能胜于同时殡葬母亲与其所生之子？此子受精心培育，习得博学教育，君望其成为晚年之依靠；却在英年之际，面颊初生细毛时即被埋葬。若仅视灾祸之性质，实无慰藉可言。然念及吾族注定必死，天命已降于斯族；苦难为常事，因生命充满此类哀伤；如此，吾人当勇敢承当所遇之事，击退绝望之攻击，且高唱那奇妙颂歌：\BibleQuote{主赐予，主亦收回；主行其所视为善者；愿主之名受赞美。}吾人尚有更多安慰之理：吾人明确受教于复活的希望，期待死者复生之时。吾知主屡次称死亡为睡眠。若吾人信靠救主之言——诚然如此——则不当为已眠者哀悼，纵其睡眠稍长于常。吾人当等候复活。当记念宇宙之主以其智慧，既知现在亦知未来，导引万事为吾人获益。一智者深知此事，论及此类死亡时曰：\BibleQuote{诚然，他迅速被接去，免得邪恶改变他的悟性。}\par

我恳求君，让我们服从那智慧的万有治理者，服从祂的法令。无论它们令人愉悦或令人忧伤，都是美好而有益的，使人明智；为忍耐者预备冠冕。\par

\SourceRecord{Translated by Blomfield Jackson.；Nicene and Post-Nicene Fathers, Second Series；Vol. 3.；Edited by Philip Schaff and Henry Wace.；Buffalo, NY: Christian Literature Publishing Co.,；1892.；<http://www.newadvent.org/fathers/2707136.htm>.}

% structure: p0001 source=h1 target=section reason=page-title

\section{\WorkTitle{第一三七号书信}}

\Emph{致掌院若望。}\par

有福的\Person{达味}曾陷入数种过失，\Person{天主}以智慧安排万物，使它们记录下来，以益于后来的人。但弑父者、凶手、不虔敬、全然卑劣的\Person{阿贝沙隆}，并非因这些过失而发动疯狂战争反对他的父亲。他发起这最不义斗争的原因，是他贪图王权。然而，神圣的\Person{达味}在这些事发生时，开始记起自己所做的错事。我也自觉内心有许多过失的罪疚，但我保持了宗徒教义训导的无玷。而那些践踏一切人神法律、在我缺席时定我罪的人，并非因我所行的错事而判决我，因为我的隐秘行为并未向他们显明；他们却捏造伪证和诽谤反对我，更确切地说，因他们公开攻击宗徒的教义，因我服从他们而将我放逐。\BibleQuote{上主如同睡醒的人，又如勇士饮酒后清醒，他击退了他的敌人，使他们永远蒙羞受辱。}伪造和虚假的教义，祂已驱散如风，并为那些在神圣福音中传授给我们的教义提供了自由宣讲。这对我而言已足享完全的喜乐。我甚至不渴望那座我终身辛劳度过的城市；我所渴望的，是看见福音真理的确立。如今，主已满足了这个渴望。因此，我非常喜乐和幸福，我歌颂我们慷慨的主，并邀请尊驾与我同乐，并在我们的赞美中，献上恳切的祈祷：愿那些现在说一套、随后又说一套、如变色龙随叶色而变的人们，因主的慈爱而坚强，建立在磐石上，并因祂的仁慈，使他们以最高敬意尊崇真理。\par

\SourceRecord{Translated by Blomfield Jackson.；Nicene and Post-Nicene Fathers, Second Series；Vol. 3.；Edited by Philip Schaff and Henry Wace.；Buffalo, NY: Christian Literature Publishing Co.,；1892.；<http://www.newadvent.org/fathers/2707137.htm>.}

% structure: p0001 source=h1 target=section reason=page-title

\section{第一三八封信}

\Emph{致贵族阿纳托利乌斯}\par

我衷心欢迎分给我的休息，并收获其有益而愉快的结果。我们爱基督的皇帝，以他真诚的虔诚为收获帝国的果实，将暴风雨般动荡的教会的平静、被侵犯的信德的胜利、福音教义的胜利，作为他王权的初果，献给了赐予他的那一位。在这些之上，他加上了纠正对我的冤枉。如此重大、如此性质的冤枉，谁曾听说过？哪个凶手在缺席时被定罪？哪个奸淫者未经审讯被判决？哪个窃贼、盗墓者、巫师、抢劫教堂者，或其他任何不法行为的作恶者，在渴望上诉法律时被阻止，并在远离时被法官的判决处死？在他们身上，从未听说过这样的事。因为，按照我们的法律，原告和被告奉命在法官面前面对面站着，而法官必须等待明明白白的真理由现，然后，唯有那时，才或释放被告为无辜，或惩罚他，因他被起诉书所指控。而在我的案件中，所采取的程序正好相反。皇帝的信禁止我参加那个著名的大公会议，而最公义的法官们在我不在场时定了我的罪，不是在公正审判之后，而是在对用来归咎我的文件过度颂扬之后。无论是天主的法律，还是人的羞愧，都未能阻止这流血的行为。主席颁布命令，把真理抛到九霄云外，讨好当下的权势。那些与我思想相同、学说与我的学说一致、并曾表达对我及我的学说钦佩的人，服从了他。那一天仍然定了某些人背信之罪，某些人怯懦之罪；而对我来说，我为真理所受的苦难给了我自信的根据。而且，我们的主基督赐予我恩惠，\Quote{不但信祂，而且也为祂受苦}。因为所有恩宠的恩赐中最大的是为主所受的苦难，神圣的宗徒甚至将它们置于伟大奇事之前。\par

在这些恩惠中，我也自夸，尽管我卑微且微不足道，没有其他夸耀的理由。我恳求阁下您，为我这卑微之人，向爱基督的皇帝及最虔诚的、亲爱的天主、善良的导师\Person{奥古斯塔}，表达感恩之情，因她以如此恩赐回报了我们慷慨的主，并使自己为真宗教的热忱成为她统治的基础和根基。此外，恳求他们的神圣陛下完成这已良好规划的工作，召开一次会议，不像上次那样由喧嚣的乌合之众组成，而是——完全避开这些人——由那些决定并高度重视神圣事物、认为所有人事都不及真理重要的人组成。如果陛下们希望为教会带来古老的和平——我确信他们希望——恳求他们的虔诚恩典参与议事，以使他们的临在震慑那些反对的人，真理可能无人反驳，而她自己凭自己的力量审查事务的状况和宗徒学说的性质。\par

我向阁下您提出这一请求，并非因为我渴望再见\Place{库鲁斯}，因为大人您知道那是一个多么孤寂的城镇，以及我如何以各种建筑的巨大支出设法掩盖了它的丑陋，而是为了使我所宣讲的能与宗徒学说一致，而我的对手的发明则是假冒和卑鄙的。一旦此事实现——但愿以上主之助说出——我将在我余下的日子里愉快满足地度过，无论主命我居住何处。对于你，在真宗教中长大、被赋予良善财富的人，努力做出这一努力是合适的，并通过你急切的劝告，使我们最虔诚的皇帝和爱基督的\Person{奥古斯塔}更加热忱——他们已经热忱地以可称赞和正当的精力巩固他们光荣的帝国。\par

\SourceRecord{Translated by Blomfield Jackson.；Nicene and Post-Nicene Fathers, Second Series；Vol. 3.；Edited by Philip Schaff and Henry Wace.；Buffalo, NY: Christian Literature Publishing Co.,；1892.；<http://www.newadvent.org/fathers/2707138.htm>.}

% structure: p0001 source=h1 target=section reason=page-title

\section{\WorkTitle{第一三九封书信}}

\Emph{致执政官兼贵族\Person{Aspar}}。\par

在阁下其他善行之外，尚需添上一件：您已使我虔诚并极其信奉基督的皇帝（天主的恩宠已任命他为臣民的福祉）知晓了对我所行的极大不公，并且您已通过一道正义的谕令，废除了一道并非正义的谕令。在天主眷顾的支持下，我视他们所视为惩罚之事为善的手段，并欣悦地迎接了我的安息；然而我确被错谬且非法地对待，尽管在真理之敌诽谤加于我的过失上，我毫无半点罪责，却被处以了最大罪犯所应受的刑罚。不，我的命运比他们的更为严酷：我未经审判而被定罪；我在缺席时被宣判；当奉皇帝之命禁止我前往厄弗所时，我竟从圣洁的法官们那里领受了至为正义的判决。如今，因阁下积极斡旋，最仁慈的陛下已撤销这一切。我觉得若保持沉默、不向您致谢是有亏于理，于是便以此信恳请阁下，在我所里既为那位得胜而虔诚的皇帝、也为那位极其敬神且虔诚的奥古斯塔面前大力为我进言。我代他们以我所能的最大热忱向我们的善主祈求，求主保佑他们的帝国安全稳固，恩赐帝国对臣民成为慈爱的保护，对敌人成为恐怖，并为众人建立尊贵的和平。恳请阁下能请求他们完全终止教会的动荡，并下令召开公议会；不是像上一次那样，由惯于扰乱、使会议陷入混乱的人组成，而是要在和平与安静中，由通晓天主之事、惯于确认宗徒谕令、摒弃那些虚假且与真理相悖之事的成员组成。我表达此希望，为使阁下能收获这种行径所能产生的善果。\par

\SourceRecord{Translated by Blomfield Jackson.；Nicene and Post-Nicene Fathers, Second Series；Vol. 3.；Edited by Philip Schaff and Henry Wace.；Buffalo, NY: Christian Literature Publishing Co.,；1892.；<http://www.newadvent.org/fathers/2707139.htm>.}

% structure: p0001 source=h1 target=section reason=page-title

\section{第一四〇封书信}

\Emph{致\Person{文科马鲁斯}大人}\par

我深感惊讶地得知，阁下虽与我素不相识，只知我所受的冤屈，却挺身为我辩护，不遗余力地消除那针对我的阴谋所造成的一切后果。然而，阁下必将从我们慷慨的主那里得到酬报，因为那位曾应许赐予一杯凉水之赏的主，无疑会将更大的赏赐赐予那些奉献更大恩惠的人。\par

我所忍受的苦难，实为古往今来之人，或至少极少数人所经历，这苦难不仅来自公开的仇敌，而且，依我看，也来自我真正的朋友。前者攻击我，后者背叛了我。\par

世上谁曾听说过这样的审判？谁曾下令将一名罪犯拘押在三十五站之外，在其缺席的情况下进行审判？有哪位法官如此残忍不仁，不仅审判那些他从未听过其声音的人，而且还以最野蛮、最不人道的方式定他们的罪？主曾命令，对那不听劝告的迷途弟兄，在第一次、第二次、第三次劝诫之后，要待他如同\BibleQuote{外邦人和税吏}。然而，这些极其公正公义的法官，甚至对那些与己同信仰的人，也不曾像对待外邦人和税吏那样对待他们。他们确实见到外邦人和税吏，偶尔还与之交谈，并在他们显示出身份尊贵时，以一切尊敬和顺从相待。但他们却下令将我逐出家门，断绝水源，剥夺一切。他们就是以这种方式，希望成为我们在天之父的效法者，\BibleQuote{祂使太阳升起，光照恶人，也光照善人；降雨给义人，也给不义的人}。关于这些人，我不再多言。主的审判台就在眼前，那里需要的不是舞台上的伪装，而是真实的生活。如今我恳请阁下，向热爱基督并得胜的皇帝，以及极其虔诚而敬主的皇后，表达我的谢忱，因为他们使真宗教成为其虔诚帝国坚固的根基；并恳请陛下们命令召开一次大公会议，不是由暴烈之徒组成——他们会扰乱讨论——而是由热爱真理之人组成，他们能巩固宗徒的训诲，并摒弃这新奇的、伪造的异端，从而使教会的和平得以稳固。我祈求，因着这些可敬的辛劳，您能从我慈爱之主的手中收获果实。\par

\SourceRecord{Translated by Blomfield Jackson.；Nicene and Post-Nicene Fathers, Second Series；Vol. 3.；Edited by Philip Schaff and Henry Wace.；Buffalo, NY: Christian Literature Publishing Co.,；1892.；<http://www.newadvent.org/fathers/2707140.htm>.}

% structure: p0001 source=h1 target=section reason=page-title

\section{第一四一书信}

\Emph{致阿科梅特\OriginalTerm{隐修院}{Acoemetæ}\OriginalTerm{院长}{Archimandrite}\Person{马塞鲁斯}。}\par

您的圣德因您圣善的生活而熠熠生辉，在地上显示出天使交往的图像，但因您对宗徒信仰的热忱而更加明亮。如同船之龙骨、房之基石，对于选择度虔敬生活的人，福音教义的真理也是如此。为了这真理，当它受到攻击时，您勇敢地战斗，并非试图保护它如同它软弱，而是彰显您的虔敬心志；因为我们的导师基督的教导赋有稳固和力量，按同一位救主的应许：\BibleQuote{阴间的门决不能战胜她。}是那仁爱而慷慨的主认为我为了这道理也应当受辱并被杀。因为我们实在将羞辱视为荣誉，将死亡视为生命。我们听到了宗徒的话：\BibleQuote{因为天主赐给我们的，不仅是信祂，而且也要为祂受苦。}但是主如同睡醒的人起来，阻止了那些对天主口出亵渎、对我不义之人的口。他却使虔敬者的舌头涌出他们惯常话语的泉源。而我，却收获着休息的美果；当我看到教会的动荡时，我悲伤，但我也因摆脱烦扰而欢喜快乐。我一直因您令人钦佩的虔敬而感到欣慰，但至今未曾写信，并非缺乏对爱德命令的尊重，而是因为我在等待合适的时机。刚才，我遇到了由您的圣德派遣办理其他事务的最虔敬而明智的隐修士们，便不失时机地实现了我愿望。我向您的虔敬致敬。我首先恳请您以祈祷支持我，并进一步通过书信使我欣慰，因为天主的恩宠，我为了福音的缘故而受攻击。\par

\SourceRecord{Translated by Blomfield Jackson.；Nicene and Post-Nicene Fathers, Second Series；Vol. 3.；Edited by Philip Schaff and Henry Wace.；Buffalo, NY: Christian Literature Publishing Co.,；1892.；<http://www.newadvent.org/fathers/2707141.htm>.}

% structure: p0001 source=h1 target=section reason=page-title

\section{\WorkTitle{第一四二号书信}}

\Emph{致同一人。}\par

今已另函寄呈尊座，并托付于贵敬爱之弟兄。兹复再致函于圣座。此举既因您堪赞之生活，亦因您为宗徒信仰所表现之可嘉热忱——既不畏皇权，亦不惧主教之结盟。盖全会虽多数迫于压力，然毕竟以署名确认此新造之异端。然圣座不为所动，坚守古训，即主藉先知及宗徒训导教会所持守者。切望我能持守此等信条，至终保持对唯一圣父、唯一圣子、唯一圣神之信德与宣认。盖独生者之降生成人并未增加天主圣三之位数。既降生后，圣三仍是圣三。此乃我初所受之教理；此乃我之信德；我以此受洗；我以此宣讲；我以此施洗；我至今持守。对于妄言父者，主曾言：\Quote{他说谎时，是出于自己}，盖师之所言，合于弟子。是故，以谎言攻击我者，所言出于自己，而非描述我之实情。我以主之言自慰：\BibleQuote{几时人因我辱骂迫害你们，捏造一切恶言毁谤你们，你们是有福的。你们欢喜踊跃罢！因为你们在天上的赏报是丰厚的。}\par

恳请尊虔代祷，使我不要与行不义者同分，而与为福音真理受不义者共分。\par

\SourceRecord{Translated by Blomfield Jackson.；Nicene and Post-Nicene Fathers, Second Series；Vol. 3.；Edited by Philip Schaff and Henry Wace.；Buffalo, NY: Christian Literature Publishing Co.,；1892.；<http://www.newadvent.org/fathers/2707142.htm>.}

% structure: p0001 source=h1 target=section reason=page-title

\section{\WorkTitle{书信第一四三号}}

\Emph{致\Person{安德肋}，\Place{君士坦丁堡}隐修士。}\par

我从未见过你的虔诚，也从未以书信往来，但我已对你深深眷恋。究竟是什么造成了这魅力，并继续点燃它呢？乃是那些品尝你蜜糖的人一致传来的报告。众人都赞叹你信仰的正统、生活的光明、灵魂的坚毅、品德的和谐、交往的吸引力和甘甜，以及真正哲学养子的所有其他特征。因这一切，我依恋你的敬虔，我的渴望甚至促使我开始通信；但是，亲爱的先生，请尽快赐我所愿，让我收到你的书信。因为当朋友相隔遥远时，书信的往来能给与极大的安慰。你将写信给一个并非持异端见解的人，而是给一个在宗徒教导中成长、宣讲天主圣三而非四位一体的宣讲者。因为实际上，我看不出那些大胆试图将独生子的两性合为一的人，与那些试图将我们的主耶稣基督、永生天主之子、降生成人的天主圣言分为两个儿子的人，在亵渎上有何区别；如果真有这样的人的话；我不能这样想；但亚略派、欧诺米派和阿波林派一直无耻地捏造这毁谤来反对教会，而且，勤勉的学生很容易看出，我们备受赞誉的教父、教会的明灯，曾因这指控——如今那些新派异端的最优秀拥护者正以此攻击我——而遭受真理仇敌的劳苦。我们智慧的主已揭露了他们的不敬，因祂不能以容忍来确认这不圣的异端。\par

因此，先生，请确信你将写信给一个与你自己怀有相同感情的人；从我的大量著作中，你很容易确信这一点。\par

那么，请给我回信，愿你的信再次蒙天主许可，点燃爱心。在写信之前，请以你的祈祷帮助我，并恳求我们良善的主引导我的脚步走上正路，使我能在祂的律法下度过余生。你这因无玷生活而获得进身之阶的人，必能轻易说服那乐意赐我们美好恩赐的祂。\par

\SourceRecord{Translated by Blomfield Jackson.；Nicene and Post-Nicene Fathers, Second Series；Vol. 3.；Edited by Philip Schaff and Henry Wace.；Buffalo, NY: Christian Literature Publishing Co.,；1892.；<http://www.newadvent.org/fathers/2707143.htm>.}

% structure: p0001 source=h1 target=section reason=page-title

\section{书翰一四四}

\Emph{致士兵们}\par

人性处处相同，但人生追求繁多且多样。有人偏爱航海生涯，有人偏爱军旅；有人成为运动员，有人成为农夫；有人操持一业，有人从事他业。暂且略过其他差异不谈：有些人热心勤勉于神圣事物，让自己受教于宗徒教义的精确教导；而另一些人则相反，成了肚腹的奴隶，以为卑下快乐的享受就是幸福。还有另外一些人，居于这两种极端之间，既不表现这值得称赞的热忱，也不拥抱放荡的生活，但仍尊崇信德的纯朴。我认为，那些攻击\Quote{有些事对天主而言全然不可能}这一论断的人，不应被归入热心且精通神圣事物者之列，而应归入要么对宗徒教义没有确切认识的人，要么归入被快乐所奴役、随一时任性飘忽不定、时而提出此论时而提出彼论的人。\par

你曾请我就这些点写信。我此刻宁愿保持沉默。但为遵从主的诫命\BibleQuote{凡求你的，就给他}，我不得不简复。\par

那么我说，宇宙的天主能做一切事，但在\Quote{一切}一词中，只包含正直和良善的事，因为祂本性既智慧又良善，不接受任何相反本性的事，而只接受符合其本性的事。如有反对者反驳这一说法，就问他们：宇宙的天主，真理的立法者，能否说谎。若他们说天主可能说谎，就驱逐他们离开你的团体，视为不敬和亵渎。若他们同意说谎对宇宙的天主不可能，就其次问他们：那正义的泉源能否变得不义。若他们承认这对万有的天主也不可能，你还要再问：那不可测度的智慧之渊能否变为不明智，天主能否不再是天主，主能否不再是主，造物主不再是造物主，善不再是善而变为恶，真光不再是光而变为反面。若他们承认所有这些以及类似的事对天主是不可能的，那么你必须对他们说：因此许多事对天主是不可能的；而这些事不可能，非但并非无能的证据，反而表明最大的能力。\par

即使在我们自己灵魂的情况下，当我们说它不能死时，我们并非断言其软弱，而是宣告其不朽的能力。同样，当我们宣认天主的不变性、无受苦性、不死性时，我们不能将变化、受苦或死亡归于天主性。假设他们主张天主能随心所欲做任何事，你必须回答他们：祂不愿做任何不符合其本性的事；祂本性是善的，因此不愿任何邪恶的事；祂本性是公义的，因此不愿任何不义的事；祂本性是真实的，因此憎恶谎言；祂本性是不变的，因此不容许变化；既然祂不容许变化，祂就始终处于同一状态和状况。这为祂自己透过先知所断言：\BibleQuote{我是主，决不改变。}而蒙福的达味说：\BibleQuote{祢却常存，祢的岁月没有穷尽。}若祂是同一的，祂就不经历变化。若祂本性超越变化和变异，祂就不会从不死变为会死，从不受苦变为会受苦，因为若这可能，祂就不会取了我们的本性。但由于祂有不死的本性，祂取了一个能受苦的身体，并连同身体取了一个人的灵魂。祂使这两者都保持未受罪的玷污，并为那曾犯罪之灵魂的缘故交出祂的灵魂，为那曾死之身体的缘故交出祂的身体。而且由于所取的身体被称为天主的独生子自己的身体，祂把身体的苦难归于自己。但四位福音作者作证：被钉在十字架上的并非天主性，而是身体；他们都异口同声教导说，\Person{阿黎玛特雅人若瑟}来到\Person{比拉多}面前，求耶稣的身体；他把耶稣的身体从木架上取下来，用细麻布裹好，安放在自己的新坟墓里；\Person{玛利亚玛达肋纳}来到坟墓寻找耶稣的身体，跑去告诉祂的门徒们，当她找不到耶稣的身体时就报告了这些事。\par

这是福音作者们一致同意的教导。但如果你们的对手强调天使曾说过\BibleQuote{请来看主安放的地方}，让这些愚昧的人知道，圣经关于得胜的斯德望也说\BibleQuote{虔诚的人将他埋葬了。}然而，只有身体才被认为适合埋葬，灵魂并未与身体一同埋葬；尽管如此，仅有身体被用共同的名称来称呼。同样，蒙福的\Person{雅各伯}对他的儿子们说：\BibleQuote{将我与我祖先同葬。}他没有说\BibleQuote{埋葬我的身体。}然后他接着说：\BibleQuote{在那里埋葬了亚巴郎和他的妻子撒辣；在那里埋葬了依撒格和他的妻子黎贝加；我也在那里埋葬了肋阿。}他没有说\BibleQuote{他们的身体。}名称是身体或灵魂共用的，但尽管如此，他只用共同的名称称呼身体。同样，我们也经常描述圣宗徒、先知和殉道者的圣龛，或许一个是\Person{狄约尼削}的，另一个是\Person{儒利安}的，再一个是\Person{哥斯玛}的。然而我们知道，只有身体的残骸躺在那里，而灵魂在更神圣的区域安息。同样的习俗也见于日常用法，我们说某人死了，某人躺在此处；虽然我们知道灵魂是不死的，不与身体共享坟墓。正因如此，天使说\BibleQuote{请来看主安放的地方}，不是因为他把天主性关在坟墓里，而是因为他用了主的名字来称呼主的身体。\par

為證明這是聖教父們的看法，且讓他們注意亞大納削——至尊的亞歷山大里亞總主教，他以宣認裝飾了其主教職——的話。他呼喊說：\Quote{生命不能死，反而使死者復活。}\par

且讓他們也聽聞遠近馳名的羅馬主教達瑪蘇的話：\Quote{若有人聲稱在十字架上，天主性承受了痛苦，而不是那承受了祂的、僕人形體——祂所取的完整形體——的靈魂連同身體受苦，則此人應受絕罰。}\par

且讓他們也聽聞羅馬教會極其神聖的主教——主良——的話，他現今寫道：\Quote{天主子按照祂能受苦的方式受苦，不是按照那承受的本性，而是按照那被承受的。因為不能受苦的本性\OriginalTerm{不能受苦的本性}{impassible nature}承受了能受苦的身體\OriginalTerm{能受苦的身體}{passible body}，並為我們交出它，為使祂能完成我們的救恩，同時保持祂自己的本性不能受苦。}\par

又說：\Quote{因為祂來，不是為毀滅祂自己的本性，而是為拯救我們的。}\par

因此，若他們責備我們說：天主能行祂所願意的，但祂願意合乎祂自己本性的事，而不合乎祂本性的，祂既非願意，也不能行；那麼，也讓他們責備這些聖人以及所有持此立場的其他人吧。甚至讓他們責備宗徒，他說：\BibleQuote{藉著兩件不可更改的事，在這些事上，天主不可能說謊。}又說：\BibleQuote{我們若不信，祂仍是信實的：祂不能否認自己。}\par

將這些段落重複給你的對手聽；若他們被說服，就為此讚美良善的主，因祂藉著你的熱誠使他們獲益。若他們仍不被說服，就不要與他們爭論教義，因為神聖的宗徒禁止\BibleQuote{在言詞上爭辯，毫無益處，反而敗壞聽者}。但你務要保存福音的教導純正無誤，以致在祂顯現的日子，你能將所託付你的，連同應得的利息，呈獻給正義的審判者，並能聽到那渴望已久的話：\BibleQuote{做得好！良善忠信的僕人，你在小事上忠信，我要委派你管理大事；進入你主人的福樂吧！}\par

\SourceRecord{Translated by Blomfield Jackson.；Nicene and Post-Nicene Fathers, Second Series；Vol. 3.；Edited by Philip Schaff and Henry Wace.；Buffalo, NY: Christian Literature Publishing Co.,；1892.；<http://www.newadvent.org/fathers/2707144.htm>.}

% structure: p0001 source=h1 target=section reason=page-title

\section{第一四五封信}

\Emph{致君士坦丁堡的隐修士们}\par

那些以舌头为武器攻击我们的天主和救主的人，竟也将他们谎言的箭矢射向祂的正直仆人，这并没有什么新奇或可惊之处。仆人因主人所受的凌辱而深感悲痛，自然也必须分担这凌辱。主早已亲自警告过他们，祂安慰祂的圣门徒说：\BibleQuote{如果人们迫害了我，也要迫害你们。} \BibleQuote{若人们称家主为贝耳则步，对他的家人更该怎样呢？}接着，祂又鼓励他们，指出诽谤容易被揭露，因为祂继续说：\BibleQuote{没有掩盖的事，将来不被揭露的；也没有隐藏的事，将来不被知道的。}我常看到这神圣预言的真理，但现在看得格外清楚。那些对我进行诽谤、花大钱买我毁灭的人，已明显看出陷入了瓦伦提努斯和巴尔德萨内斯的谬误。他们原希望，只要用谎言的磨石磨快舌头来伤害我，就能掩盖自己的邪恶。因为自从我看到这早已熄灭的异端被这些人重新燃起，我就不停地高声呼喊，在私下和公开场合，在社交聚会和天主的圣殿中，都作见证，努力驳斥他们反对信仰的阴谋。于是，他们将侮辱倾泻在我头上，声称我宣讲两个儿子。但他们本应当面指证我，而不是背后诽谤。他们却完全相反。他们通过皇帝敕令在居鲁士将我手脚捆绑；他们迫使那些极正义的审判官不经审判就定我的罪，并对一个远在三十五站之外的人作了最公正的判决。任何被控行巫术或掘坟盗墓、杀人或通奸的罪犯，都从未遭受过这种对待。但目前我要放过这些审判官，因为主已临近，\BibleQuote{他以正义审判普世，以他的真理审判万民。}祂不仅追究言语和行为，甚至连恶念也要清算。但我认为应当驳斥所捏造的虚假指控。他们有什么证据证明我断言两个儿子？如果我是一个沉默寡言的人，或许有怀疑的余地；但我所做的正是为宗徒法令辩护，为主羊群提供教导的牧场，为此我写了三十五卷解释\WorkTitle{圣经}的著作，证明异端的虚假。这些人编造的谎言因此很容易被驳倒。成千上万的听众证明我教导了\WorkTitle{福音}教义的真理，任何愿意检验的人都可以看到我的著作摆在世人面前。我不是为两个儿子辩护，而是为天主的独生子辩护，对抗外教人、对抗犹太人、对抗染上亚略和欧诺弥乌斯瘟疫的人、对抗支持阿波利纳里乌斯疯狂的人、对抗受到马尔基翁堕落毒害的人——我从未停止奋斗：试图说服外教人，永生天主的永恒圣子本身就是宇宙的创造者；说服犹太人，先知们关于祂的预言；说服亚略派和欧诺弥乌斯派，祂与圣父同体、同尊、同能；说服马尔基翁的疯狂追随者，祂不仅是善的，而且是正义的，并且祂是救主，不是如他们杜撰的那样拯救别人的作品，而是拯救自己的作品。总而言之，在与每一种异端作战时，我命令人们俯伏朝拜独一的圣子。\par

何须多言，用几句话就能驳斥谎言？我们规定，每年前来领受圣洗的人，要仔细学习由神圣教父们在尼西亚所订立的信经；并按所受的命令，以圣父、圣子及圣神之名给他们授洗，逐一念诵每个名字。此外，在举行神圣礼仪时，无论清晨黄昏，还是将一日分为三时，我们都光荣圣父、圣子及圣神。如果如诽谤者所说，我们宣讲两个儿子，那么我们光荣哪一个，不朝拜哪一个？相信有两个儿子，却只将光荣颂归于一个，那是极端的愚昧。又有谁如此迷乱，听见圣保禄的话\BibleQuote{只有一个主，一个信德，一个洗礼；}，以及\BibleQuote{只有一个主耶稣基督，万物藉他而有；}，却制定与圣神教导相悖的律法，把独一者一分为二？但我是在无谓多言，因为这些人从小在谎言中长大，甚至不敢断言他们曾听我说过此类的话；但他们断言我宣讲两个儿子，因为我承认我们的主基督的两性。他们不愿看到，每个人类都有不朽的灵魂和可朽的身体；但至今没有人因为人兼有灵魂和身体而称保禄为两个保禄，也没有人称伯多禄为两个伯多禄，亚巴辣罕或亚当亦然。每个人都承认本性的区别，却不称一个人为两个保禄。同样，当我们称我们的主耶稣基督为天主的独生子、降生成人的天主圣言、既是天主子又是人子时——如\WorkTitle{圣经}所教导我们的——我们并不断言两个儿子，而是承认天主性和人性的独特属性。但是，那些否认祂所取于我们的人性的人，听到这些话就无法不恼怒。\par

我理当指出他们从哪些来源得到了这种不敬。\Person{西满}、\Person{默南德尔}、\Person{刻尔多}与\Person{玛尔基翁}绝对否认降生成人，并称童贞女的生育为神话。然而，\Person{瓦伦提诺}、\Person{巴西里德}、\Person{巴尔得撒内斯}与\Person{赫尔摩尼俄斯}及其追随者，接受童贞女的怀孕与生育，但他们否认天主圣言从童贞女取得任何东西，而是犹如经过一个管道那样经过她，仅以幻像向人类显现，看似是人，如同祂曾被\Person{亚巴郎}及其他一些古人看见一样。\Person{阿尔略}与\Person{欧诺弥俄斯}则相反，认为祂取了一个身体，但天主性扮演了灵魂的角色，以便他们将祂言行中卑微的部分归之于天主性。\Person{阿波利纳里}确实断言祂带着身体取了灵魂，但不是理性灵魂，而是被称为动物灵魂或\OriginalTerm{植物灵魂}{phytic}的灵魂。他们争辩说，天主性取代了心灵的位置。他从外邦哲学家那里学到了灵魂与心灵的区别，而神圣圣经说人由灵魂与身体构成。因为我们读到：\BibleQuote{上主天主用地上的灰土形成了人，在他鼻孔内吹了一口气，人就成了一个有灵的生物。}主在神圣福音中对祂的宗徒说：\BibleQuote{你们不要害怕那杀害肉身，而不能杀害灵魂的。}\par

这些教义之间的分歧如此之大。这些人现已竭尽全力，在亵渎方面超越\Person{阿波利纳里}、\Person{阿尔略}与\Person{欧诺弥俄斯}，并竭力重新栽种那由\Person{瓦伦提诺}与\Person{巴尔得撒内斯}旧日播下、而后被最卓越的农人拔除的异端。如同\Person{瓦伦提诺}与\Person{巴尔得撒内斯}一样，他们否认我们主的身体是从我们的本性取来的。但教会追随宗徒的足迹，在主基督内默观完美的天主性与完美的人性。因为正如祂取了身体，并非祂需要身体，而是藉此使所有身体获得不朽；同样，祂也取了灵魂——身体的向导，使每个灵魂藉此分享祂的不变性。因为即使灵魂是不朽的，却并非不变；它们经历许多频繁的变化，时而因这个对象，时而因那个对象体验快乐。由此，当我们被改变并倾向于更坏时，我们就犯了错误。但在复活之后，我们的身体享有不朽与不可朽坏，我们的灵魂则享有不受苦与不变。为此，天主的独生子取了身体和灵魂，使它们免受一切责备，并为人类献上了祭献。这就是为什么祂被称为我们的大司祭；祂被称为大司祭不是作为天主，而是作为人。祂作为人献祭，作为天主与圣父及圣神一同接受祭献。如果只有亚当的身体犯了罪，那么唯独身体应受益于治疗。但由于灵魂不仅参与了罪，而且是罪中的首位，因为首先是思想形成罪的图像，然后藉由身体实行，我以为，灵魂也应得到医治。然而，当神圣圣经清楚地宣告这些道理时，也许用推理来证明它们是多此一举。这教义由圣\Person{达味}和非常神圣的\Person{伯多禄}明确教导，一位从远古预言，另一位解释他的预言。宗徒之长的话是：\BibleQuote{达味既是先知，也知道天主曾以誓词向他起了誓，要从他的后裔中立一位，坐在他的宝座上；他预见了这事，就讲论基督的复活说：祂的灵魂没有被遗弃在阴府，祂的肉身也没有见到腐朽。}\par

现在，祂用这几句话给了我们很多关于同一问题的教导。首先祂指出，\OriginalTerm{被取的人性}{assumed nature}起源于大卫的腰；其次，祂不仅取了身体，也取了不朽的灵魂，第三，祂将身体和灵魂交付于死亡，并在重新取回后，按自己的意愿使它们复活。祂自己的话是：\BibleQuote{拆毁这座圣殿，三天之内我要把它重建起来。}但是我们知道天主性是不朽的。受苦的是\OriginalTerm{可受难的}{passible}，而\OriginalTerm{不受苦难的}{impassible}保持不受苦难。因为天主圣言成为人，不是要使不受苦难的本性成为可受难的，而是要借\OriginalTerm{受难}{Passion}将不受苦难的恩赐赋予可受难的本性。主自己在神圣的福音中曾说：\BibleQuote{我有权舍掉我的性命，也有权再取回它；谁也不能从我夺去，而是我自愿舍掉它；}\BibleQuote{为再取回它。}又说：\BibleQuote{为此父爱我，因为我为羊舍掉我的性命。}又说：\BibleQuote{现在我的心灵烦乱}\BibleQuote{我的心灵忧闷得要死。}关于祂的身体，祂说：\BibleQuote{我所要赐给的食粮就是我的肉，是为世界的生命而赐给的。}当祂传授神圣奥迹，擘开象征并分给大家时，祂说：\BibleQuote{这是我的身体，为你们而\Emph{擘开}的，以赦免罪过。}又说：\BibleQuote{这是我的血，为多人流出来，以赦免罪过。}又说：\BibleQuote{你们若不吃人子的肉，不喝祂的血，在你们内便没有生命}和\BibleQuote{谁吃我的肉，并喝我的血，必有永生}祂加上说\BibleQuote{在他内}。可以引述无数相同性质的段落，无论是在旧约还是新约，都指出身体与灵魂的被取，以及它们出于亚巴郎和大卫。阿黎玛特雅的\Person{若瑟}来到\Person{比拉多}那里，求取\Person{耶稣}的身体；神圣福音的四重权威告诉我们，他如何领取了身体，用细麻布包裹好，安放在坟墓中。我确实悲伤哀叹，因错谬的攻击，我被迫对那些自认为与我同信一信仰的人，引用我曾用以反对\Person{马尔西翁}瘟疫受害者的论据——其中，靠天主的恩宠，我已使超过一万人悔改，并引领他们领受\OriginalTerm{圣洗圣事}{Holy Baptism}。教会的儿女有谁曾对这些点有过疑惑？谁没有引用过圣教父们的这个教导？伟大的\Person{巴西略}的著作充满此道；同样，他的战友\Person{额我略}和\Person{安斐罗奇}，以及西方那些杰出的恩宠导师——大\Place{罗马}的主教\Person{达玛苏}、\Place{米兰}的\Person{盎博罗削}，还有\Place{迦太基}的\Person{西彼廉}，他为这些教义赢得了殉道者的荣冠。著名的\Person{亚大纳削}曾五次被逐出他的羊群，被迫流亡；他的老师\Person{亚历山大}也为此教义奋斗。\Person{厄斯达丢}、\Person{默肋提}、\Person{弗拉维雅诺}，东方的明灯；\Person{厄弗冷}，圣神的竖琴，每日以恩宠的溪流灌溉\Place{叙利亚}人民；\Person{若望}和\Person{阿提古}，真理的响亮宣讲者；以及比他们更早的人物：\Person{依纳爵}、\Person{波利卡普}、\Person{依肋内}、\Person{犹斯定}、\Person{希波律}，他们中大多数不仅在主教行列中闪耀，也装饰了殉道者的队伍。\par

那位现今统治着伟大的罗马、从西方四面八方传播正统教导光芒的至圣的\Person{良}，也在他自己的书信中向我表达了信仰的这一独特标志。这一切都清楚地教导：天主的独生子、永恒的天主，由圣父不可言说地所生，是独一圣子；并且在降生成人后，祂被称为人子和人，不是因为祂改变了成为人，因为祂的本性是不变的，而是因为祂取了属于我们的（人性）。他们也教导：祂作为天主是不受苦的、不朽的，作为人是可死的、可受苦的；但在复活之后，即使就祂的人性而言，祂也领受了不受苦和不朽，因为，虽然身体仍是身体，却是不受苦的、不朽的，确是一个神性的身体，并因神性的光荣而受光荣。蒙福的\Person{保禄}用这些话清楚地告诉了我们：\BibleQuote{因为我们的家乡是在天上，我们也等待救主，主耶稣基督，祂必要改变我们卑微的身体，使之相似祂光荣的身体。}祂没有说\Quote{祂的光荣}，而是说\Quote{祂光荣的身体}；而且主自己曾对祂的宗徒说过：\BibleQuote{站在这里的人中，有些人在未尝到死亡之前，必要看见人子带着祂父的光荣来临。}之后，过了六天，祂带他们上了一座高山，并在他们面前变了容貌，祂的面貌如同太阳，祂的衣服洁白如光。藉此，祂显示了第二次来临的方式。祂教导，所取的人性不是不受限的（因为这是天主性独有的特性），而是要发出神性光荣的闪光，放射出超越视觉能力的光芒。带着这光荣，祂被接升天；带着这光荣，天使们说祂要再来；因为他们的话是：\BibleQuote{那位从你们中被接升天的人，必要如你们所见祂升天一样，同样降来。}而且，当祂复活后被神圣的宗徒们看见时，祂给他们看了祂的手和脚；又给\Person{多默}看了祂的肋旁和钉痕与枪伤。因为为了那些坚决否认取了肉体的人，以及另外那些声称复活后身体的本性改变成了天主性本性的人，祂保留了钉痕和枪伤的印记，不变。并且，当祂使所有其他身体都从任何残缺中复活时，在祂自己的身体上留下了受苦的印记，为的是使否认取了身体的人，能藉祂的受苦而被指出错误；而那些主张祂的身体改变成了另一种本性的人，能藉钉痕得知它仍保持自己本有的特质。如果有人以为，主在复活后身体并非依然是身体，其证据是祂在门户紧闭时进入门徒们中间，那么让这样的人记住：当祂的身体尚属有死时，祂怎样在海面上行走；怎样在保持童贞封缄完整后诞生；以及当被图谋反对祂的人包围时，祂又怎样时常从他们手中逃脱。但我何须提及主呢？祂不仅是一个人，而且是在万世之前的天主，祂要做什么都容易。让他们说说\Person{哈巴谷}是怎样在顷刻之间从\Place{犹太}被迁移到\Place{巴比伦}，穿过兽穴的盖子，把食物带给\Person{达尼尔}，然后又返回，而没有破坏兽穴的封缄。追问主奇迹的方式是纯粹愚蠢的；但除了已经说过的，还应当知道：复活后，我们的身体也将是不朽的和永生的，且从属地之物解脱后，将变得轻盈而如以太。此外，神圣的\Person{保禄}用这些话明确地教导了我们：\BibleQuote{播种的是可朽坏的，复活起来的是不朽坏的；播种的是软弱的，复活起来的是强壮的；播种的是可耻的，复活起来的是光荣的；播种的是属生灵的身体，复活起来的是属神的身体。}并且在另一处：\BibleQuote{我们要被提到云彩中，在空中迎接主。}那么，如果圣人们的身体变得轻盈如以太，并且容易在空中穿行，我们就不能惊讶，主与独生子的天主性相结合的身体，在复活后成为不朽时，在门户紧闭时进入了。\par

从宗徒和先知那里，不难举出无数其他证据。但以上所述已足以表明我教导的要旨。我信一个圣父、一个圣子、一个圣神；并宣认一个天主性、一个主性、一个性体及\OriginalTerm{三个位格}{three hypostases}。因为独生子的降生成人并未增加天主圣三的数目，使天主圣三成为\OriginalTerm{四一性}{quaternity}，而是即使在降生成人之后，天主圣三仍然是圣三。当我宣认天主的独生子成为人时，我并不否认祂所取的本性，而是如我所说，既宣认那取的本性，也宣认那被取的本性。结合并未混淆两性的属性。因为如果空气通过其各部分接受光并不停止为空气，同时也不消灭光的本性——因为我们用眼睛看见光，用触觉认识空气，正如我们遭遇冷、热、潮湿或干燥——那么称天主性与人性的结合为混淆，实属愚昧。倘若被造的本性，既拥有从属且暂时的存在，当结合并在某种意义上混合时，仍然保持完好无损，而当光离去时，空气的本性仍独存，那么我认为，那创造万物的本性，当与祂从我们这里所取的本性结合并联合时，更应被承认继续保持其纯洁，并以同样方式使祂所取的本性保持完好无损。同样，金子与火接触时，既参与火的颜色和力量，却不失去自己的本性，而同时仍为金子并具有火的活跃特性。同样，主的身体是一个身体，但它是主的不受苦难、不朽、永存的身体，是属神且为神圣光荣所荣耀的。它既未与天主性分离，也不是其他任何人的，而只是天主的独生子自己的。因为它并不向我们显示另一位格，而是独生子自己披戴了我们的本性。\par

这是我不断宣讲的道理。另一方面，那些否认为我们而行的降生成人的人称我为异端者，他们所采取的行径有点像不贞的妇女：她们一边出卖自己的姿色，一边用她们行当的侮辱攻击诚实的妇女，并将属于她们自己淫荡的语言用在那些厌恶如此淫荡的妇女身上。埃及正是如此行事。她自愿陷入卑劣欲望的奴役。她将奴颜婢膝的谄媚倾注在一个贞洁的人身上。然后，当她的诡计未能引诱他，她的肉欲情欲的网罗也未能捕捉他时，她便称那忠于纯洁的人为犯奸淫者。\par

但这些人将因他们反对信德的诡计，以及他们为我设下的陷阱，而受天主的审问。我只嘱咐那些受了针对我的不实指控影响的人：要留一只耳朵给被告，而不要把两只都给原告。如此，他们将履行神圣的法律，这法律规定：\BibleQuote{不可传播谣言}，和\BibleQuote{在人与人之间、兄弟之间要秉公审判}。这些话中，神圣的法律告诫我们：不要相信针对不在场者所说的毁谤，而要当面审判被告。\par

\SourceRecord{Translated by Blomfield Jackson.；Nicene and Post-Nicene Fathers, Second Series；Vol. 3.；Edited by Philip Schaff and Henry Wace.；Buffalo, NY: Christian Literature Publishing Co.,；1892.；<http://www.newadvent.org/fathers/2707145.htm>.}

% structure: p0001 source=h1 target=section reason=page-title

\section{第一四六号信}

\Emph{致管家若望}\par

安宁与无忧无虑的生活对我而言十分可喜。因此，我封闭了隐修院的大门，并谢绝与朋友的交往。\par

但我得到消息，有人正对福音的信仰发动新的攻击，因此断定我的沉默可能带来危险。当对某位可朽的君王行不义时，不仅是犯罪的肇事者，连那些站在一旁、未努力驱散攻击者的人，也面临惩罚的危险：那么，那些敢于轻看对我们天主和救主所发的亵渎言语的人，该受怎样的惩罚呢？正是这种恐惧促使我现在写信，揭露我所听闻的革新。\par

据说城中有传闻，某些司铎祈祷后按惯例结束，有的说：\Quote{因为荣耀归于你、你的基督和圣神}，有的说\Quote{借着你基督的恩宠和慈爱，荣耀归于你和你的圣神}，那位非常聪明的总执事禁止使用\Quote{基督}一词，并说应当荣耀\Quote{独生子}。若是真的，这种不敬虔已无可超越。因为他要么将独一的主耶稣基督分裂为两个儿子，视独生子为合法而本性的，视基督为收养的、非婚生的，因此不配在光荣颂中被尊崇；要么是在支持那如今以狂乱喧嚣侵入我们的异端。若一场猛烈的暴风雨正压迫我们，任何人都会以为这亵渎者是因恐惧异端起源者的势力，而按时势的需要来配合自己的亵渎。但如今，被亵渎的那位斥责了风和海，并以平静祝福了颠簸的教会，而宗徒的宣讲在陆地与海洋各处传播，哪里还有亵渎的余地？就连最近卑鄙地将肉身与神性视为同一天性纳入教会教义的那些人，也从未禁止向主基督献上赞美。这事实易从从那里回来的人得知。一个在教会等级中占据首位的人，本应知晓圣经，并从中学习到：正如真理的使者将独生子与圣父同列，同样他们也使用\Quote{基督}这一名号代替\Quote{儿子}，时而将他与圣父同列，时而与圣神同列；因为基督不是别的，正是天主的独生子。我们可以引用神圣的保禄写给格林多人、却教导世界的经文：\Quote{只有一个天主父，万物出于祂……并有一位主耶稣基督，万物藉祂而有。}这样，他称同一人为基督、耶稣、主和万物的创造者。而写给得撒洛尼人时，他说：\Quote{愿天主父和我们的主耶稣基督亲自指引我们的路到你们那里。}在他的第二封给同一群人的信中，他将基督置于父前，并非颠倒顺序，而是教导：名称的顺序并不表明尊严和本性的区别。他的话是：\Quote{愿我们的主耶稣基督自己和天主、我们的父，那爱了我们、并藉恩宠赐予我们永久的安慰和美好的希望，安慰你们的心，并在一切善言善事上坚固你们。}在《罗马书》的结尾，经过某些劝勉后，他补充说：\Quote{弟兄们，我为我们的主耶稣基督的缘故，并为圣神的爱，恳求你们。}现在，如果他曾知道基督是不同于儿子的另一位，他就不会把他置于圣神之前。写给格林多人时，在书信的开头，他仅以基督的名字就足以影响信友。\Quote{弟兄们，我因我们主耶稣基督的名，恳求你们大家言谈一致。}而第二次写信给他们时，他这样结束：\Quote{愿我们主耶稣基督的平安、天主父的爱和圣神的共融与你们众人同在。}在这里，他不仅将基督的名置于圣神之前，也置于父之前；这在所有教会中，是奥迹礼仪的开端。\par

因此，按照这一特殊规定，我们天主和救主耶稣基督的尊名应从奥秘著作中删去。但关于这一点，不必多说。神圣宗徒的每一封信的开头都以这一称呼为标记。有时是\Quote{保禄，耶稣基督的仆人，蒙召作宗徒}；有时是\Quote{保禄，蒙召作耶稣基督的宗徒}；有时是\Quote{保禄，天主的仆人，耶稣基督的宗徒}。而他将祝福语与开端相配，源于同一源头，并将圣子的称号与天主父相联系，说：\Quote{愿恩宠与平安，由天主我们的父和主耶稣基督赐予你们。}最后，他以祝福结束他的信件：\Quote{愿我们主耶稣基督的恩宠与你们众人同在，阿们。}\par

还可以找到大量额外证据，易于得知我们的主耶稣基督不是别人，正是那组成天主圣三的圣子。因为同一位置者在万世之前就是独生子和天主圣言，复活后被称为耶稣和基督，从事实中接受这些名称。耶稣意为救主；\Quote{你要称祂名为耶稣，因为祂要把自己的民族从罪恶中拯救出来。}\par

祂被称为基督，是因为祂作为人受了圣神的傅油，并被称作我们的大司祭、宗徒、先知和君王。古时，天主的梅瑟曾呼号说：\BibleQuote{上主你的天主要从你中间、从你的弟兄中，给你兴起一位像我一样的先知。}天主的达味也呼喊说：\BibleQuote{上主发了誓，决不反悔：你照默基瑟德的品位，永为司祭。}这预言由天主的宗徒所证实。又说：\BibleQuote{我们既然有一位伟大的、进入了诸天的大司祭，天主子耶稣，就让我们坚持所明认的信仰罢。}\par

祂作为天主，在万世之前就是君王，这是先知诗歌所教导我们的，它说：\BibleQuote{天主，你的宝座永远常存；你王国的权杖是正直的权杖。}\par

祂作为人的尊威也显示给我们。因为祂作为天主和造物主，拥有万物的主权，祂作为人也接受了这尊威，因此接着又说：\BibleQuote{你爱慕正义，憎恨邪恶；因此天主，你的天主，用喜乐的油傅你，胜过你的同伴。}在第二篇圣咏中，受傅者亲自说：\BibleQuote{我已由祂被立为君王，在熙雍祂的圣山上，我要传报上主的命令：上主对我说：\InnerQuote{你是我的儿子，我今日生了你。你向我请求，我就将外邦人赐你作产业，将地极赐你作基业。}}祂这样说，是作为人说的，因为作为人，祂接受了祂作为天主所拥有的一切。而在圣咏的开始，预言的恩赐将祂与天主圣父并列，说：\BibleQuote{外邦人为何暴怒，万民为何妄想？地上的君王起来，首长聚在一起，反对上主和祂的受傅者。}\par

因此，谁也不可愚昧地以为基督是唯一天主子以外的另一位。我们不要自以为比圣神的恩赐更聪明。让我们听伟大的伯多禄的话：\BibleQuote{你是基督，永生天主之子。}让我们听主基督确认这宣认，因为祂说：\BibleQuote{在这磐石上，我要建立我的教会，阴间的门不能战胜她。}因此，明智的保禄，教会最卓越的工程师，也再没有立定别的根基。他说：\BibleQuote{我好像一个明智的建筑师，立定了根基，别人在上面建筑；但各人应留心怎样在上面建筑，因为除已立定的根基，即耶稣基督外，任何人不能再立别的根基。}那么，他们怎能想到别的根基呢？因为他们受命不是要立根基，而是在已立的根基上建造。天主的作者承认基督是根基，并以此称号为荣，如同他说：\BibleQuote{我已被钉在十字架上同基督一起，所以，我生活已不是我生活，而是基督在我内生活。}又说：\BibleQuote{在我看来，生活原是基督，死亡乃是福利。}又说：\BibleQuote{因为我曾决定，在你们中不知道别的，只知道耶稣基督，这被钉在十字架上的基督。}不久之前他说：\BibleQuote{我们所宣讲的，却是被钉在十字架上的基督：这为犹太人固然是绊脚石，为外邦人是愚妄，但为蒙召的人，不拘是犹太人或希腊人，基督却是天主的德能和天主的智慧。}他在致迦拉达人书中写道：\BibleQuote{但是，从母胎中已选拔我，以恩宠召叫我的天主，却决意将祂的圣子启示给我，叫我在外邦人中传扬祂。}但在致格林多人书中，他没有说\Quote{圣子}，而是说\BibleQuote{被钉在十字架上的基督}，这样并没有违反他的使命，而是承认同一位是耶稣、基督、主、独生者和天主的圣言。同样，在他致罗马人书的开头，他称自己为\Quote{耶稣基督的仆人}，并描述自己为\Quote{被选拔传天主的福音，这福音是天主先前借祂的先知在圣经上所预许的，论及祂的圣子，我们的主耶稣基督，祂按肉身是生于达味的后裔，按圣德的灵，借复活从死者中，被立为具有大能的天主子}等等。他称同一位既是耶稣基督，又是达味之子，又是天主子，作为天主和万有的主，然而在书信中间，在提到犹太人之后，他接着说：\BibleQuote{他们的先祖也是他们的，并且按肉身说，基督也是从他们来的，祂在万有之上，天主永远受赞美，阿们。}在这里，他说那位按肉身从犹太人而来的，是永恒的天主，并由心思正确者赞颂为万有之主。同样的教导也见于宗徒致卓越的弟铎书：\BibleQuote{期待所希望的幸福，和我们伟大的天主及救主耶稣基督光荣的显现。}在这里，他称同一位既是救主，又是伟大的天主，又是耶稣基督。在另一处他写道：\Quote{在基督和天主的国内。}此外，天使的合唱团向牧童们宣告：\BibleQuote{今天在达味城中，为你们诞生了一位救世者，他是主——基督。}\par

但对那些昼夜默思天主法律的人来说，实在无须写下所有此类证明；上述内容足以说服最顽固的反对者，不要分割神性的称号。然而，有一点我无法容忍略过。据说他曾说，有许多\OriginalTerm{诸基督}{Christs}，但只有一位\OriginalTerm{子}{Son}。我认为他是因无知而陷入此错误的。因为他若读过神圣圣经，便会知道我们慷慨的主也将\OriginalTerm{子}{Son}的称号赐予了许多人。立法者梅瑟，这位古代历史的作者，说：\BibleQuote{天主的儿子们看见人的女儿美貌，就随意挑选，娶来为妻}（\BibleRef{创 6:2}）；万有的天主亲自对这先知说：\BibleQuote{你要对法郎说：以色列是我的儿子，我的长子}（\BibleRef{出 4:22}）；在伟大的歌中他说：\BibleQuote{外邦人哪，你们当与祂的百姓一同欢乐；愿天主的众子在祂里面都刚强}（\BibleRef{申 32:43}）；又通过依撒意亚先知的口说：\BibleQuote{我养育儿女，将他们养大，他们竟悖逆了我}（\BibleRef{依 1:2}）；并通过三倍有福的达味说：\BibleQuote{我曾说：你们是神，都是至高者的儿子}（\BibleRef{咏 82:6}）；智慧的保禄致罗马人书中如此写道：\BibleQuote{因为凡被天主的圣神引导的，都是天主的儿子。你们所受的不是奴仆的心，仍旧害怕；所受的乃是儿子的心，因此我们呼叫：阿爸，父！圣神亲自与我们的心一同证明我们是天主的儿女；既是儿女，便是后嗣，就是天主的后嗣，和基督同作后嗣。如果我们和祂一同受苦，也必和祂一同得荣耀}（\BibleRef{罗 8:14-17}）；他致迦拉达人书中写道：\BibleQuote{你们既是儿子，天主就差遣祂儿子的灵进入你们的心，呼叫：阿爸，父！可见，从此以后，你不是奴仆，乃是儿子了；既是儿子，就靠着耶稣基督成为天主的后嗣}（\BibleRef{迦 4:6-7}）；他教导厄弗所人的是：\BibleQuote{在爱中预定我们借着耶稣基督得儿子的名分，归于祂自己}（\BibleRef{弗 1:5}）。\par

那么，若因基督这名通用，我们就不该将基督尊为天主，我们同样会因祂是子而不敢敬拜祂，因为这也是赐予许多人的称号。我为何说子呢？连天主这名本身也是天主赐予许多人的。\BibleQuote{上主，诸神之神，发言召呼大地}（\BibleRef{咏 50:1}）。又说：\BibleQuote{我曾说：你们是神}（\BibleRef{咏 82:6}），以及\BibleQuote{不可毁谤神}（\BibleRef{出 22:28}）。许多人也擅自将此名归于自己。欺骗人类的\OriginalTerm{邪魔}{dæmons}将此称号给了偶像；因此耶肋米亚喊道：\BibleQuote{那不曾创造天地的神，必从地上、从天下消灭}（\BibleRef{耶 10:11}）；又说：\Quote{他们为自己铸造了银像和金像}；依撒意亚先知嘲笑制造偶像后说：\BibleQuote{他用一部分在火中焚烧；用一部分烤肉吃，使自己暖和，说：啊哈，我暖和了，我看到了火光；剩下的他做成神，向它叩拜，说：求你拯救我，因为你是我的神}（\BibleRef{依 44:16-17}）；于是先知为他们悲叹说：\BibleQuote{他们要知道他们的心是灰烬}（\BibleRef{依 44:20}）。达味圣咏者教导我们歌唱：\BibleQuote{因为万民的神都是偶像，但上主创造了诸天}（\BibleRef{咏 96:5}）。\par

但这类称号的通用用法并不使那些受教于真宗教的人感到困扰。我们知道，邪魔将神性之名妄加于自己和偶像身上，而圣人们则是白白领受这一尊荣。\par

实际上，按其本质，万有的天主及其独生子和圣神乃是天主。这是由可钦的\Person{保禄}清楚教导我们的，他说：\BibleQuote{因为虽有称为神的，或在天上，或在地上，就如那许多的神，许多的主，可是我们只有一位天主，就是圣父，万物都出于祂，我们也在祂内；也只有一位主，就是\Person{耶稣基督}，万物都藉着祂而有，我们也藉着祂而有。}圣神被称为天主的神，人的灵魂也是如此，因为经上记载：\BibleQuote{祂的气息发出，}\BibleQuote{义人的灵魂和神体，请赞美上主。}圣咏作者\Person{达味}也称天使为神体：\BibleQuote{祂以风为使者，以火焰为仆役。}我何必提及天使和人的灵魂呢？连魔鬼也被主这样称呼：\BibleQuote{他要另外带去七个比自己更恶的魔鬼，进去住在那里，那人末后的情况就比先前更坏了。}然而，这个称呼的运用并不冒犯虔敬的读者，因为圣父、祂的独生子及祂的圣神，按其本质是一位天主；而成了人的圣言，我们的主\Person{耶稣基督}，按其本质是圣父的独生子；那完成天主圣三的护慰者是一位圣神。因此，虽然有许多被称为\Quote{父}，我们却敬拜一位父，即万世之前的父，祂自己将这称号赐予人，正如宗徒所说：\BibleQuote{因此，我在天主父面前屈膝，上天下地的各家都是从祂得名。}那么，不要因为别人被称为\Quote{受傅者}，我们就剥夺了对主\Person{耶稣基督}的敬拜。因为正如虽然有许多被称为神和父，却只有一位天主和万有之父，是在万世之前；虽然有许多被称为子，却只有一位真实而本性的子；虽然有许多被称为神体，却只有一位圣神；同样，虽然有许多被称为受傅者，却只有一位主\Person{耶稣基督}，万物藉着祂而有。教会极恰当地持守这个名称，因为她听见了\Person{保禄}——新娘的护送者——呼喊：\BibleQuote{我以妒爱将你们许配给一个丈夫，将你们如同贞洁的童女献给基督。}又说：\BibleQuote{你们作丈夫的，要爱你们的妻子，如同基督爱了教会。}又说：\BibleQuote{为此，人要离开父母，依附自己的妻子，两人成为一体。这奥秘真是伟大！但我是指基督和教会说的。}请听他所说：\BibleQuote{基督为我们成了诅咒，赎我们脱离法律的诅咒。}在别处又说：\BibleQuote{难道你们不知道，我们受过洗归于\Person{耶稣基督}的人，就是受洗归于祂的死吗？}又在另一处：\BibleQuote{凡受洗归于基督的人，就是穿上了基督。}又说：\BibleQuote{穿上主\Person{耶稣基督}，不要为肉体安排，去满足私欲。}\par

那些因天主恩赐而得福的人，既已学会认识这些段落及其他类似经文，心中燃起对慷慨之主的温暖爱情，便常将这最亲爱的名挂在唇边，以\WorkTitle{雅歌}的话呼喊：\BibleQuote{我的爱人属于我，我属于他；}\BibleQuote{我坐在他的荫下，满心欢喜，他的果实甘甜可口。}此外，我们所深爱的这个名，乃是从基督之名而来。我们被称为基督徒。\par

关于这名，万有的主说：\BibleQuote{上主天主必以另一名称呼自己的仆人，这名称在地上必受祝福。}以下是教会特别持守此名的原因。当天主的独生子成为人时，祂便被称为基督，那时人性接受了理智之光的光线；那时真理的使者向世界放射光芒。然而，教会的导师们常常不加区别地使用独生子的各种名称；有时他们光荣圣父、圣子及圣神；有时光荣圣父与基督及圣神；但就含义而言，这并无差异。因此，主命令以圣父、圣子及圣神之名施洗之后，有福的\Person{伯多禄}对接受他讲道并问该做什么的人说：\BibleQuote{你们每人要相信并因我们的主耶稣基督之名受洗。}仿佛这名包含了神圣命令的全部效力。同样的教导也由伟大的\Person{巴西略}——\Place{卡帕多细亚}的明灯，更可以说是世界的明灯——清楚地给了我们。他的话是：\Quote{基督之名是全体的宣认。}它同时指出圣父（傅油者）、圣子（受傅者）及圣神（藉以傅油者）。此外，在\Place{尼西亚}召开大公会议的三倍有福的教父们，在说明我们必须相信一位天主圣父之后，又加上：\Quote{也相信一位主耶稣基督，天主的独生子。}他们以此教导：主耶稣基督本身就是天主的独生子。\par

此外，还必须补充：我们不可断言升天之后，主基督不再是基督，而只是独生子。我们知道，神圣的福音、\BibleBook{宗徒}{大事录}和那位宗徒的书信都是在升天之后写成的。正是在升天之后，神圣的\Person{保禄}高呼：\BibleQuote{我们既然有一位经过了诸天的大司祭，天主子耶稣，就让我们坚守所宣认的信仰。}\BibleRef{希 4:14}又说：\BibleQuote{因为基督并非进入人手所造的圣所，那是真实圣所的模型；而是进入了天堂本身，如今为我们出现在天主面前。}\BibleRef{希 9:24}在谈及我们在天主内的望德后他又加上：\BibleQuote{我们拿这望德当作灵魂的锚，既坚固又安全，且深入幔内；那前驱耶稣已为我们进入那里，按照默基瑟德的品位永为大司祭。}\BibleRef{希 6:19-20}当给真福\Person{弟铎}写信论及第二次来临时他说：\BibleQuote{期待那有福的望德，和我们伟大的天主与救主耶稣基督光荣的显现。}\BibleRef{铎 2:13}致\Place{得撒洛尼}人信中也写类似的话：\BibleQuote{因为他们自己向我们传述，我们怎样进了你们中间，你们怎样离弃偶像归向天主，以事奉永生的真天主，并期待祂的圣子从天上降下，就是从死人中复活的耶稣，祂拯救我们脱离将要来临的愤怒。}\BibleRef{得前 1:9-10}又说：\BibleQuote{愿主使你们彼此之间的爱德和向众人的爱德增长满溢，就像我们爱你们一样，好使你们的心在我们主耶稣基督与众圣人来临的时候，在天主我们的父面前，在圣洁上无可指摘。}\BibleRef{得前 3:12-13}当第二次写信给同一教会时他说：\BibleQuote{弟兄们，现在我们因我们主耶稣基督的来临，和我们的聚集到祂面前，请求你们。}\BibleRef{得后 2:1}稍后预言敌基督的毁灭时他加上：\BibleQuote{主要用口中的气消灭他，并以祂来临的光辉毁灭他。}\BibleRef{得后 2:8}当劝勉罗马人和睦时他说：\BibleQuote{但你怎么判断你的弟兄？或者，你怎么轻视你的弟兄？因为我们众人都要站在基督的审判台前，因为经上记载说：我指着我的永生起誓，万膝必向我跪拜，万口必承认天主。}\BibleRef{罗 14:10-11}主自己宣布祂第二次来临时，除其他事外也说了这句话：\BibleQuote{那时，若有人对你们说：看，基督在这里，或在那里，你们不要信。因为如同闪电从东方发出，直照到西方，人子的来临也将这样。}\BibleRef{玛 24:23,27}\par

在祂的身体成为不死和不朽之后，祂称自己为人子，以那可见的本性为自己命名，因为天主性对天使来说确实是不可见的，正如主自己曾说的：\BibleQuote{从来没有人见过天主。}\BibleRef{若 1:18}又对伟大的\Person{梅瑟}说：\BibleQuote{没有人看见了我还能活着。}\BibleRef{出 33:20}\par

\BibleQuote{所以从今以后，我们不再按肉性认识谁了；即便我们曾按肉性认识过基督，如今却不再这样认识祂了。}\BibleRef{格后 5:16}这些话不是神圣的\Person{保禄}宗徒写来废除那所取的本性，而是为了确证我们将来自身的不朽、不死和属神生命。\par

因此\Person{保禄}宗徒继续说：\BibleQuote{所以谁若在基督内，他就是一个新受造物；旧的已成过去，看，一切都成了新的。}\BibleRef{格后 5:17}祂说的是将来的事，好像已经成就。我们尚未被赐予不死，但我们将被赐予；当被赐予时，我们不会成为无身体的，而是穿上不死。\BibleQuote{因为}，神圣的\Person{保禄}宗徒说，\BibleQuote{我们不愿赤身，而愿穿上衣服，为使那有死的被生命吞灭。}\BibleRef{格后 5:4}又说：\BibleQuote{因为这可朽坏的必须穿上不可朽坏的，这可死的必须穿上不死的。}\BibleRef{格前 15:53}如此，他没有说主是无身体的，而是教导我们相信就连这可见的本性也是不朽坏的，并以天主性的光荣受到光荣。在致\Place{斐理伯}人的书信中，他将这教导赐给我们更加清楚：\BibleQuote{因为我们的家乡是在天上；从中我们也等待救主，主耶稣基督，祂将改变我们卑贱的身体，使之相似祂光荣的身体。}\BibleRef{斐 3:20-21}藉着这些话，他清楚地教导我们，主的身体是一个身体，但是一个天主性身体，并以天主性的光荣受到光荣。\par

那么，我们不要回避那使我们获得救恩、并使万物更新的名字，正如我们的导师本人在他的\WorkTitle{厄弗所书}中所说的：\BibleQuote{按祂在祂自己内所定下的善意，为在时期圆满时，将天上和地上的万有，总归于基督元首。}我们倒不如从这蒙福的言语学习，我们该如何藉着将基督的名字与我们的天主圣父相连，来光荣我们的恩主。在\WorkTitle{罗马书}中，宗徒说：\BibleQuote{我的福音和耶稣基督的宣讲，是按照那永世以来隐秘、如今却彰显出来的奥秘的启示，并藉着先知们的经书，依照永生天主的命令，晓谕万民，使他们信德的服从；愿光荣藉着耶稣基督归于唯一的天主，直到永远。阿们。}在写给\WorkTitle{厄弗所书}时，他如此赞美：\BibleQuote{愿光荣归于那能按着在我们内运行的大能，充充足足地成就一切，超过我们所求所想的，愿祂在教会中，藉着基督耶稣，得着光荣，直到世世代代，永永远远。阿们。}稍早他又说：\BibleQuote{因此，我在我们主耶稣基督的父面前屈膝，天上地上各家，都是从祂得名。}再往后他说：\BibleQuote{凡事要奉我们主耶稣基督的名，常常感谢父天主。}当他以祝福回报斐理伯人的慷慨时，他说：\BibleQuote{我的天主必照祂在基督耶稣内的光荣的丰富，使你们一切所需用的都充足。}他为希伯来人祈祷说：\BibleQuote{但愿赐平安的天主，就是那凭着永久的盟约之血、使群羊的大牧人、我们主耶稣从死里复活的天主，在各样善事上成全你们，叫你们遵行祂的旨意，又藉着耶稣基督在你们心里行祂所喜悦的事。愿光荣归于祂，直到永永远远。阿们。}宗徒不仅在光荣时，也在劝勉和警告时，将基督与天主圣父连结。他对蒙福的弟茂德宣告说：\BibleQuote{我在天主和主耶稣基督面前嘱咐你。}又说：\BibleQuote{我在那赐万物生活的天主面前，并在那在本丢·彼拉多面前作过美好宣认的基督耶稣面前嘱咐你：要保守这命令，毫无玷污，无可指责，直到我们的主耶稣基督显现；到了日期，那可称颂、独有权能的万王之王、万主之主，就是那独一不死、住人不能靠近的光内、是人未曾看见也是不能看见的，要将祂显明出来。愿尊贵和永远的权能归给祂。阿们。}\par

这些是我们从神圣的宗徒们所学得的教训；这是若望和玛窦——福音讯息的两条浩荡江河——给予我们的教导。后者说：\BibleQuote{亚巴郎之子，达味之子耶稣基督的族谱。}前者在揭示那些万世之前的事物时写道：\BibleQuote{在起初已有圣言，圣言与天主同在，圣言就是天主。祂在起初就与天主同在。万物是藉着祂而造成的。}\par

\SourceRecord{Translated by Blomfield Jackson.；Nicene and Post-Nicene Fathers, Second Series；Vol. 3.；Edited by Philip Schaff and Henry Wace.；Buffalo, NY: Christian Literature Publishing Co.,；1892.；<http://www.newadvent.org/fathers/2707146.htm>.}

% structure: p0001 source=h1 target=section reason=page-title

\section{\WorkTitle{第一四七封信}}

\Emph{\Person{若望}，\Place{杰尔马尼西亚}的主教。}\par

收到您尊座的前一封信后，我立即回复了。关于当前的事态，不可能抱有任何好的希望。我担心这是普遍背道的开始。因为当我们看到那些为厄弗所所行之事（如他们所说，是出于暴力）而哀叹的人，毫无悔改的迹象，反而固守他们的不法行为，并在不义和不敬之上继续建造；当我们看到其余的人没有采取一致行动否认他们的作为，也不拒绝与那些固守不法行为的人共融，我们还能抱有什么好的希望呢？倘若他们赞赏所发生的事，仿佛一切都妥善正确，那么他们坚持自己所称赞的，倒也合情合理。但倘若他们（如他们所说）为所行之事哀叹，并声称那是出于强迫和暴力，那么他们为何不弃绝那不法的行为？为何这短暂易逝的当下，反倒被置于那必定到来的未来之前？他们为何公然说谎，否认教义中引入了任何新事？我因何种谋杀和巫术而被驱逐？此人犯了什么奸淫？他打开了什么坟墓？甚至外人也清楚，我和其余人是因教义而被驱逐。主多姆诺斯也是如此，因为他拒绝接受\Quote{《章节》}，就被这些优秀的人废黜，而他们却称赞这些章节，并承认坚守它们。我读过他们的提议，他们视我为异端的头目而拒绝了我，并以同样理由驱逐了其他人。\par

所发生的事足以证明，他们以为救主制定实践德性的法则，更多地是为哈马克索比人，而不是为他们。当有人指控彼西底的坎迪迪安，控告他多次奸淫及其他恶行时，据说议会主席说道：\Quote{倘若你们指控教义问题，我们接受你的控告；但我们来此不是为了裁决奸淫之事。}因此，被东方主教会议驱逐的阿特尼乌与阿塔纳修，竟被命令返回他们自己的教会；仿佛我们的救主没有制定关于行为的法则，而只命令我们遵守教义——而这些最聪明的人却带头败坏了教义。那么让他们停止嘲笑吧；让他们不再试图掩盖他们用言语和暴力所确认的不敬。倘若并非如此，就让他们告诉我们大屠杀的原因；让他们以书面形式承认我们救主两个本性的区分，以及结合是毫无混淆的；让他们宣告，结合后天主性和人性都保持完整。\BibleQuote{天主是轻慢不得的。}让他们否认那些他们曾多次弃绝、如今又在厄弗所批准的章节。不要让他们用谎言欺骗您的尊座。他们曾称赞我在安提约基雅的言论，那时我们本是弟兄，当他们被选为读经者、被祝圣为执事、司铎和主教时，在我讲论结束时，他们曾拥抱我，亲吻我的头、胸和手；有人甚至抱住我的膝盖，称我的教导为宗徒的教导——正是如今他们所谴责、所咒骂的教导。他们曾称我为明灯，不仅照亮东方，更照亮全世界，而如今我却被剥夺了公权，甚至（在他们能力所及）连面包也没有。他们甚至咒骂所有与我交谈的人。而那个不久之前被他们废黜、称为瓦伦提尼安与阿波利纳里的人，如今却被尊为信仰的殉道者，他们匍匐在他脚下，请求宽恕，称他为神师。难道木虱也会随石变色，变色龙也会随叶换皮，就像这些人随时代改变心意吗？我向他们放弃尊荣、地位、品级以及此生的一切奢华。为了宗徒的教义，我等待着他们认为可怕的那些灾祸，从主的审判中找到足够的安慰。因为我希望，为了这冤屈，主会赦免我的许多罪过。\par

现在我恳求您的尊座，小心不义的共融，并坚持让他们弃绝所行之事。若他们拒绝，就避开他们，视之为信仰的叛徒。阁下稍微等待，看暴风雨是否过去，我们并未认为是可责难的。但在东方首席主教被祝圣之后，每个人的心意必会显明。先生，请屈尊为我祈祷。此时我极需要这一帮助，好能抵挡所策划的一切恶谋。\par

\SourceRecord{Translated by Blomfield Jackson.；Nicene and Post-Nicene Fathers, Second Series；Vol. 3.；Edited by Philip Schaff and Henry Wace.；Buffalo, NY: Christian Literature Publishing Co.,；1892.；<http://www.newadvent.org/fathers/2707147.htm>.}

% structure: p0001 source=h1 target=section reason=page-title

\section{第一四八书}

\Emph{在\Person{加纳里乌斯}版本中。}\par

这是\Quote{最神圣的主教\Person{济利禄}的备忘录，由他派遣\Person{波西多尼乌斯}带往\Place{罗马}，关于\Person{聂斯托利}之事。} (Cyrill. Ep. XI. tom. lxxvii. 85.)\par

\SourceRecord{Translated by Blomfield Jackson.；Nicene and Post-Nicene Fathers, Second Series；Vol. 3.；Edited by Philip Schaff and Henry Wace.；Buffalo, NY: Christian Literature Publishing Co.,；1892.；<http://www.newadvent.org/fathers/2707148.htm>.}

% structure: p0001 source=h1 target=section reason=page-title

\section{第一四九封书信}

\Emph{是\Quote{\Person{若望}、\Place{安提约基雅}主教，致\Person{聂斯多略}书信的抄本。}}\par

这封信有时被认为是\Person{狄奥多莱特}真实撰写的。\par

\SourceRecord{Translated by Blomfield Jackson.；Nicene and Post-Nicene Fathers, Second Series；Vol. 3.；Edited by Philip Schaff and Henry Wace.；Buffalo, NY: Christian Literature Publishing Co.,；1892.；<http://www.newadvent.org/fathers/2707149.htm>.}

% structure: p0001 source=h1 target=section reason=page-title

\section{第一五〇书信}

\Emph{\Place{库罗司}主教\Person{塞奥多雷}致\Place{安提约基雅}主教\Person{若望}书}\par

我读到你所寄来、请求我以书面驳斥并向众人揭露其异端含义的绝罚条文时，深感悲痛。我忧心的是，一位被委任为牧者、受托照管如此庞大的羊群、并受命医治群羊中病弱者的人，自身竟也病入膏肓，且企图将他的疾病传染给羔羊，对待圈中羊群的残忍远胜野兽。野兽确实撕咬那些散离羊群的羊；但他却身在羊群之中，当被视为救主和保护者时，竟将隐秘的错误引入那些信任他的受害者之中。对于公开的攻击，尚可防备；但当攻击以友谊为伪装时，受害者便会疏于防范，伤害便轻易造成。因此，从内部发动战争的敌人远比从外部进攻的危险得多。\par

更令我痛心的是，他竟然以真宗教的名义、以牧者的尊严，发表异端亵渎之词，重拾\Person{阿波里纳}早已被消除的空虚邪恶的教导。除此之外，他不仅支持这些观点，甚至还敢于绝罚那些拒绝参与他亵渎的人——如果这些作品确实出自他手，而非某个真理的敌人冒用他的名字所写，正如古老的传说所载，将纷争的苹果抛入中间，从而煽风点火。\par

但无论这篇作品出自他本人还是别人冒他之名，我本人，在圣神之光的助佑下，按所赐予我的能力限度，尽我所能审查了这一异端腐朽的意见并驳斥了它们。我以福音作者和宗徒们的教导来面对它们。我揭露了这教义的怪异之处，并证明了它与神圣真理有着多么巨大的分歧。我通过将其与圣神的话语相比较，并指出它与神圣者之间存在着何等离奇刺耳的不协调，来完成此事。\par

至于这种绝罚的狂妄，我要说：保罗，真理清晰声音的宣告者，绝罚了那些败坏福音和宗徒教导的人，并且大胆地针对天使而非那些遵守神学家所定律法的人；他以下述祝福坚固他们：\BibleQuote{凡以此准则而行的人，愿平安与怜悯降在他们身上，并及于天主的以色列。} 那么，让这些作品的作者从宗徒的咒骂中收获他劳苦的应得酬报和他异端种子的果实吧。我们将持守圣教父们的教导。\par

我在本信之后附上了我的反驳论证，好使您阅读后能判断我是否有效摧毁了异端命题。我将每个绝罚条文分别列出，并附上反驳陈述，以便读者易于理解，并使对诸信理的驳斥清晰明了。\par

\SourceRecord{Translated by Blomfield Jackson.；Nicene and Post-Nicene Fathers, Second Series；Vol. 3.；Edited by Philip Schaff and Henry Wace.；Buffalo, NY: Christian Literature Publishing Co.,；1892.；<http://www.newadvent.org/fathers/2707150.htm>.}

% structure: p0001 source=h1 target=section reason=page-title

\section{第一五一封信}

\Emph{\Person{狄奥多勒}致\Place{幼发拉底西斯}、\Place{奥斯若恩}、\Place{叙利亚}、\Place{腓尼基}和\Place{西利西亚}诸隐修士的书信或讲话。}\par

当我默观教会在当前危机中的状况——最近袭击圣船的风暴、狂飙、波涛汹涌、深夜的深沉黑暗，以及除此之外，船员的争吵、桨手间的斗争、领航员的醉酒，还有最后坏人的不合时宜行动——我便想起耶肋米亚的哀歌，与他一同哭喊：\BibleQuote{我的五脏六腑，我的五脏六腑！我心中剧痛，我的心在我内骚动，}为藉我眼中的泪水驱散沮丧的浓云，我求助于泪泉。在如此狂暴的风暴中，领航员们应当清醒，与风暴战斗，并关注船只的安全；水手们应当停止争吵，努力通过祈祷和技巧消除危险；船员们应当保持和平，既不与彼此争吵，也不与领航员争吵，而要恳求海的主用祂的杖驱散黑暗。如今没有人愿意做任何这类事；正如夜战中所发生的那样，我们彼此无法辨认，我们放过敌人，却将武器浪费在自己人身上；我们为敌人伤害自己的同伴，而旁观者则嘲笑我们醉酒的愚蠢，享受我们的灾难，并乐于看到我们自相残杀。这一切的责任在于那些试图败坏宗徒信仰的人，他们胆敢在福音的教导之外添加一个怪异的教义；在于那些接受了他们发出并带有绝罚令送到帝都的亵渎\OriginalTerm{章节}{Chapters}，并如同他们所想象的凭自己的签名确认了它们的人。但这些章节无疑是从\Person{阿波利拿里}的酸根中长出的；它们沾染了\Person{亚略}和\Person{欧诺米}的错误；仔细审视，你会发现它们并未摆脱\Person{摩尼}和\Person{华伦提努}的不敬。\par

在他的第一章中，他拒绝了为我们所做的经世，教导说天主圣言并没有接受人性，而是祂自己变成了血肉，从而主张降生成人不是真实的，而是表象和外表。这是\Person{马尔西翁}、\Person{摩尼}和\Person{华伦提努}不敬的结果。\par

在他的第二章和第三章中，他似乎完全忘记了他在前言中所说的，引入了位格结合和本性结合的相遇，通过这些术语，他描绘了神性本性和仆人形态之间发生了某种混合和混淆。这是\Person{阿波利拿里}异端的创新。\par

在第四章中，他否认了传福音者与宗徒用语之间的区别，并拒绝接受正统教父教导所允许的，即将神圣尊荣的术语理解为指代天主性，而将按人性方式所说谦卑的术语用于所取的本性；由此，心地正直的人不难察觉其与不敬的亲近关系。因为亚略和欧诺米曾断言天主的独生子是受造物、从无中制成、且是仆人，他们竟敢将那些以卑微和人性方式所说的话应用于祂的天主性；借此手段来确立本质的差异与不相似。此外，简而言之，他还论证说，那不可受难、不可改变的基督之天主性本身受了苦、被钉十字架、死了并被埋葬。这甚至超过了亚略和欧诺米的疯狂，因为就连那些胆敢称宇宙的创造者和造物主为受造物的人，也未曾达到这种不敬的程度。更进一步，他亵渎圣神，否认祂发自父，正如主的话所说，反而坚持祂来自子。这是阿波里纳派种子的果实；我们接近了马其顿纽斯的邪恶耕作。这些是那埃及人的后代，卑劣父亲的更卑劣子女。这种生长——那些受托医治灵魂的人本应在它尚在母腹时就使之流产，或一出生就将其毁灭，因为它对人类危险致命——却被这些可敬之人珍视，并大力推广，既导致他们自己的毁灭，也导致所有听从他们之人的毁灭。相反，我们热切希望保持我们的传承不受触动；我们努力保存我们所领受并赖以受洗及施洗的信德完好无损、纯洁无瑕。我们宣认我们的主耶稣基督，完美天主和完美的人，拥有理性灵魂和身体，在万世之前由父所生，就天主性而言；在末期为了我们人类和我们的救恩，由童贞玛利亚所生；同一主就天主性而言与父同性同体，就人性而言与我们同性同体。因为两性结合了。因此我们承认一位基督、一位子、一位主；但我们不破坏结合；我们相信这种结合是在不混淆的情况下完成的，顺从主对犹太人的话：\Quote{你们拆毁这座圣殿，三天之内我要把它重建起来。} 反之，如果存在混合和混淆，且从两者中生出了一性，祂就应该说：\Quote{拆毁我，三天之内我就要被复活。} 但如今，为表明按本性的天主与圣殿之间的区别，并且两者是一基督，祂的话是：\Quote{你们拆毁这座圣殿，三天之内我要把它重建起来，} 清楚教导说，不是天主在被拆毁，而是圣殿。后者的本性是可朽坏的，而前者的权能复活了那正在被拆毁的。此外，正是顺从神圣的圣经，我们承认基督是天主和人。我们的主耶稣基督是天主，这由蒙福的传福音者若望所肯定：\BibleQuote{在起初已有圣言，圣言与天主同在，圣言就是天主。他起初就与天主同在。万物是藉着他造的；凡受造的，没有一样不是藉着他造的。} 又说：\BibleQuote{那普照每个人来到世上的真光。} 主自己也明确教导说：\BibleQuote{看见了我，就是看见了父。} 以及\BibleQuote{我与父原是一体}和\BibleQuote{我在父内，父也在我内}（\BibleRef{若 10:38}），蒙福的保禄在\WorkTitle{致希伯来人书}中说：\BibleQuote{他是天主光荣的反映，是他本体的真像，以自己大能的话支撑万有}，在\WorkTitle{致斐理伯人书}中说：\BibleQuote{你们该怀有基督耶稣所怀有的心情；他虽具有天主的形体，却没有以自己与天主同等为应当把持不舍的，反而空虚自己，取了奴仆的形体}。在\WorkTitle{致罗马人书}中：\BibleQuote{他们的祖先也是他们的祖先；按血统说，基督也出自他们，他是万有的天主，当受赞美直到永远。阿们。} 在\WorkTitle{致弟铎书}中：\BibleQuote{期待所希望的幸福，和伟大的天主及我们的救主耶稣基督光荣的显现。} 依撒意亚高呼：\BibleQuote{有一个婴孩为我们诞生了，有一个儿子赐给了我们；他肩上担负着王权；他的名字要称为神奇的谋士、强有力的天主、永远之父、和平之王。} 又说：\BibleQuote{他们带着锁链前来归顺你，向你俯伏恳求说：惟独天主在你中间，以外没有别的神，没有神。你实在是隐藏的天主，以色列的天主、救主。} 然而，\Quote{厄玛奴耳}之名指示了天主和人，因为在福音中它被解释为\Quote{天主与我们同在}，也就是说\Quote{天主在人内}，天主在我们的本性内。神圣的耶肋米亚也发出预言：\BibleQuote{这就是我们的天主，没有别的可与他相比。他找出了全部知识的道路，将它赐给祂的仆人雅各伯和祂的爱徒以色列；此后，他在地上显现，与人往来。} 此外，还可以找到无数其他段落，无论是在神圣福音和宗徒著作中，还是在先知预言中，都表明我们的主耶稣基督是真天主。\par

主自己教导我们，在降生成人之后，祂被称为人，正如祂对犹太人说：\Quote{你们为什么图谋杀害我？} \Quote{一个把真理告诉你们的人。} 在 \BibleBook{格林多}{前书} 中，蒙福的\Person{保禄}写道：\Quote{因为死亡是藉着一个人来的，复活也是藉着一个人来的，}为指明他所说的是谁，他又解释自己的话并说：\Quote{就如\Person{亚当}内众人死了，照样在\Person{基督}内众人也要复活。}在写给\Person{弟茂德}的信中，他说：\Quote{因为只有一个天主，在天主与人之间的中保也只有一人，就是\Person{基督耶稣}，祂是人。}在 \BibleBook{宗徒}{大事录} 中，他在\Place{雅典}的演讲中这样说：\Quote{这些蒙昧无知的时代，天主都置之不问；但现在祂却命令各处的人都要悔改，因为祂已定了一个日子，要按正义藉着祂所指定的人审判天下，并且赐给众人凭据，因为祂已使祂从死人中复活了。}蒙福的\Person{伯多禄}在向犹太人讲道时说：\Quote{诸位以色列人，请听这些话：纳匝肋人耶稣是天主用祂在你们中间藉着祂所行的奇迹、奇事和征兆所证明的人，正如你们所知道的。}先知\Person{依撒意亚}在预言主\Person{基督}的苦难时，在稍前称祂为天主之后，称祂为人，经文说：\Quote{一个忧苦的人，熟悉病痛。} \Quote{祂确实背负了我们的病痛，承担了我们的忧苦。}我本可收集其他一致的圣经经文并写在这封信里，只因我知道你熟习神圣的启示，正如在 \BibleBook{}{圣咏} 中被称为有福的人所应有的。现在我将收集证据的工作留给你自己的勤勉，继续我的主题。\par

因此，我们宣认我们的主\Person{耶稣基督}是真实的天主和真实的人。我们不将一位\Person{基督}分成两个位格，而是相信两个本性无混淆地联合。这样，我们就能毫不费力地驳斥异端者多方面的亵渎：因为那些背叛真理之人的错误繁多且多样，我们将逐一指出。\Person{马西翁}和\Person{玛乃斯}否认天主圣言取了人性，也不相信我们的主\Person{耶稣基督}是由童贞女所生。他们说天主圣言本身被塑造成人形，只是外表像人，而非真实的人。\par

\Person{瓦伦廷}和\Person{巴尔德撒纳斯}承认祂的诞生，但否认祂取了我们的本性，并声称天主之子使用童贞女如同一个单纯的管道。\par

\Person{撒伯里乌}，利比亚人；\Person{佛提努}，加辣提亚人\Person{马塞鲁}；以及\Person{撒摩沙塔的保禄}，说一个纯粹的人由童贞女所生，但公开否认永恒的\Person{基督}是天主。\par

\Person{亚略}和\Person{欧诺米}主张天主圣言只从童贞女取了身体。\par

\Person{阿波林}在身体之外加上了无理性的灵魂，仿佛天主圣言的降生成人不是为有理性的存在，而是为无理性的存在；而宗徒的教导是：完全的天主取了完全的人，正如经上所证：\Quote{祂虽具有天主的形体，却取了奴仆的形体；}因为\OriginalTerm{形体}{form}代替\OriginalTerm{本性}{nature}和\OriginalTerm{实体}{substance}，表明拥有天主的本性，祂取了奴仆的本性。\par

因此，当我们与最早发明不敬之事的\Person{马西翁}、\Person{玛乃斯}和\Person{瓦伦廷}辩论时，我们努力从神圣的圣经中证明主\Person{基督}不仅是天主，也是人。\par

然而，当我们向无知者证明\Person{亚略}、\Person{欧诺米}和\Person{阿波林}关于\OriginalTerm{救恩计划}{œconomy}的教导不完整时，我们藉着圣神的神圣启示表明所取的本性是完美的。\par

\Person{撒伯里乌}、\Person{佛提努}、\Person{马塞鲁}和\Person{撒摩沙塔的保禄}的不敬之处，我们通过神圣圣经的证据来驳斥，证明主\Person{基督}不仅是人，而且是永恒的天主，与圣父同一性体。祂取了有理性的灵魂，这由主自己所说的话表明：\Quote{现在我心神烦乱，我该说什么呢？父啊！救我脱离这时辰吧；但我正是为此才来到这时辰。}又说：\Quote{我的心神忧闷得要死。}在另一处：\Quote{我有权舍掉我的灵魂，也有权再取回它；没有人能夺去我的灵魂。} \BibleRef{若 10:18} 天使对\Person{若瑟}说：\Quote{起来，带着婴孩和祂的母亲，到以色列地去；因为那些想要谋害婴孩性命的人已经死了。}福音作者说：\Quote{耶稣在智慧和身量上，在天主和人前的喜爱上，日益增长。}现在，在身量和智慧上增长的，不是永远完美的天主性，而是那在时间中形成、成长并达于完美的人性。\par

因此，主基督的一切人性特征——即饥饿、口渴、疲倦、睡眠、恐惧、出汗、祈祷、无知等——我们都确认是属于我们的人性，这人性是天主圣言为完成我们的救恩而取来并与自己结合的。但是，使残废者恢复行动、使死人复活、供应饼食以及所有其他奇迹，我们相信是神性能力的工作。在这个意义上，我说同一主基督既受苦又消除痛苦；即，就可见的一面而言，祂受苦；就不可言说地内住的神性而言，祂消除痛苦。这一点由圣福音作者的记载所不容置疑地证明，我们从中学到：当祂躺在马槽中、裹着襁褓时，被一颗星所宣告，被贤士朝拜，被天使歌颂。如此，我们虔敬地辨别：襁褓、没有床榻以及一切的贫穷属于人性；而贤士的旅程、星星的引导和天使的陪伴则宣告不可见者的神性。同样地，祂逃往\Place{埃及}，以逃避\Person{黑落德}的愤怒，因为祂是真人；但正如先知所说，\BibleQuote{祂震动埃及的偶像}，因为祂本来就是天主。祂受割损；祂遵守法律；并献上取洁的祭品，因为祂出自叶瑟的根。而且，作为人，祂在法律之下；后来却废除了法律并赐予新盟约，因为祂是立法者，并曾通过先知们应许祂自己要赐予它。祂受\Person{若翰}的洗礼；这显示出祂分担了我们的一切。祂由至高之处的圣父所见证，并由圣神所指明；这宣告祂是永恒的。祂饥饿；但祂用五个饼喂饱了数千人；后者是神圣的，前者是人性的。祂口渴，并要水喝；但祂是生命之泉；前者出于祂人性的软弱，后者出于祂神性的大能。祂在船上睡着了，却使海上的风暴止息；前者出于祂的人性，后者出于祂那赋予一切事物存在的能力和创造之力。祂走路疲倦；却治愈了跛子，并从坟墓中复活死人；前者出于人性的软弱，后者出于超越此世的能力。祂害怕死亡，又消灭了死亡；前者显示祂是会死的，后者显示祂是不死的，或者更确切地说，是生命的赐予者。\BibleQuote{祂被钉在十字架上}，如\Person{保禄}所说，\BibleQuote{因软弱}。但同一\Person{保禄}说：\BibleQuote{然而祂因天主的德能而活着}。让\Quote{软弱}这个词教导我们，祂并非作为全能者、\OriginalTerm{无界限、不变、不易}{the Uncircumscribed, the Immutable and Invariable}而被钉在木头上；而是那被天主的德能所赋予生命的本性，按照宗徒的教导，曾死过并被埋葬，死亡和埋葬都属于\OriginalTerm{形式的仆人}{form of the servant}的特征。\BibleQuote{祂打破了铜门，斩断了铁闩}，并消灭了死亡的权势，三天之内复活了祂自己的圣殿。这些是\OriginalTerm{形式的天主}{form of God}的证据，与主的话一致：\BibleQuote{你们拆毁这座圣殿，三天之内我要把它重建起来}。因此，在同一个基督内，通过苦难我们默观人性，通过奇迹我们领悟神性。我们不把两个本性分成两个基督，并且我们知道天主圣言是由圣父所生，而我们的人性是从\Person{亚巴郎}和\Person{达味}的后裔所取。因此蒙福的\Person{保禄}在谈论\Person{亚巴郎}时说：\BibleQuote{并不说和许多子孙，而是说和一个子孙，那就是基督}。他写给\Person{弟茂德}的书信中说：\BibleQuote{你当记住耶稣基督，达味的后裔，依据我的福音，从死人中复活了}。他写给罗马人的书信中说：\BibleQuote{关于祂的儿子耶稣基督……按肉身说是从达味的后裔生的}。又说：\BibleQuote{他们是祖先，基督按肉身说也是从他们而来的}。圣史写道：\BibleQuote{耶稣基督，达味之子，亚巴郎之子的族谱}。蒙福的\Person{伯多禄}在\WorkTitle{宗徒大事录}中说\Person{达味}：\BibleQuote{既是先知，也知道天主曾向他起誓，要从他的后裔中立一位坐在他的宝座上；他预先看见了，就讲论基督的复活}。天主对\Person{亚巴郎}说：\BibleQuote{地上万民都要因你的后裔获得祝福}。\Person{依撒意亚}说：\BibleQuote{从叶瑟的树干将生出一根枝条，从他的根上将长出一个嫩芽；上主的神，智慧和聪敏的神，超见和刚毅的神，明达和孝爱之神，以及敬畏上主的神将住在他内}。稍后又说：\BibleQuote{到那一天，叶瑟的根将成为万民的旗帜；外邦人要寻求它，祂安息之所将充满光荣}。\par

從這些引文可以清楚地看出：按肉身說，基督是亞巴郎和達味的後裔，與他們有同一本性；按天主性說，祂是永恆的聖子和聖言，不可言喻地、超乎人類方式地由聖父所生，並作為光輝和印跡及聖言與聖父同永恆。因為正如聖言與理智、光輝與光不可分離地相連一樣，獨生子與祂自己的聖父也是如此。所以我們斷言，我們的主耶穌基督是天主的獨生子和首生子；在降生成人之前和之後都是獨生子，但在由童貞女誕生後才成為首生子。因為\Quote{首生子}這個名稱在某種意義上似乎與\Quote{獨生子}相對：由某人所生的獨生子被稱為\Quote{獨生子}，而眾兄弟中的最年長者被稱為\Quote{首生子}。神聖的聖經宣稱唯獨天主聖言是由聖父所生；但獨生子藉著由童貞女攝取我們的本性，並屈尊稱那些信靠祂的人為兄弟，也成為首生子；因此，同一位基督，作為天主是獨生子，作為人則是首生子。這樣，我們承認兩個性體，欽崇同一位基督，並向祂獻上同一的欽崇，因為我們相信結合從在聖母聖胎中成孕的那一刻就開始了。因此，我們也稱神聖的童貞女為天主之母和人之母，因為主基督自己在神聖的聖經中被稱為天主和人。\Quote{厄瑪奴耳}這個名字宣告了兩個性體的結合。如果我們承認基督既是天主又是人，並如此稱呼祂，那麼誰會如此缺乏感覺，以至於不敢在使用\Quote{天主之母}的同時也使用\Quote{人之母}這個詞呢？因為我們對主基督使用這兩個稱謂。為此，童貞女受到尊敬並被稱為\Quote{充滿恩寵的}。那麼，有什麼有理智的人會反對按照救主的稱號來稱呼童貞女呢？因為她因祂的緣故受到信徒的尊敬。從她所生的那一位並不是因她而被敬拜，而是她因從她所生的那一位而受到最高的尊崇。\par

假設基督只是天主，並且祂的存在起源於童貞女，那麼就稱呼童貞女只為\Quote{天主之母}，因為她生育了一位本性為天主的存在者。但如果基督既是天主又是人，並且祂從永遠就是天主（因為祂並非開始存在，而是與生祂的聖父同永恆），而在這些最後的日子裡，祂按人性誕生為人，那麼，願意從兩個方向界定教義的人，就應當為童貞女制定稱謂，並解釋哪一個適合本性，哪一個適合結合。但如果有人想要發表讚詞、創作聖詩、重複讚美，並且自然希望使用最尊貴的名號，那麼，不應像前面那樣制定教義，而是以修辭性的讚頌，表達對奧秘大能的一切欽佩，讓他滿足心中的願望，使用崇高的名號，讓他讚美，讓他驚嘆。在正統教師的著作中可以找到許多這類例子。但在所有場合都應尊重節制。願一切讚美歸於那位說\Quote{節制是最好的}的人，儘管他不屬於我們的羊群。\par

這是教會信仰的宣認；這是福音作者和宗徒們所教導的教義。為了這個信仰，靠天主的恩寵，我甘願死多次。我們努力將這個信仰傳授給那些現在迷失的人，一再向他們挑戰討論，渴望向他們顯示真理，但沒有成功。由於懷疑自己可能被明顯駁斥，他們迴避了交鋒；因為虛假確實是腐敗的，是黑暗的伴侶。經上記載：\Quote{凡作惡的，都不到光明裡來，免得他的行為被暴露}在光明中。\par

因此，經過多次努力，我未能說服他們認識真理，便回到自己的教會，同時充滿憂傷和喜樂：喜樂是因為我自己脫離了謬誤，憂傷是因為我的肢體不健康。所以我懇求你盡全力向我們慈愛的主祈禱，向祂呼求說：\Quote{主啊，求祢憐憫祢的百姓，不要使祢的產業受辱}。\InnerQuote{求祢牧養我們，主啊，使我們不再像起初那樣，當祢沒有治理我們，祢的名也沒有被呼求幫助我們的時候。}\Quote{我們成了鄰居的羞辱，四周之人的譏誚和嘲笑}，因為邪惡的教義進入了祢的產業。他們褻瀆了祢的聖殿，因為外邦人的女子因我們的苦難而歡樂。不久之前我們還同心同語，現在卻分裂成許多種語言。但是，主我們的天主，求祢賜給我們平安，這是我們因輕視祢的誡命而失去的。主啊，我們不認識別的祢。我們呼求祢的名。\Quote{使兩者合一，拆毀中間隔斷的牆}，即興起的罪惡。求祢將我們一個一個聚集起來，成為祢的新以色列，重建耶路撒冷，聚集以色列被驅散的人。願我們再次成為一棧，都受祢的餵養；因為祢是善牧，\Quote{為羊捨命}。\Quote{醒來吧，主啊，祢為什麼睡覺？起來，不要永遠丟棄我們。}斥責風和海；賜給祢的教會平靜和安全，脫離波浪。\par

这些话语以及类似的话语，我恳求你对万有的天主述说；因为祂是良善的、充满慈爱的，并且永远成全敬畏祂之人的意愿。因此，祂必垂听你的祈祷，驱散这比\Place{埃及}的灾祸更深沉的黑暗。祂必赐给你祂自己的爱之平安，聚集分散的人，欢迎被驱逐的人。那时，在义人的帐幕中必听到\Quote{欢呼与救恩的声音}。那时，我们要向祂呼求：\Quote{我们因祢使我们受苦的日子，和我们遭难的年岁，而欢乐。}而你，当你的祈祷蒙应允时，将用这些话赞美祂：\Quote{天主是应当称颂的，祂没有转面不听我的祈祷，也没有从祂那里收回祂的仁慈。}\par

% structure: p0022 source=h2 target=subsection reason=relative-heading-rank; base=h2

\subsection{证明降生成人之后，我们的主耶稣基督是一个圣子。}

编造诽谤我的人指控我把唯一的主耶稣基督分成两个儿子。但我绝无此意，反而认为所有胆敢如此说的人都是不虔敬的。因为我受了神圣圣经的教导，要钦崇一个圣子，即我们的主耶稣基督，天主的独生子，降生成人的天主圣言。因为我们宣认同一位基督既是永恒的天主，又在末日为了人的救恩而成为人；但成为人并非由于神性的改变，而是由于人性的摄取。因为神性的本性是不变不易的，正如在万世之前生了祂的圣父的本性一样。凡关于圣父本质所能理解的，也必完全存在于独生子的本质中；因为祂是由那本质所生的。我们的主教导了这一点，当祂对\Person{斐理伯}说\Quote{看见了我，就是看见了父}；又在另一处说\Quote{凡父所有的，都是我的}；此外还说\Quote{我与父原是一体}，还有许多别的章节可以引述，表明本质的同一性。\par

由此可知，祂并没有成为天主：祂本是天主。\Quote{在起初已有圣言，圣言与天主同在，圣言就是天主。}祂并不是人：祂成了人，祂这样成为人，是由于取了我们的本性。正如蒙福的\Person{保禄}所说：\Quote{祂虽具有天主的形体，却不以自己与天主同等为僭越，反而空虚自己，取了奴仆的形体。}又说：\Quote{祂没有取得天使的本性，而是取得了\Person{亚巴郎}的后裔。}又说：既然孩子们都同有血肉，祂也照样取了同样的血肉。因此祂既是可受苦的，又是不可受苦的；既是可朽的，又是不朽的；作为人，祂是可受苦的、可朽的；作为天主，祂是不可受苦的、不朽的。作为天主，祂复活了自己已死的肉体——正如祂自己的话所宣告的：\Quote{你们拆毁这座圣殿，三天之内我要把它重建起来。}作为人，直到受难时，祂是可受苦的、可朽的。因为复活之后，即使作为人，祂也是不可受苦的、不朽的、不腐坏的；祂发出神圣的闪电，但并非按肉体被改变成神性的本性，而是仍然保存着人性的特征。祂的身体也并非无限的，因为这只是神性本性的特性，而是保持在它原有的界限之内。这一点，祂在复活后对门徒所说的话中教导我们：\Quote{你们看我的手和脚，实在是我；你们摸摸我，看看！因为鬼神没有肉和骨头，像我这样你们是有的。}就在这样被看见的时候，祂升了天；祂应许要这样再来，这样被相信祂的人和钉死祂的人看见，因为经上记载：\Quote{他们要仰望他们所刺透的那一位。}因此我们钦崇圣子，但在祂身上我们默观两种本性的完美，包括那取了和那被取的本性，一种来自天主，一种来自\Person{达味}。为此，祂既被称为生活天主之子，又被称为\Person{达味}之子；每一种本性都得到其应有的称号。因此，神圣圣经称祂为天主和人，而蒙福的\Person{保禄}高呼：\Quote{只有一个天主，在天主与人之间的中保也只有一人，就是降生成人的基督耶稣；祂奉献了自己，为众人作赎价。}但在这里被称为人的那一位，他又在另一处描述为天主，因为他说：\Quote{期待所希望的幸福，和伟大的天主及我们的救主耶稣基督的光荣显现。}而在另一处，他同时使用了两个名称：\Quote{按血统说，基督也是从他们来的；祂是在万有之上，永远受赞美的天主。阿们。}\par

这样，他说明了同一位基督，按肉体说是从犹太人来的，按天主性说是在万有之上的天主。同样，先知\Person{依撒意亚}写道：\Quote{祂是一个受苦的人，熟悉病痛……祂实在承担了我们的疾病，背负了我们的痛苦。}不久之后他说：\Quote{谁能述说祂的世代？}这不是指人说的，而是指天主。同样，天主通过\Person{米该亚}说：\Quote{你，\Place{白冷}，\Place{犹大}地啊，你在\Place{犹大}的首领中并不是最小的，因为将有首领从你那里出来，牧养我的百姓以色列；祂的根源从亘古、从太初就有。}说\Quote{从你那里将出来一位首领}，显明了降生成人的\OriginalTerm{经济}{œconomy}；补充说\Quote{祂的根源从亘古、从太初就有}，则宣告了在万世之前由圣父所生的天主性。\par

既然我们受了神圣圣经这样的教导，又发现历代在教会中杰出的教师们也有同样的意见，我们就尽力保持我们的传承不受侵犯；钦崇一个圣子、一个天主圣父、一个圣神；同时承认肉体与神性之间的区别。正如我们断言那些把我们的主耶稣基督分成两个儿子的人，偏离了圣宗徒所走过的道路，同样我们宣告那些主张独生子的神性与人性成为一体的人，是陷入了相反的山谷。这些道理我们持守；这些道理我们宣讲；为了这些道理我们战斗。\par

诽谤者诬蔑我敬拜两个儿子的毁谤已被明显的事实所驳斥。我教导所有前来领受圣洗圣事的人相信尼西亚所颁布的信德；当我举行\OriginalTerm{重生的圣事}{sacrament of regeneration}时，我为那些宣认信德的人施洗，因圣父、圣子及圣神之名，逐一念诵每位称号。当我在教会中举行礼仪时，我习惯将光荣归于圣父、圣子及圣神；不是众子，乃是子。那么，如果我坚持两个儿子，哪一个被我光荣，哪一个受冷落？因为我还不至于糊涂到承认两个儿子而让其中一个不受尊崇。由此可见，诽谤已显为是诽谤——因为我只敬拜一位独生子，即降生成人的天主圣言。我称童贞玛利亚为\Quote{天主之母}，因为她生了厄玛努尔，即\Quote{天主与我们同在}。但预言厄玛努尔的先知稍后论及他写道：\BibleQuote{有一婴孩为我们诞生，有一子赐给了我们；权柄必担在他的肩头上；他的名字称为大谋士的天使、奇妙的、策士、全能的天主、大能者、和平的君王、永世之父。}既然从童贞女所生的婴孩被称为\Quote{全能的天主}，那么母亲被称为\Quote{天主之母}就是合理的。因为母亲分享子嗣的尊荣，童贞玛利亚既是主基督为人的母亲，又是祂为主、为造物主、为天主的婢女。\par

由于这一称谓的差异，神圣的保禄说他\Quote{无父、无母、无族谱，既无生之始，亦无命之终}。就人性而言，他是无父的；因为作为人，他独由母亲而生。就天主性而言，他是无母的；因为他从永远独由圣父所生。再者，作为天主，他是无族谱的；而作为人，他则有族谱。因为经上记载：\Quote{耶稣基督的族谱，他是达味之子，亚巴郎之子。}神圣的路加也记载了他的族谱。同样，作为天主，他无生之始，因为他生于诸世之前；亦无命之终，因为他的本性是不朽且\OriginalTerm{不受苦难}{impassible}的。但作为人，他既有生之始——生于奥古斯都凯撒年间，也有命之终——死于提庇留凯撒年间。然而，如今正如我已说过的，他的人性也已不朽；并且，正如他升天，他还要再来，正如天使所说：\BibleQuote{这离开你们被接到天上去的耶稣，你们见他怎样往天上去，他也要怎样来临。}\par

这便是神圣的先知传给我们的道理；这便是众圣宗徒的道理；这便是东方与西方伟大圣人的道理：著名的依纳爵，他由伟大的伯多禄的右手接受了\OriginalTerm{大司祭职}{archpriesthood}，并为基督的宣认而被野兽吞噬；还有伟大的尤斯塔提乌斯，他曾主持大公会议，因对真道的热诚而被放逐。这道理由杰出的默肋提乌斯宣讲，他同样付出不少辛劳，为宗徒的教义三次被逐离羊群；由弗拉维雅努斯——帝国京城的荣耀；由令人钦佩的厄弗辣因——天主恩宠的器皿，他用叙利亚文为我们留下了美好的笔文化遗产；由西彼廉——迦太基及整个利比亚的杰出牧者，他为基督的缘故在火中殉道；由伟大的罗马主教达玛苏，以及米兰的荣耀安波罗修，他们用罗马的语言宣讲并书写了这道理。\par

同样的道理由亚历山大城的伟大明灯亚历山大和亚大纳削所教导，他们同心同德，经历了举世闻名的苦难。这是帝国京城伟大教师们牧养羊群的草场：额我略——真理光辉的支持者；若望——世界的导师；阿提库斯——他们在教座和思想上的继承者。凭借这些道理，巴西略——真理的伟大明灯、同父母的额我略，以及从他手中领受大司祭职的盎非罗基，教导了同时代人，并在著作中留给我们作为美好遗产。要述说波利卡普、依勒内、美多迪乌斯、希波利特以及教会其他教师，我时间不够了。总之，我坚持我既跟随\OriginalTerm{神圣的圣言}{divine oracles}，也跟随所有这些圣人。借着圣神的恩宠，他们深入了天主默感的圣经深处，既自己领悟其内涵，并使所有愿意学习的人明白。言语的不同并未造成教义的差异，因为他们是同一源泉所流出的圣神恩宠的管道。\par

\SourceRecord{Translated by Blomfield Jackson.；Nicene and Post-Nicene Fathers, Second Series；Vol. 3.；Edited by Philip Schaff and Henry Wace.；Buffalo, NY: Christian Literature Publishing Co.,；1892.；<http://www.newadvent.org/fathers/2707151.htm>.}

% structure: p0001 source=h1 target=section reason=page-title

\section{第一五二封信}

\Emph{东方主教们向皇帝的报告，说明他们的经过，并解释安提约基雅主教迟迟未到的原因。}\par

我们遵照陛下虔诚来信的命令，已前往厄弗所大都会。在那里，我们发现教会事务混乱，且被内争所扰。其原因在于：亚历山大的济利禄和厄弗所的默农联合起来，召集了大群乡民，并禁止了五旬节大庆的庆祝，以及晚间和早晨的日课。\par

他们关闭了圣堂和殉道圣所；他们与受他们欺骗的人另行集会；他们行了无数不义之事，将圣教父们的\OriginalTerm{教规}{canons}和陛下的谕令都踩在脚下。而且，这些行动是在最尊贵的伯爵、爱基督的陛下之使者坎狄狄安，以书面和口头方式下达命令后进行的，命令大公会议必须等待来自帝国各地的至圣主教们到达，然后才按照陛下的命令正式召开。此外，亚历山大的济利禄在他们会议召开前两天曾写信给我，安提约基雅主教，说整个大公会议都在等我到达。因此，我们罢免了上述两人，济利禄和默农，并禁止他们参与教会的一切礼仪。其余参与他们不义之人，我们处以绝罚，直到他们弃绝并诅咒济利禄所颁布的、充满\OriginalTerm{欧诺弥异端}{Eunomian heresy}和\OriginalTerm{亚略异端}{Arian heresy}的\OriginalTerm{章句}{Chapters}，并服从陛下的命令，与我们一同集会，并与我们一起有条不紊且精确地审查所争议的问题，确认圣教父们的虔诚教义。\par

至于我本人迟迟未到之事，请陛下明鉴：因考虑到陆路路途遥远（我们正是走陆路），我已非常快速地赶来，连续走了四十个驿站，途中未作停留；爱基督的陛下可从沿途城镇的居民那里得知。此外，我还在安提约基雅被当地的饥荒、民众每日的骚动以及异常的雨季耽搁了多日，雨季导致山洪暴涨，威胁城镇安全。\par

\SourceRecord{Translated by Blomfield Jackson.；Nicene and Post-Nicene Fathers, Second Series；Vol. 3.；Edited by Philip Schaff and Henry Wace.；Buffalo, NY: Christian Literature Publishing Co.,；1892.；<http://www.newadvent.org/fathers/2707152.htm>.}

% structure: p0001 source=h1 target=section reason=page-title

\section{第一五三封信}

\Emph{同一人致皇后\Person{普尔凯利亚}与\Person{欧多克西亚}之报告。}\par

我们原希望能以不同的言词向你们虔敬的陛下们报告，但如今我们被迫向你们陈述以下事实，因为我们受到亚历山大城的\Person{济利禄}和厄弗所的\Person{默农}暴虐专权的逼迫。按照教会的法律及你们虔敬的陛下们的命令，应当的做法是等候在途中的敬虔主教们到来，与他们一同查考关于真信仰的争议问题，研讨所提出的论点，并在仔细调查后确认宗徒的教义。他们曾写信给我说会等候我们的到来。他们听说我们距此仅三天路程。然后他们独自召集了一次不合法的会议，竟敢采取不义、非法且荒谬绝伦的行动。他们这样做，尽管由你们虔敬的基督徒陛下们为维持良好秩序所派遣的最尊贵的伯爵\Person{康狄迪亚努斯}曾以书面和口头明确告诫他们等候已召集的敬虔主教们到来，不得对真信仰作任何革新，而应坚持我们敬虔君主的指示。但他们听到了皇帝的谕旨和最尊贵的伯爵\Person{康狄迪亚努斯}的劝告后，仍然无视应有的秩序。正如先知所说：\BibleQuote{他们孵毒蛇蛋，织蜘蛛网；谁要吃它们的蛋，打破时发现腐烂，里面有一条毒蛇。}因此我们满怀信心地呼喊：\BibleQuote{他们的网不能成为衣服，他们也不能用他们的作品遮蔽自己。}\par

他们封闭了教堂和殉道者圣堂；禁止庆祝神圣的五旬节；此外，他们还派遣他们暴虐统治的爪牙进入主教们的私宅，发出骇人的威胁，强迫他们在非法文件上签名。因此，我们鉴于他们所有荒谬的行为，已将上述\Person{济利禄}和\Person{默农}撤职，剥夺他们的主教职务。至于他们那些不法的同伙，无论是出于阿谀还是恐惧，我们均予以绝罚，直到他们认识到自己的创伤，诚心悔改，诅咒\Person{济利禄}那些沾染了\Person{阿波林}、\Person{亚略}和\Person{欧诺弥}异端邪说的异端章节，重获在\Place{尼西亚}大公会议上诸教父的信德，并服从我们基督徒君主虔敬的命令，和平地、毫无骚乱地集合召开会议，愿意仔细审查提交给他们的议题，并诚实地保护福音信仰的纯洁。\par

\SourceRecord{Translated by Blomfield Jackson.；Nicene and Post-Nicene Fathers, Second Series；Vol. 3.；Edited by Philip Schaff and Henry Wace.；Buffalo, NY: Christian Literature Publishing Co.,；1892.；<http://www.newadvent.org/fathers/2707153.htm>.}

% structure: p0001 source=h1 target=section reason=page-title

\section{书信154}

\Emph{同一位给\Place{君士坦丁堡}元老院的报告。}\par

\SourceRecord{Translated by Blomfield Jackson.；Nicene and Post-Nicene Fathers, Second Series；Vol. 3.；Edited by Philip Schaff and Henry Wace.；Buffalo, NY: Christian Literature Publishing Co.,；1892.；<http://www.newadvent.org/fathers/2707154.htm>.}

% structure: p0001 source=h1 target=section reason=page-title

\section{第一五五封信}

\Emph{安提约基雅主教\Person{若望}及其支持者致君士坦丁堡圣职人员书}\par

\SourceRecord{Translated by Blomfield Jackson.；Nicene and Post-Nicene Fathers, Second Series；Vol. 3.；Edited by Philip Schaff and Henry Wace.；Buffalo, NY: Christian Literature Publishing Co.,；1892.；<http://www.newadvent.org/fathers/2707155.htm>.}

% structure: p0001 source=h1 target=section reason=page-title

\section{\WorkTitle{第一五六封书信}}

\Emph{同一人致\Place{君士坦丁堡}人民的信。}\par

\SourceRecord{Translated by Blomfield Jackson.；Nicene and Post-Nicene Fathers, Second Series；Vol. 3.；Edited by Philip Schaff and Henry Wace.；Buffalo, NY: Christian Literature Publishing Co.,；1892.；<http://www.newadvent.org/fathers/2707156.htm>.}

% structure: p0001 source=h1 target=section reason=page-title

\section{第一五七函}

\Emph{东方（主教）会议向胜利的皇帝呈报二次宣布开除济利禄和默默隆的报告。}\par

陛下为帝国和天主教会之益而照耀的虔敬，曾命我等聚集于厄弗所，旨在为教会带来和平与裨益，而非制造混乱与纷扰。陛下的谕旨清晰明确地表明了对基督的教会之虔敬与和平的意愿。然而，亚历山大的济利禄，一个似乎生来就为教会之祸害的人，与厄弗所的默默隆的胆大妄为结伙后，首先违反了陛下安定与虔敬的谕令，并显出其普遍的败坏。陛下曾下令就信仰问题进行调查与审慎考查，且须经众人同意与和谐。济利禄受到挑战，更确切地说，因他最近送往帝都的书信及其所附绝罚条款而自行定罪，其中他被证明与不敬神且属异端的阿波利纳里持有相同观点，却不顾此情势，仿佛我们生活在无皇帝统治之下，而走向各种无法无天。他本应为自己关于主耶稣基督的不健全意见负责，但他篡取了既非法典亦非陛下谕令所授予的权柄，贸然陷入种种混乱与非法行为。\par

因上述理由，拒绝接受他对信仰之损害伎俩的圣主教会议，基于上述原因，将他罢免。它亦罢免了默默隆，此人自始至终是他的谋士与同谋，不断煽动反对极其神圣的主教们，因他们拒绝赞同其有害的异端学说；他关闭了教堂及所有祈祷之所，仿佛我们生活在异教徒和天主的敌人当中；他引入了厄弗所的暴民，以致我们每日处于极大危险之中，我们并不顾及防卫，而只关心真宗教的正确教义。因为清除这些人就等于确立正统信仰。\par

从其亲自撰写的章节中，陛下不难看出他不敬神的心志。他被证实，可以说，试图将死于异端的不敬神的阿波利纳里从阴府中召上来，并攻击教会与正统信仰。他的出版书显示，他同时诅咒了福音作者、宗徒以及作为教会先祖的继承者，他们不是出于自己的想象，而是藉圣神感动，宣讲了真信仰，传报了福音；这一信仰和福音确实与这个人所持守和教导的相反，而他通过灌输这些，希望使自己的私人邪恶主宰世界。由于我们无法容忍此点，我们遵循了正当程序，同时倚靠天主的恩宠与陛下的善意。\par

我们知道，陛下对任何事物都不比对神圣信仰更为尊崇，陛下与您极其有福的先祖皆在此信仰中成长。从他们那里，您继承了永久的帝国权杖，不断镇压反对宗徒教义者。上述的济利禄正是这样一个反对者，他借助默默隆之助，如攻占堡垒般占据了厄弗所，并理所当然地与他的盟友一同承受罢免判决。理所当然：因为除上述一切外，他们胆敢以各种袭击与暴力手段对待我们，而我等为遵从陛下谕旨聚集公议会，已不顾家、国与自身的一切要求。\par

除非陛下之虔敬干预，命我等在附近另一地点聚集，以便能根据圣经和教父著作，无可辩驳地驳斥济利禄及其伎俩的受害者，否则我们如今正遭受暴政之害。我们已仁慈地将这些人逐出共融，并给予他们若悔改便有希望得救的提示；然而，仿佛置身于未开化士兵的战役中，他们至今仍为其非法行为提供手段。有些人是早已被罢免，又被济利禄恢复职位；有些人是被自己的都主教绝罚，又被济利禄重新接纳共融；其他人则因各种指控而被处以刑罚，却被他提升至尊荣。自始至终，他行动的主要动机是试图凭借人数众多来实现其异端目的，因为他并未正确认识到，在关乎真宗教之事上，所需要的不是人数，而是教义的正确性和宗徒教义的真理。需要的是那些不仅以胆大妄为和专横自恃，而是藉虔诚地运用宗徒的证言与榜样来确立这些要点的人。\par

基于所有这些理由，我们恳求并呼求陛下速来救助遭受攻击的真理，并不延迟地纠正这些人的专横鲁莽；因为它如飓风般将我等中较不节制者席卷进入有害的异端。陛下的虔敬曾关心波斯及蛮族中的教会；理当不忽视那些在罗马帝国边界内被风暴颠簸的教会。\par

\SourceRecord{Translated by Blomfield Jackson.；Nicene and Post-Nicene Fathers, Second Series；Vol. 3.；Edited by Philip Schaff and Henry Wace.；Buffalo, NY: Christian Literature Publishing Co.,；1892.；<http://www.newadvent.org/fathers/2707157.htm>.}

% structure: p0001 source=h1 target=section reason=page-title

\section{Letter 158}

\Emph{东方（主教们）向最虔诚的皇帝呈递的报告，该报告与前述呈递给最尊贵的伯爵\Person{依肋内}的报告一同递交。}\par

在接到陛下虔诚的书信时，我们曾希望最近打击天主教会的那场埃及风暴会被驱散。但我们失望了。那些人因疯狂而变得更加胆大妄为；他们毫不顾及公正且适时地对他们作出的撤职权判决，也未曾因陛下的斥责而变得稍加节制。他们践踏了陛下虔诚的法律和圣教父们的教规，其中一些人已被撤职，另一些人已被绝罚，却仍在祈祷堂里举行庆典、共融。我们，正如我们已经告知陛下爱基督的威严那样，在收到陛下仁慈的信函后，虽只愿在宗徒教堂内祈祷，却不仅被阻止，而且实际上被投石追赶了相当远的一段距离，以至于我们不得不全速逃跑以求安全。相反，我们的对手认为他们可以为所欲为。他们拒绝调查争议问题，拒绝为\Person{济利禄}的异端章节辩护，拒绝承认其中包含的不敬神的明显证据。他们只是由于无耻而变得厚颜无耻，而对当前问题的检查需要的不是无耻，而是冷静、知识和在教义问题上的技巧。\par

在这种情况下，我们不得不派遣最尊贵的伯爵\Person{依肋内}先行一步，前去觐见陛下，并说明事态的状况。他对所发生的一切有着准确的信息，并从我们这里学到了许多补救方法，通过这些方法或许有可能使天主圣教会恢复平静。我们恳求陛下仁慈地耐心听他陈词，并下令迅速执行任何陛下认为有益的措施，以免我们在此被压垮而无法忍受。\par

\SourceRecord{Translated by Blomfield Jackson.；Nicene and Post-Nicene Fathers, Second Series；Vol. 3.；Edited by Philip Schaff and Henry Wace.；Buffalo, NY: Christian Literature Publishing Co.,；1892.；<http://www.newadvent.org/fathers/2707158.htm>.}

% structure: p0001 source=h1 target=section reason=page-title

\section{第一五九书}

\Emph{同一人致长官与院长书。}\par

\SourceRecord{Translated by Blomfield Jackson.；Nicene and Post-Nicene Fathers, Second Series；Vol. 3.；Edited by Philip Schaff and Henry Wace.；Buffalo, NY: Christian Literature Publishing Co.,；1892.；<http://www.newadvent.org/fathers/2707159.htm>.}

% structure: p0001 source=h1 target=section reason=page-title

\section{第一六〇书}

\Emph{同一人致\Person{总督}与\Person{斯科拉斯提库斯}书。}\par

\SourceRecord{Translated by Blomfield Jackson.；Nicene and Post-Nicene Fathers, Second Series；Vol. 3.；Edited by Philip Schaff and Henry Wace.；Buffalo, NY: Christian Literature Publishing Co.,；1892.；<http://www.newadvent.org/fathers/2707160.htm>.}

% structure: p0001 source=h1 target=section reason=page-title

\section{\WorkTitle{第一六一封书信}}

\Emph{由\Place{安提约基雅}总主教\Person{若望}及其支持者通过\Person{帕拉狄乌斯·玛基斯特里亚努斯}呈递给皇帝的陈情书。}\par

\SourceRecord{Translated by Blomfield Jackson.；Nicene and Post-Nicene Fathers, Second Series；Vol. 3.；Edited by Philip Schaff and Henry Wace.；Buffalo, NY: Christian Literature Publishing Co.,；1892.；<http://www.newadvent.org/fathers/2707161.htm>.}

% structure: p0001 source=h1 target=section reason=page-title

\section{书信一百六十二}

\Emph{\Person{狄奥多勒}致\Place{萨莫萨塔}主教\Person{安德肋}的书信，写自\Place{厄弗所}。}\par

我自厄弗所致书问候你的圣德，我因你的病弱而祝贺你，认为你蒙天主珍爱，因为你通过报告而非亲身经验知道了这里发生的恶行。真是邪恶！它们超越所有想象和历史事件；它们迫使眼泪不断流淌。教会的身体有被肢解的危险；——不，我可以说它已经被切开了第一刀；——除非智慧的医生恢复并重新连接那些不健康和断裂的肢体。埃及人再次对天主发怒，与祂的仆人梅瑟和亚郎争战，以色列的大部分人站在敌人一边；因为愿意为真宗教受苦的健康之人实在太少。古老的原则被践踏。被撤职的人履行司铎职责，而撤职他们的人却坐在家中叹息。被相同判决逐出教会的人自行解除了被撤职者的撤职。这就是由埃及人、巴勒斯坦人、来自本都和亚细亚教区的人以及与他们为伴的西方人所举行的会议的笑话。\par

在异教时代，有什么哑剧演员，甚至任何闹剧，如此嘲笑过宗教？实际上，有什么闹剧作家曾上演过这样的戏剧？有什么剧作家写过如此悲惨的悲剧？如此巨大而众多的烦恼困扰着天主的教会，我只叙述了其中极小一部分。\par

\SourceRecord{Translated by Blomfield Jackson.；Nicene and Post-Nicene Fathers, Second Series；Vol. 3.；Edited by Philip Schaff and Henry Wace.；Buffalo, NY: Christian Literature Publishing Co.,；1892.；<http://www.newadvent.org/fathers/2707162.htm>.}

% structure: p0001 source=h1 target=section reason=page-title

\section{\WorkTitle{书信163}}

\Emph{东方专员的第一封信，寄往加采东，其中包含\Person{德奥多勒}。}\par

我们抵达加采东时（因为我们自己和我们的对手都未获准进入君士坦丁堡，由于卓越的隐修士们的骚乱），我们听说，在我们出现之前八天（看哪，至虔敬君王的荣耀！），\Person{聂斯多略}阁下已被从厄弗所解职，自由地去往他愿意去的地方；对此我们深感忧虑，因为非法且非正式的行为如今似乎有了某种效力。然而，请尊驾确信，我们将热切地为信德而战，并愿意战斗至死。今天，果尔比亚月十一日，我们期待我们至虔敬的皇帝前往鲁菲尼亚努姆，在那里审理案件。\par

因此我们恳求尊驾向主基督祈祷，帮助我们能够确认圣教父们的信德，并将这些损害教会的章节连根拔除。我们恳求尊驾与我们同心同德，并持守你对正统信德随时准备的热忱。写此信时，\Person{希梅里乌斯}阁下尚未与我们见面，或许在路上受阻。但不要因此烦恼。唯愿你的虔敬大力支持我们，我们相信阴霾将消散，真理将照耀。\par

\SourceRecord{Translated by Blomfield Jackson.；Nicene and Post-Nicene Fathers, Second Series；Vol. 3.；Edited by Philip Schaff and Henry Wace.；Buffalo, NY: Christian Literature Publishing Co.,；1892.；<http://www.newadvent.org/fathers/2707163.htm>.}

% structure: p0001 source=h1 target=section reason=page-title

\section{第一六四封信}

\Emph{同一作者致同一收信人的第二封信，表达在胜利中过早的凯旋。}\par

靠您的圣座的祈祷，我们最虔诚的君王恩准我们觐见；藉天主的恩宠，我们胜过对手，因为我们所有的观点都被最爱基督的皇帝接受了。其他人的报告被宣读，他拒绝了那些他觉得不适合接受且无足轻重的内容。那些报告充满了对\Person{济利禄}的攻击，并请求召他来解释自己的行为。迄今为止他们尚未成功，反而听到了关于真宗教的演讲，即关于信仰体系的演讲，并确认了诸位真福教父们的信仰。我们进一步驳斥了\Person{阿加齐}，因为他在注释中写道天主性是\OriginalTerm{可受难的}{passible}。我们的虔诚君王对如此严重的亵渎感到震惊，以至于扔掉披风，后退几步。我们知道整场集会都欢迎我们作为真宗教的捍卫者。\par

我们最虔诚的君王认为好：任何人都可以解释自己的观点，并报告给他的陛下。我们回答说，除了真福教父们在\Place{尼西亚}所作的阐述外，我们不可能作出任何其他阐述，陛下对此感到满意。因此我们提交了由圣座签署的信式。此外，\Place{君士坦丁堡}全体民众不断出来见我们，恳求我们为信仰勇敢作战。我们尽力约束他们，以免冒犯对手。我们寄了一份阐述的副本，以便制作两份副本，并请圣座在两份上签字。\par

\SourceRecord{Translated by Blomfield Jackson.；Nicene and Post-Nicene Fathers, Second Series；Vol. 3.；Edited by Philip Schaff and Henry Wace.；Buffalo, NY: Christian Literature Publishing Co.,；1892.；<http://www.newadvent.org/fathers/2707164.htm>.}

% structure: p0001 source=h1 target=section reason=page-title

\section{第一六五封信}

\Emph{同一位写给相同的人的信。}\par

致如今在\Place{厄弗所}的各位最虔诚的主教：\Person{约翰内斯}、\Person{希梅里乌斯}、\Person{保禄}、\Person{阿普林吉乌斯}、\Person{狄奥多雷}，问候。\newline{}我们已被准予第五次觐见。我们详细讨论了异端\WorkTitle{十二章}的问题，并一再向最虔诚的皇帝发誓，我们不可能与对手共融，除非他们弃绝那些\WorkTitle{十二章}。此外，我们指出，即使\Person{济利禄}弃绝了他的\WorkTitle{十二章}，我们也无法接纳他，因为他已成为如此不虔敬之异端的异端魁首。尽管如此，我们仍未能取得进展，因为我们的对手急不可耐，他们的听众既无法遏制他们傲慢的企图，也无法强迫他们进行探究和辩论。他们就这样避开了对\WorkTitle{十二章}的审查，不允许对此进行任何讨论。然而，正如你们所恳求的，我们准备坚持到底，视死如归。我们拒绝接纳\Person{济利禄}及其\WorkTitle{十二章}；我们不会让这些人加入共融，除非那些不当的信仰附加物被弃绝。因此，我们恳请你们的圣德继续表明我们的心意和努力。这场战斗是为了真正的宗教；为了我们唯一的希望——借此我们期盼在来世享受救世主的慈爱。至于最虔诚而神圣的主教\Person{聂斯多略}，请你们的虔敬知晓：我们曾试图提及他，但迄今未能成功，因为所有人对他都心怀恶意。尽管如此，我们仍会尽力而为，虽则如此，也要利用任何可能的机会和听众的善意，在天主助佑下实现这一目标。但为使你们的圣德也不忽略这一点，要知道：我们眼见\Person{济利禄}的同党通过专制、欺骗、谄媚和贿赂欺骗了所有人，曾多次恳求最虔诚的皇帝和最尊贵的王子们，既将我们遣返东方，也让你们的圣德回家。因为我们开始认识到，我们是在徒然浪费时间，毫无终点，因为\Person{济利禄}处处逃避讨论，他确信他\WorkTitle{十二章}中发布的亵渎之言可以公开驳斥。最虔诚的皇帝在许多劝诫后决定，我们都各自返回自己的家中，并且此外，埃及人和\Place{厄弗所}的\Person{默农}应留在各自的地方。这样，埃及人将能继续通过贿赂蒙蔽人。那个犯下数不胜数罪行的人将返回他的教区。另一个无辜之人，仅被勉强允许回家。我们和这里所有人向你们和与你们同在的所有弟兄致意。\par

\SourceRecord{Translated by Blomfield Jackson.；Nicene and Post-Nicene Fathers, Second Series；Vol. 3.；Edited by Philip Schaff and Henry Wace.；Buffalo, NY: Christian Literature Publishing Co.,；1892.；<http://www.newadvent.org/fathers/2707165.htm>.}

% structure: p0001 source=h1 target=section reason=page-title

\section{第一六六封信}

\Emph{委员们的首次请愿书，发自\Place{加采东}，呈递给皇帝。}\par

本当希望真宗教之言不被荒谬的解释所玷污，尤其不被那些获得司铎职和在教会中身居高职的人所玷污；他们不知怎地受野心、权力欲和某些可怜的许诺所诱，藐视基督的一切诫命。他们唯一的动机是巴结一个胆大妄为的人，自信他及其同谋者能成功操控整个事务；我指的是\Person{亚历山大里亚的济利禄}。他凭其轻浮侵入了天主的圣教会，提出异端学说，并自以为能用论证支持这些学说。他指望单靠\Person{梅农}和上述阴谋的主教们的帮助，逃脱对罪人的惩罚。\par

我们爱好沉默；通常我们劝人采取哲学的作法。然而，现在我们意识到，保持沉默并修养哲学将意味着抛弃信德，我们便向您——在至高美善之下，世界唯一的保存者——恳求。我们知道，特别属于您的职责是对真宗教的关切，因为您迄今一直保护它，也因它而受保护。\par

我们恳求您接纳这篇论述，如同我们的辩护将要在至圣天主面前陈述；并非因为我们在此神圣事业中不够活跃，而是因为我们忠于真宗教，并为其发言。因为在基督教时代，神职人员的本分莫过于在如此忠信的君王面前作证，尽管我们本已准备好奉献身体，并为信德在战斗中千次舍命。因此，我们藉着洞察万有的天主、将按公义审判万人的主耶稣基督、以及藉着圣神——您因祂的恩宠而执掌帝国——并藉着蒙选的天使——他们是您的守护者，有朝一日您将看见他们站在可畏的宝座前，不停息地向天主献上那如今试图被腐化的可畏的赞颂——我们恳求您的虔敬，您如今正被某些人的诡计所围困，他们阻止您接见，并支持将异端章节引入信德中，这些章节完全违背健全的教义，染有异端色彩。请下令所有赞同这些章节或同意它们，并在您允诺宽恕后仍欲争论的人，都出来服从教会的纪律。先生，没有比为真理而战更配得上皇帝了，为此您曾急忙与波斯人及其他蛮族作战，基督因您向祂的热忱而赐予您辉煌的胜利。我们恳求您将这些争议的问题以书面形式呈于您的虔敬面前，因为这样更容易看出其意图，而违法者将永远被定罪。但若有人不顾他要为之负责的言论，企图以其教导压倒正确的信德，那么您的正义和判断将考虑，那些不愿为他们所引进的教义承担风险、拒绝服从您的命令以逃避因做错事而被定罪的人，是否已将教师的名称丢弃；也不认为他们值得反驳，以免他们的共谋被证明无果。因为现在很清楚，从那些由他们祝圣的人看来，其中一些人为了回报这种不敬，已考虑通过授予尊位来讨好某些人，并设计了其他手段。这将变得更加明朗；您的虔敬将很快看到，他们将像分发基督信德的战利品一样分发他们背叛的酬报。\par

然而我们——其中一些人很久以前由非常虔诚的\Person{尤文纳尔}（\Place{耶路撒冷}主教）祝圣——一直保持沉默，尽管我们应当为教规争辩，以免显得为自己的名声烦恼。我们现在完全知道他在\Place{第二腓尼基}和\Place{阿拉比亚}的活跃诡计。我们实在没有时间顾及这些事。我们宁愿被剥夺托付给我们职位的那些地方，甚至生命，也不愿失去我们为信德的即发热忱。对那些人的企图，我们将以天主的判决和您的虔敬来反对。\par

现在我们也恳求，真宗教可能是您唯一和首要的关切，那在神圣的\Person{君士坦丁}时代艰难地燃起、由您有福的祖父和父亲维持、并由陛下在波斯人和其他蛮族中扩展的正教之光明，不要让它在您皇宫的最内院里变得昏暗，或在您安宁的日子里被驱散。\par

先生，您不会将分裂的基督教送进\Place{波斯}；在这里，我们自己也不会有什么伟大的事，因为我们因争论而苦恼，而且他们那边也没有人存在来解决争论；没有人会参与分裂的圣言和圣事；没有人会在不失信德的情况下与那些从未被谴责的著名教父和圣人断绝关系。没有帝国成功会被允许给一个内部纷争的民族；从真宗教的敌人那里会激起一阵嘲笑；他们恶意争论的所有其他有害后果多得数不胜数。\par

如果有人轻视神学科学，那么让他成为世上任何一个人，而不是主已将世界的最高政府交给他的人。我们的请求是，您的虔敬将作出判断，因为天主将引导您的智力达到精确的理解。最后，如果这不可行（我们无法知道您虔敬的所有事务），我们恳求您安宁地允许我们平安旅行回家。我们知道，如此长期的耽搁给托付给我们的教区带来了冒犯的理由，因为那些人甚至在神圣事务中寻找引起不和的时机，而这些时机不能带来任何益处。\par

\SourceRecord{Translated by Blomfield Jackson.；Nicene and Post-Nicene Fathers, Second Series；Vol. 3.；Edited by Philip Schaff and Henry Wace.；Buffalo, NY: Christian Literature Publishing Co.,；1892.；<http://www.newadvent.org/fathers/2707166.htm>.}

% structure: p0001 source=h1 target=section reason=page-title

\section{\WorkTitle{Letter 167}}

\Emph{同一恳求的第二道上书，从\Place{加采东}呈递给\Person{德奥多西奥·奥古斯都}。}\par

陛下屡次亲自从我们本人以及我们所遣使者处得悉，真信仰的教义似乎面临被败坏的危险，教会的身体显然正被那些颠倒一切、践踏全部教会秩序和所有帝国法律、制造混乱以图巩固\Person{亚力山大的济利禄}所倡异端的人撕裂。因为当陛下首次召集我们前往\Place{厄弗所}，欲查究所生问题并确认由圣教父们所制定的福音与宗徒信仰时，在所有被召集的主教尚未抵达之前，那些自行召开私人会议的人便以书面方式确认了那些与\Person{亚略}、\Person{欧诺米}及\Person{亚波利纳里}不敬神行为一脉相承的异端章节。他们欺骗了一些人，恐吓了一些人，又接纳了早已背负异端指控者进入共融，并以贿赂诱使那些未曾与他们共融的人加入；另一些人或是以不当职位的希望被煽动。于是这些人在身边聚集了一大群追随者，仿佛他们全然不知真宗教并非由人数，而是由真理所彰显。\par

最尊敬的\Person{坎迪狄安}伯爵第二次宣读了陛下的诏书，吩咐以温和、兄弟般的方式审查新近提出的问题。然而，当所有虔诚的主教正在聚集时，这次宣读毫无效果。\par

随后，高贵的\Person{帕拉迪乌斯·玛吉斯特里亚努斯}带来陛下另一道诏书，命令所有私下、单独通过的决议必须撤销；必须重新召集会议并确认真道；但一如往常，这道圣旨被这些毫无顾忌的人所藐视。\par

接着，最尊贵的\Person{若望}大人——时任\Quote{Comes Largitionum}——到来，又带来一道圣谕，声称已决定废黜那三位（主教），要铲除所滋生的罪行，并令全体确认由神圣而蒙福的教父们在\Place{尼西亚}所制定的信仰。这些普世嘲弄者一如往常，也违反了这项法令。\par

因为他们在听了这道信件之后，并未改变行径；他们与被废黜者共融，仍称他们为主教，并拒绝摒弃那些已被提出、有损并败坏虔诚信仰的章节——尽管我们曾多次召唤他们前来讨论。因为我们手头已有对这些异端章节的清晰驳斥。\par

作为这些说法的证据，我们有最尊贵的\Person{若望}大人；当他第三次、第四次传唤双方，却不敢以他们的不顺服为借口来开脱这种行为时，他认为值得召唤我们到此处来。\par

我们立刻前来；抵达后我们一刻不曾歇息，在陛下和显赫的议会面前呈上我们的恳求，请求他们为这些章节展开争论并加以讨论，或者相反地，因其违背正道信仰而予以拒绝，坚守由蒙福教父们在\Place{尼西亚}会议上所制定的信仰。\par

他们拒绝做任何此类事；他们坚持其异端行径；然而他们却被允许进入圣堂并履行司祭职。而我们，无论在\Place{厄弗所}还是在此处，早已被剥夺了共融；在那里和这里，我们都经历了数不尽的危险；当被装扮成隐修士的奴隶投石险些丧命时，我们把这一切视为最好的——因为我们为了真理而甘愿承受这样的待遇。\par

后来，陛下认为我们与反对派应当再次聚集，以便强制那些抗拒者审查教义。当我们等候此事发生之时，陛下起驾前往京城，并下令那些被指控异端、因而其中一些人已被我们废黜、另一些被开除教籍并随后须受教会纪律处分的人，进入京城履行司祭职并进行圣秩授任。然而我们这些为了真宗教进行如此巨大斗争、在为正统教义而战中不避任何危险的人，既未被召入京城以服务于危殆的信仰并力争正统，也未被允许返家；而是一直困在\Place{加采东}，为教会受分裂之苦而悲伤叹息。\par

因此，由于我们没有收到任何答复，我们认为有必要以本信函在天主、基督和圣神面前禀告陛下：若有人在正统教义解决之前，由这些持有异端意见的人之中任何一位被授任，他必然应从整个教会中被剪除——无论是圣职界还是与之不合的平信徒都如此。因为没有一个虔诚之人会容忍与异端者共融，从而废止自己的救恩。\par

当此事发生时，陛下将被迫违背自己的意愿行事。因为分裂将超乎预期地扩大，从而使真宗教的捍卫者忧心忡忡，无法忍受自己灵魂的失落以及\Person{济利禄}那些不敬神的教义——即争辩者所渴望捍卫的——的确立。\par

事实上，许多真宗教的支持者绝不接受\Person{济利禄}的教义；我们——包括整个东方教省、本都教省、亚细亚教省、色雷斯教省、伊利里库姆教省以及意大利诸教省的信众，还有曾向陛下呈上最蒙福的\Person{盎博罗削}为反对这种新生迷信所写论著的那些人——也将绝不允许。\par

为了避免这一切以及进一步烦扰陛下，我们恳求、祈求并哀求您颁布诏令：在因陛下爱基督之崇高而被召集的正统信仰解决之前，不得举行任何圣秩授任。\par

\SourceRecord{Translated by Blomfield Jackson.；Nicene and Post-Nicene Fathers, Second Series；Vol. 3.；Edited by Philip Schaff and Henry Wace.；Buffalo, NY: Christian Literature Publishing Co.,；1892.；<http://www.newadvent.org/fathers/2707167.htm>.}

% structure: p0001 source=h1 target=section reason=page-title

\section{书信第168号}

\Emph{来自同一人的第三次请求，发自加采东，致君王。}\par

我们从未期望您虔诚的召集会有这样的结果。我们被尊荣地召集，作为司铎受君王之召；我们被召集来批准圣教父们的信德。因此，出于对虔敬君主的适当服从，我们来了。到达之后，我们对教会同样忠实，对您的谕令同样恭敬。从我们到达厄弗所的那一天直到此刻，我们无间断地遵从了您的命令。\par

然而，看来在这时候，我们的节制对我们毫无益处；非但如此，就我们所见的，反而大大阻碍了我们。我们这些如此行事的人，至今仍被拘留在加采东；而现在我们被告知可以回家。但那些把一切搅乱、使世界充满骚动、企图将教会撕裂为二、公开攻击真宗教的人，却执行司铎职务，挤满教堂，并按他们所想拥有祝圣权——其实这权柄是被他们非法窃取的——在教会中煽动叛乱，而且本应用于穷人的钱财，他们却花在他们的打手上。\par

但您不但是他们的皇帝，也是我们的皇帝。因为您帝国中不小的部分是东方，那里正信始终照耀着，此外还有其他我们被召集来的省份和教区。\par

愿陛下不要轻视那正在被败坏的信德，您和您的先祖们就是在这信德中受洗的；教会的根基是建立在这信德上的；最圣的殉道者们正是为这信德而喜乐地承受了无数种死亡；借着这信德，您战胜了蛮族，消灭了暴君；如今在您征服阿非利加的战争中，您也需要这信德。因为如果您为祂神圣的教义而奋战，并禁止教会身体的肢解，那么万有的天主必会与您并肩作战；因为如果基利罗引入教会、其他异端者加以确认的意见得势，教会就必被肢解。\par

对于这些真理，我们在厄弗所和此地已多次在天主面前作证。我已向您的圣善提供了信息，如同在万有的天主面前陈明。因为这是我们被要求的，正如在神圣圣经中，先知和宗徒们都教导的；正如真福保罗所说：\BibleQuote{我在天主面前，并在那将要复活死者的主耶稣基督面前，以及那在本丢·彼拉多面前作了美好见证的主耶稣基督面前，嘱咐你：}又如天主命令厄则克耳向百姓宣告，并加上警告说：\BibleQuote{如果你不警告他，……我必向你追讨他的血。}\par

因着这句话的敬畏，我们再次禀告陛下：那些被允许占有教堂、并教导阿波利纳里、亚略和欧诺弥学说的人，违规范地、违反教规地举行一切神圣礼仪，并毁灭所有接近他们者的灵魂——如果真有人愿意听他们的话。因为靠着天主的恩宠，祂眷顾一切，并愿意所有的人都得救，大多数人民是纯正的，并热切依恋虔敬的教义。正是为了这些人，我们忧伤。\par

在我们的痛苦和恐惧中，唯恐那瘟疫渐渐蔓延，攻击更多的人，使邪恶变成普遍，我们如此教导您的安宁，并继续劝诫您；我们恳求陛下听从我们的祈祷，禁止对在尼西亚开会的圣教父们的信德作任何添加。\par

如果在我们恳求之后，您的虔诚拒绝这在天主面前赐下的教义，我们就要向你们跺掉脚上的尘土，并和真福保罗一同呼喊：\BibleQuote{我们的手是洁净的，不沾你们的血。}因为我们从到达这次著名会议的那一刻起，日夜不停地向君王、贵族、士兵、司铎和百姓作证：我们坚守圣教父们传交给我们的信德。\par

\SourceRecord{Translated by Blomfield Jackson.；Nicene and Post-Nicene Fathers, Second Series；Vol. 3.；Edited by Philip Schaff and Henry Wace.；Buffalo, NY: Christian Literature Publishing Co.,；1892.；<http://www.newadvent.org/fathers/2707168.htm>.}

% structure: p0001 source=h1 target=section reason=page-title

\section{第一六九封书信}

\Person{狄奥多勒特斯}，\Place{居鲁士}主教，自\Place{迦克墩}致\Place{希拉波里}的\Person{亚历山大}。\par

我们已用尽一切手段：礼遇、严厉、恳求、雄辩，向最虔敬的皇帝和显赫的议会陈情，在那洞察万有的天主、将以正义审判世界的主耶稣基督、圣神及其选民天使面前作证，唯恐这信德如今被那些维护并大胆签署异端教义者所败坏，并请求它以与尼西亚相同的措辞被订立，且摒弃那引入进来、导致真正宗教之损失与毁灭的异端。然而迄今未产生丝毫效果，因听众时而倾向一方，时而倾向另一方。\par

尽管如此，所有这些困难并未能阻止我坚持己见，而是靠天主的恩宠，我继续奋进。我甚至向最虔敬的皇帝起誓断言，\Person{西利尔}和\Person{梅农}与我绝无和好可能，且我们绝不会与任何不预先弃绝异端章节的人共融。这是我们的心意。那些\Quote{各人只求自己的事，并不求耶稣基督的事}的人，其目的是要违背我们的意愿与他们和好。但这与我无关，因为天主衡量我们的动机，考验我们的品格，且不会对违背我们意愿的行为施以惩罚。敬告圣座：若我曾当着最虔敬的皇帝或显赫的议会直言关于我们这位朋友的事，便立即被指责为叛乱者。宫廷党对他深恶痛绝。这最令人烦恼。尤其最虔敬的皇帝，简直无法忍受听到他的名字，且公开说：\Quote{谁也不可向我提及此人。}有一次他曾向我举例说明。然而，只要我在此地，我不会停止为我们这位神父的利益服务，因我知道不敬虔者冤枉了他。\par

我渴望圣座和我自己都能摆脱此事。从中无善可期，因所有法官都信靠金钱，并主张天主神性及人性本质为一。\par

然而全体民众靠天主的恩宠情况良好，不断出来见我们。我已开始向他们讲道，并举行了人数众多的共融礼。\par

第四次我详细论述了信德，他们如此喜悦地聆听，竟至第七时辰才离去，甚至忍受了正午的炎热。在一个带有四个柱廊的大庭院中聚集了庞大的人群，我从靠近屋顶的平台上方讲道。\par

相反，所有圣职人员连同杰出的隐修士们都极力反对我，以致当我们从\Place{鲁菲尼亚努姆}返回，即最虔敬的皇帝来访之后，投石开始，我的许多同伴被民众和假隐修士打伤。\par

最虔敬的皇帝知道暴民聚集反对我，便独自走向我，说：\Quote{我知道你们在不当聚集。}我于是说：\Quote{既然您允许我说话，请仁慈地听我说。被绝罚的异端者得以在教堂中履行职责，而我，为信德奋战，因此被他人排除在共融之外，却不得进入教堂，这公平吗？}他答道：\Quote{我该怎么办？}我说：\Quote{您的\OriginalTerm{财税长官}{comes largitionum}在以弗所所做的：当他发现有些人在聚集，而我们没有聚集时，他阻止他们说：\InnerQuote{如果你们不和平，我不会允许任何一方聚集。}圣座也本应指示这里的主教，在我们集会向众人宣示您公正的判决之前，禁止对立派和我们双方聚集。}他对此答道：\Quote{这不是我该命令主教的。}我回答：\Quote{您也不可命令我们，我们将占一座教堂，并聚集。圣座将会发现我们这边的人数比他们多得多。}此外，我指出我们既没有圣经诵读，也没有献祭，只有\Quote{为信德和陛下的祈祷，以及虔诚的交谈。}于是他同意了，不再禁止。结果，更多的人群涌向我们，并欣然聆听我们的教导。因此我恳请圣座祈祷，愿我们的事有一个悦乐天主的结局。我每日处于危险中，怀疑隐修士和圣职人员的诡计，因为我亲眼目睹他们的影响力和疏忽。\par

\SourceRecord{Translated by Blomfield Jackson.；Nicene and Post-Nicene Fathers, Second Series；Vol. 3.；Edited by Philip Schaff and Henry Wace.；Buffalo, NY: Christian Literature Publishing Co.,；1892.；<http://www.newadvent.org/fathers/2707169.htm>.}

% structure: p0001 source=h1 target=section reason=page-title

\section{第一七〇号书信}

\Emph{若干东方人士——被派往君士坦丁堡者——致主教鲁富斯的信。}\par

\Signature{致我们至为虔诚圣洁的同工鲁富斯主教，若望、希默里乌斯、狄奥多雷及其他同仁在主内问安。}\par

真宗教与教会的和平，我们认为，因你的圣德不在场而受损非小。倘若你身在现场，你本可止息所生的骚乱和所敢行的暴力，并可为压制那些引入正统信仰中的异端，以及那宗徒与福音的教义——该教义从先祖代代相传，最终传至我们——而与我们并肩战斗。\par

然而，议会中的许多成员，借用先知的话，\Quote{已经偏离了正道}，而且\Quote{全都变得污秽}，因为他们放弃了从圣教父们领受的信仰，并签署了\Person{亚历山大的济利禄}的\WorkTitle{十二章}，这些章节充满阿波利拿里错误，与\Person{亚略}和\Person{欧诺米}的亵渎一致，并诅咒所有不接受其不加掩饰的异端的人。对于这场打击教会的瘟疫，我们这些从东方聚集的人，以及来自不同教区的其他人，进行了有力的抵抗，目的是确保受祝福的圣教父们在\Place{尼西亚}会议上传递的信仰得到批准。因为，如你的圣德所知，其中既不缺福音教义的教导，也不缺对任何异端的驳斥。\par

为了这信仰，我们继续战斗，轻视现世生命的一切喜乐与忧愁，只求我们能完好保存这祖传遗产。为此，我们已罢免了\Person{济利禄}和\Person{梅姆农}；前者是异端的主要推动者，后者是他在一切事上的帮手和同谋，用以批准并维护那些为毁灭教会而公布的\WorkTitle{十二章}。我们也已逐出教会所有敢于签署和支持这些不敬教义的人，直到他们诅咒这些教义，并回到\Place{尼西亚}圣教父们的信仰。\par

但我们的忍耐并未给他们带来益处。直到今日，他们继续为那些有害的教义战斗，并把自己钉在教规的法律上，该法律明文规定：\Quote{如果任何主教被会议罢免，或任何司铎或执事被其本主教罢免，仍未等待会议判决而执行圣职，他就没有机会为自己辩护，即使在另一个会议上也不行；而且所有与他共融的人也都要被逐出教会。}如今，这项法律已被被罢免者和被逐出教会者所破坏。因为在得知被罢免和被逐出教会后，他们立即执行了圣职，并且继续这样做，明显不信那说过\BibleQuote{凡你们在地上所束缚的，在天上也要被束缚}的一位。\par

我们曾认为最好立即将此事告知你的圣德，但期待某种有利的改变，我们一直等待至今。然而我们失望了。他们继续为这不敬的异端战斗，完全不理会十分虔诚的皇帝的劝告。他五次会见我们，命令他们要么拒绝\Person{济利禄}的\WorkTitle{十二章}，因其违反信仰，要么愿意为他们辩护，并表明它们如何与圣教父们的信条一致。我们手头有证据，本可证明它们完全反对正统教导，且大部分与异端一致。\par

因为正是在这些\WorkTitle{十二章}中，这些有害作品的作者教导说，独生子的天主性受苦了，而不是祂为了我们的救恩所摄取的人性受苦；那寓居的天主性明显地将其自身的身体所受的苦归为己有，尽管其本性上并无受苦；此外，还认为天主性与人性成了一个本性——因为他是这样解释\BibleQuote{圣言成了血肉}，仿佛天主性经历了某种变化，变成了血肉。\par

此外，他还诅咒那些在宗徒和福音作者关于主基督的用语之间作出区分的人，将那些屈辱的用语归于人性，而将那些神圣荣耀的用语归于主基督的天主性。正是由于这些观点，亚略派和欧诺米派将屈辱的用语归于天主性，竟毫不畏惧地宣称天主圣言是被造成和被造的，属于不同实体，且与圣父不相似。\par

从这些陈述中会引出何等的亵渎，并不难看出。一种本性的混淆被引入，而将以下话语应用在天主圣言身上：\BibleQuote{我的天主，我的天主，你为什么舍弃了我？}以及\BibleQuote{父啊！若是可能，就让这杯离开我吧}，还有饥饿、口渴和天使加力；祂说：\BibleQuote{现在我心神烦乱}和\BibleQuote{我的心灵忧闷得要死}，以及所有类似属于基督人性的篇章。任何人都能看出这些陈述与\Person{亚略}和\Person{欧诺米}的亵渎如何相符；因为他们无法确立实体的差异，就把苦难和屈辱的用语，如上所述，联系到基督的天主性上。\par

请你的尊驾确信，如今在他们的教会中，亚略派教师所传讲的教义不外乎是：那些支持\OriginalTerm{同本体}{homousion}的人目前持有与\Person{亚略}相同的观点，而且经过漫长岁月后，真理如今终于被揭示出来了。\par

反之，我们持守着在\Place{尼西亞}集会的真福教父们以及他们杰出的继任者——\Place{安提约基雅}的\Person{尤斯塔提乌斯}、\Place{凯撒勒雅}的\Person{巴西略}、\Person{额我略}、\Person{若望}、\Person{亚太纳削}、\Person{德奥斐禄}、\Place{罗马}的\Person{达玛苏斯}、\Place{米兰}的\Person{盎博罗削}——的教导，并追随他们虔诚的足迹。因为所有这些人都追随宗徒的言语，为我们留下了正统信仰的精确准则，我们所有东方人都热切希望保持不变。同样的愿望也属于\Place{彼提尼亚}人、\Place{帕夫拉戈尼亚}人、\Place{第二卡帕多细亚}、\Place{丕西狄亚}、\Place{米西亚}、\Place{特萨利亚}、\Place{洛多培}以及更多不同省份的人。\Place{意大利}人显然也不会容忍这种新奇的教义；因为非常虔诚和圣洁的\Place{米兰}主教\Person{玛尔提努斯}给我们写了一封信，并将真福\Person{盎博罗削}关于主降生成人的著作呈送给了非常虔诚的皇帝，该著作的教导与这些异端章节相对立。\par

并且敬请圣座知晓：\Person{济利禄}和\Person{默农}不满足于败坏正统信仰，还践踏了所有教规。因为他们接纳了在各地区各教区被绝罚的人进入共融。另有人身负异端指控，与\Person{塞勒斯提乌斯}和\Person{佩拉纠}同一心思（因为他们是\OriginalTerm{优奇忒}{Euchitæ}或\OriginalTerm{热忱者}{Enthusiasts}），因此被其教区和教省总主教绝罚，他们却无视所有教会纪律，接纳这些人进入共融，从而从一切可能的方面壮大其追随者，显示出他们更倾向于用暴力而非虔诚来推行其教导。因为当他们在虔诚上被剥光时，他们走投无路，便发明了另一种暴力——人肉之墙，以为通过贿赂的倾泻就能战胜教父们的信仰。但只要圣座施展力量，并如常为保卫真宗教而战斗，这些伎俩将毫无用处。因此，我们恳请最神圣的阁下，警惕这异端无耻引入者的共融；并告知远近所有人：这正是三倍真福的\Person{达玛苏斯}因之而罢免异端者\Person{阿波利纳里}、\Person{维塔利}和\Person{提摩泰}的要点；并且那封作者隐藏异端并用真理的涂层粉饰的信函，不可轻率地接受。因为在该章节中，他大胆暴露了自己的不虔敬，并敢诅咒所有不同意他的人，而在信中他却卑鄙地试图伤害较为单纯的读者。\par

因此，圣座必须当心不可忽视此事，以免当你们为时已晚地看到这异端被巩固时，徒然哀叹，并因再也无法保卫真理事业而受苦。\par

我们还给您寄去了一份我们呈递给最虔诚和爱基督的皇帝的备忘录副本，其中包含了\Place{尼西亞}圣教父们的信仰。我们在其中拒绝了\Person{济利禄}新发明的异端，并判定它们与正统信仰相对立。\par

由于遵照最虔诚皇帝的命令，我们只有八人前往\Place{君士坦丁堡}，我们附上了神圣主教会议给我们的命令副本，以便您了解其中所包含的省份。圣座将从教省总主教的署名中获悉这些省份。我们问候与您在一起的弟兄们。\par

\SourceRecord{Translated by Blomfield Jackson.；Nicene and Post-Nicene Fathers, Second Series；Vol. 3.；Edited by Philip Schaff and Henry Wace.；Buffalo, NY: Christian Literature Publishing Co.,；1892.；<http://www.newadvent.org/fathers/2707170.htm>.}

% structure: p0001 source=h1 target=section reason=page-title

\section{第一七一书}

\Emph{\Person{狄奥多雷}在和解后致\Place{安提约基雅}主教\Person{若望}的书信。}\par

天主以智慧治理万有，顾念我们的一致，并关切祂子民的救恩，祂使我们聚集在一起，并表明我们众人的看法彼此相符。我们聚集一处，阅读了埃及书信；我们仔细审察了其要旨，发现其内容完全符合我们自己的陈述，且全然反对那十二章，直至今日我们仍持续与之争战，视其为违反真宗教。他们的教训是：天主圣言成了血肉之躯；有位格的合一，且合一中的结合是本性上的；天主圣言是从死者中首先复生的。他们禁止在论及我们的主时所使用的术语有任何区别，此外还包含其他与宗徒所播种的种子相悖、出于异端稗子的道理。然而，这封书信却因宗徒般的高贵起源而增色。因为其中显示我们的主耶稣基督为完美的天主和完美的人；它表明两个本性及其区别；一种不相混合的合一，不是通过掺和与混合而成，而是以不可言喻和神圣的方式，明确保存各本性的属性；天主圣言的不受苦难与不死；圣殿的可受苦难与暂时屈服于死亡，以及藉着合一的天主的能力而复活；圣神并非出自圣子，也不是由圣子获得存在，而是由圣父所发，且按理被称为属于圣子，因其为同一实体。看到这书信中的正统信仰，我们颂扬那位治愈我们口吃舌结，并使我们不和谐的声音变为美妙音乐的和谐的主。\par

\SourceRecord{Translated by Blomfield Jackson.；Nicene and Post-Nicene Fathers, Second Series；Vol. 3.；Edited by Philip Schaff and Henry Wace.；Buffalo, NY: Christian Literature Publishing Co.,；1892.；<http://www.newadvent.org/fathers/2707171.htm>.}

% structure: p0001 source=h1 target=section reason=page-title

\section{\WorkTitle{第一七二封书信}}

\Emph{狄奥多莱致聂斯多略的书信}\par

致极可敬、虔诚的主、至圣的教父\Person{聂斯多略}，主教\Person{狄奥多莱}在主内致敬。我想阁下深知，我既不为繁华世故所喜，也不关心此生的利益与名誉，更不受其他教区的吸引。即使没有从其他来源学到这一课，单单我所治理的这座城市的孤寂，已足以教导我这套哲学。这座城不仅以孤寂闻名，也多有纷扰，足以抑制那些最喜活动之人的热情。\par

因此，不要让任何人说服阁下，说我因垂涎任何教区，便闭眼接受埃及著作视为正统。实在说，实话实说，我在多次阅读并仔细审查之后，发现它们毫无异端污染，但我曾犹豫是否重视它们，尽管我确实不喜爱其作者，他是搅乱世界的骚乱的始作俑者。为此，我希望在审判之日得以免罚，因为公义的审判者审察动机。至于对阁下所做的不义不法之事，即使有人要砍断我的双手，我也决不赞同，靠天主的恩宠帮助我并扶持我的软弱。我已将这一点以书面形式告知了那些要求我这样做的人。我已将我对您来信的答复寄给阁下，以便您知道，靠着天主的恩宠，时间从未改变我，如同那些蜈蚣和变色龙，它们以周围的石头和树叶的颜色来模仿。我和所有与我同在的人，在主内问候您身边的所有弟兄。\par

\SourceRecord{Translated by Blomfield Jackson.；Nicene and Post-Nicene Fathers, Second Series；Vol. 3.；Edited by Philip Schaff and Henry Wace.；Buffalo, NY: Christian Literature Publishing Co.,；1892.；<http://www.newadvent.org/fathers/2707172.htm>.}

% structure: p0001 source=h1 target=section reason=page-title

\section{第一七三封信}

\Emph{致\Place{君士坦丁堡}隐修士\Person{安德肋}的书信。}\par

\BibleQuote{天主是忠信的，他决不许你们受试探超过你们的能力；相反，他必在试探中为你们开一条出路，使你们能承担。}并且驳斥了谎言——尽管现在已被驳倒的谎言断言得到认可——而真理的力量已经显示。因为，看，那些曾以不虔敬的推理混淆了我们救主基督的本性，竟敢宣讲一个本性，因此侮辱了至圣可敬的天主的大司祭\Person{聂斯多略}的人，他们的口被嚼环和辔头勒住，正如先知所说，从错谬转向正义，已再次学习真理，采纳了那位在真理事业中承受了战斗主力的人的声明。因为，他们现在宣认两个本性，而非一个本性，绝罚所有宣讲混合与混淆的人。他们朝拜基督不受苦难的神性；他们将苦难归于肉体；他们区分福音书中的词语，将崇高和神性归于神性，将卑微归于人性。这些就是从\Place{埃及}带来的文献。\par

\SourceRecord{Translated by Blomfield Jackson.；Nicene and Post-Nicene Fathers, Second Series；Vol. 3.；Edited by Philip Schaff and Henry Wace.；Buffalo, NY: Christian Literature Publishing Co.,；1892.；<http://www.newadvent.org/fathers/2707173.htm>.}

% structure: p0001 source=h1 target=section reason=page-title

\section{书信第一七四号}

\Emph{致\Place{尼科美底亚}主教\Person{希梅里乌斯}。}\par

我们愿通知尊座，在阅读并多次讨论从埃及带来的信函后，我们发现它与教会的道理相符。至于那十二章，我们已证明其相反之处，并且至今仍继续反对它们。因此，我们决定，如果尊座已收复了天主托付于你的各教会，你应当与埃及人、君士坦丁堡人以及曾与我们一同对抗他们的其他人共融，因为他们已宣示持守我们的信德，或更确切地说，持守宗徒的信德；但不要同意对至圣可敬的\Person{聂斯脱利}所声称的定罪。因为我们视之为不虔不义，在双方同为被告的案件中，对一人施以恩惠，却对另一人关闭悔改之门。尤其不义不虔的是将无辜者判处死刑。尊座应确信，在收复各教会之前，你不应与他们共融。因为此事不仅我，连我们教区的所有圣主教也在最近的会议上议决了。\par

\SourceRecord{Translated by Blomfield Jackson.；Nicene and Post-Nicene Fathers, Second Series；Vol. 3.；Edited by Philip Schaff and Henry Wace.；Buffalo, NY: Christian Literature Publishing Co.,；1892.；<http://www.newadvent.org/fathers/2707174.htm>.}

% structure: p0001 source=h1 target=section reason=page-title

\section{第一七五封信}

\Emph{致\Place{希拉波利}的\Person{亚历山大}}。\par

我已通知了你的圣座：若非常神圣且可敬的主教、我的主\Person{聂斯托利}的教义受到谴责，我将不与那些这样做的人共融。如果这能使你的圣座乐意被插入正在送往\Place{安提约基雅}的信中，那就这样吧。那么，我恳求你，不要再耽搁了！\par

\SourceRecord{Translated by Blomfield Jackson.；Nicene and Post-Nicene Fathers, Second Series；Vol. 3.；Edited by Philip Schaff and Henry Wace.；Buffalo, NY: Christian Literature Publishing Co.,；1892.；<http://www.newadvent.org/fathers/2707175.htm>.}

% structure: p0001 source=h1 target=section reason=page-title

\section{第一七六封书信}

\Emph{致同一\Person{亚历山大}，在他得知安提约基雅主教\Person{若翰}\OriginalTerm{绝罚了}{anathematized}\Person{聂斯多略}的学说之后。}\par

敬请尊驾明鉴：当我阅读致皇帝的书信时，我非常忧苦，因为我深知该信的作者因怀同样见解，竟不明智且不虔敬地谴责了一位从未持有或教导任何违反纯正道理之人。然而，那\OriginalTerm{绝罚}{anathema}的格式虽比他对谴责的赞同更易令读者忧伤，却仍给我些许安慰，因它并非以宽泛笼统之词陈述，而加以条件限制。因为他并未说\Quote{我们绝罚他的学说}，而是说\Quote{凡他所说或所持，超出宗徒道理所准许的以外}。\par

\SourceRecord{Translated by Blomfield Jackson.；Nicene and Post-Nicene Fathers, Second Series；Vol. 3.；Edited by Philip Schaff and Henry Wace.；Buffalo, NY: Christian Literature Publishing Co.,；1892.；<http://www.newadvent.org/fathers/2707176.htm>.}

% structure: p0001 source=h1 target=section reason=page-title

\section{书信第177封}

\Emph{致\Place{萨莫萨塔}主教\Person{安德肋}书}\par

显贵的\Person{阿里斯托劳斯}从\Place{埃及}派遣\Person{玛吉斯特里亚努斯}携\Person{济利禄}的一封信到来，他在信中绝罚\Person{亚略}、\Person{欧诺弥}、\Person{阿波林纳}及所有主张基督的天主性可受苦、并坚持两性混淆与混合的人。对此我们感到欢喜，尽管他未赞同我们的声明。他要求进一步签署已通过的谴责，并要求绝罚神圣主教\Person{聂斯多略}的道理。您的圣德深知，若有人不加区分地绝罚那位至圣可敬主教的道理，这无异于他似乎绝罚了真宗教。\par

那么，我们若被迫，必须绝罚那些称基督纯属凡人、或将我们唯一的主耶稣基督分作两个儿子并否认其天主性的人，等等。\par

\SourceRecord{Translated by Blomfield Jackson.；Nicene and Post-Nicene Fathers, Second Series；Vol. 3.；Edited by Philip Schaff and Henry Wace.；Buffalo, NY: Christian Literature Publishing Co.,；1892.；<http://www.newadvent.org/fathers/2707177.htm>.}

% structure: p0001 source=h1 target=section reason=page-title

\section{第一七八封书信}

\Emph{致\Place{希埃拉波利斯}的\Person{亚历山大}书。}\par

我认为，超过一切，至圣可敬的主教、我的主\Person{若望}必定因我拒绝同意谴责至圣可敬的主教\Person{聂斯脱利}，也拒绝违背在\Place{塔尔索}、\Place{加采东}和\Place{厄弗所}所作的承诺而感到欣慰。\par

他还记得我们离开后，在\Place{安提约基雅}经常从我们这里收到的内容。\par

因此，不要让任何人欺骗圣座，使您相信我会这样做，因为天主无疑站在我一边，并坚固我。\par

\SourceRecord{Translated by Blomfield Jackson.；Nicene and Post-Nicene Fathers, Second Series；Vol. 3.；Edited by Philip Schaff and Henry Wace.；Buffalo, NY: Christian Literature Publishing Co.,；1892.；<http://www.newadvent.org/fathers/2707178.htm>.}

% structure: p0001 source=h1 target=section reason=page-title

\section{第一七九号书信}

\Emph{济利禄致安提约基雅主教若望书，反对德奥多勒。}\par

\SourceRecord{Translated by Blomfield Jackson.；Nicene and Post-Nicene Fathers, Second Series；Vol. 3.；Edited by Philip Schaff and Henry Wace.；Buffalo, NY: Christian Literature Publishing Co.,；1892.；<http://www.newadvent.org/fathers/2707179.htm>.}

% structure: p0001 source=h1 target=section reason=page-title

\section{第一百八十号书信}

\Emph{据某些人推测，此为德奥多勒致安提约基亚主教多姆诺斯的书信，写于亚历山大主教济利禄去世之时。}\par

那恶徒终于费尽周折走了。善人和温和者逝去得太快；恶人却延长寿命多年。\par

我以为，万善的施予者使前者提前脱离人世的烦扰；祂像竞赛的胜利者一样解放他们，并将他们迁往更好的生命，那无死亡、无悲伤、无挂虑的生命，是操持德性者的奖赏。反之，那些喜爱和行恶的人，被允许稍长久地享受此世生活，或为了让他们饱尝邪恶之后学习德性的功课，或为了在此世为他们自身的邪恶道路付出惩罚，在多年此世悲伤与邪恶的波浪中颠簸。\par

然而，这个恶徒并没有像其他人那样被我们灵魂的统治者释放，以便他能更长久地拥有那些看似充满喜乐的事物。主知道这人的恶意每日增长，并损害教会肢体，便像瘟疫一样将他砍掉，\BibleQuote{从以色列身上除去了羞辱。}他的幸存者确实为他的离去而高兴。死者或许感到悲伤。有某种担忧的根据，恐怕他们如此厌烦他的陪伴，以至于把他送回给我们，或者他像路喀阿诺斯所著中库尼司科斯的暴君一样逃离他的护送者。\par

因此必须极其小心，这尤其是圣座的责任，要告诉殡殓者行会在他的坟墓上放置一块非常大且重的石头，以免他再回来，再次显露出他反复无常的心。让他把他的新教义带到阴间去吧，在那里昼夜宣讲。我们丝毫不怕他通过发表反对真宗教的公开演说，以及赋予不朽本性以死亡，来分裂他们。他不仅会被熟谙神圣法律的幽灵用石头砸死，还会被纳默罗德、法郎、散乃黑黎布，或其他任何天主的敌人用石头砸死。\par

但我是在浪费言语。这可怜人已缄默无言，无论他愿不愿意，\BibleQuote{他的气息一断，归于灰土，当日的计划尽都消亡。}他也注定要承受另一种沉默。他的行为被揭露，绑住了他的舌头，塞住他的口，遏制他的激情，使他哑口无言，并使他俯首在地。\par

我真为这可怜人感到难过。说实话，他去世的消息并未使我毫无混杂的喜悦，而是带有悲伤。看到教会摆脱这种瘟疫，我欢喜快乐；但想到这个恶徒直到死时也没有从他的罪行中得到安息，反而继续尝试更大更严重的事，我便悲伤哀恸。据说，他的想法是第二次攻击真教义，从而搅乱帝都，并指控圣座支持这些教义。但天主看见了，并未忽视。\BibleQuote{将钩子放进他的鼻子，将辔头戴在他的嘴上}，并使他归于所出之土。愿圣座的祈祷使他获得怜悯与慈悲，愿天主无边的宽仁超过他的邪恶。我恳求圣座驱散我灵魂的焦虑。许多不同的报告在四处流传，令我惊恐，预告普遍的不幸。甚至有人说，圣座在违背己意的情况下出发去朝廷，但迄今为止，我视这些报告为不实而轻视。但发现每个人都在重复同一个故事，我认为应当尽力从圣座那里了解真相，以便若为假，我可以嘲笑这些说法；若为真，则不至于无故悲伤。\par

\SourceRecord{Translated by Blomfield Jackson.；Nicene and Post-Nicene Fathers, Second Series；Vol. 3.；Edited by Philip Schaff and Henry Wace.；Buffalo, NY: Christian Literature Publishing Co.,；1892.；<http://www.newadvent.org/fathers/2707180.htm>.}

% structure: p0001 source=h1 target=section reason=page-title

\section{第一八一封书函}

\Emph{致\Place{科莫}主教\Person{阿本狄}书}\par

给我的亲爱的主内至圣兄弟\Person{阿本狄}，\Person{德奥多勒}在主内致候。我已发现你的虔敬虔诚地保存了真实和宗徒传承的信德；我感谢全能天主，因为濒临危险的真理由你的圣德得以更新并显扬。\par

古时，洪水之后，\Person{诺厄}和他的儿子们被留下作人类的种子。同样，在我们的时代，西方的教父们被保存下来，好使东方的圣教会能藉着他们保存那被一个邪恶的新异端所威胁、快要遭蹂躏和毁灭的真宗教。我们很可以引用先知的话：\BibleQuote{万军的上主若不是给我们留下了少许残余，我们早已像\Place{索多玛}，早已像\Place{哈摩辣}了。}这样，从这邪恶异端而来的天主的愤怒像洪水猛兽般落在我们身上。\par

现在我们承认我们的救主带着人性身体，以及天主子一位，祂有完美的天主性和完美的人性。我们不将我们唯一的主耶稣基督分成两个儿子，因为祂是一位；但我们认出天主与人的区别；我们知道一位出自圣父，另一位出自\Person{达味}和\Person{亚巴郎}的后裔，按照圣经，并且天主性不受苦难，那先前受苦的身体如今也脱离了苦难；因为复活后，它显然已脱离一切苦难。\par

这是我们从至圣且虔诚的总主教我主\Person{良}那里学到的。因为我们读了他写给神圣而蒙福的\Person{弗拉维安}的信，并感谢主的慈爱，因为我们找到了一位真理的辩护者和捍卫者。对这封信，我表示赞成，并将它的副本附在我现在的信函后，我也签了名，从而证明我遵守宗徒规条，即真实的教义；我至今持守它们，并因它们而受苦。\par

我的主\Person{依巴斯}和我的主\Person{阿奎利诺}也表示同意；新异端的创立者曾借助帝国权力武装起来反对他们。\par

剩下的就是，你与你的至圣同僚们救助圣教会，驱散威胁她的战争。驱逐那被鼓动起来反对真理的邪恶党派；将古老的和平归还各教会；这样，你们将从上主——祂应许赐予这恩惠——那里得到你们宗徒劳作的果实。\par

所有至为虔诚和敬主的司铎、可敬的执事和弟兄们，我藉你的圣德问候他们；我和我这里的全体人员问候你的尊威。\par

\SourceRecord{Translated by Blomfield Jackson.；Nicene and Post-Nicene Fathers, Second Series；Vol. 3.；Edited by Philip Schaff and Henry Wace.；Buffalo, NY: Christian Literature Publishing Co.,；1892.；<http://www.newadvent.org/fathers/2707181.htm>.}
